%%%%%%%%%%%%%%%%%%%%%%%%%%%%%%%%%%%%%%%%%%%%%%%%%%%%%%%%%%%%%%%%%%%%%%
% Lớp 12 | Nguyên hàm, tích phân | Ứng dụng thực tế và hình học của tích phân
% Dạng 2D4?3-2: Bài toán thực tế sử dụng diện tích hình phẳng
% Số câu: 94  (sinh tự động bởi tach_cau.py; từ nay thầy sửa trực tiếp file này)
%%%%%%%%%%%%%%%%%%%%%%%%%%%%%%%%%%%%%%%%%%%%%%%%%%%%%%%%%%%%%%%%%%%%%%

%%%%%%%%%%%%%%%% TRẮC NGHIỆM 4 PHƯƠNG ÁN %%%%%%%%%%%%%%%%

% nguồn: Phan1_Lop_12_Dai_so_Chuong_4
\begin{ex}%[2D4N3-2]%[ma:7bf51c79]
	Diện tích hình phẳng giới hạn bởi đồ thị hàm số $y=x^3-x, y=3 x$ và hai đường thẳng $x=0, x=3$  bằng
	\choice 
		{$\dfrac{17}{4}$} 
		{\True  $\dfrac{41}{4}$}
		{$\dfrac{9}{4}$} 
		{$\dfrac{57}{4}$} 
	 \loigiai{
		Ta có $S=\displaystyle\int_0^3\left|x^3-4 x\right|\mathrm{\,d}x=\displaystyle\int_0^2\left(-x^3+4 x\right)\mathrm{\,d}x+\displaystyle\int_2^3\left(x^3-4  x\right)\mathrm{\,d}x=\dfrac{41}{4}$.
	} 
\end{ex}

% nguồn: Phan1_Lop_12_Dai_so_Chuong_4
\begin{ex}%[2D4H3-2]%[ma:26f94c36]
	Tốc độ chuyển động $v\,\mathrm{(m/s)}$ của ca nô trong khoảng thời gian $40$ giây được thể  hiện như Hình $1$. Quãng đường đi được của ca nô trong khoảng thời gian này là
	 \choice
		{$400 m$} 
		{$350 m$} 
		{\True$310 m$} 
		 {$200 m$}
		\begin{center}
			\begin{tikzpicture}
					\draw [->](-.5,0) node[below right]{$O$}--(5,0) node[above]{$t (s)$}; 
					\draw[->](0,-0.5)--(0,1.5) node[right]{$v  (m/s)$};
					\foreach \i  /\t/\g in {(0,0),(0.8,0)/8/-90,(3,0)/30/-90,(4,0)/40/-90,(3,1)/ /0,(0.8,1)/ /0,(0,1)/10/180}{\fill \i circle (1pt);}
					\foreach \i  /\t/\g in {(0.8,0)/8/-90,(3,0)/30/-90,(4,0)/40/-90,(3,1)/ /0,(0.8,1)/ /0}{\path \i node[below]{$\t$};}
					\path (0,1)  node[left]{$10$};
					 \draw (0,0)--(0.8,1)--(3,1)--(4,0);
					\draw[dashed](0,1)--(0.8,1)--(0.8,0) 
						 (3,0)--(3,1);
				\end{tikzpicture}
	\end{center}
	\loigiai{ 
		Quãng  đường đi được của ca nô trong $40$ giây là\\
		 $S=\displaystyle\int_{0}^{8} \dfrac{10}{8}x\mathrm{\,d}x+\displaystyle\int_{8}^{30} 10\mathrm{\,d}x+\displaystyle\int_{30}^{40} (-x+40)\mathrm{\,d}x=\dfrac{5x^2}{8}\bigg|_{0}^{8}+10x\bigg|_{8}^{30}+\left(-\dfrac{x^2}{2}+40x\right)\bigg|_{30}^{40}=310m$
	 }
\end{ex}

% nguồn: Phan1_Lop_12_Dai_so_Chuong_4
\begin{ex}%[2D4V3-2]%[ma:307df547]
	\immini{ 
	Trong một công viên, người ta xây dựng một hồ nước với hình dạng giới hạn bởi hai đường cong $y=\sqrt{x}$ và $y=\dfrac{1}{\sqrt{27}}x^2$. Hãy tính diện tích mặt hồ (kết quả làm tròn đến hàng phần chục).}{\begin{tikzpicture}[scale=0.7,font=\footnotesize,line join=round,
			 line cap=round,>=stealth]
		 \def \xmin{-2}
		 \def \xmax{4}
		\def \ymin{-1} 
		\def \ymax{3.5} 
		\draw[->]  (\xmin,0)--(\xmax,0) node[shift=(-110:0.2)] {$x$};
		\draw[->] (0,\ymin)--(0,\ymax+1) node[shift=(-150:0.2)]  {$y$};
		\fill (0,0) circle  (1pt) node[shift=(-135:0.25)]{$O$};
		\draw (1.5,2)  node{$y=\sqrt{x}$}
			(3.4,0.6) node{$y=\dfrac{1}{\sqrt{27}}x^2$}; 
		\clip (\xmin,\ymin) rectangle  (\xmax,\ymax);
		\draw[smooth,samples=100,domain=0:\xmax] plot  (\x,{sqrt (\x)});
		\draw[smooth,samples=100,domain=0:\xmax] plot  (\x,{sqrt (1/27)*(\x)^2});
		\foreach \x/\g in  {1/-90,2/-90,3/-90}
		 \draw (\x,1pt)--(\x,-1pt) + (\g:0.4) node {$\x$};
		 \foreach \y/\g in {2/180,1/180,3/180}
		\draw (1pt,\y) --(-1pt,\y) + (\g:0.4)  node {$\y$};
	\end{tikzpicture}
	 }
	\choice 
		{\True $1{,}7$} 
		{$1{,}5$} 
		{$2{,}5$} 
		{$2{,}3$} 
	\loigiai{ 
		Phương trình hoành độ giao  điểm
		\[\sqrt{x}=\dfrac{1}{\sqrt{27}}x^2  \Leftrightarrow \hoac{&x=0\\&x=3.}\]
		Diện tích mặt hồ  là
	\[S=\displaystyle  \int_0^3 \left(\sqrt{x} - \dfrac{1}{\sqrt{27}}x^2\right) \mathrm{\,d} x =\sqrt{3} \approx 1{,}7.\]}
\end{ex}

% nguồn: Phan1_Lop_12_Dai_so_Chuong_4
\begin{ex}%[2D4V3-2]%[ma:e9ef8518]
	Một viên gạch hoa hình vuông cạnh $60\ \mathrm{cm}$. Người thiết kế đã sử dụng bốn đường parabol có chung đỉnh tại tâm viên gạch để tạo ra bốn cánh  hoa (được tô đen như hình).
	\begin{center}
		\begin{tikzpicture}[scale=0.1,declare function={
					f (\x)  = \x*\x/30;
					 g (\x) = -\x*\x/30;
					h (\y) = \y*\y/30; 
					k (\y) =  -\y*\y/30;
					 }]
				\draw (-30,-30) rectangle (30,30); 
				 \fill[gray!30,domain=-30:30,variable=\x]
					plot ({\x},{f  (\x)}) -- plot ({\x},{g (\x)}) -- cycle;
				 \fill[gray!30,domain=-30:30,variable=\y]
					plot ({h  (\y)},{\y}) -- plot ({k (\y)},{\y}) -- cycle;
				\draw[domain=-30:30,smooth,variable=\x] plot  ({\x},{f (\x)});
				\draw[domain=-30:30,smooth,variable=\x] plot ({\x},{g  (\x)});
				\draw[domain=-30:30,smooth,variable=\y] plot ({h  (\y)},{\y});
				 \draw[domain=-30:30,smooth,variable=\y] plot ({k (\y)},{\y});
			\end{tikzpicture}
	\end{center}
	Diện  tích mỗi cánh hoa của viên gạch bằng
	\choice 
		{$\dfrac{400}{3}\  \mathrm{cm}^2$}
		{\True $300\ \mathrm{cm}^2$} 
		 {$800\ \mathrm{cm}^2$}
		 {$\dfrac{800}{3}\ \mathrm{cm}^2$}
	 \loigiai{
		Chọn hệ trục tọa độ $Oxy$ như hình  vẽ bên dưới
		\begin{center}
			\begin{tikzpicture}[scale=0.1,declare function={
						 f (\x) = \x*\x/30;
						 g (\x) = -\x*\x/30;
						h  (\y) = \y*\y/30;
						k (\y) = -\y*\y/30; 
						 }]
					\draw (-30,-30) rectangle  (30,30);
					\fill[gray!50,domain=0:30,variable=\x] 
						plot ({\x},{f (\x)}) -- 
						plot[domain=30:0,variable=\y]({h  (\y)},{\y}) -- cycle;
					\filldraw (30,30) circle (2pt) node[above left] {$A$}; 
					 \draw[domain=-30:30,smooth,variable=\x] plot ({\x},{f (\x)});
					\draw[domain=-30:30,smooth,variable=\x]  plot ({\x},{g (\x)});
					\draw[domain=-30:30,smooth,variable=\y] plot  ({h (\y)},{\y});
					\draw[domain=-30:30,smooth,variable=\y] plot ({k  (\y)},{\y});
					\draw[->] (-35,0) --  (35,0) node[above left] {$x$};
					\draw[->]  (0,-35) -- (0,35) node[below right] {$y$};
					\node at (2,-2) {$O$}; 
					\draw[red,dashed] (-35,-35) -- (35,35) node[above right] {$y = x$}; 
				\end{tikzpicture}
		\end{center}
		 Giả sử parabol $(P)\colon y=ax^2$ đi qua điểm $A(30;30)$ $\Rightarrow a= \dfrac{1}{30}$ $\Rightarrow (P)\colon y= \dfrac{1}{30}x^2$.\\
		Khi đó, diện  tích một cánh hoa là
		$$S=2 \cdot \displaystyle \int_0^{30} \left(x-\dfrac{1}{30}x^2\right)\,\mathrm{d}x =300 \,  \mathrm{cm}^2$$
		Vậy mỗi  cánh hoa có diện tích $300\ \mathrm{cm}^2$.
	} 
\end{ex}


%%%%%%%%%%%%%%%% ĐÚNG/SAI %%%%%%%%%%%%%%%%

% nguồn: Phan1_Lop_12_Dai_so_Chuong_4
\begin{ex}%[2D4H3-2]%[ma:38dfd1f4]
	Cho hình phẳng $(H)$ giới hạn bởi đồ thị hàm số $(C)\colon y=f(x)=6\sqrt{x}-5$, đường thẳng $y=1$ và hai đường thẳng $x=1$, $x=4$. 
	\choiceTF 
		 {\True $\displaystyle\int f(x) \mathrm{\,d}x=4 x \sqrt{x}-5 x+C$}
		{\True $\displaystyle\int[f(x)-1] \mathrm{\,d}  x=4 x \sqrt{x}-6 x+C$}
		{$\displaystyle\int_1^4[f(x)-1]  \mathrm{\,d} x=6$}
		{\True Diện tích của hình phẳng $(H)$ bằng $10$} 
	 \loigiai{
		\begin{itemchoice}
			 \itemch Ta có
			$$\displaystyle\int f(x) \mathrm{\,d} x=\displaystyle\int\left( 6 \sqrt{x}-5 \right) \mathrm{\,d} x=4 x \sqrt{x}-5 x+C.$$ 
			\itemch Ta  có
			$$\displaystyle\int[f(x)-1] \mathrm{\,d} x=\displaystyle\int\left(6 \sqrt{x}-6\right) \mathrm{\,d} x=4 x \sqrt{x}-6 x+C.$$ 
			 \itemch Ta có
			$$\displaystyle\int_1^4[f(x)-1] \mathrm{\,d} x=\displaystyle\int_1^4\left(6 \sqrt{x}-6\right) \mathrm{\,d} x=\left(4 x \sqrt{x}-6 x\right)\Big|_1 ^4=(32-24)-(4-6)=10.$$ 
			\itemch 
			$$ 
			S=\displaystyle\int_1^4\left|f(x)-1\right| \mathrm{\,d} x=\displaystyle\int_1^4\left|6 \sqrt{x}-6\right| \mathrm{\,d} x=\left|\left(4 x \sqrt{x}-6  x\right)\Big|_1^4 \right|=10.
			$$ 
		\end{itemchoice}
	} 
\end{ex}

% nguồn: Phan1_Lop_12_Dai_so_Chuong_4
\begin{ex}%[2D4H3-2]%[ma:0db6ac3e]
	Các nhà kinh tế sử dụng đường cong Lorenz để minh họa sự phân phối thu nhập trong một quốc gia. Gọi $x$ là đại diện cho  phần trăm số gia đình trong một quốc gia và $y$ là phần trăm tổng thu nhập, mô hình $y=x$ sẽ đại diện cho một quốc gia mà các gia đình có thu nhập như nhau. Đường cong Lorenz $y=f(x)$, biểu thị sự phân phối thu nhập thực tế. Diện tích giữa hai mô hình này, với $0 \le x \le 100$, biểu thị "sự bất bình đẳng về thu nhập" của một quốc gia. Năm $2005$, đường cong Lorenz của Hoa Kỳ có thể được mô hình hóa bởi hàm số
	$$ y = (0{,}00061x^2+0{,}0218x+1{,}723)^2,\hspace{0.1cm} 0 \le x \le  100.$$
	Trong đó $x$ được tính từ các gia đình nghèo  nhất đến giàu có nhất.\\
	 \textit{(Theo R.Larson, Brief Calculus: An Applied Approach, $8^{\text {th }}$ edition, Cengage Learning, $2009$)}
	\choiceTF 
		{\True Theo thứ tự  từ các gia đình nghèo nhất đến giàu nhất, tổng thu nhập thực tế của $60\%$ các gia đình đầu tiên chiếm chưa đến $30\%$ so với tổng thu nhập của toàn bộ các gia đình}
		{Nếu sắp xếp các gia đình theo thứ tự từ nghèo nhất đến giàu nhất, rồi chia thành $10$ nhóm bằng nhau từ $1$ đến $10$, tổng thu nhập của các gia đình trong nhóm 3 chiếm khoảng $8{,}56\%$ tổng thu nhập của toàn bộ các gia đình} 
		{Sự bất bình đẳng về  thu nhập của Hoa Kì năm $2005$ được xác định bởi công thức
		 $\displaystyle\int_0^{100} \left[x - (0{,}00061x^2+0{,}0218x+1{,}723)^2\right]\mathrm{d}x$}
		{\True Sự bất bình  đẳng về thu nhập của Hoa Kỳ năm $2005$ đã vượt quá $2000$}
	 \loigiai{
		Đặt  $L(x)=(0{,}00061x^2+0{,}0218x+1{,}723)^2$.
		\begin{enumerate}
			\item  \textbf{Đúng.}
			\begin{align*}
				L(60) &= (0{,}00061\cdot 60^2 + 0{,}0218\cdot 60 + 1{,}723)^2  = (2{,}196 + 1{,}308 + 1{,}723)^2 \\ &= (5{,}227)^2 \approx 27{,}32
			\end{align*}
			Vì $27{,}32 < 30$ nên mệnh đề  đúng.
			\item \textbf{Sai.} "Nhóm 3" tương ứng với các gia đình từ $20\%$ đến $30\%$. Thu nhập của nhóm này được tính bằng  $L(30)-L(20)$.
			\begin{itemize}
				\item $L(30) = (0{,}00061\cdot 30^2+0{,}0218\cdot 30+1{,}723)^2 = (0{,}549+0{,}654+1{,}723)^2  = (2{,}926)^2 \approx 8{,}56$.
				\item $L(20) = (0{,}00061\cdot 20^2+0{,}0218\cdot 20+1{,}723)^2 = (0{,}244+0{,}436+1{,}723)^2 = (2{,}403)^2 \approx  5{,}77$.
			\end{itemize}
			Thu nhập của nhóm 3 là $8{,}56\% - 5{,}77\% = 2{,}79\%$.  Con số $8{,}56\%$ là tổng thu nhập của 3 nhóm đầu tiên.
			\item \textbf{Đúng.} Sự bất bình  đẳng về thu nhập (hệ số Gini nhân với diện tích tam giác dưới đường $y=x$) được định nghĩa là diện tích giữa đường cong bình đẳng tuyệt đối $y=x$ và đường cong Lorenz $y=L(x)$. Công thức tính diện tích này chính là $\displaystyle\int_0^{100} [x - L(x)]\mathrm{d}x$.
			\item \textbf{Đúng.} Ta cần tính giá trị của tích phân  ở câu c).
			$A  = \displaystyle\int_0^{100} x \mathrm{d}x - \int_0^{100} (0{,}00061x^2+0{,}0218x+1{,}723)^2 \mathrm{d}x \approx 2060{,}37$.
			Vậy mệnh đề này  đúng.
		\end{enumerate}
	 }
\end{ex}

% nguồn: Phan1_Lop_12_Dai_so_Chuong_4
\begin{ex}%[2D4H3-2]%[ma:c941403f]
	Cho một viên gạch  men có dạng hình vuông $OABC$ như hình vẽ. Sau khi tọa độ hóa, ta có $O(0;0)$, $A(0;1)$, $B(1;1)$, $C(1;0)$ và hai đường cong lần lượt là đồ thị hàm số $y=x^3$ và $y=\sqrt[3]{x}$.
	\begin{center}
		\begin{tikzpicture}[scale=3,font=\footnotesize,line join=round,line cap=round,>=stealth]
				\draw[->,line width = 1pt] (-.2,0)--(0,0)  node[below right]{$O$}--(1.5,0) node[below]{$x$};
				\draw[->,line width = 1pt]  (0,-.2) --(0,1.5) node[right]{$y$};
				\foreach  \x in {}{
					\draw (\x,0)  node[below]{$\x$} circle (1pt);
				 }
				\draw [domain=0:1.1,samples=100] 
					plot  (\x,{(\x)^3}) node[above]{$y=x^3$};
				\draw [domain=0:1.3,samples=100] 
					plot  (\x,{(\x)^.333333}) node[right]{$y=\sqrt[3]{x}$};
				\draw [dashed] (0,1)--(1,1)--(1,0); 
				\draw (1,0) node[below]{$ C  $};
				\draw (0,1)  node[left]{$ A $};
				\draw  (1,1) node[below right]{$ B $};
				\draw[pattern,opacity=.2,line width = 1pt]  (0,0)
					plot[domain=0:1] (\x,{(\x)^0.33})--(1,1)--cycle; 
				 \draw[pattern,opacity=.2,line width = 1pt] (0,0) plot[domain=0:1] (\x,{(\x)^3})--(1,1)--cycle;
			\end{tikzpicture}
	\end{center}
	 \choiceTF
		{\True Diện tích hình phẳng giới hạn bởi đồ thị  hàm số $y=\sqrt[3]{x}$, trục $Ox$, đường thẳng $x=0$ và đường thẳng $x=1$ được tính bằng công thức $S=\displaystyle \int_0^1 \left| \sqrt[3]{x}\right| \mathrm{d}x$}
		{Diện tích hình phẳng giới hạn bởi đồ thị hàm số $y=x^3$, trục $Ox$, đường thẳng $x=0$ và đường thẳng $x=1$ có giá trị bằng $\dfrac{3}{4}$  (đvdt)}
		 {Diện tích hình phẳng giới hạn bởi đồ thị hàm số $y=x^3$ và $y=\sqrt[3]{x}$, đường thẳng $x=0$ và đường thẳng $x=1$ được tính bằng công thức $S=\displaystyle \int_0^1 (x^3 - \sqrt[3]{x})\mathrm{d}x$}
		{\True Diện tích phần không được tô đậm  trên viên gạch men có giá trị bằng $\dfrac{1}{2}$ (đvdt)}
	\loigiai{ 
		Trên  đoạn $[0;1]$, ta có $\sqrt[3]{x} \ge x^3$.
		\begin{enumerate}
			\item \textbf{Đúng.} 
			Công thức đúng là $S = \displaystyle\int_0^1 \left| \sqrt[3]{x}  \right| \mathrm{d}x$.
			\item  \textbf{Sai.}
			Diện tích giới hạn bởi $y=x^3$, trục $Ox$  từ $x=0$ đến $x=1$ là
			$S = \displaystyle\int_0^1 x^3  \mathrm{d}x = \dfrac{1}{4}$.
			\item  \textbf{Sai.}
			Diện tích  hình phẳng giới hạn bởi hai đồ thị được tính bằng công thức tích phân của (hàm trên trừ hàm dưới).
			 Công thức đúng là $S = \displaystyle\int_0^1 (\sqrt[3]{x} - x^3)\mathrm{d}x$.
			 \item \textbf{Đúng.}
			 Diện tích phần tô đậm (giới hạn bởi hai đường cong) là\\
			$S_{\text{ tô đậm}} 
			= \displaystyle\int_0^1 (\sqrt[3]{x}-x^3)\mathrm{d}x = \dfrac{3}{4}-\dfrac{1}{4}=\dfrac{1}{2}$.\\ 
			Diện tích hình vuông $OABC$ là $1 \times  1 = 1$.\\
			Diện tích phần không  tô đậm = Diện tích hình vuông - Diện tích phần tô đậm $= 1 - \dfrac{1}{2}=\dfrac{1}{2}$.
		\end{enumerate}
	 }
\end{ex}

% nguồn: Phan1_Lop_12_Dai_so_Chuong_4
\begin{ex}%[2D4V3-2]%[ma:439beeaf]
	Vịnh Hạ Long là một địa  danh du lịch được nhiều người biết đến trên thế giới, nơi đây vẫn còn nhiều quần thể đảo lớn nhỏ chưa được khám phá. Một công ty du lịch quyết định khai thác khu vực có một số đảo nhỏ với hình dáng đặc biệt nếu nhìn từ trên xuống, trong số đó có hai hòn đảo mà phần giới hạn đất liền của nó được mô phỏng như hai đồ thị hàm số trên hình. Với hệ trục tọa độ $Oxy$ thích hợp, đơn~vị trên mỗi trục là $100$ mét, đường cong mô tả cho hòn đảo thứ nhất có dạng $y=\log_a x$ đi qua điểm có tọa độ $(3;1)$.
	\begin{center}
		\begin{tikzpicture}[line cap=butt,line join=miter,>=stealth]
				 \tikzset{
					declare function={ 
					f  (\x) = ln (\x)/ln (3);
					g  (\x) = 1 + 3*3^(\x);
					h (\x) =  \x + 1;
					 }}
				\begin{scope}
					\clip (-3,-2) rectangle (5,5); 
					 \path[fill=green!65!cyan,smooth,samples=100,fill opacity=0.3,draw]
						plot[domain=0.05:5.1]  (\x,{f (\x)})--(5,-2)--cycle
						plot[domain=-3.1:1.2] (\x,{g  (\x)})--(-3.1,5)--cycle;
					\draw[smooth,samples=100] plot[domain=-3:4] (\x,{h  (\x)});
				\end{scope}
				 \draw[->] (-3,0)--(5,0)
					 node[shift={(-100:7pt)},font=\normalsize]{$x$};
				\draw[->]  (0,-2)--(0,5)
					node[shift={(180:7pt)},font=\normalsize]{$y$}; 
				\draw (5pt,0)  |- (0,5pt)
					(0,0)  node[shift={(225:9pt)},font=\normalsize]{$O$};
				\foreach \x  in {1,3,4}{
					\draw  (\x,1.5pt)--(\x,-1.5pt) node[font=\scriptsize,shift={(-50:5pt)}]{\x};
				 }
				\foreach \y in {-1,1}{ 
					\draw (1.5pt,\y)--(-1.5pt,\y) node[font=\scriptsize,shift={(180:5pt)}]{\y}; 
				 }
				\draw[dashed] (0,1)  -| (3,0)
					(0,-1)--(4,-1)--(4,1.25); 
				\path (1/3,-1) node[below right]{$A$} 
					(4,-1)  node[below right]{$B$}
					(4,1.26)  node[below right]{$C$}
					(3.5,1.1) node[above left]{$y=\log_a x$} 
					(0,1)--(3,4) node[above,sloped,pos=0.8]{$y=x+1$}; 
				\fill (1/3,-1) circle  (1pt)
					(4,-1) circle  (1pt)
					(4,1.26)  circle (1pt)
					(3,1)  circle (1pt);
			\end{tikzpicture}
	\end{center}
	 \choiceTF
		{Điểm có tọa  độ $(9;3)$ thuộc đường cong $y=\log_a x$}
		{Chủ dự  án muốn xây dựng một nơi trực tiếp nhìn ra biển để du khách tham quan, ăn uống...Họ đã lựa chọn khu vực tam giác cong $ABC$ như trong hình (đường cong $AC$ tiếp giáp biển); diện tích khu vực này là $536~\mathrm{m}^2$ (làm tròn đến hàng đơn~vị)}
		{\True Chủ dự án đã thuê một số kỹ sư rất giỏi toán (\textbf{đặc biệt giỏi  về hàm số mũ-log}) đi khảo sát khu vực này và họ nhận thấy có thể bồi đắp thêm cho hòn đảo thứ hai để đường cong giáp biển $y=g(x)$ của nó đối xứng với đường cong $y=\log_a x$ qua đường thẳng $y=x+1$. Khi đó đường cong $g(x)=1+3\cdot 3^x$}
		{Chủ dự án định xây một cây cầu nối  liền hai hòn đảo, khoảng cách ngắn nhất theo đường chim bay của cây cầu bằng $285~\mathrm{m}$ (làm tròn đến hàng đơn~vị của mét)}
	\loigiai{ 
		\begin{itemchoice}
			\itemch\textbf{Mệnh đề  sai}.\\
			Đường cong $y=\log_a x$  đi qua điểm $(3;1)$ nên $1=\log_a 3\Rightarrow a=3$.\\
			Khi đó hàm số trở thành $y=\log_3 x$, đường  cong này không đi qua điểm $(9; 3)$.
			\itemch \textbf{Mệnh  đề sai}.\\
			Điểm $A\left(x_A;-1\right)$ thuộc đồ thị hàm số $y=\log_3  x\Rightarrow \log_3 x=-1\Rightarrow x=3^{-x}=\dfrac{1}{3}$.\\
			Diện tích tam giác cong $ABC$ là phần hình phẳng được giới hạn bởi hai đồ thị $y=\log_3 x$; $y=-1$ cùng các đường thẳng $x=\dfrac{1}{3}$; $x=4$.\\ 
			Do đó diện tích cần  tính là $S=100\times \displaystyle\int_{\frac{1}{3}}^4\left|\log_3 x-(-1)\right|\,\mathrm{d}x \approx 571~\mathrm{m}^2$.
			 \itemch \textbf{Mệnh đề đúng}.\\
			Gọi $M\left(x_M;  y_M \right)\in \left(C_1 \right)\colon y=\log_3 x$ và $N(x; y)\in \left(C_2 \right)\colon y=g(x)$.\\
			$M$, $N$ đối  xứng qua $x-y+1=0$ nên ta có
			\begin{eqnarray*}
				\heva{&\dfrac{x+x_M}{2}-\dfrac{y+y_M}{2}+1=0 \\&1\cdot \left(x-x_M \right)+1\cdot \left(y-y_M \right)=0}\Leftrightarrow \heva{&x-y+x_M-y_M+2=0 \\  &x+y-x_M-y_M=0}\Leftrightarrow \heva{&y_M=x+1 \\ &x_M=y-1.}
			\end{eqnarray*} hay $M(y-1; x+1)$.\\
			Vì $M\in \left(C_1 \right)$ nên $x+1=\log_3 (y-1)\Rightarrow  y-1=3^{x+1} \Rightarrow y=3^{x+1}+1$ hay $y=g(x)=1+3\cdot 3^x$.
			\itemch \textbf{Mệnh đề  sai}.\\
			Xét tiếp tuyến của đường cong $y=\log_3 x$ biết tiếp  tuyến song song với đường thẳng $y=x+1$.\\
			 Hệ số góc tiếp tuyến là $k=1$; gọi $M\left(x_0; y_0 \right)$ là tiếp điểm thì
			\begin{eqnarray*}
				 f'\left(x_0 \right)=k\Rightarrow \dfrac{1}{x_0 \ln 3}=1\Rightarrow x_0=\dfrac{1}{\ln 3}; y_0=\log_3 \dfrac{1}{\ln 3}.
			\end{eqnarray*}
			Độ dài ngắn nhất cây cầu (theo đường chim bay) bằng hai lần khoảng cách từ $M\left(\dfrac{1}{\ln 3}; \log_3 \dfrac{1}{\ln 3} \right)$ đến đường  thẳng $y=x+1$.\\
			Ta  có $d_{\min}=2\cdot \dfrac{\left|\dfrac{1}{\ln 3}-\log_3 \dfrac{1}{\ln 3}+1\right|}{\sqrt{2}} \times 100\approx 282~\mathrm{m}$.
		\end{itemchoice}
	} 
\end{ex}

% nguồn: Phan1_Lop_12_Dai_so_Chuong_4
\begin{ex}%[2D4V3-2]%[ma:351b6cf9]
	Vịnh Hạ Long là một địa danh du lịch được nhiều người biết đến trên thế giới, nơi đây vẫn còn nhiều đảo lớn nhỏ chưa được khám phá. Một công ty du lịch quyết định khai thác khu vực có một số đảo  nhỏ với hình dáng đặc biệt nếu nhìn từ trên xuống, trong đó có hai hòn đảo mà phần giới hạn đất liền của nó được mô phỏng như hai đồ thị hàm số trên hình. Với hệ trục $Oxy$ thích hợp, đơn vị trên mỗi trục là $100$\,m, đường cong mô tả hòn đảo thứ nhất có dạng $y=\log_a x$ đi qua điểm có tọa độ $(3;1)$.
	\begin{center}
		\begin{tikzpicture}[scale=1,font=\footnotesize,line join=round,line cap=round,>=stealth]
				\def\f (\x){ln (\x)/ln  (3)}
				\def\g (\x){1+3^((\x)+1)} 
				\pgfmathsetmacro{\c}{ln (4)/ln  (3)}
				\pgfmathsetmacro{\a}{1/3} 
				\draw[teal,->] (-2,0)--(0,0) node[below left]{$O$}--(5,0) node[above  right]{$x$};
				\foreach \x in {3,4}{\ifnum \x=0 \else \draw[teal,thin] (\x,2pt)--(\x,-2pt) node[below right]{$\x$};  \fi}
				\draw[teal,->] (0,-2)--(0,5)  node[right]{$y$};
				\foreach \y in  {-1,1}{\ifnum \y=0 \else \draw[teal,thin] (2pt,\y)--(-2pt,\y) node[left]{$\y$}; \fi}
				\clip (-2,-2) rectangle  (5,5);
				\fill[gray!30!,opacity =.7] 
					plot[domain = 0.15:5,samples  = 200,variable=\x]({\x},{\f (\x)})
					 --plot[domain = 5:0.15,samples = 200,variable=\x]({\x},{-2});
				\fill[gray!30!,opacity =.7] 
					 plot[domain = -2:0.5,samples = 200,variable=\x]({\x},{\g (\x)})
					--plot[domain = 0.5:-2,samples =  200,variable=\x]({\x},{5});
				\draw[teal,domain=0.05:5,samples=100,smooth] 
					 plot (\x,{\f (\x)}) node[smooth,above left]{$y=\log_ax$};
				 \draw[teal,domain=-2:0.5,samples=100,smooth]
					plot  (\x,{\g (\x)});
				\draw[teal,dashed](3,0)|-(0,1) 
					 (4,0)|-(0,-1)
					 (4,0)--(4,\c) node[below right]{$C$}
					 (4,-1) node[below]{$B$}
					(\a,-1)  node[below right]{$A$};
				 \foreach \x/\y in{3/1,4/-1,4/\c,\a/-1}
					{ 
				\fill (\x,\y) circle  (1pt);
				} 
			\end{tikzpicture}
	\end{center}
	\choiceTF 
		 {Điểm có tọa độ $(9;3)$ thuộc đường cong $y=\log_a x$}
		{Chủ dự án muốn xây dựng một nơi trực tiếp nhìn ra biển để du khách tham quan, ăn uống. Họ đã lựa chọn khu vực tam giác cong $ABC$ như trong hình (đường cong $AC$ tiếp giáp biển);  diện tích của khi vực này là $536$\,m$^2$ (kết quả làm tròn đến hàng đơn vị)}
		{Chủ dự án đã thuê một số kỹ sư giỏi Toán (đặc biệt giỏi về hàm số mũ- logarit) đi khảo sát khu vực này và họ nhận thấy có thể bồi đắp thêm cho hòn đảo thứ hai để đường cong giáp biển $y=g(x)$ của nó đối xứng với đường cong $y=\log_a x$ qua  đường thẳng $y=x+1$. Khi đó đường cong $g(x)=1+3\cdot 3^x$}
		{Chủ dự án định xây một cây cầu nối liền hai đảo, khoảng cách ngắn nhất theo đường chim bay của cây  cầu bằng $285$\,m (kết quả làm trong đến hàng đơn vị)}
	 \loigiai{
		\begin{itemchoice}
			\itemch Đường cong $y=\log_ax$  đi qua điểm $(3;1)$ nên $\log_a3=1\Leftrightarrow a=3$.\\
			Khi đó,  hàm số trở thành $y=\log_3x$.\\
			Với $x=9$  thì $y=\log_39=2$ nên đường cong này đi qua điểm $(9;2)$.
			\itemch Điểm $A\left(x_A;-1\right)$ thuộc đồ thị hàm số $y=\log_3x$ nên $\log_3x_A=-1\Leftrightarrow x_A=\dfrac{1}{3}$.\\ 
			Khi đó, tam giác cong $ABC$ là phần hình phẳng được giới hạn bởi hai đồ  thị hàm số $y=\log_3x$; $y=-1$ cùng các đường thẳng $x=\dfrac{1}{3}$ và $x=4$.\\
			 Do đó, diện tích khu vực này là $S=100^2\cdot\displaystyle\int_{\tfrac{1}{3}}^4\left|\log_3x+1\right|\mathrm{\,d}x\approx 571\,000$\,m$^2$.
			 \itemch
		 \immini{
			 Gọi $M\left(x_M;y_M\right)\in \left(C_1\right)\colon y=\log_3x$ và $N(x;y)\in \left(C_2\right)\colon y=g(x)$.\\
			$M$, $N$ đối xứng nhau qua đường thẳng $d\colon x-y+1=0$ khi và chỉ khi trung điểm của đoạn thẳng $MN$ sẽ thuộc đường thẳng $d$ và $MN$ vuông góc với $d$, nghĩa là 
			 \allowdisplaybreaks
			\begin{eqnarray*}
				&&\heva{&\dfrac{x+x_M}{2}-\dfrac{y+y_M}{2}+1=0\\&1\cdot (x-x_M)+1\cdot  (y-y_M)=0}\\
				&\Leftrightarrow&\heva{&x-y+x_M-y_M+2=0\\&x+y-x_M-y_M=0}\\ 
				 &\Leftrightarrow&\heva{&x_M-y_M=-x+y-2\\&x_M+y_M=x+y}\\
				&\Leftrightarrow&\heva{&x_M=y-1\\&y_M=x+1}\Rightarrow M(y-1;x+1). 
			\end{eqnarray*}
			} 
			{ 
			\begin{tikzpicture}[scale=.85,font=\footnotesize,line join=round,line cap=round,>=stealth]
					\def\f (\x){ln (\x)/ln  (3)}
					\def\g (\x){1+3^((\x)+1)} 
					 \pgfmathsetmacro{\c}{ln (4)/ln (3)}
					\pgfmathsetmacro{\a}{1/3} 
					 \draw[teal,->] (-2,0)--(0,0) node[below left]{$O$}--(5,0) node[above right]{$x$};
					\foreach \x in {3,4}{\ifnum \x=0 \else \draw[teal,thin] (\x,2pt)--(\x,-2pt) node[below right]{$\x$};  \fi}
					\draw[teal,->]  (0,-2)--(0,5) node[right]{$y$};
					\foreach \y in {-1,1}{\ifnum \y=0  \else \draw[teal,thin] (2pt,\y)--(-2pt,\y) node[left]{$\y$}; \fi}
					\clip (-2,-2) rectangle  (5,5);
					\fill[gray!30!,opacity  =.7]
						 plot[domain = 0.15:5,samples = 200,variable=\x]({\x},{\f (\x)})
						--plot[domain = 5:0.15,samples = 200,variable=\x]({\x},{-2}); 
					\fill[gray!30!,opacity =.7] 
						plot[domain =  -2:0.5,samples = 200,variable=\x]({\x},{\g (\x)})
						--plot[domain  = 0.5:-2,samples = 200,variable=\x]({\x},{5});
					\draw[teal,domain=0.05:5,samples=100,smooth] 
						plot (\x,{\f (\x)})  node[smooth,above left]{$y=\log_ax$};
					\draw[teal,domain=-2:0.5,samples=100,smooth] 
						plot (\x,{\g  (\x)});
					\draw[teal] (-2,-1)  -- node[midway,above,sloped] {$y=x+1$}
						 (5,6);
					 \draw[teal,dashed](3,0)|-(0,1)
						 (4,0)|-(0,-1)
						(4,0)--(4,\c) node[below  right]{$C$}
						(4,-1)  node[below]{$B$}
						(\a,-1) node[below right]{$A$}; 
					\foreach  \x/\y in{3/1,4/-1,4/\c,\a/-1}
						{ 
					\fill  (\x,\y) circle (1pt);
					 }
				\end{tikzpicture}
			 }
			 \noindent
			Vì $M\in \left(C_1\right)$  nên $\log_3(y-1)=x+1\Leftrightarrow y-1=3^{x+1}\Leftrightarrow y=1+3\cdot 3^x$.\\
			Vậy $y=g(x)=1+3\cdot 3^x$. 
			\itemch Xét tiếp tuyến của đường cong $y=\log_3x$ biết tiếp tuyến song  song với đường thẳng $y=x+1$.\\
			Hệ số góc  của tiếp tuyến là $k=1$.\\
			Gọi  $M\left(x_0;y_0\right)$ là tiếp điểm.\\
			Đạo hàm $y'=\dfrac{1}{x\ln  3}$, khi đó hệ số góc của tiếp tuyến tại $M$ là
			\allowdisplaybreaks 
			\begin{eqnarray*}
				y'(x_0)=k&\Leftrightarrow&\dfrac{1}{x_0\ln  3}=1\\
				&\Leftrightarrow&x_0=\dfrac{1}{\ln 3}. 
			\end{eqnarray*}
			Suy  ra $y_0=\log_3\left(\dfrac{1}{\ln 3}\right)$.\\
			Độ dài ngắn nhất của cây cầu (theo đường chim bay) sẽ bằng hai lần khoảng cách từ điểm $M\left(\dfrac{1}{\ln 3};\log_3\left(\dfrac{1}{\ln 3}\right)\right)$ đến đường thẳng $d\colon x-y+1=0$.\\ 
			Ta có $\mathrm{d}(M,d)=\dfrac{\left|\dfrac{1}{\ln 3}-\log_3\left(\dfrac{1}{\ln  3}\right)+1\right|}{\sqrt{1^2+(-1)^2}}\approx 1{,}41$.\\
			Vậy độ dài ngắn nhất của cây cầu gần bằng $2\cdot 100\cdot 1{,}41=282$\,m. 
		\end{itemchoice}
	 }
\end{ex}

% nguồn: Phan1_Lop_12_Dai_so_Chuong_4
\begin{ex}%[2D4V3-2]%[ma:b373f74c]
	Nhằm điều tiết mực nước của hồ chứa trong những ngày mưa lớn, một hồ thủy điện thông báo mở cửa xả lũ trong $5$ giờ, bắt đầu từ thời điểm $0$h. Lưu lượng nước xả lũ là lượng nước hồ thủy điện xả về hạ  lưu trên mỗi đơn vị thời gian, được ước tính bởi hàm số $f(t)=at^3+bt^2+1$ (đơn vị nghìn m$^3$/s), trong đó $t$ được tính bằng đơn vị giờ, là thời gian từ $0$h đến $5$h. Tại thời điểm $0$h, lưu lượng nước xả về hạ lưu là $1$ nghìn m$^3$/s và bắt đầu tăng cho đến khi đạt cực đại bằng $3$ nghìn m$^3$/s thì lưu lượng nước giảm dần, đến thời điểm $5$h thì trở về lại $1$ nghìn m$^3$/s. Hàm số $y=f(t)$ có đồ thị như hình vẽ dưới đây.
	\begin{center}
		\begin{tikzpicture}[line join = round,line cap=round,>=stealth,font=\footnotesize,scale=1]
				\draw[teal,->]  (-0.6,0) -- (6,0) node[below]{$t$(giờ)};
				\draw[teal,->] (0,-0.6) -- (0,3.6)  node[right]{$y$(nghìn m$^3$/s)};
				\draw[teal] (0,0) node[below left]{$O$}; 
				\draw[teal,smooth,blue,samples=100,domain=0:5] plot  (\x,{-0.108*(\x)^3 + 0.54*(\x)^2 + 1});
				 \draw[teal,dashed] (0,3) node[left]{$3$}circle (1pt)--(3.3,3)--(3.3,0);
				\draw[teal,dashed] (0,1) node[left]{$1$}circle (1pt)--(5,1)--(5,0)  node[below]{$5$};
			\end{tikzpicture}
	\end{center}
	 \choiceTF
		 {\True $f(5)=1$}
		{Giá trị của  $a$ là $-0{,}2$}
		{\True Lưu lượng  nước xả lũ đạt cực đại tại thời điểm $3$h$20$}
		{Tổng lượng nước được hồ chứa xả  về hạ lưu từ $0$h đến $5$h là $20{,}25$ triệu m$^3$}
	 \loigiai{
		 Xét $f(t)=at^3+bt^2+1$.\\
		 Vì điểm $(5;1)$ thuộc đồ thị hàm số $f(t)$ nên
		$$ f(5) = 1 \Leftrightarrow a\cdot 5^3 + b\cdot 5^2 + 1 = 1 \Leftrightarrow 125a + 25b  = 0 \Leftrightarrow b = -5a.$$
		Ta có $f'(t) = 3at^2 + 2bt$.\\ 
		Suy  ra
		\allowdisplaybreaks 
		\begin{eqnarray*}
			 f'(t) = 0 \Leftrightarrow 3at^2 + 2bt = 0 \Leftrightarrow 3at^2 - 10at = 0 &\Leftrightarrow& at(3t-10) = 0\\
			&\Leftrightarrow&  \hoac{&at = 0\\&3t-10=0} \\
			&\Leftrightarrow& \hoac{&t=0\\&t=\dfrac{10}{3}}  \text{ (loại $t=0$ vì $0<t<5$).}
		\end{eqnarray*}
		Do đó, lưu  lượng xả đạt cực đại tại $t = \dfrac{10}{3}$.\\
		Suy ra 
		$$ f\left(\dfrac{10}{3}\right) = 3 \Leftrightarrow \left(\dfrac{10}{3}\right)^3a + \left(\dfrac{10}{3}\right)^2b + 1 = 3 \Leftrightarrow \dfrac{1\,000}{27}a  + \dfrac{100}{9}\cdot (-5a) = 2 \Leftrightarrow a = -\dfrac{27}{250}.$$
		Do đó  $b = -5a = -5\cdot \left(-\dfrac{27}{250}\right) = \dfrac{27}{50}$.\\
		Vậy $f(t) = -\dfrac{27}{250}t^3  + \dfrac{27}{50}t^2 + 1$.
		\begin{itemchoice}
			\itemch Dựa vào đồ thị, ta có  $f(5) = 1$.
			 \itemch Dựa vào kết quả trên, ta có $a = -\dfrac{27}{250} = -0{,}108$.
			\itemch Lưu lượng nước xả lũ đạt cực đại tại thời điểm $t =  \dfrac{10}{3} = 3 + \dfrac{1}{3}$ (giờ) hay $3$h$20$.
			 \itemch Tổng lượng nước được hồ chứa xả về từ $0$h đến $5$h là
			\[\displaystyle\int_0^5 f(t)\mathrm{\,d}t = 3\,600\displaystyle\int_0^5\left(-\dfrac{27}{250}t^3 + \dfrac{27}{50}t^2 + 1\right) \mathrm{\,d}t = 119{,}250  \text{ (nghìn m$^3$/giờ)} = 119{,}25 \text{ (triệu m$^3$/giờ)}.\]
		\end{itemchoice}
	} 
\end{ex}

% nguồn: Phan1_Lop_12_Dai_so_Chuong_4
\begin{ex}%[2D4V3-2]%[ma:abc0e61f]
	 \immini
	 {Chất điểm chuyển động theo quy luật vận tốc $v(t)$ (m/s) có dạng đường thẳng khi $0\leq t\leq3(\mathrm{~s})$ và $8\leq t\leq15(s)$ và $v(t)$ có dạng đường Parabol khi $3\leq t\leq8(s)$ (như hình vẽ bên).
	 \choiceTF
		{\True Quãng đường chất điểm đi được trong khoảng thời gian từ $8$  giây đến $15$ giây bằng $73{,}5$ (m)}
		{\True  Quãng đường chất điểm di chuyển được trong 3 giây đầu tiên là~$S_3=\displaystyle\int_0^{3}11\mathrm{\,d}t$ (m)}
		{Vận tốc trung bình $v_{tb}$ của chất  điểm trong khoảng thời gian từ $3$ giây đến $8$ giây thoản mãn $v_{tb}<7$ (m/s)}
		{\True Vận tốc của chất điểm tại  thời điểm $t=3$ là~$v(3)=11$ (m/s)}
		} 
		{\begin{tikzpicture}[line join=round,line cap=round,>=stealth,y=0.25cm,x=0.25cm]
		\draw[->](-1,0)--(18,0) node[below]{$t$(s)}; 
		\draw[->](0,-1)--(0,25) node[right]{$v$(m/s)}; 
		 \draw (0,11)--(3,11)
			(8,21)--(15,0); 
		 \draw[samples=100,domain=3:8]plot (\x,{2*(\x)^2-20*(\x)+53});
		 \draw[dashed](3,0)--(3,11)
			(5,0)--(5,3) 
			 (8,0)--(8,21)--(0,21)
			(0,0) node[below left]{$O$}; 
		\foreach \x in{3,5,8,15}{\draw(\x,0) circle  (1pt) node[below]{$\x$};}
		\foreach  \x in{3,11,21}{\draw(0,\x) circle (1pt) node[left]{$\x$};}
	\end{tikzpicture}
		 }
	\loigiai{ 
		\begin{itemchoice}
			\itemch Trong khoảng thời gian từ $8$ giây đến $15$ giây, đồ thị $v(t)$ là~một đường thẳng đi qua hai điểm có tọa độ $(8;21)$ và  $(15;0)$.\\
			Gọi $v(t)=at+b$. Từ  giả thiết ta có $\heva{&8a+b=21\\
			 &15a+b=0}\Leftrightarrow\heva{&a=-3\\
			&b=45.}$\\ 
			Do đó  $v(t)=-3t+45(8\leq t\leq15)$.\\
			Quãng đường chất điểm đi được trong quãng thời gian này  là~$\displaystyle\int_8^{15}(-3t+45)\mathrm{\,d}t=73{,}5$ m.
			\itemch Trong khoảng thời gian từ $0$ đến $3$ giây, đồ thị $v(t)$ là~một đường  thẳng đi qua hai điểm có tọa độ $(0;11)$ và $(3;11)$ nên $v(t)=11$ m/s.\\
			Quãng đường chất điểm di  chuyển được trong $3$ giây đầu tiên là~$\displaystyle\int_0^3{11}\mathrm{\,d}t$ m.
			\itemch Trong khoảng thời gian từ $3$ đến $8$ giây, đồ thị  $v(t)$ là~một Parabol đi qua ba điểm có tọa độ $(3;11)$, $(5;3)$ và $(8;21)$.\\
			Gọi Parabol có phương trình  dạng $v(t)=at^2+bt+c$.\\
			Từ giả thiết ta có $\heva{&9a+3b+c=11\\ 
			&25a+5b+c=3\\ 
			&64a+8b+c=21}\Leftrightarrow\heva{&a=2\\ 
			 &b=-20\\
			 &c=53.}$\\
			 Dó đó $v(t)=2t^2-20t+53$ với $(3\leq t\leq8)$.\\
			Quãng đường  chất điểm đi được từ $3$ đến $8$ giây là~$S=\displaystyle\int_3^8\left(2t^2-20t+53\right)\mathrm{\,d}t=\dfrac{115}{3}$ (m).\\
			Vận tốc trung bình của chất điểm chuyển động  từ giây thứ $3$ đến giây thứ $8$ là
			\[v_{tb}=\dfrac{S}{5}=\dfrac{115}{3\cdot5}=\frac{23}{3}\approx7{,}67\quad>7\quad\text{  (m/s)}.\]
			\itemch Dựa vào đồ thị $v(t)$ (m/s) thì vận tốc của chất điểm tại thời điểm  $t=3$ là~$v(3)=11$ (m/s).
		\end{itemchoice}
	 }
\end{ex}


%%%%%%%%%%%%%%%% TRẢ LỜI NGẮN %%%%%%%%%%%%%%%%

% nguồn: Phan2_Lop_12_Dai_so_Chuong_4
\begin{bt}%[2D4H3-2]%[ma:4c346818]
	Ông A  làm một cái cửa nhà hình parabol có chiều cao từ mặt đất đến đỉnh là $2,25$ mét, chiều rộng tiếp giáp với mặt đất là $3$ mét. Giá thuê mỗi mét vuông là $150$ nghìn đồng. Vậy số tiền ông A phải trả là bao~nhiêu nghìn đồng?
	\shortans{675} 
	\loigiai{ 
	\immini{Do tính đối xứng của parabol nên ta có thể chọn hệ trục tọa độ $ Oxy$ sao cho $(P)$ có đỉnh $ I\in Oy$ như hình vẽ (đơn  vị trên trục là mét).\\
		Gọi parabol $(P)\colon y=a x^2+bx+c$.\\ 
		Ta xác  định được phương trình của parabol là $(P)\colon y=-x^2+\dfrac{9}{4}$.\\
		 Diện tích cửa là:\\
		 $S=\displaystyle\int_{-\tfrac{3}{2}}^{\tfrac{3}{2}}{\left(-x^2+\dfrac{9}{4}\right)\mathrm{\,d}x}=2\displaystyle\int_0^{\tfrac{3}{2}}{\left(-x^2+\dfrac{9}{4}\right)\mathrm{\,d}x}=\dfrac{9}{2}$ m$^2$.
		 }
		{\begin{tikzpicture}>=stealth,line join=round,line cap=round]
		\def\a{-1} \def\b{0}  \def\c{9/4}
		 \def\xmin{-2} \def\xmax{2.2}
		\def\ymin{-1}  \def\ymax{3}
		\draw[->] (\xmin,0)--(\xmax,0) node  [below]{$x$};
		\draw[->] (0,\ymin)--(0,\ymax) node [left]{$y$}; 
		\node at (0,0) [below left]{$O$}; 
		\fill (-1.5,0) circle  (1pt) node [below] {$-\dfrac{3}{2}$}
			(1.5,0) circle (1pt) node  [below] {$\dfrac{3}{2}$}
			 (0,2.25) circle (1pt) node [above right] {$\dfrac{9}{4}$}
			 (0,1) circle (1pt) node [left] {$1$};
		\clip (\xmin+0.1,\ymin+0.1)  rectangle (\xmax-0.1,\ymax-0.1);
		\draw[smooth,samples=300,domain=-1.5:1.5] plot (\x,{\a*(\x)^2+\b*(\x)+\c}); 
		\fill[pattern=north east lines,opacity=0.8] plot[domain=-1.5:1.5](\x,{\a*(\x)^2+\b*(\x)+\c})--cycle; 
	\end{tikzpicture}
		 }
		 Số tiền phải trả là $T=\dfrac{9}{2}\cdot 150=675$ nghìn đồng.
	} 
\end{bt}

% nguồn: Phan2_Lop_12_Dai_so_Chuong_4
\begin{bt}%[NBV2526]%[2D4H3-2]%[ma:511c5fdf]
	Bác Bình muốn nhờ thợ trang trí một bức tường hình chữ nhật $ADEF$ với kích thước $EF=6$\,m,  $DE=3$\,m sao cho cân xứng hai nửa.
	Phần gạch chéo là hình giới hạn bởi  đường gấp khúc $AFED$ và nửa đường tròn đường kính $AD$, được thuê sơn với đơn giá $250\,000$ đồng mỗi mét vuông.
	Phần màu trắng giới hạn bởi nửa đường tròn  đường kính $AD$ và một đường parabol (có đỉnh $H$ cách đường thẳng $AB$ một khoảng bằng $2$\,m và đi qua hai điểm $B$, $C$ nằm trên cạnh $AD$ thỏa mãn $AB=CD=1{,}5$\,m) được thuê trang trí bằng bức phù điêu đắp bằng xi măng với đơn giá $1\,950\,000$ đồng mỗi mét vuông (tham khảo hình vẽ).
	Hỏi bác Bình phải trả bao nhiêu triệu đồng  để trang trí bức tường như vậy (kết quả làm tròn đến hàng phần mười)?
	\begin{center}
		\begin{tikzpicture}[line join=round,font=\footnotesize,scale=1]
				\tikzset{declare function={f  (\x)=-8/9*(\x)^2+2;},smooth,samples=150}
				\path (-3,0)  coordinate (A)
					(-1.5,0) coordinate (B) 
					(1.5,0) coordinate  (C)
					(3,0) coordinate  (D)
					(3,3) coordinate  (E)
					(-3,3) coordinate  (F)
					(0,{f (0)})  coordinate (H);
				 \draw[fill=gray!20] plot[domain=-1.5:1.5] (\x,{f (\x)});
				 \draw[pattern=north east lines,pattern color=gray] (D) arc (0:180:3)--(F)--node[above]{$6$~m}
					(E)--node[right]{$3$~m}cycle--node[above]{$1{,}5$~m} 
					(C)--(A); 
				\draw[dashed]  (0,0)--(0,3);
				\foreach \i/\j in {A/-135,B/-110,C/-70,D/-45,E/45,F/135,H/35}\fill (\i) circle (1pt) node[shift={(\j:.2)}]{$\i$}; 
			\end{tikzpicture}
	\end{center}
	\shortans{20,7} 
	\loigiai{ 
		Ta  có $S_{ADEF}=6 \cdot 3=18$\, (m$^2$).\\
		Diện tích của hình tròn đường kính $AD$ là $S_1=\dfrac{1}{2}\pi \cdot  3^2 = \dfrac{9\pi}{2}$\, (m$^2$).\\
		Chi phí cho phong cách $F_1 = \left(18-\dfrac{9\pi}{2}\right) \cdot 0{,}25$  (triệu đồng).\\
		Xét hệ trục  tọa độ $Oxy$ sao cho $D(3;0)$, $E(3;3) \Rightarrow H(0;2)$, $C(1{,}5;0)$.\\
		Parabol nhận đỉnh là $H$ và đi qua $B$, $C$ có  phương trình dạng $y=ax^2+2$ thỏa mãn
		\[a\left(\dfrac{3}{2}\right)^2+2=0 \Leftrightarrow a =  -\dfrac{8}{9}. \]
		Suy ra phương trình parabol  $(P)$ là $y=-\dfrac{8}{9}x^2+2$.\\
		Diện tích  phần giới hạn bởi parabol và $AD$ là $S_2= \displaystyle \int_{-1{,}5}^{1{,}5}\left(-\dfrac{8}{9}x^2+2\right)\mathrm{d}x=4$\,(m$^2$).\\
		Chi phí cho phần đắp phù điêu là $F_2=(S_1-S_2) \cdot 1{,}95=\left(\dfrac{9\pi}{2}-4\right) \cdot  1{,}95$ (triệu đồng).\\
		Vậy  tổng chi phí là $F_1+F_2=\left(18-\dfrac{9\pi}{2}\right) \cdot 0{,}25 + \left(\dfrac{9\pi}{2}-4\right) \cdot 1{,}95 \approx 20{,}7$ (triệu đồng).
	 }
\end{bt}

% nguồn: Phan2_Lop_12_Dai_so_Chuong_4
\begin{bt}%[2D4H3-2]%[ma:ac35bdae]
	Một doanh nghiệp sản xuất một mặt hàng với chi phí cận biên  được mô tả bởi hàm số:
	$$ f(x) = \frac{1}{10}(x^2-16x+93) \text{ , } x \text{ là số sản  phẩm sản xuất}. $$
	Giả sử rằng doanh nghiệp bán được hết số lượng sản phẩm sản xuất được. Biết rằng doanh thu cận biên được mô tả bởi hàm số $g(x)=\left(\dfrac{4}{5}\right)^{x-8}+5$, với $x$ là số lượng sản phẩm được bán ra. Giả sử rằng tổng chi phí khi chưa sản xuất sản phẩm nào là 0 đồng và tổng doanh thu khi chưa bán được sản phẩm nào là $0$ đồng. Hỏi khi sản xuất $8$ sản phẩm và bán hết thì doanh nghiệp thu được lợi nhuận là bao nhiêu? (kết  quả làm tròn đến hàng đơn vị)
	 \shortans{22}
	 \loigiai{
		Hàm tổng chi phí $F(x)$ là nguyên hàm của hàm chi phí cận biên $f(x)$. Ta có $F(x)=\displaystyle \int f(x) \mathrm{d} x=\int \dfrac{1}{10}\left(x^2-16 x+93\right) \mathrm{d}x$ và  $F(x)=\dfrac{1}{10}\left(\dfrac{x^3}{3}-8 x^2+93 x\right)+C.$ \\
		Với điều kiện ban đầu $F(0)=0$ ta có $\dfrac{1}{10}\left(\dfrac{0^3}{3}-8  \cdot 0^2+93 \cdot 0\right)+C=0 \Rightarrow C=0.$\\
		 Do đó $F(x)=\dfrac{1}{10}\left(\dfrac{x^3}{3}-8 x^2+93 x\right).$\\
		 Hàm tổng doanh thu $G(x)$ là nguyên hàm của hàm doanh thu cận biên $g(x)$. \\
		Ta có $G(x)=\displaystyle \int g(x) \mathrm{d} x=\int\left(\left(\frac{4}{5}\right)^{x-8}+5\right) \mathrm{d} x$ và $G(x)=\dfrac{1}{\ln\dfrac{4}{5}}\left(\dfrac{4}{5}\right)^{x-8}+5 x+C.$\\ 
		Với điều kiện  ban đầu $G(0)=0$, suy ra $\dfrac{1}{\ln\dfrac{4}{5}}\left(\dfrac{4}{5}\right)^{-8}+5.0+C=0 \Rightarrow C=-\dfrac{1}{\ln\dfrac{4}{5}}\left(\dfrac{4}{5}\right)^{-8}.$\\
		Do  đó $G(x)=\dfrac{1}{\ln\dfrac{4}{5}}\left(\dfrac{4}{5}\right)^{x-8}+5 x-\dfrac{1}{\ln\dfrac{4}{5}}\left(\dfrac{4}{5}\right)^{-8}.$\\
		Khi đó  $F(8)=\dfrac{1}{10}\left(\dfrac{8^3}{3}-8.8^2+93.8\right)$ và $G(8)=\dfrac{1}{\ln\dfrac{4}{5}}\left(\dfrac{4}{5}\right)^{8-8}+5.8-\dfrac{1}{\ln\frac{4}{5}}\left(\dfrac{4}{5}\right)^{-8}.$\\
		Lợi  nhuận khi sản xuất và bán hết 8 sản phẩm là\\
	Lợi nhuận $=G(8)-F(8) \approx 21{,}96$ đồng.} 
\end{bt}

% nguồn: Phan2_Lop_12_Dai_so_Chuong_4
\begin{bt}%[2D4H3-2]%[ma:73826f8b]
	Để trang trí cho một phòng trong một tòa nhà, người ta vẽ lên tường một hình như sau: trên mỗi cạnh của hình lục giác đều có cạnh bằng $2$dm có một cánh hoa hình parabol, đỉnh của parabol cách cạnh $3$dm và nằm phía ngoài hình lục giác, đường parabol đó đi qua hai đầu mút của mỗi cạnh (xem  hình sau).
	\begin{center}
		\begin{tikzpicture}[x=0.75pt,y=0.75pt,yscale=-1,xscale=1]
				 \draw (290.17,179.33) -- (257.98,160.75) -- (257.98,123.58) -- (290.17,105) -- (322.35,123.58) -- (322.35,160.75) -- cycle;
				\draw (257.98,123.58).. controls (179.32,136.14) and (179.32,148.7)..  (257.98,161.25);
				\draw (290.17,105).. controls (239.96,43.15)  and (229.09,49.42).. (257.55,123.83);
				\draw  (257.98,161.25).. controls (229.52,235.65) and (240.4,241.93).. (290.6,180.08);
				 \draw (290.17,179.33).. controls (340.37,241.18) and (351.24,234.91).. (322.79,160.5);
				\draw (322.79,123.83).. controls (351.25,49.42) and  (340.38,43.15).. (290.17,105);
				\draw (322.35,160.75).. controls  (401.02,148.19) and (401.02,135.64).. (322.35,123.08);
			\end{tikzpicture}
	\end{center}
	Hãy tính diện tích của hình nói trên (kể cả hình lục giác đều) để mua sơn trang trí cho phù hợp  (kết quả làm tròn đến chữ số thập phân thứ nhất sau dấu phẩy).
	\shortans{34,4} 
	 \loigiai{
		Tổng diện tích $S$ của hình cần  tính bằng tổng diện tích của hình lục giác đều ($S_{lg}$) và diện tích của 6 cánh hoa hình parabol ($S_p$).
		$$ S = S_{lg}  + 6 \cdot S_p .$$
		Diện tích của của hình lục giác đều là 
		$$ S_{lg} = 6 \cdot \dfrac{2^2\sqrt{3}}{4} = 6\sqrt{3}  \text{ (dm}^2\text{ )} .$$
		Mỗi cánh hoa là một hình parabol. Áp dụng công thức tính diện tích hình parabol, diện tích bằng $\frac{2}{3}$ diện tích hình chữ nhật ngoại  tiếp.
		$$ S_p = \frac{2}{3} \times (\text{chiều rộng}) \times (\text{chiều cao})  = \frac{2}{3} \times 2 \times 3 = 4 \text{ (dm}^2\text{ )}. $$
		Tổng diện  tích là
		$$ S = S_{lg} + 6 \cdot S_p = 6\sqrt{3} + 6 \cdot 4 = 6\sqrt{3} + 24 \approx 34{,}3923  \text{ (dm}^2\text{ )}.$$
		Vậy, diện tích của hình nón trên là khoảng  \textbf{34{,}39 dm$^2$}.
	 }
\end{bt}

% nguồn: Phan2_Lop_12_Dai_so_Chuong_4
\begin{bt}%[2D4H3-2]%[ma:6b6db5aa]
	Các nhà kinh tế sử dụng đường cong Lorenz để minh họa sự phân phối thu nhập trong một quốc gia. Gọi $x$ là đại diện cho phần trăm số gia đình trong một quốc gia và $y$ là phần trăm tổng thu nhập, mô hình $y=x$ sẽ đại diện cho một quốc gia mà các gia đình có  thu nhập như nhau. Đường cong Lorenz $y=f(x)$, biểu thị sự phân phối thu nhập thực tế. Diện tích giữa hai mô hình này, với $0\leq x\leq 100$, biểu thị \lq\lq sự bất bình đẳng về thu nhập\rq\rq\ của một quốc gia. Năm 2009, đường cong Lorenz của Hoa Kỳ có thể được mô hình hóa bởi hàm số
	\[y=\left(0{,}00061x^2+0{,}0224x+1{,}666\right)^2,0\leq x\leq  100.\]
	Trong đó $x$ được tính  từ các gia đình nghèo nhất đến giàu có nhất. Sự bất bình đẳng về thu nhập của Hoa Kì vào năm 2009 bằng bao nhiêu (kết quả làm tròn đến chữ số hàng đơn vị)?
	\shortans{2086} 
	 \loigiai{
		Sự bất bình đẳng về thu nhập của Hoa Kì vào năm 2009 là diện tích hình phẳng $S$ giới hạn bởi các đường  $y=x$, $y=\left(0{,}00061x^2+0{,}0224x+1{,}666\right)^2$, $x=0$, $x=100$.\\
		Ta có  \[S=\int_0^{100}\left|\left(0{,}00061x^2+0{,}0224x+1{,}666\right)^2-x\right|\mathrm{\,d}x\approx 2\,086.\]
	 }
\end{bt}

% nguồn: Phan2_Lop_12_Dai_so_Chuong_4
\begin{bt}%[2D4H3-2]%[ma:a329dbf1]
	Tại một nhà máy ta gọi $C(x)$ là tổng chi phí (tính theo triệu đồng) để sản xuất $x$ tấn sản phẩm A trong một tháng. Khi đó đạo hàm $C'(x)$ gọi là chi phí cận biên, cho biết tốc độ tăng tổng chi  phí theo sản phẩm được sản xuất. Giả sử chi phí cận biên (tính theo triệu đồng trên tấn) của nhà máy được ước lượng bởi công thức:
	$$ C'(x)=5-0{,}06x+0{,}00072x^2 \text{ với } 0  \le x \le 150. $$
	Biết rằng $C(0)=30$ triệu đồng, gọi là chi phí cố định. Tính tổng chi phí khi nhà máy sản xuất  $100$ tấn sản phẩm A trong tháng.
	\shortans{470} 
	\loigiai{ 
		Ta  có
		\begin{align*}
			C(100)-C(0)&=\displaystyle \int_0^{100} C^{\prime}(x) \mathrm{d} x=\displaystyle \int_0^{100}\left(5-0{,}06 x+0{,}00072 x^2\right) \mathrm{d} x\\ 
			&=5 \int_0^{100}  \mathrm{~d} x-0.06 \int_0^{100} x \mathrm{~d} x+0.00072 \int_0^{100} x^2 \mathrm{~d} x=440.
		\end{align*}
		Vậy, tổng chi phí khi nhà máy sản xuất $100$ tấn sản phẩm trong tháng là  \textbf{$440+30=470$} triệu đồng.
	} 
\end{bt}

% nguồn: Phan2_Lop_12_Dai_so_Chuong_4
\begin{bt}%[2D4H3-2]%[ma:e914b9c3]
	Tính diện tích hình phẳng giới hạn bởi đồ thị hàm số $f(x)=3^x$, trục $O x$, đường thẳng $x=0$; $x=1$. (\textit{Làm tròn kết quả  đến hàng phần mười})
	 \shortans{1,8}
	\loigiai{Ta có 
		$$S=\displaystyle\int_0^1  3^x \mathrm{\,d} x=\dfrac{3^x}{\ln 3}\Big|_0 ^1=\dfrac{3}{\ln 3}-\dfrac{1}{\ln 3}=\dfrac{2}{\ln 3} \approx 1{,}8.$$
	} 
\end{bt}

% nguồn: Phan2_Lop_12_Dai_so_Chuong_4
\begin{bt}%[2D4H3-2]%[ma:8c1b0904]
	Một viên gạch hoa hình vuông cạnh 60 cm. Người thiết kế đã sử dụng bốn đường parabol có chung đỉnh tại tâm viên gạch để tạo ra bốn cánh hoa (được tô màu  sẫm như hình vẽ bên).
	\begin{center}
		\begin{tikzpicture}
				\draw  (2,2)--(2,-2)--(-2,-2)--(-2,2)--cycle;
				 \draw[fill=cyan,smooth,samples=200] plot[domain=-2:2,variable=\y]({0.5*(\y)^2},{\y})--plot[domain=2:0,variable=\x]({\x},{0.5*(\x)*(\x)})--plot[domain=0:2,variable=\x]({\x},{(-0.5)*(\x)*(\x)});
				\draw[fill=cyan,smooth,samples=200] plot[domain=-2:2,variable=\y]({-0.5*(\y)^2},{\y})--plot[domain=-2:0,variable=\x]({\x},{0.5*(\x)*(\x)})--plot[domain=0:-2,variable=\x]({\x},{(-0.5)*(\x)*(\x)}); 
			\end{tikzpicture}
	\end{center}
	Tính diện  tích phần cánh hoa của viên gạch.
	 \shortans{1200}
	\loigiai{ 
		\begin{center}
			\begin{tikzpicture}[scale=1,>=stealth]
					\draw[->]  (-1,0)--(3,0) node[below] {$x$};
					\draw[->] (0,-1)--(0,3) node[left]  {$y$};
					\draw[fill] (0,0) circle node[below left=-2pt]  {$O$};
					 \draw[dashed] (2,0)--(2,2)--(0,2);
					 \draw[fill] (2,0) circle (1pt) node[below] {$2$};
					\draw[fill] (0,2) circle  (1pt) node[left] {$2$};
					\draw[fill=cyan,smooth,samples=200] plot[domain=0:2,variable=\y]({0.5*(\y)^2},{\y})--plot[domain=2:0,variable=\x]({\x},{0.5*(\x)*(\x)}); 
				\end{tikzpicture}
		\end{center}
		Đặt hệ trục tọa độ $Oxy$ với gốc $O$ tại tâm của viên gạch. Khi đó, các  đỉnh của hình vuông có tọa độ là $(\pm 30, \pm 30)$.
		Cánh hoa trong góc phần tư thứ nhất bị giới hạn bởi  hai đường parabol
		\begin{itemize}
			\item Parabol có đỉnh tại $O(0{,}0)$, trục đối xứng là $Oy$ và đi qua điểm  $(30, 30)$.
			Phương trình có dạng $y=ax^2$. Thay tọa độ điểm  $(30{,}30)$ vào, ta có: $30 = a \cdot 30^2 \Rightarrow a=\dfrac{1}{30}$.
			 Vậy, phương trình parabol dưới là $y=\dfrac{1}{30}x^2$.
			\item Parabol có đỉnh tại $O(0{,}0)$,  trục đối xứng là $Ox$ và đi qua điểm $(30, 30)$.
			Phương trình có dạng $x=ay^2$. Tương tự, ta có $a=\dfrac{1}{30}$. 
			Vậy, phương trình parabol trên (xét theo biến $x$) là  $x=\dfrac{1}{30}y^2 \Rightarrow y = \sqrt{30x}$ (vì ta xét trong góc phần tư thứ nhất nên $y \ge 0$).
		\end{itemize}
		Hai  parabol này cắt nhau tại $O(0{,}0)$ và $(30{,}30)$.
		Diện tích một cánh hoa ($S_1$) trong góc phần tư thứ nhất được tính bằng tích phân của  hiệu hai hàm số từ $x=0$ đến $x=30$:
		$$ S_1 = \int_{0}^{30} \left(\sqrt{30x}  - \frac{1}{30}x^2\right) \mathrm{d}x = 300 \text{ (cm}^2\text{ )} .$$
		Tổng diện tích của bốn  cánh hoa là
		 $$ S = 4 \cdot S_1 = 4 \cdot 300 = 1200 \text{ (cm}^2\text{ )}. $$
		 Vậy, diện tích phần cánh hoa của viên gạch là \textbf{1200 cm$^2$}.
	} 
\end{bt}

% nguồn: Phan2_Lop_12_Dai_so_Chuong_4
\begin{bt}%[2D4H3-2]%[ma:0e9a1c15]
	 \immini{
	Trong cuộc thi thiết kế logo cho câu lạc bộ NBK ART, một bạn học sinh đã sử dụng phần mềm GeoGebra để thiết kế ra một mẫu logo và gửi đi dự thi. Logo được tạo bằng  cách vẽ một hình tròn bán kính bằng $\sqrt{29}$ cm, sau đó cắt bốn góc đối xứng nhau bởi bốn parabol giống nhau như hình vẽ dưới đây. Phần hình phẳng còn lại sau khi cắt được tô màu, chính là hình dạng cuối cùng của logo dùng để dự thi. Tính diện tích (đơn vị: cm$^2$) của logo đó (làm tròn kết quả đến hàng đơn vị).
	}{\begin{tikzpicture}[scale=0.5,font=\footnotesize,line join=round,line cap=round,>=stealth]
		 \pgfmathsetmacro\a{sqrt (29)}
		\begin{scope}
			\clip (-5,-5) rectangle  (5,5);
			\draw (0,0) circle  (\a);
		\end{scope}
		\draw[<->] (2,5)--(-2,5) node[pos=0.5,above]{$4$ cm}; 
		\draw[<->] (0,1)--(0,5) node[pos=0.5,right]{$4$ cm}; 
		\draw[smooth,samples=100,domain=-2:2] plot (\x,{(\x)^2+1}); 
		 \draw[smooth,samples=100,domain=-2:2,rotate=-90] plot (\x,{(\x)^2+1});
		\draw[smooth,samples=100,domain=-2:2,rotate=90] plot (\x,{(\x)^2+1}); 
		\draw[smooth,samples=100,domain=-2:2,rotate=180] plot  (\x,{(\x)^2+1});
	\end{tikzpicture}}
	 \shortans {$44$}
	\loigiai{ 
		\begin{center}
			\begin{tikzpicture}[scale=0.5,font=\footnotesize,line join=round,line cap=round,>=stealth]
					\pgfmathsetmacro\a{sqrt  (29)}
					\begin{scope}
						\draw (0,0) circle  (\a);
					\end{scope}
					\draw[->] (-6,0)--(7,0) node[above] {$x$}; 
					\draw[->]  (0,-6)--(0,7) node[left] {$y$};
					\draw (0,0) circle (1pt) +  (-135:0.5) node{$O$};
					\draw  (0,0)--(2,5) circle (1pt) +(90:0.5) node{$A$} node[sloped,above,pos=0.5] {$\sqrt{29}$ cm};
					\draw[dashed] (0,5)--(2,5)--(2,0); 
					 \draw[smooth,samples=100,domain=-2:2] plot (\x,{(\x)^2+1});
					\draw[smooth,samples=100,domain=-2:2,rotate=-90] plot (\x,{(\x)^2+1}); 
					 \draw[smooth,samples=100,domain=-2:2,rotate=90] plot (\x,{(\x)^2+1});
					\draw[smooth,samples=100,domain=-2:2,rotate=180]  plot (\x,{(\x)^2+1});
					\foreach  \x/\g in {2/-90}
					\draw (\x,1pt)--(\x,-1pt) + (\g:0.4) node {$\x$}; 
					\foreach  \y/\g in {5/180}
					\draw (1pt,\y) --(-1pt,\y) + (\g:0.4)  node {$\y$};
				\end{tikzpicture}
		\end{center}
		Từ đồ thị hàm số, ta  có $(P) \colon y =x^2+1$.\\
		Phương trình đường tròn là 
		\[x^2+y^2=29 \Leftrightarrow y=  \sqrt{29-x^2} \text{ (Nửa dường tròn phía trên trục hoành)}.\]
		Diện tích phần giới hạn  của parabol và đường tròn là
		\[V'=\displaystyle \int_{-2}^2 \left(\sqrt{29-x^2}-x^2-1\right) \mathrm{\,d}x=11{,}7.\] 
		Diện tích  của logo là
		 \[V=\pi \sqrt{29}^2-4 \cdot 11{,}7 \approx 44.\]
	 }
\end{bt}

% nguồn: Phan2_Lop_12_Dai_so_Chuong_4
\begin{bt}%[2D4H3-2]%[ma:676a18c5]
	Một mảnh vườn hoa dạng hình  tròn có bán kính bằng $6$ m. Phần đất trồng hoa là phần tô màu trong hình vẽ bên. Biết hai hình chữ nhật $ABCD$, $MNPQ$ có $AB=MQ=6$ m và kinh phí trồng hoa là $100\,000$ đồng/m$^2$. Tính số tiền cần để trồng hoa? (đơn vị: triệu đồng, kết quả làm tròn đến hàng phần chục).
	\begin{center}
		\begin{tikzpicture}[scale=0.7,font=\footnotesize,line join=round,line cap=round,>=stealth]
				\coordinate  (O) at (0,0);
				\coordinate (A1) at  (3,3);
				\coordinate (A2)  at (-3,3);
				\coordinate (A3)  at (-3,-3);
				\coordinate (A4)  at (3,-3);
				 \draw[name path=(C)] (O) circle (6 cm);
				\coordinate (N)  at ($(O) + (30:6)$);
				\coordinate  (B) at ($(O) + (60:6)$);
				 \coordinate (A) at ($(O) + (120:6)$);
				\coordinate  (M) at ($(O) + (150:6)$);
				\coordinate  (Q) at ($(O) + (210:6)$);
				\coordinate (D) at ($(O)  + (240:6)$);
				\coordinate (C) at ($(O) +  (300:6)$);
				\coordinate (P) at  ($(O) + (330:6)$);
				\draw  (N)--(A1)--(B)
					 (A)--(A2)--(M)
					(Q)--(A3)--(D) 
					 (C)--(A4)--(P);
				\fill[pattern=dots]plot[samples=200,domain=-3:3,smooth,variable=\x] (\x,{(36-(\x)^2)^(1/2)})--plot[samples=200,domain=3:-3,smooth,variable=\x]  (\x,{-(36-(\x)^2)^(1/2)})--cycle;
				 \fill[pattern=dots]plot[samples=200,domain=5.196:6,smooth,variable=\x] (\x,{(36-(\x)^2)^(1/2)})--plot[samples=200,domain=6:5.196,smooth,variable=\x] (\x,{-(36-(\x)^2)^(1/2)})--cycle;
				 \fill[pattern=dots]plot[samples=200,domain=-6:-5.196,smooth,variable=\x] (\x,{(36-(\x)^2)^(1/2)})--plot[samples=200,domain=-5.196:-6,smooth,variable=\x] (\x,{-(36-(\x)^2)^(1/2)})--cycle;
				 \fill[pattern=dots](A)--(B)--(A1)--(N)--(P)--(A4)--(C)--(D)--(A3)--(Q)--(M)--(A2)--(A)--cycle;
				\foreach \i/\g in  {N/0,B/90,A/90,M/180,Q/180,D/-90,C/-90,P/0}{\draw(\i) circle (1pt)
					($(\i)+(\g:3mm)$)  node[scale=1]{$\i$};}
			\end{tikzpicture}
	\end{center}
	 \shortans{10,2}
	 \loigiai{
		Chọn hệ trục tọa  độ $Oxy$ sao cho như hình vẽ.
		\begin{center}
			\begin{tikzpicture}[scale=0.7,font=\footnotesize,line join=round,line cap=round,>=stealth]
					\tikzset{every  node/.style={scale=1.1}}
					\draw[->,red] (-7,0)--(7,0)  node[below left] {$x$};
					\draw[->,red] (0,-7)--(0,7) node[below left]  {$y$};
					\draw (0,0) node [below left] {$O$}; 
					\foreach \x/\nx in {-3/-3,3/3} 
					\draw[thin] (\x,1pt)--(\x,-1pt) node [below] {$\nx$}; 
					 \coordinate (O) at (0,0);
					\coordinate (A1) at  (3,3);
					 \coordinate (A2) at (-3,3);
					\coordinate (A3) at (-3,-3); 
					\coordinate  (A4) at (3,-3);
					 \draw[name path=(C)] (O) circle (6 cm);
					\coordinate  (N) at ($(O) + (30:6)$);
					\coordinate (B) at  ($(O) + (60:6)$);
					\coordinate (A) at ($(O) + (120:6)$); 
					\coordinate (M) at ($(O) +  (150:6)$);
					\coordinate (Q) at  ($(O) + (210:6)$);
					\coordinate (D) at  ($(O) + (240:6)$);
					\coordinate  (C) at ($(O) + (300:6)$);
					\coordinate (P) at ($(O) + (330:6)$); 
					 \draw (N)--(A1)--(B)
						 (A)--(A2)--(M)
						 (Q)--(A3)--(D)
						 (C)--(A4)--(P);
					\fill[pattern=dots,thin]plot[samples=200,domain=-3:3,smooth,variable=\x]  (\x,{(36-(\x)^2)^(1/2)})--plot[samples=200,domain=3:-3,smooth,variable=\x] (\x,{-(36-(\x)^2)^(1/2)})--cycle;
					\fill[pattern=dots]plot[samples=200,domain=5.196:6,smooth,variable=\x]  (\x,{(36-(\x)^2)^(1/2)})--plot[samples=200,domain=6:5.196,smooth,variable=\x] (\x,{-(36-(\x)^2)^(1/2)})--cycle;
					 \fill[pattern=dots]plot[samples=200,domain=-6:-5.196,smooth,variable=\x] (\x,{(36-(\x)^2)^(1/2)})--plot[samples=200,domain=-5.196:-6,smooth,variable=\x] (\x,{-(36-(\x)^2)^(1/2)})--cycle;
					\fill[pattern=dots](A)--(B)--(A1)--(N)--(P)--(A4)--(C)--(D)--(A3)--(Q)--(M)--(A2)--(A)--cycle; 
					\foreach  \i/\g in {N/0,B/90,A/90,M/180,Q/180,D/-90,C/-90,P/0}{\draw(\i) circle (1pt)
						($(\i)+(\g:3mm)$)  node[scale=1]{$\i$};}
					\draw(5.196,0) circle (1pt)  node[below]{$3\sqrt{3}$};
					\draw(0,5.196)  circle (1pt) node[left]{$3\sqrt{3}$};
					\draw(0,3)  circle (1pt) node[left]{$3$};
					\draw[dashed](0,3)--(3,3)--(3,0) 
						(5.196,3)--(5.196,0); 
				\end{tikzpicture}
		\end{center}
		Khi đó phương trình đường  tròn tâm $O$ bán kính bằng $6$ là $x^2+y^2=36$.\\
		 Diện tích phần trồng hoa là
		 \[S=\pi\cdot 6^2-4\displaystyle\int ^{3\sqrt{3}}_{3}\left(\sqrt{36-x^2}-3\right)\mathrm{\,d}x\approx 101{,}75\ (\text{m}^2).\]
		Đổi $100\,000$ đồng  $=0{,}1$ triệu đồng.\\
		Số tiến cần  để trồng hoa là
		\[S\cdot  0{,}1\approx 10{,2}\ (\text{triệu đồng}).\]
	 }
\end{bt}

% nguồn: Phan2_Lop_12_Dai_so_Chuong_4
\begin{bt}%[2D4H3-2]%[ma:3d607054]
	Người ta sử dụng đường cong Lorenz để minh họa sự phân phối thu  nhập trong một quốc gia. Gọi $x$ là đại diện cho phần trăm số gia đình trong một quốc gia và $y$ là phần trăm tổng thu nhập, mô hình $y=x$ sẽ đại diện cho một quốc gia mà các gia đình có thu nhập như nhau. Đường cong Lorenz $y=f(x)$, biểu thị sự phân phối thu nhập thực tế. Diện tích giữa hai mô hình này, với $0 \le x \le 100$, biểu thị "sự bất bình đẳng về thu nhập" của một quốc gia. Năm $2009$, đường cong Lorenz của Hoa Kỳ có thể được mô hình hóa bởi hàm số
	$$ y=(0{,}00061x^2+0{,}0224x+1{,}666)^2, 0 \le x \le 100.  $$
	Trong đó, $x$ được tính từ các gia đình nghèo nhất đến giàu có nhất. Sự bất bình đẳng về thu nhập của Hoa Kì vào năm 2009 là bao nhiêu? (kết  quả làm tròn đến hàng đơn vị)
	\shortans{2086} 
	\loigiai{Sự bất bình đẳng về thu nhập của Hoa Kì vào năm $2009$ là diện tích hình phẳng $S$  giới hạn bởi
		$\heva{y&=x  \\
		y&=\left(0{,}00061 x^2+0{,}0224  x+1{,}666\right)^2 \\
		 x&=0 \\ x&=100}$ \\
		Diện tích hình phẳng là 
		$S=\displaystyle \int_0^{100}\left|\left(0{,}00061 x^2+0{,}0224  x+1{,}666\right)^2-x\right| \mathrm{d} x \approx 2085{,}5.$
	 }
\end{bt}

% nguồn: Phan2_Lop_12_Dai_so_Chuong_4
\begin{bt}%[2D4V3-2]%[ma:88b28bda]
	\immini[thm]{Nhà bác Ba có tất cả $8$ cánh cửa sắt hình chữ nhật với chiều dài $2$\,m và chiều rộng $1$\,m. Hai mặt của mỗi cánh cửa được thiết kế như hình vẽ dưới đây. Trong đó,  phần được tô đậm được sơn màu xanh, phần còn lại được sơn màu trắng. Mỗi phần sơn màu trắng có đường biên cong là một phần của parabol có đỉnh nằm trên cạnh của hình chữ nhật. Biết rằng chi phí để sơn màu xanh là $120$ nghìn đồng/m$^2$ và chi phí sơn màu trắng là $110$ nghìn đồng/m$^2$. Hỏi để sơn toàn bộ số cửa sắt trên, bác Ba phải tốn bao nhiêu triệu đồng (làm tròn kết quả đến hàng phần trăm theo đơn vị triệu đồng).}
	{\begin{tikzpicture}[scale=0.45,line join=round,line cap=round,>=stealth,x=4cm,y=4cm]
		\tikzset{every node/.style={scale=0.8}} 
		\def \ha{4*(\x)^2} 
		\draw [fill=cyan!50] plot[samples=50,domain=0:0.5,smooth,variable=\x] (\x,{\ha}) 
			 --plot[shift={(1,0)},samples=50,domain=-0.5:0,smooth,variable=\x] (\x,{\ha})
			--plot[shift={(1,0)},yscale=-1,samples=50,domain=0:-0.5,smooth,variable=\x] (\x,{\ha}) 
			 --plot[yscale=-1,samples=50,domain=0.5:0,smooth,variable=\x] (\x,{\ha});
		\draw (0,-1) rectangle  (1,1);
		\draw [dashed,<->](-0.1,0)--(-0.1,1) node[pos=0.5,sloped,above]{$1$ m}; 
		\draw [dashed,<->](0,1.1)--(0.5,1.1) node[pos=0.5,sloped,above]{$0{,}5$ m}; 
	\end{tikzpicture}}
	 \shortans{3,63}
	 \loigiai{
		\begin{center}
			\begin{tikzpicture}[scale=0.7,line join=round,line cap=round,>=stealth,x=4cm,y=4cm]
					\tikzset{every  node/.style={scale=0.8}}
					\def \xmin{-0.7} 
					\def \xmax{0.7} 
					\def  \ymin{-0.2}
					 \def \ymax{1.5}
					\def  \ha{4*(\x)^2}
					\fill  [pattern=dots](0,0)--(0.5,0)--(0.5,1)-- plot[samples=50,domain=0.5:0,smooth,variable=\x] (\x,{\ha});
					 \draw[->] (\xmin,0)--(\xmax,0) node[above] {$x$};
					\draw[->] (0,\ymin)--(0,\ymax) node[left] {$y$}; 
					\draw  (0,0) node [below right] {$O$};
					\foreach  \x/\g/\t in {-0.5/-90/{-0{,}5},0.5/-90/{0{,}5}}
					\draw[thin] (\x,1pt)--(\x,-1pt)+(\g:0.1) node[scale=1]{$\t$}; 
					\foreach \y/\g  in {1/135}
					\draw[thin]  (1pt,\y)--(-1pt,\y)+(\g:0.1) node[scale=1]{$\y$};
					\begin{scope}
						\clip (\xmin+0.01,\ymin+0.01) rectangle  (\xmax-0.01,\ymax-0.01);
						 \draw[samples=50,domain=-0.7:0.7,smooth,variable=\x] plot (\x,{\ha});
					\end{scope}
					\foreach \x/\y in  {0/0,0.5/0,0.5/1}
					\fill (\x,\y)  circle (1pt);
					\foreach \x/\y in  {-0.5/1,0.5/1}
					\draw  [dashed] (\x,0) |-(0,\y);
					\path (0.35,0.2) node  [circle,fill=white,inner sep=0] {$S_1$};
					\begin{scope}[shift={(-2.5,0.5)}]
						\draw [fill=cyan!50] plot[samples=50,domain=0:0.5,smooth,variable=\x] (\x,{\ha}) 
						--plot[shift={(1,0)},samples=50,domain=-0.5:0,smooth,variable=\x] (\x,{\ha}) 
						--plot[shift={(1,0)},yscale=-1,samples=50,domain=0:-0.5,smooth,variable=\x] (\x,{\ha}) 
						--plot[yscale=-1,samples=50,domain=0.5:0,smooth,variable=\x]  (\x,{\ha});
						\draw (0,-1) rectangle (1,1); 
						 \draw[red,->] (0,0)--(1.2,0) node[above] {$x$};
						\draw[red,->]  (0,0)--(0,1.2) node[left] {$y$};
						 \draw [red](0,0) node [below left] {$O$};
						\draw [dashed,<->](-0.1,0)--(-0.1,1) node[pos=0.5,sloped,above]{$1$ m}; 
						 \draw [dashed,<->](0,1.1)--(0.5,1.1) node[pos=0.5,sloped,above]{$0{,}5$ m};
					\end{scope}
				\end{tikzpicture}
		\end{center}
	\immini{Xét đồ thị parabol trên  hệ trục tọa độ $Oxy$ như sau\\
		 Gọi parabol có phương trình là $y=a x^2+b x+c$.\\
		Parabol đi qua các điểm có tọa  độ $(-0{,}5; 1)$, $(0{,}5; 1)$, $(0; 0)$, nên ta có hệ sau \[\heva{&c=0\\
		&a \cdot 0{,}5^2+b \cdot 0{,}5=1\\ 
		&a \cdot\left(-0{,}5\right)^2+b \cdot(-0{,}5)=1}  \Leftrightarrow\heva{&c=0\\&a=4\\&b=0.}\]
		Vậy ta có phương  trình parabol $y=4x^2$.
		} 
		{\begin{tikzpicture}[line cap=butt,line join=miter,>=stealth,font=\normalsize]
		\tikzset{declare function={xmin=-1.5; xmax=1.5; ymin=-1; ymax=2.5; 
			f  (\x)=4*(\x)^2;
			}, 
			 smooth,samples=450
		 }
		 \draw[dashed](-.5,0) node[below,scale=.7]{$-\dfrac{1}{2}$}|-(0,1) node[above left]{$1$}
			 (.5,0) node[below,scale=.7]{$\dfrac{1}{2}$}|-(0,1);
		\draw[->] (xmin,0)--(xmax,0) node[shift={(-100:7pt)}]{$ x $}; 
		\draw[->] (0,ymin)--(0,ymax) node[shift={(190:7pt)}]{$ y $}; 
		 \fill (0,0) node[shift={(-45:9pt)}]{$ O $};
		\clip (xmin,ymin) rectangle  (xmax,ymax);
		\draw  plot[domain=xmin:xmax] (\x,{f (\x)});
	\end{tikzpicture}}
	\immini{Diện tích  hình phẳng giới hạn bởi $y=4x^2$, $y=0$, $x=0$, $x=0{,}5$ là
		 \[S_1=\displaystyle\int_0^{0{,}5} 4x^2\mathrm{\,d}x=\dfrac{1}{6}.\]
		Diện tích của phần tô đậm  được sơn màu xanh của mỗi mặt cánh cửa là \[S_2=4\cdot S_1=4\cdot \dfrac{1}{6}=\dfrac{2}{3}.\]}
		{\begin{tikzpicture}[line cap=butt,line join=miter,>=stealth,font=\normalsize]
		\tikzset{declare function={xmin=-1.5; xmax=1.5; ymin=-1;  ymax=2.5;
			 f (\x)=4*(\x)^2;
			}, 
			 smooth,samples=450
		} 
		 \draw[dashed](-.5,0) node[below,scale=.7]{$-\dfrac{1}{2}$}|-(0,1) node[above left]{$1$}
			(.5,0) node[below,scale=.7]{$\dfrac{1}{2}$}|-(0,1); 
		\fill[pattern=north west lines](0,0)--plot[domain=0:.5]  (\x,{f (\x)})--(.5,1)--(.5,0);
		\draw[->] (xmin,0)--(xmax,0) node[shift={(-100:7pt)}]{$  x $};
		\draw[->] (0,ymin)--(0,ymax) node[shift={(190:7pt)}]{$ y  $};
		\fill (0,0) node[shift={(225:10pt)}]{$ O $}; 
		\clip (xmin,ymin)  rectangle (xmax,ymax);
		\draw  plot[domain=xmin:xmax] (\x,{f (\x)});
	\end{tikzpicture}}
		\noindent Diện tích của phần không tô đậm tô  đậm được sơn màu trắng của mỗi mặt cánh cửa là \[S_3=1\cdot 2-S_2=2-\dfrac{2}{3}=\dfrac{4}{3}.\]
		Có  tất cả $8$ cánh cửa sắt hình chữ nhật, mỗi cánh có $2$ mặt, vậy có tất cả $16$ mặt như hình vẽ.\\
		Vậy 
		\begin{itemize}
			\item Tổng diện tích phần tô đậm được sơn màu  xanh là $16S_2=16\cdot \dfrac{2}{3}=\dfrac{32}{3}~\left(\mathrm{m}^2\right)$.
			\item Tổng diện tích phần không tô đậm  được sơn màu trắng là $16S_3=16\cdot \dfrac{4}{3}=\dfrac{64}{3}~\left(\mathrm{m}^2\right)$.
			\item Tổng chi phí để sơn màu xanh là $\dfrac{32}{3} \cdot 120=1\,280$  (nghìn đồng) $=1{,}28$ (triệu đồng).
			\item Tổng chi phí để sơn màu trắng là $\dfrac{64}{3}  \cdot 110\approx 2\,346$ (nghìn đồng) $\approx 2{,}35$ (triệu đồng).
			\item Tổng chi  phí sơn là $1{,}28+2{,}35=3{,}63$ (triệu đồng).
		\end{itemize}
	 }
\end{bt}

% nguồn: Phan2_Lop_12_Dai_so_Chuong_4
\begin{bt}%[2D4V3-2]%[ma:79a611a3]
	Một nhóm học sinh lớp $12$ đã lên bản thiết kế mẫu  hoa văn cho một loại gạch men lát nền nhà. Các em đã vẽ $4$ đường cong như hình, từ đó tạo thành một hình $(H)$ khép kín ở giữa viên gạch để tạo điểm nhấn. Cụ thể cách dựng hình được thực hiện như sau
	\begin{itemize}
		\item Dựng hệ trục $Oxy$ với điểm $O$ là tâm của viên gạch, tia $Ox$ hướng sang phải và tia $Oy$ hướng lên trên,  đơn~vị trên mỗi trục là $5$ cm.
		\item Các em  lấy $O$ là tâm viên gạch và $A$ là trung điểm một cạnh viên gạch, xác định được điểm $B$ thỏa mãn $\overrightarrow{OB}=\dfrac{5}{6} \overrightarrow{OA}$.
		\item Dựng đường thẳng $\Delta\colon 5x-9=0$. Đường cong $(L_1)$ là tập hợp các  điểm $M$ thỏa mãn $3MB=5\mathrm{d}(M, \Delta)$.
		\item Lấy đối xứng đường cong $(L_1)$ qua tâm $O$ và qua các đường chéo của viên gạch thì được các đường  cong còn lại.
	\end{itemize}
	Biết viên gạch là hình vuông có kích thước $60$ cm; hỏi diện tích hình $(H)$ là bao nhiêu centimét  vuông (làm tròn đến hàng đơn~vị)?
	\begin{center}
		\begin{tikzpicture}[scale=0.3]
				\draw (0,0) rectangle  (12,12);
				\coordinate  (O) at (6,6);
				\fill (O)  circle (2pt);
				 \draw[domain=9:12,smooth,variable=\x,samples=200]
					plot  ({\x},{6 + 4*sqrt ((\x - 6)^2/9 - 1)});
				 \draw[domain=9:12,smooth,variable=\x,samples=200]
					plot ({\x},{6 - 4*sqrt ((\x - 6)^2/9  - 1)});
				\draw[domain=9:12,smooth,variable=\x,samples=200] 
					 plot ({6 + 4*sqrt ((\x - 6)^2/9 - 1)},{\x});
				 \draw[domain=9:12,smooth,variable=\x,samples=200]
					plot ({6 - 4*sqrt ((\x - 6)^2/9  - 1)},{\x});
				 \draw[domain=9:12,smooth,variable=\x,samples=200]
					plot ({12 - \x},{12 - (6 + 4*sqrt ((\x -  6)^2/9 - 1))});
				 \draw[domain=9:12,smooth,variable=\x,samples=200]
					plot ({12 - \x},{12 - (6 - 4*sqrt ((\x -  6)^2/9 - 1))});
				 \draw[domain=9:12,smooth,variable=\x,samples=200]
					plot ({12 - (6 + 4*sqrt  ((\x - 6)^2/9 - 1))},{12 - \x});
				 \draw[domain=9:12,smooth,variable=\x,samples=200]
					plot ({12 - (6 - 4*sqrt ((\x - 6)^2/9 -  1))},{12 - \x});
			\end{tikzpicture}
		 \hspace{0.5 cm}
		\begin{tikzpicture}[scale=0.3]
				\path  (6,6) coordinate (O)
					 (12,6) coordinate (A)
					($(O)!5/6!(A)$) coordinate  (B)
					(9.5,{6 + 4*sqrt ((9.5  - 6)^2/9 - 1)}) coordinate (M);
				\draw  (0,0) rectangle (12,12);
				 \draw[domain=9:12,smooth,variable=\x,samples=200]
					plot ({\x},{6 + 4*sqrt ((\x - 6)^2/9 - 1)}); 
				\draw[domain=9:12,smooth,variable=\x,samples=200] 
					plot ({\x},{6 - 4*sqrt ((\x -  6)^2/9 - 1)});
				\node  at (10.8,3.5) {$(L_1)$};
				\draw[domain=9:12,smooth,variable=\x,samples=200] 
					plot ({6 + 4*sqrt ((\x - 6)^2/9  - 1)},{\x});
				\draw[domain=9:12,smooth,variable=\x,samples=200] 
					plot ({6 - 4*sqrt ((\x  - 6)^2/9 - 1)},{\x});
				\node at  (6,10) {$(L_2)$};
				 \draw[domain=9:12,smooth,variable=\x,samples=200]
					plot ({12 - \x},{12  - (6 + 4*sqrt ((\x - 6)^2/9 - 1))});
				 \draw[domain=9:12,smooth,variable=\x,samples=200]
					plot ({12 - \x},{12  - (6 - 4*sqrt ((\x - 6)^2/9 - 1))});
				\node at  (1.5,7.5) {$(L_3)$};
				 \draw[domain=9:12,smooth,variable=\x,samples=200]
					plot ({12 - (6 + 4*sqrt ((\x - 6)^2/9  - 1))},{12 - \x});
				 \draw[domain=9:12,smooth,variable=\x,samples=200]
					plot ({12 - (6  - 4*sqrt ((\x - 6)^2/9 - 1))},{12 - \x});
				\node at  (4,1.5) {$(L_4)$};
				\draw  (O)--(A);
				\foreach \t/\g in  {A/0,B/-90,O/-90,M/0}{
					\draw (\t) circle (1pt) node[shift={(\g:7pt)},font=\scriptsize]{$ \t  $};
				 }
			\end{tikzpicture}
	\end{center}
	 \shortans{1168}
	\loigiai{ 
		\begin{center}
			\begin{tikzpicture}[scale=0.3]
					\path (6,6) coordinate  (O)
						(12,6) coordinate (A) 
						 ($(O)!5/6!(A)$) coordinate (B)
						(9.5,{6 + 4*sqrt ((9.5 - 6)^2/9 -  1)}) coordinate (M);
					\draw[->] (-1,6)--(14,6) node[below]{$x$}; 
					\draw[->]  (6,-1)--(6,14) node[left]{$y$};
					\draw (0,0)  rectangle (12,12);
					 \draw[domain=9:12,smooth,variable=\x,samples=200]
						plot ({\x},{6 + 4*sqrt  ((\x - 6)^2/9 - 1)});
					 \draw[domain=9:12,smooth,variable=\x,samples=200]
						plot ({\x},{6 - 4*sqrt  ((\x - 6)^2/9 - 1)});
					\node at (10.8,3.5)  {$(L_1)$};
					 \draw[domain=9:12,smooth,variable=\x,samples=200]
						plot ({6  + 4*sqrt ((\x - 6)^2/9 - 1)},{\x});
					\draw[domain=9:12,smooth,variable=\x,samples=200] 
						plot ({6 - 4*sqrt ((\x - 6)^2/9 -  1)},{\x});
					\node  at (7.5,10) {$(L_2)$};
					\draw[domain=9:12,smooth,variable=\x,samples=200] 
						plot ({12 - \x},{12 - (6 + 4*sqrt ((\x - 6)^2/9 - 1))}); 
					\draw[domain=9:12,smooth,variable=\x,samples=200] 
						 plot ({12 - \x},{12 - (6 - 4*sqrt ((\x - 6)^2/9 - 1))});
					\node  at (1.5,7.5) {$(L_3)$};
					\draw[domain=9:12,smooth,variable=\x,samples=200] 
						 plot ({12 - (6 + 4*sqrt ((\x - 6)^2/9 - 1))},{12 - \x});
					 \draw[domain=9:12,smooth,variable=\x,samples=200]
						plot ({12 - (6 - 4*sqrt ((\x - 6)^2/9  - 1))},{12 - \x});
					\node at (4,1.5) {$(L_4)$}; 
					 \draw (O)--(A);
					\foreach \t/\g in  {A/-60,B/-90,O/230,M/0}{
						\draw (\t) circle (1pt) node[shift={(\g:7pt)},font=\scriptsize]{$  \t $};
					 }
				\end{tikzpicture}
		\end{center}
		Ta chọn hệ trục tọa độ như  hình vẽ với mỗi đơn~vị trên trục bằng $5$ cm.\\
		Đường cong $(L_1)$ là tập hợp điểm $M$ thỏa $3MB=5\mathrm{d}(M, \Delta)$ hay \[\dfrac{MB}{\mathrm{d}(M, \Delta)}=\dfrac{5}{3}.\qquad  (1)\]
		Phương trình $\Delta\colon  x=\dfrac{9}{5}$.\qquad (2)\\
		Từ (1) và (2) ta thấy $M$ thuộc một nhánh của hyperbol  $(L_1)\colon \dfrac{x^2}{a^2}-\dfrac{y^2}{b^2}=1$; trong đó điểm $B(5;0)$ là một trong hai tiểu điểm của $(L_1)$ nên $c=5$.\\
		Đường chuẩn $x=\dfrac{a}{\mathrm{e}}=\dfrac{a^2}{c}=\dfrac{9}{5} \Rightarrow a=3$.\\ (thử lại ta thấy  $\dfrac{MB}{d(M, \Delta)}=\dfrac{5}{3}=e$ (hợp lí).\\
		Do  $c^2=a^2+b^2\Rightarrow b=\sqrt{c^2-a^2}=4$.\\
		Phương trình  $(L_1)\colon \dfrac{x^2}{9}-\dfrac{y^2}{16}=1~(x > 3)$.\\
		 Xét giao điểm của $(L_1)$ với đương thẳng $y=x$, ta có $\dfrac{x^2}{9}-\dfrac{x^2}{16}=1~(x > 3)\\
		\Rightarrow x=\dfrac{12\sqrt{7}}{7}  \approx 4{,}54$.\\
		Từ phương trình $\dfrac{x^2}{9}-\dfrac{y^2}{16}=1\Rightarrow \dfrac{y^2}{16}=\dfrac{x^2}{9}-1\Rightarrow y=4\sqrt{\dfrac{x^2}{9}-1}~(x > 3)$.\\ 
		Diện tích cần tính là \[S=8\left(\displaystyle\int _0^3x\mathrm{\,d}x+\displaystyle\int_3^{\tfrac{12\sqrt{7}}{7}}\left(x-4\sqrt{\dfrac{x^2}{9}-1} \right)\mathrm{\,d}x \right)\cdot 25\approx 1\,168~ \mathrm{cm}^2.\] 
	 }
\end{bt}

% nguồn: Phan2_Lop_12_Dai_so_Chuong_4
\begin{bt}%[2D4V3-2]%[ma:1a6277dd]
	Họa sĩ thiết kế một loại gạch trang trí có dạng như \textbf{Hình 1}, gạch có dạng hình vuông cạnh $8\,$dm. Khi đặt bản vẽ trong hệ tọa độ $(Oxy)$ với đơn vị của mỗi trục là $1\,$dm thì mỗi nét cong phía trong thuộc một trong hai đường hypebol $y=-\dfrac{4}{x}$ và $y=\dfrac{4}{x}$ (\textbf{Hình 2}), các cạnh của viên gạch lần lượt thuộc $4$ đường thẳng $x=-4$, $x=4$, $y=-4$, $y=4$. Người ta sơn màu hồng vào phần hình được gạch chéo như \textbf{Hình 3}. Diện tích phần sơn  màu hồng là bao nhiêu dm$^2$ (làm tròn kết quả đến hàng phần mười).\\
	\begin{center}
		\resizebox{!}{5cm}{\begin{tikzpicture}
		\draw[thin,opacity=.5] (-4,-4)  grid (4,4);
		\draw[smooth]  plot[domain=1:4](\x,{4/(\x)})
			 plot[domain=-4:-1](\x,{4/(\x)})
			 plot[domain=-4:-1](\x,{-4/(\x)})
			plot[domain=1:4](\x,{-4/(\x)}); 
		\draw (-4,-4) rectangle  (4,4);
		 \draw (0,-4) node[below,scale=2]{\textbf{Hình 1}};
	\end{tikzpicture}}
		\resizebox{!}{5cm}{\begin{tikzpicture}[line cap=butt,line join=miter,>=stealth]
		\draw[thin,opacity=.5] (-4,-4)  grid (4,4);
		\draw[smooth]  plot[domain=1:4](\x,{4/(\x)})
			plot[domain=-4:-1](\x,{4/(\x)}) 
			 plot[domain=-4:-1](\x,{-4/(\x)})
			 plot[domain=1:4](\x,{-4/(\x)});
		\draw (-4,-4)  rectangle (4,4);
		\draw[->] (-4.5,0)--(4.5,0) 
			 node[shift={(-100:7pt)},font=\normalsize]{$x$};
		\draw[->]  (0,-4.5)--(0,4.5)
			 node[shift={(180:7pt)},font=\normalsize]{$y$};
		\draw (5pt,0) |- (0,5pt) 
			(0,0) node[shift={(225:9pt)},font=\normalsize]{$O$}; 
		\foreach  \x in {-4,-3,-2,-1,1,2,3,4}{
			\fill (\x,0) circle (1pt)  node[shift={(-90:9pt)},fill=white]{$\x$};
		} 
		 \foreach \y in {-4,-3,-2,-1,1,2,3,4}{
			\fill (0,\y) circle (1pt)  node[shift={(180:9pt)},anchor=east,fill=white,inner sep=0pt]{$\y$};
		} 
		\draw  (0,-4) node[below,scale=2]{\textbf{Hình 2}};
	\end{tikzpicture}}
		\resizebox{!}{5cm}{\begin{tikzpicture}
		 \draw[thin,opacity=.5] (-4,-4) grid (4,4);
		\draw[smooth,pattern = north  west lines] plot[domain=1:4](\x,{4/(\x)})
			-- plot[domain=4:1](\x,{-4/(\x)})-- plot[domain=-1:-4](\x,{4/(\x)})-- 
			 plot[domain=-4:-1](\x,{-4/(\x)})--cycle;
		\draw (-4,-4) rectangle (4,4); 
		\draw (0,-4)  node[below,scale=2]{\textbf{Hình 3}};
	\end{tikzpicture}}
	\end{center}
	 \shortans{38,2}
	\loigiai{Ta chia phần sơn màu  hồng thành $5$ phần diện tích $S_1$, $S_2$, $S_3$, $S_4$, $S_5$ như hình dưới đây.
		\begin{center}
			\resizebox{!}{5cm}{\begin{tikzpicture}[>=stealth]
		\draw[thin,opacity=.5]  (-5.5,-5.5) grid
			(5.5,5.5); 
		\draw[->] (-5.5,0) --  (5.5,0) node[below] {$x$};
		\draw[->] (0,-5.5)  -- (0,5.5) node[left]{$y$};
		\draw (-0.1,0) node[above  right] {$O$};
		\foreach \x in  {1,2,3,4}
		 \draw[shift={(\x,0)}] (0pt,2pt) -- (0pt,-2pt) node[below right]{$\x$};
		\foreach \x in  {-4,-3,-2,-1}
		\draw[shift={(\x,0)}] (0pt,2pt) -- (0pt,-2pt)  node[above left]{$\x$};
		\foreach  \y in {1,2,3,4}
		\draw[shift={(0,\y)}] (2pt,0pt)  -- (-2pt,0pt) node[above right]{$\y$};
		\foreach  \y in {-1,-2,-3,-4}
		\draw[shift={(0,\y)}] (2pt,0pt) -- (-2pt,0pt) node[below right]{$\y$}; 

		\draw[smooth] plot[domain=0.727:5.5](\x,{4/(\x)})  node[above]{$y=\dfrac{4}{x}$};
		\draw[smooth] plot[domain=-0.727:-5.5](\x,{4/(\x)})  node[below]{$y=\dfrac{4}{x}$};
		 \draw[smooth] plot[domain=-0.727:-5.5](\x,{-4/(\x)}) node[above]{$y=-\dfrac{4}{x}$};
		 \draw[smooth] plot[domain=0.727:5.5](\x,{-4/(\x)}) node[below]{$y=-\dfrac{4}{x}$};
		 \draw[smooth] plot[domain=-5.5:5.5]({-4,(\x)});
		 \draw[smooth] plot[domain=-5.5:5.5]({4,(\x)});
		\draw[smooth]  plot[domain=-5.5:5.5](\x,{-4});
		\draw[smooth]  plot[domain=-5.5:5.5](\x,{4});
		 \draw (0,-5.5) node[below,scale=2]{\textbf{Hình 2}};
	\end{tikzpicture}}~
			\resizebox{!}{5cm}{\begin{tikzpicture}
		 \draw[thin,opacity=.5] (-4,-4) grid (4,4);
		\draw[smooth]  plot[domain=1:4](\x,{4/(\x)})
			 plot[domain=-4:-1](\x,{4/(\x)})
			plot[domain=-4:-1](\x,{-4/(\x)}) 
			plot[domain=1:4](\x,{-4/(\x)}); 
		\draw 
			 (-4,0)--(-1,0)
			(4,0)--(1,0) 
			 (-1,4)--(-1,-4)
			 (1,4)--(1,-4)
			(-4,-4)  rectangle (4,4);
		 \path (-2,1) node {$S_1$}
			(2,1) node  {$S_2$}
			(2,-1) node {$S_3$} 
			(-2,-1)  node {$S_4$}
			 (0,0) node {$S_5$};
		 \draw (0,-4) node[below,scale=1.5]{\textbf{Hình 1}};
	\end{tikzpicture}}
		\end{center}
		Ta nhận thấy $S_1$, $S_2$,  $S_3$, $S_4$ có diện tích bằng nhau nên ta chỉ cần tính một trong bốn diện tích trên.\\
		Ta có $S_1$ là diện tích hình phẳng được giới hạn bởi đường  thẳng $x=-4$, $x=-1$, đồ thị $y=-\dfrac{4}{x}$ và trục hoành $Ox$ nên
		 \[S_1=\displaystyle \int_{-4}^{-1} \left|\dfrac{-4}{x} \right| \mathrm{\,d}x =4 \ln 4 .\]
		$S_5$ là diện tích hình chữ nhật có hai  kích thước lần lượt là $2$ dm và $8$ dm nên $S_5=2 \cdot 8=16$ dm$^2$.\\
		Vậy diện tích phần sơn màu hồng (làm tròn kết quả đến hàng phần mười) là \\$S_1+S_2+S_3+S_4+S_5=4 \cdot S_1 +S_5  =16 \ln 4+16 \approx 38{,}2\,$ dm$^2$.
	} 
\end{bt}

% nguồn: Phan2_Lop_12_Dai_so_Chuong_4
\begin{bt}%[2D4V3-2]%[ma:d04eec49]
	Một bức tường hình chữ nhật $ABCD$ có $AB=6\,m$, $BC=4\, m$ được bạn An trang trí bằng cách vẽ một Parabol có trục đối xứng là trục đối xứng của hình chữ nhật và đỉnh của Parabol cách $CD$ một đoạn là $1\, m$, các nhánh Parabol chia hình chữ  nhật thành ba phần $H_1$, $H_2$, $H_3$ như hình vẽ bên dưới ($HK=2\, m$).\\
	\begin{center}
		\begin{tikzpicture}[line join=round,line cap=round,>=stealth,font=\footnotesize,scale=1]
				 \path
					(-3,4) coordinate (A) 
					(3,4)  coordinate (B)
					 (-3,0) coordinate (D)
					(3,0) coordinate  (C)
					({sqrt (5)},4) coordinate (I) 
					({-sqrt (5)},4)  coordinate (P)
					(1,0) coordinate  (K)
					(-1,0)  coordinate (H);
				\begin{scope}
					\clip (-3,0) rectangle  (3,4);
					 \fill[yellow!30] (-3,0) rectangle (3,4);
					\fill[cyan!30] 
						(A) -- (D) --  (H)
						-- plot[domain=-1:-3,samples=100] (\x,{(\x)^2-1}) 
						-- (P) -- cycle; 
					 \fill[green!50]
						(B) -- (C) -- (K) 
						--  plot[domain=1:3,samples=100] (\x,{(\x)^2-1})
						-- (I)  -- cycle;
				\end{scope}
				\begin{scope}
					\clip (-3,0) rectangle (3,4); 
					\draw[samples=200,domain=-3:3,smooth,variable=\x] plot (\x,{1*(\x)^2-1}); 
				\end{scope}
				\draw  (A)--(B) node[above,midway,sloped]{$6$\,m}--(C) node[right,midway]{$4$\,m}--(K)--(H) node[below,midway,sloped]{$2$\,m}--(D)--(A)
					(-2.2,2)  node{$H_1$}
					(0,2) node{$H_2$} 
					 (2.2,2) node{$H_3$};
				\foreach \i/\g in  {A/135,B/45,C/-45,D/-135,K/-90,H/-90}{\draw(\i) circle (1pt)
					($(\i)+(\g:3mm)$)  node[scale=1]{$\i$};}
			\end{tikzpicture}
	\end{center}
	Phần $H_1$ được sơn màu xanh da trời, phần $H_2$ được sơn màu vàng, phần $H_3$ được sơn màu xanh lá cây. Biết rằng mỗi hộp sơn các màu chỉ sơn được $3\,(m^2)$ tường, đồng thời giá  của hộp sơn màu xanh da trời là $120\,000$ đồng/hộp, hộp sơn màu vàng là $130\,000$ đồng/hộp, hộp sơn màu xanh lá cây là $110\,000$ đồng/hộp. Tính giá tiền bạn An phải trả để sơn bức tường này. Biết rằng cửa hàng sơn chỉ bán số nguyên hộp (\textit{đơn vị là triệu đồng và làm tròn kết quả đến hàng phần trăm}).
	\shortans{1,11} 
	\loigiai{ 
		\begin{center}
			\begin{tikzpicture}[line join=round,line cap=round,>=stealth,font=\footnotesize,scale=1]
					 \path
						(0,0) coordinate (O) 
						 (-3,4) coordinate (A)
						(3,4)  coordinate (B)
						 (-3,0) coordinate (D)
						(3,0) coordinate (C) 
						(0,-1) coordinate  (E)
						({sqrt (5)},4) coordinate (I) 
						({sqrt (5)},0)  coordinate (J)
						({-sqrt  (5)},4) coordinate (P)
						(1,0) coordinate  (K)
						(-1,0)  coordinate (H);
					\begin{scope}
						\clip (-3,0) rectangle (3,4); 
						\fill[yellow!30] (-3,0) rectangle  (3,4);
						 \fill[cyan!30]
						 (A) -- (D) -- (H)
						-- plot[domain=-1:-3,samples=100] (\x,{(\x)^2-1}) 
						-- (P) --  cycle;
						\fill[green!50] 
						(B) --  (C) -- (K)
						-- plot[domain=1:3,samples=100] (\x,{(\x)^2-1}) 
						--  (I) -- cycle;
					\end{scope}
					\begin{scope}
						\clip  (-3,-2) rectangle (3,4);
						\draw[samples=200,domain=-3:3,smooth,variable=\x] plot (\x,{1*(\x)^2-1}); 
					\end{scope}
					 \draw[->] (-4,0)--(4,0) node[below left] {$x$};
					\draw[->] (0,-1.5)--(0,5.1) node[below left]  {$y$};
					\draw  (A)--(B)--(C)--(K)--(H)--(D)--(A);
					\draw[dashed]  (I)--(J);
					\fill (P) circle (1pt); 
					\foreach \i/\g in {A/135,B/45,C/-45,D/-135,K/-80,H/-120,I/90,J/-90,E/-120,O/-45}{\draw(\i) circle (1pt) 
						($(\i)+(\g:3mm)$) node[scale=1]{$\i$};} 
				\end{tikzpicture}
		\end{center}
		Ta thiết lập hệ trục tọa độ $Oxy$ với $E(0;-1)$  là đỉnh của Parabol như hình vẽ.\\
		Gọi Parabol cần tìm có dạng $y=ax^2+c$  với ($a\neq 0$).\\
		 Do Parabol đi qua điểm $K(1;0)$ và có đỉnh $E(0;-1)$ nên suy ra $a=1$ và $c=-1$.\\
		Do đó  Parabol là $y=x^2-1$.\\
		Gọi giao điểm của Parabol với đường thẳng $AB$ là điểm $I$ $ (IB<IA)$.\\ 
		Suy  ra điểm $I\left(\sqrt{5};4 \right)$.\\
		Gọi hình chiếu vuông góc của điểm $I$ lên $DC$ là $J$.\\ 
		Ta tính diện tích $S_0$ của hình phẳng giới  hạn bởi các đường thẳng $x=\sqrt{5}$, $y=0$ và parabol\\
		$$S_0=\displaystyle\int_1^{\sqrt{5}} \left( x^2-1\right)  \mathrm{\,d}x=\left( \dfrac{x^3}{3}-x\right) \bigg|_1^{\sqrt{5}}=\dfrac{2\sqrt{5}+2}{3}\,(m^2).$$
		Suy ra diện tích  hình $H_3$ và $H_1$ là\\
		 $$S_{H_3}=S_{H_1}=S_0+S_{IJCB}=\dfrac{2\sqrt{5}+2}{3}+4\cdot \left(3-\sqrt{5} \right)\approx 5{,}21\,\,(m^2).$$
		Vậy với mỗi diện tích $H_1$  và $H_3$ bạn An cần mua $2$ hộp sơn mỗi loại.\\
		Diện tích  hình $H_2$ là
		$$S_{H_2}=S_{ABCD}-2S_{H_3}\approx 6\cdot4-2\cdot 5{,}21=13{,}58\,\,(m^2).$$ 
		Vậy với diện tích hình $H_3$ bạn An cần mua  $5$ hộp sơn màu vàng.\\
		Tổng số tiền bạn An cần để  mua sơn là
		\begin{align*}
			T=2\cdot 120\,000 + 2\cdot 110\,000 + 5\cdot 130\,000=1\,\,110\,\,000\,\,\,\, (\text{đồng})=1{,}11\,\,\,\, (\text{triệu đồng}). 
		\end{align*}
	 }
\end{bt}

% nguồn: Phan2_Lop_12_Dai_so_Chuong_4
\begin{bt}%[2D4V3-2]%[ma:67a337a3]
	Một ô tô đang chạy với vận tốc $10~\text{ m/s}$ thì gặp chướng ngại vật, người lái xe đạp phanh. Từ thời điểm đó, ô tô chuyển động chậm dần đều với gia tốc $a=-2~\text{ m/s}^2$, trong đó $t$ là thời gian được tính bằng giây kể từ lúc bắt đầu đạp phanh. Hỏi ô tô đi chuyển được bao nhiêu mét kể từ  lúc đạp phanh đến khi dừng hẳn?
	 \shortans{25}
	\loigiai{ 
		 Phương trình vận tốc của xe là
		$$ 
		v(t) = v_0 + at = 10-2t. 
		 $$
		 Xe dừng lại khi $v(t)=0$ hay
		$$ 
		 10-2t=0
		 \Leftrightarrow
		t=5. 
		$$ 
		Trong khoảng thời gian từ khi đạp thắng ($t=0$) tới khi dừng lại $(t=5)$, xe luôn đi tới (không đổi chiều chuyển động) nên quãng đường nó  đi được là
		$$ 
		s =  \displaystyle\int_{0}^{5} v(t) \;\mathrm{d}t
		=  \displaystyle\int_{0}^{5} (10-2t) \;\mathrm{d}t
		= 25. 
		$$ 
	} 
\end{bt}

% nguồn: Phan2_Lop_12_Dai_so_Chuong_4
\begin{bt}%[2D4V3-2]%[ma:2463d52f]
	\immini{Người ta thiết kế một mẫu gạch lát nền nhà có dạng hình vuông, cạnh $4$ dm. Bốn góc viên gạch màu trắng, phần ở giữa màu xanh (tham khảo hình vẽ). Đường viền của phần màu xanh bao gồm $4$ đoạn thẳng  nằm trên các cạnh của hình vuông và bốn đường cong có tính chất: Tích khoảng cách từ một điểm bất kỳ thuộc đường cong đó đến hai trục đối xứng của viên gạch (hai đường thẳng đi qua tâm viên gạch và song song với hai cạnh vuông góc) bằng $2$ dm$^2$.}{
	\resizebox{!}{4cm}{\begin{tikzpicture}[declare function={a=2; b=1;}]
		\path (-b,a)  coordinate (A)
			(b,a)  coordinate (B)
			 (a,b) coordinate (C)
			 (a,-b) coordinate (D)
			 (b,-a) coordinate (E)
			 (-b,-a) coordinate (G)
			 (-a,-b) coordinate (H)
			(-a,b)  coordinate (K)
			(-a,a)  coordinate (M)
			(a,a) coordinate  (N)
			(a,-a) coordinate (P) 
			 (-a,-a) coordinate (Q);
		\begin{scope}
			\clip  (N) rectangle (Q);
			\draw[fill=cyan!80!blue] (A)--(B)--plot[domain=0.1:a](\x,{2/(\x)}) 
				 --(D)--plot[domain=a:0.1](\x,{-2/(\x)})--(G)--plot[domain=0.1:a](-\x,{-2/(\x)})
				--(K)--plot[domain=a:0.1](-\x,{2/(\x)})--cycle; 
		\end{scope}
		\draw (N) rectangle  (Q);
	\end{tikzpicture}}}
	\noindent Hãy cho biết phần màu xanh có diện tích bằng bao  nhiêu decimét vuông (làm tròn kết quả đến hàng phần mười)?
	\shortans{13,5} 
	 \loigiai{
		\begin{center}
			\begin{tikzpicture}[declare function={a=2; b=1;},>=stealth]
					\path  (-b,a) coordinate (A)
						 (b,a) coordinate (B)
						(a,b) coordinate  (C)
						(a,-b)  coordinate (D)
						(b,-a)  coordinate (E)
						(-b,-a) coordinate (G) 
						(-a,-b) coordinate  (H)
						(-a,b) coordinate (K) 
						(-a,a) coordinate (M) 
						 (a,a) coordinate (N)
						(a,-a)  coordinate (P)
						(-a,-a)  coordinate (Q);
					\begin{scope}
						\clip (N) rectangle (Q); 
						 \draw[fill=cyan!15] (A)--(B)--plot[domain=0.1:a](\x,{2/(\x)})
						 --(D)--plot[domain=a:0.1](\x,{-2/(\x)})--(G)--plot[domain=0.1:a](-\x,{-2/(\x)})
						--(K)--plot[domain=a:0.1](-\x,{2/(\x)})--cycle; 
					\end{scope}
					\draw (N) rectangle  (Q);
					\fill (1,0) circle  (1pt) node[below right]{$1$}
						 (-1,0) circle (1pt) node[below left]{$-1$}
						(2,0) circle (1pt)  node[below right]{$2$}
						(-2,0)  circle (1pt) node[below left]{$-2$}
						 (0,2) circle (1pt) node[above right]{$2$}
						(0,-2)  circle (1pt) node[below right]{$-2$}
						(0,1)  circle (1pt) node[above right]{$1$}
						 (0,-1) circle (1pt) node[below right]{$-1$};
					\draw[dashed] (A)--(G) 
						 (B)--(E)
						 (K)--(C)
						(H)--(D); 
					\foreach  \x/\g in {A/90,B/90,C/0,D/0,E/-90,G/-90,H/180,K/180}
					\fill (\x) circle (1pt) node[shift={(\g:9pt)}]{$\x$}; 
					\draw[->]  (-2.5,0)--(3,0) node[below]{$x$};
					\draw[->]  (0,-2.5)--(0,3) node[left]{$y$};
					\fill (0,0) circle  (1pt) node[below left]{$O$};
				\end{tikzpicture}
		\end{center}
		Gắn hệ trục tọa độ $O xy$ vào viên gạch sao cho  hai trục trùng với hai đường đối xứng, gốc $O$ ở tâm hình vuông (xem hình).\\
		Giả sử tọa độ một  điểm nằm trên đường viền cong là $\left(x; y \right)$. \\
		 Theo giả thiết, ta có $|x y| =2 \Rightarrow \hoac{& y = \dfrac{2}{x} \\ & y = - \dfrac{2}{x}.}$ \\
		Các đường viền $AK$, $DE$ là một phần của đồ thị hàm số $y = - \dfrac{2}{x}$.  \\
		Các đường viền  $BC$, $GH$ là một phần của đồ thị hàm số $y = \dfrac{2}{x}$. \\
		Vậy,  diện tích phần màu xanh là
		\allowdisplaybreaks 
		\begin{eqnarray*}
			S &=&\displaystyle\int_{-2}^{-1} \left|-\dfrac{2}{x} - \dfrac{2}{x}\right| \mathrm{\, d}x + \displaystyle\int_{1}^{2} \left|\dfrac{2}{x} - \dfrac{- 2}{x}\right| \mathrm{\, d}x + S_{ABEG}\\ 
			&=&\displaystyle\int_{-2}^{-1} \left(-\dfrac{2}{x} - \dfrac{2}{x}\right) \mathrm{\, d}x + \displaystyle\int_{1}^{2} \left( \dfrac{2}{x} + \dfrac{2}{x}\right) \mathrm{\, d}x + 8  \\
			&=& - 4  \ln |x| \Bigr|^{-1}_{-2} + 4 \ln |x| \Bigr|^{2}_{1} +8 \approx 13{,}5\, \left( \mathrm{dm}^2 \right).
		\end{eqnarray*}
	 }
\end{bt}

% nguồn: Phan2_Lop_12_Dai_so_Chuong_4
\begin{bt}%[2D4V3-2]%[ma:35cbed8d]
	Trong lễ kỷ niệm 30 năm thành lập trường THPT Chu Văn An – Biên Hòa (1994 - 2024), nhà trường có cho dựng 1 cổng chào hình parabol  với điểm cao nhất là $6m$ và chiều rộng giữa hai chân cổng chào là $8m$. Hỏi diện tích của cổng chào là bao~nhiêu mét vuông?
	 \shortans{32}
	\loigiai{ 
		Đặt hệ quy chiếu $Oxy$, vào 1 điểm chân cổng chào.\\ 
		Khi đó,  ta cần xác định hàm số $y=ax^2+bx+c$, biết parabol có đỉnh $I(4;6)$ và đi qua gốc tọa độ $O(0;0)$.\\
		Thay tọa độ điểm $O$ vào hàm  số, suy ra: $c=0$\\
		Ta có hệ phương trình: $\heva{&16a+4b=6  \\ &8a+b=0} \Leftrightarrow \heva{&a=\dfrac{-3}{8} \\ &b=3}$\\
		Suy ra hàm  số $y = \dfrac{-3}{8}x^2+3x$\\
		Diện tích của cổng chào: $S  = \displaystyle \int_0^8\left|\dfrac{-3}{8}x^2+3x \, dx \right| = 32 \, (m^2)$
	 }
\end{bt}

% nguồn: Phan2_Lop_12_Dai_so_Chuong_4
\begin{bt}%[2D4V3-2]%[ma:8ce1d338]
	Khu vực trung tâm một quảng trường (xem hình vẽ bên dưới) có dạng hình tròn đường kính $AB$ bằng $12$ m. Người ta trang trí khu vực này bằng hai đường parabol đối xứng nhau qua $AB$, nằm trong hình tròn, đi qua các điểm $A$, $B$ và có đỉnh cách  mép hình tròn $1$ m. Phần giới hạn bởi $2$ parabol được trồng hoa với chi phí $300$ nghìn đồng mỗi mét vuông, phần còn lại được lát gốm sứ với chi phí $900$ nghìn đồng mỗi mét vuông. Tổng chi phí để hoàn thành khu vực này là bao nhiêu triệu đồng (làm tròn kết quả đến hàng đơn vị)?
	\begin{center}
		\begin{tikzpicture}[line cap=butt,line join=miter,>=stealth]
				\path (0,0) coordinate (O); 
				\tikzset{declare function={xmin=-4; xmax=4; ymin=-3; ymax=4; 
					 f (\x)=-0.22*(\x)^2+2;
					 g (\x)=0.22*(\x)^2-2;
					 },
					 smooth,samples=450
				 }
				\fill  (0,0) node[shift={(-90:7pt)},font=\normalsize]{$ O $};
				\fill  (0,0) circle (1pt);
				\path (0,2) node[below]{$I$} 
					(0,3) node[above]{$J$} 
					 (0,2)--(0,3) node[right,pos=.5]{$1$ m}
					(-3,0)  node[left]{$A$}
					 (3,0) node[right]{$B$};
				\clip  (xmin,ymin) rectangle (xmax,ymax);
				 \draw plot[domain=-3:3] (\x,{f (\x)});
				 \draw plot[domain=-3:3] (\x,{g (\x)});
				\draw 
					(-3,0)--(3,0) 
					 (0,2)--(0,3)
					(O)  circle (3);
			\end{tikzpicture}
	\end{center}
	 \shortans{54}
	 \loigiai{
		Ta  có diện tích phần parabol giới hạn bởi parabol đi qua $A$, $I$, $B$ và đường thẳng $AB$ bằng \[\dfrac{2}{3} AB\cdot OI=\dfrac{2}{3} \cdot 12\cdot 5=40\left(\mathrm{m}^2\right).\]
		Do đó diện tích phần trồng hoa bằng  $40\cdot 2=80\left(\mathrm{m}^2\right)$.\\
		Suy ra diện tích phần lát gốm sứ bằng $\pi  \cdot 6^2-80=36\pi-80\left(\mathrm{m}^2\right)$.\\
		 Vậy tổng chi phí bằng $80\cdot 0{,}3+0{,}9\cdot (36\pi-80) \approx 54$ (triệu đồng).
	} 
\end{bt}

% nguồn: Phan2_Lop_12_Dai_so_Chuong_4
\begin{bt}%[2D4V3-2]%[ma:3bc0e3dd]
	Một khoảng sân của một ngôi nhà có hình dạng là một phần của parabol, nếu chọn hệ trục tọa độ $Oxy$ như hình vẽ bên dưới thì parabol có phương trình là $y=x^2$. Chủ nhà muốn thiết kế một hồ nuôi cá cảnh có dạng hình tròn bán kính $1$ m, hình tròn được thiết kế tiếp xúc với $2$ nhánh  của parabol, phía trên của hình tròn ngăn bởi một bức tường mỏng tiếp xúc với hình tròn. Phần đất còn lại của khoảng sân giới hạn bởi đường tròn, $2$ nhánh của parabol và bức tường được thiết kế trồng hoa (phần gạch sọc trong hình). Tính diện tích mà chủ nhà thiết kế để trồng hoa, biết rằng đơn vị trên mỗi trục tọa độ là mét. (Kết quả làm tròn đến hàng phần trăm).
	\begin{center}
		\begin{tikzpicture}[scale=1.5,font=\footnotesize,line join=round,line cap=round,>=stealth]
				 \draw[->] (-2,0) -- (2,0) node[below]{$x$};
				\draw (0,0) node[below left]{$O$}circle (1pt); 
				\draw (1.5,2.25)  node[right]{$A$}circle (1pt);
				 \draw (-1.5,2.25) node[left]{$B$}circle (1pt);
				\draw (-1.5,2.25)--(1.5,2.25); 
				\clip (-2,3.7) rectangle  (2.7,-0.5);
				\draw[smooth,samples=100,domain=-1.7:1.7]  plot (\x,{(\x)^2});
				 \fill[pattern=crosshatch] plot[domain=-1.5:1.5](\x,{(\x)^2})--cycle;
				 \draw[fill=white] (0,1.25) circle (1);
				\draw (0,1.25) node[left]{$I$}circle (1pt); 
				\draw (1.6,2.7) node[right]{$y=x^2$}; 
				\draw[->] (0,-0.5) -- (0,3)  node[left]{$y$};
				 \draw (0,1.25)--(0.87,0.75);
			\end{tikzpicture}
	\end{center}
	\shortans{1,36} 
	\loigiai{ 
		Gọi $I(0;y_0)$ là tâm của đường  tròn suy ra
		$$(C)\colon x^2+(y-y_0)^2=1\Rightarrow \hoac{&y=y_0+\sqrt{1-x^2}\\&y=y_0-\sqrt{1-x^2.}}$$ 
		Do Parabol luôn  tiếp xúc nửa dưới của đường tròn nên ta có
		 $$\heva{&y_0-\sqrt{1-x^2}\\&\left(y_0-\sqrt{1-x^2}\right)'=(x^2)'}\Rightarrow x^2=\dfrac{3}{4}\Rightarrow y_0=\dfrac{5}{4}.$$
		Đường thẳng  $AB$ có phương trình $y=\dfrac{9}{4}$ và điểm $A\left(\dfrac{3}{2};\dfrac{9}{4}\right)$, điểm $B\left(-\dfrac{3}{2};\dfrac{9}{4}\right)$.\\
		Diện tích cần tìm  là $S=\displaystyle\int_{-\frac{3}{2}}^{\frac{3}{2}} \left(\dfrac{9}{4}-x^2\right) \mathrm{\,d}x - 1^2\cdot \pi\approx 1{,}36$.
	 }
\end{bt}

% nguồn: Phan2_Lop_12_Dai_so_Chuong_4
\begin{bt}%[2D4V3-2]%[ma:e315d747]
	Một sảnh ga  có thông số như hình dưới, được thông gió bằng hai quạt, mỗi quạt có lưu lượng
	 khoảng $80$m$^3$\,/phút.
	\begin{center}
		\begin{tikzpicture}[line join=round,line cap=round,scale=0.7,transform shape,>=stealth,declare function={a=4; b=a; h=2; f (\x)=-0.5* (\x)^2+2*(\x)+h; g (\x)=f (\x-5)+3;}]
				\definecolor{arsenic}{rgb}{0.23,0.27,0.29} 
				\definecolor{burntorange}{rgb}{1.0,0.4,0.0} 
				\tikzset{fan/.pic={ 
				\draw (0,.8) circle  (2.25cm);
				\draw (0,.8) circle (2.15cm); 
				 \draw (0,.8) circle (1.42cm);
				\draw (0,.8) circle  (1.48cm);
				\draw (0,.8) circle  (6mm);
				\draw[fill=arsenic] (0,.8) circle (5mm); 
				 \def\N{
					 (0,.8)
					..  controls +(145:1.3) and +(170:1).. (0,2.8)
					..  controls +(-10:1.4) and +(-20:1).. (.6,1.76)
					.. controls +(160:.4) and  +(100:.4).. (.2,.9) -- cycle
				 }
				\foreach \i/\j/\k in {0/0/0,120/.7/-1.2,240/-.7/-1.2} 
					{ 
				\draw[rotate=\i,shift={(\j,\k)}]  \N;
				\fill[burntorange,rotate=\i,shift={(\j,\k)}] \N; 
				 }
					}} 
				 \path
					(0,0) coordinate (A) 
					(a,0)  coordinate (D)
					(5,3) coordinate (B) 
					($(B)+(D)-(A)$)  coordinate (C)
				\foreach \p in {A,B,C,D}{ 
					($(\p)+(90:h)$)  coordinate (\p')};
				\fill[red!15]  (1.8,{f (1.8)}) --(6,{g (6)}) plot[domain=6:9] (\x,{g (\x)})--(4,2)--plot[domain=4:1] (\x,{f (\x)})--cycle;
				\draw[gray!70](1.8,{f (1.8)}) --(6,{g  (6)});
				\draw[red,domain=0:4] plot (\x,{f (\x)}); 
				\draw[red,domain=6:9] plot (\x,{g (\x)}); 
				\draw[dashed,domain=5:6] plot  (\x,{g (\x)});
				 \draw[dashed] (A')--(B')--(C')
					(B')--(B)--(C) 
					 (A)--(B);
				\draw (A')--(D')--(D)--(A) node[pos=0.5,above]{$20$\,m}--cyclenode[pos=0.5,below,sloped]{$8$\,m} 
					 (D')--(C')--(C)--(D) node[pos=0.5,above,sloped]{$60$\,m};
				\draw[|<->|](A')--++(0,h) node[pos=0.5,above,sloped]{$10$\,m}; 
				\foreach  \x/\goc in {A/160,B/180,C/0,D/0,A'/160,B'/180,C'/-30,D'/0}{
					\fill (\x)  circle (1pt);
				} 
				 \path (1.2,2.5) pic[scale=.25]{fan}
					 (2.8,2.5) pic[scale=.25]{fan};
			\end{tikzpicture}
	\end{center}
	Biết mái sảnh là kết cấu gỗ dạng parabol. Hỏi cần bao nhiêu phút để trao đổi toàn bộ không khí trong sảnh?  (Kết quả làm tròn đến hàng đơn vị).
	\shortans{110} 
	\loigiai{ 
		Chọn hệ trục tọa độ  $Oxy$ như hình (trục $Ox$ ngang, trục $Oy$ thẳng đứng).
		\begin{center}
			\begin{tikzpicture}[scale=1,line join=round,line cap=round,>=stealth]
					 \tikzset{declare function={a=-0.5; b=0; c=3; f (\x)=a*(\x)^2+b*(\x)+c;
						xI=-b/2*a; yI=f (xI);  xmin=xI-2.5; xmax=xI+2.5;
						ymin=yI-0.5; ymax=yI+1;}, 
						 smooth,samples=100
					 }
					 \draw[->] (xmin-0.5,0)--(0,0) node[below left]{$O$} --(xmax+1,0) node[below left] {$x$};
					 \draw[->] (0,-0.5)--(0,ymax) node[below right] {$y$};
					\fill  (0,0) circle (1pt);
					 \draw[red,domain=-2.45:2.45] plot (\x,{f (\x)});
					\fill (-2.45,0) circle (1pt)  node[shift={(-90:7pt)}]{$-10$}
						(2.45,0) circle (1pt) node[shift={(-90:7pt)}]{$10$} 
						 (0,3) circle (1pt) node[shift={(35:7pt)}]{$10$};
				\end{tikzpicture}
		\end{center}
		 Gọi $f(x)=a x^2+c$ là phương trình đường parabol của mái nhà ga.\\
		Ta có $\heva{&f(0)=10\\&f(10)=0}\Rightarrow\heva{&c=10\\&10^2a+c=0}\Rightarrow\heva{&a=-\dfrac{1}{10}\\&c=10.}$\\ 
		Do đó $f(x)=-\dfrac{1}{10}x^2+10$. 
		Khi đó diện tích thiết diện ngang  của sảnh gồm hai phần
		\begin{itemize}
			\item Phần hình  hộp chữ nhật có diên tích đáy là $S_1=20\times 8=160$\,m$^2$.
			\item Phần diện tích giới hạn bởi parabol và trục hoành có diện tích  là\\ $S_2=\int_{-10}^{10}f(x)\mathrm{d}x=2\int_{0}^{10}\left(-\dfrac{1}{10}x^2+10\right)\mathrm{d}x\approx 133{,}3$\,m$^2$.
		\end{itemize}
		 Thể tích không khí trong sảnh là $V=\left(S_1+S_2\right)\cdot 60\approx 17\,600$\,m$^3$.\\
		Vậy  thời gian trao đổi hết không khí trong sảnh là $\dfrac{17\,600}{2\cdot 80}=110$\, phút.
	 }
\end{bt}

% nguồn: Phan2_Lop_12_Dai_so_Chuong_4
\begin{bt}%[2D4V3-2]%[ma:e92ec6a3]
	Ở bể bơi  thành phố sẽ xây một cầu trượt nước mới. Cầu trượt dài $5$ m và có thể trượt từ cả hai phía. Những người thận trọng trượt về bên trái, nơi mép trượt được cho bởi $f(x)=\mathrm{e}^x$; những người ưa mạo hiểm trượt về bên phải, nơi mép trượt cao và dốc hơn, mô hình bởi $g(x)=-\dfrac{4}{9}x^2+x+1$.
	Biết rằng tại $x=0$ hai đường  cong nối trơn với nhau (không gãy, không góc nhọn), đơn vị $x$, $y$ tính bằng mét và hai thành bên của cầu trượt được ốp nhựa. Cần bao nhiêu m$^2$ nhựa? (Kết quả làm tròn đến hàng phần trăm).
	\begin{center}
		\begin{tikzpicture}[scale=1,font=\footnotesize,line join=round,line cap=round,>=stealth]
				\def\a{-4/9} \def\b{1} \def\c{1} 
				\def\xmin{-3} \def\xmax{4} 
				\def\ymin{-.5} \def\ymax{2.5} 
				\draw[->] (\xmin,0)--(\xmax,0) node [below]{$x$}; 
				 \draw[->] (0,\ymin)--(0,\ymax) node [right]{$y$};
				\fill  (0,0) circle (1.0pt) node[below left]{$O$};
				\clip (\xmin+0.1,\ymin+0.1) rectangle  (\xmax-0.1,\ymax-0.1);
				 \draw[smooth,samples=300,domain=-2:0] plot (\x,{e^(\x)});
				\draw[smooth,samples=300,domain=0:3] plot (\x,{\a*(\x)^2+\b*(\x)+\c}); 
				 \fill (0,1) circle (1.0pt) node[left]{$1$}
					 (-2,0) circle (1.0pt) node[below right]{$-2$}
					(3,0)  circle (1.0pt) node[below]{$3$};
				\draw[dashed] (-2,0)--(-2,0.15); 
			\end{tikzpicture}
	\end{center}
	\shortans{8,72} 
	 \loigiai{
		Diện tích ốp nhựa  (một bên).\\
		Trái $A_f=\displaystyle\int_{-2}^0{\mathrm{e}^x\mathrm{\,d}x}=\mathrm{e}^0-\mathrm{e}^{-2}\approx 0{,}86\,  \left(\text{m}^2\right)$.\\
		Phải $A_g=\displaystyle\int_{0}^3{\left( -\dfrac{4}{9}x^2+x+1 \right)\mathrm{\,d}x}=\left[-\dfrac{4}{27}x^3+\dfrac{1}{2}x^2+x\right]  \Bigg|_0^3=3{,}5\, \left(\text{m}^2\right)$.\\
		Tổng  diện tích một bên $A_1=0,86+3,5=4{,}36\, \left(\text{m}^2\right)$.\\
		Vậy tổng diện  tích hai bên $A_2=2\cdot 4{,}36=8{,}72\, \left(\text{m}^2\right)$.
	} 
\end{bt}

% nguồn: Phan2_Lop_12_Dai_so_Chuong_4
\begin{bt}%[2D4V3-2]%[ma:02c5aa33]
	Một công ty đang thiết kế một bảng quảng cáo hình chữ nhật $ABCD$ có kích thước $AB=12$ m và $AD=8$ m. Phần trung tâm của bảng sẽ được in nội dung quảng cáo, được mô  tả là phần tô đậm (xem hình minh họa). Hai đường cong trong hình là một phần của đồ thị hàm số có dạng $y=\dfrac{ax+b}{cx+d}$, đường tiệm cận đứng và tiệm cận ngang của đồ thị hàm số này đều cách điểm $A$ một khoảng bằng $5$ m. Đồ thị giao với cạnh $AB$ tại điểm $E$ thỏa mãn $\dfrac{AE}{AB}=\dfrac{7}{15}$. Diện tích phần in nội dung quảng cáo là bao nhiêu mét vuông (làm tròn kết quả đến hàng phần mười)?
	\begin{center}
		\begin{tikzpicture}[join = round,cap = round,>=stealth,scale = 0.7,declare function={a=3;},font=\footnotesize]
				\path  (-5,-5) coordinate (A)
					(7,-5) coordinate (B) 
					(0.6,-5)  coordinate (E)
					(7,3) coordinate (C) 
					(-5,3)  coordinate (D)
					(7,-a/7) coordinate (F); 
				\fill[pink] (-1,3)--(C)--plot[domain=7:0.6] (\x,{(0*(\x)+-3)/(1*(\x)+0)})--(E)--(A)--(-5,0.6)--plot[domain=-5:-1]  (\x,{(0*(\x)+-3)/(1*(\x)+0)})--cycle;
				\draw[teal,dashed]  (-5,0)--(7,0);
				\draw[teal,dashed] (0,-5)--(0,3); 
				\begin{scope}
					\clip (-5,-5) rectangle  (7,3);
					\draw[teal,samples=200,domain=-5:-0.01,smooth,variable=\x]  plot (\x,{(0*(\x)+-3)/(1*(\x)+0)});
					\draw[teal,samples=200,domain=0.01:7,smooth,variable=\x]  plot (\x,{(0*(\x)+-3)/(1*(\x)+0)});
					 \draw[teal,dashed,thin] (-5,0/1)--(7,0/1);
				\end{scope}
				 \draw[teal] (A)--(B)--(C)--(D)--cycle;
				\foreach \t/\g  in {A/-90,B/-90,C/90,D/90,E/-90,F/0}{
					\draw[teal] (\t) circle (1pt) node[shift={(\g:7pt)}]{$ \t $}; 
				} 
			\end{tikzpicture}
	\end{center}
	 \shortans{64,2}
	\loigiai{ 
		\begin{center}
			\begin{tikzpicture}[line join=round,line cap=round,>=stealth,scale=0.7,declare function={a=3;},font=\footnotesize]
					\tikzset{every node/.style={scale=0.9}} 
					\path  (-5,-5) coordinate (A)
						 (7,-5) coordinate (B)
						(0.6,-5) coordinate  (E)
						(7,3)  coordinate (C)
						(-5,3) coordinate (D) 
						(7,-a/7) coordinate (F); 
					\fill[pink] (-1,3)--(C)--plot[domain=7:0.6]  (\x,{(0*(\x)+-3)/(1*(\x)+0)})--(E)--(A)--(-5,0.6)--plot[domain=-5:-1](\x,{(0*(\x)+-3)/(1*(\x)+0)})--cycle;
					\draw[teal,dashed] (-5,0)--(7,0); 
					\draw[teal,dashed] (0,-5)--(0,3); 
					\draw[teal,->]  (-6,0)--(8,0) node[right] {$x$};
					\draw[teal,->] (0,-6)--(0,4.5)  node[above] {$y$};
					\draw[teal] (0,0) node [above left] {$O$}; 
					\begin{scope}
						\clip (-5,-5) rectangle  (7,3);
						\draw[teal,samples=200,domain=-5:-0.01,smooth,variable=\x] plot  (\x,{(0*(\x)+-3)/(1*(\x)+0)});
						\draw[teal,samples=200,domain=0.01:7,smooth,variable=\x] plot (\x,{(0*(\x)+-3)/(1*(\x)+0)}); 
					\end{scope}
					\draw[teal]  (A)--(B)--(C)--(D)--cycle;
					\draw[teal,dashed] (-1,0) node[below]{$-1$}--(-1,3) 
						(0.6,0) node[above]{$0,6$}--(0.6,-5); 
					\foreach \t/\g in  {A/-90,B/-90,C/90,D/90,E/-90,F/0}{
						 \draw[teal] (\t) circle (1pt) node[shift={(\g:7pt)}]{$ \t $};
					 }
					\path  (-5,0) node[above left]{$-5$}
						 (7,0) node[above right]{$7$}
						(0,3) node[above  left]{$3$}
						 (0,-5) node[below left]{$-5$};
				\end{tikzpicture}
		\end{center}
		Chọn hệ trục toạ độ  $Oxy$ trùng với hai tiệm cận như hình vẽ.\\
		 Suy ra $A(-5;-5)$; $B(7;-5)$; $C(7;3)$.\\
		Tiệm cận đứng  $x=-\dfrac{d}{c}=0\Rightarrow d=0$.\\
		Tiệm cận ngang $y=\dfrac{a}{c}=0\Rightarrow a=0\Rightarrow y=\dfrac{b}{cx}$  $(C)$.\\
		Ta có 
		 \[
		\dfrac{AE}{AB}=\dfrac{x_E-x_A}{12}=\dfrac{x_E+5}{12}=\dfrac{7}{15}\Rightarrow  x_E=\dfrac{3}{5}=0{,}6.\]
		Điểm $E(0{,}6;-5)$ thuộc $(C)$ nên  ta được $-5=\dfrac{b}{c\cdot 0{,}6}\Rightarrow b=-3c$. Khi đó $y=-\dfrac{3}{x}$.\\
		Gọi $K$ là giao điểm của đồ thị $(C)$ với đường thẳng  $y=3$ suy ra $K(-1;3)$.\\
		Từ  hình vẽ ta có
		 \[
		S_1=\displaystyle\int_{-5}^{-1}\left(3+\dfrac{3}{x}\right) \mathrm{\,d}x=\left(3x+3\ln|x|\right)\Big|_{-5}^{-1}=12-3\ln 5. 
		 \]
		\[ 
		 S_2=\displaystyle\int_{0{,}6}^{7}\left(-\dfrac{3}{x}+5\right)\mathrm{\,d}x=\left(-3\ln|x|+5x\right)\Big|_{0{,}6}^{7}=32+3\ln\dfrac{3}{35}.
		\] 
		Diện tích  phần in nội dung quảng cáo (\textbf{phần tô đậm}) là
		\[ 
		 S=S_{\text{ ABCD}}-S_1-S_2=12\cdot 8-12+3\ln 5-32+3\ln\dfrac{3}{35}=52+3\ln\dfrac{175}{3}\approx 64{,}2\ \mathrm{m}^2.
		 \]
	} 
\end{bt}

% nguồn: Phan2_Lop_12_Dai_so_Chuong_4
\begin{bt}%[2D4V3-2]%[ma:ce24c526]
	 \immini
	{Hình vẽ bên cạnh minh họa mặt cắt đứng của mặt trước một ngôi nhà có dạng hình chữ  nhật $5 \mathrm{m} \times 5{,}5 \mathrm{m}$ với các kích thước được cho như trong hình bên dưới. Phần tô đậm trong hình vẽ ứng với phần kính được lắp đặt, phần còn lại lót gạch men. Phần lắp kính được chia làm 2 miền $\left(H_1\right)$ và $\left(H_2\right)$, trong đó miền $\left(H_1\right)$ có đường biên cong là một phần của parabol. Chi phí lắp kính là $1{,}8$ triệu đồng cho $1 \, \mathrm{m}^2$ và chi phí lót gạch men là $0{,}25$ triệu đồng cho $1 \, \mathrm{m}^2$. Tổng chi phí cần hoàn thành mặt trước ngôi nhà là bao nhiêu triệu đồng? (\textit{Làm tròn kết quả đến hàng phần mười}).
	 }
	{\begin{tikzpicture}[>=stealth,line join=round,line cap=round,color=blue,font=\footnotesize,scale=0.7]
		\path (0,0) coordinate  (O)
			(5,0) coordinate (A) 
			 (0,5.5) coordinate (C)
			(5,5.5)  coordinate (B)
			(5,2.5)  coordinate (D)
			(0,2.5)  coordinate (E)
			(5,3.5)  coordinate (F)
			(0,3.5)  coordinate (G);
		\fill[pattern =  north east lines,pattern color=cyan,smooth] (O)--(A)-- (D)--(E)--cycle;
		\fill[pattern = north east lines,pattern color=red,smooth] (0,3.5) --plot[domain=0:5.0]  (\x,{-0.32*(\x)^2+1.6*(\x)+3.5})-- (5.0,3.5) -- cycle;
		\draw (O)--(A)--(B)--(C)--cycle  (D)--(E)
			(F)--(G); 
		\draw[violet,samples=150,smooth,domain=0:5] plot  (\x,{-0.32*(\x)^2+1.6*(\x)+3.5});
		\draw (2.5,3.6)  node[above]{$(H_1)$};
		\draw (2.5,1.2) node[above]{$(H_2)$}; 
		\draw[<->] (0,-1)--(5,-1)  node[above,midway]{$5$ m};
		\draw[<->] (-1,0)--(-1,5.5) node[above,midway,left]{$5{,}5$  m};
		\draw[<->] (5.5,0)--(5.5,2.5) node[above,midway,right]{$2{,}5$ m}; 
		\draw[<->]  (5.5,3.5)--(5.5,5.5) node[above,midway,right]{$2$ m};
	\end{tikzpicture}
	} 
	\shortans{36,6} 
	 \loigiai{
		Diện tích miền $\left(H_1\right)$ bằng $S_1 = \dfrac{2}{3} \cdot 5 \cdot  2 = \dfrac{20}{3} \left(\mathrm{m}^2\right)$.\\
		Diện tích miền $\left(H_2\right)$ bằng $S_2 = 5 \cdot 2{,}5 =  12{,}5 \left(\mathrm{m}^2\right)$.\\
		Chi phí lắp kính bằng $1{,}8 \cdot \left(\dfrac{20}{3} + 12{,}5\right) = \dfrac{69}{2}$ (triệu đồng).\\ 
		Chi phí lát gạch men bằng $0{,}25  \cdot\left(5 \cdot 5{,}5 - 12{,}5 - \dfrac{20}{3}\right) = \dfrac{25}{12}$ (triệu đồng).\\
		Tổng chi phí bằng $\dfrac{69}{2} + \dfrac{25}{12} = \dfrac{439}{12} \approx 36{,}6$ (triệu  đồng).
	} 
\end{bt}

% nguồn: Phan2_Lop_12_Dai_so_Chuong_4
\begin{bt}%[2D4V3-2]%[ma:43348582]
	Hưởng ứng chính sách hiện đại hoá nông thôn, người dân ở khu phố $A$ đồng lòng cùng nhau góp tiền để bê tông một đường đi trong khu phố (phần được tô  đậm trong hình vẽ).
	\begin{center}
		\begin{tikzpicture}[line join = round,line cap=round,>=stealth,font=\footnotesize,scale=1]
				\draw[->]  (-3,0) -- (8,0) node[below]{$x$};
				\draw[->] (0,-1) -- (0,10) node[left]{$y$}; 
				\draw (0,0) node[below left]{$O$}circle  (1pt);
				\draw[dashed]  (-2,0)--(-2,2)
					(7,0)--(7,9.59); 
				\clip  (-2.5,9.8) rectangle (8,-1);
				\draw[smooth,samples=100,domain=-2.1:7.2] plot (\x,{0.03125*(\x)^3  - 0.375*(\x)+1.5});
				\draw[smooth,blue,samples=100,domain=-2.1:7.5] plot (\x,{0.03125*(\x)^3  - 0.625*(\x)+0.5});
				 \fill (-2,1.5) node[below left]{$C$} circle (1pt);
				\fill (-2,2) node[above]{$A$}  circle (1pt);
				\fill (7,9.59) node[right]{$B$} circle (1pt); 
				\fill (7,6.84) node[right]{$D$} circle  (1pt);
				\fill (-2,0) node[below]{$-2$}  circle (1pt);
				\fill  (7,0) node[below]{$7$} circle (1pt);
			\end{tikzpicture}
	\end{center}
	Biết rằng khi chọn hệ trục tọa độ $Oxy$ với đơn vị độ dài trên mỗi trục tọa độ là $10$ m, các đường cong $AB$, $CD$ là mép đường được cho bởi đồ thị hàm số $f(x) = \dfrac{1}{32}x^3 -  \dfrac{3}{8}x + \dfrac{3}{2}$ và $g(x)=\dfrac{1}{32}x^3 - \dfrac{5}{8}x + \dfrac{1}{2}$, đồng thời lớp bê tông được đổ dày $16\,$cm và giá tiền $1$ m$^3$ bê tông là $1\,080\,000$ đồng. Tính số tiền (đơn vị triệu đồng) cần dùng để đổ bê tông con đường đó (làm tròn kết quả đến hàng đơn vị).
	\shortans{253} 
	\loigiai{ 
		Diện tích của hình phẳng giới hạn bởi  hai đồ thị hàm số đã cho là
		 $$
		S=\displaystyle\int_{-2}^{7} \left[f(x) - g(x)  \right] \mathrm{\,d}x
		= \displaystyle\int_{-2}^{7} \left[\left(\dfrac{1}{32}x^3 - \dfrac{3}{8}x  + \dfrac{3}{2}\right)
		- \left(\dfrac{1}{32}x^3  - \dfrac{5}{8}x + \dfrac{1}{2}\right)\right]\mathrm{\,d}x
		= \dfrac{117}{8} \text{ (đvdt).} 
		$$ 
		Do đơn vị độ dài trên mỗi trục tọa độ là $10\,$ m nên diện tích thực tế của  con đường là
		$$S_1 = \dfrac{117}{8} \cdot 10^2 = \dfrac{2925}{2}\, (\text{m}^2).$$ 
		Thể tích khối bê tông cần đổ là $V=S_1 \cdot h =  \dfrac{2925}{2} \cdot 0{,}16 = 234 \,(\text{m}^3)$.\\
		Số tiền cần  dùng để đổ bê tông con đường đó là $234 \cdot 1{,}08 \approx 253$ (triệu đồng).
	 }
\end{bt}

% nguồn: Phan2_Lop_12_Dai_so_Chuong_4
\begin{bt}%[2D4V3-2]%[ma:b917a841]
	\immini{Ông A có mảnh đất hình chữ nhật chiều dài $16$ m, chiều rộng $6$ m. Ông A muốn thiết kế mảnh đất đó thành ba phần bởi đường cong hypebol (như hình vẽ). Ông A dự định trồng hoa ở hai phần đất được giới hạn bởi đường cong hypebol và chiều dài mảnh đất, phần đất còn lại ông lát đá. Biết chi phí trồng hoa là $150\,000$ đồng $/ $1$ \mathrm{~m}^2$ và lát đá là $300\,000$ đồng $/ $1$ \mathrm{~m}^2$. Tính chi phí mà ông A bỏ ra để hoàn thành dự định của mình, (\textit{đơn vị triệu} đồng, \textit{kết quả làm tròn đến hàng phần chục}).}{\begin{tikzpicture}[scale=0.7,line join=round,line cap=round,>=stealth]
		\tikzset{every node/.style={scale=1}} 
		\def\xmin{-4.5}\def\xmax{4.5}\def\ymin{-11}\def\ymax{9} 
		\clip  (\xmin,\ymin) rectangle (\xmax,\ymax);
		\draw[smooth,samples=200,domain=1:3] plot (\x,{sqrt  (8*(\x)^(2.0)-8)});
		\draw[smooth,samples=200,domain=1:3] plot (\x,{-sqrt  (8*(\x)^(2.0)-8)});
		\foreach \x in {-1}\draw (\x,0.1)--(\x,-0.1); 
		 \draw[smooth,samples=50,domain=-3:-1] plot (\x,{sqrt (8*(\x)^(2.0)-8)});
		\draw[smooth,samples=50,domain=-3:-1] plot (\x,{-sqrt  (8*(\x)^(2.0)-8)});
		\coordinate (B) at  (-3,8);
		 \coordinate (A) at (-3,-8);
		\coordinate (C)  at (3,-8);
		\coordinate (D)  at (3,8);
		\coordinate (P) at (1,0); \foreach \diem in {A,B,C,D,P} \fill  (\diem) circle (1pt);
		\draw[smooth](A)--(B)--(D)--(C)--(A); 
		\draw[dashed]($(A)!.5 ! (B)$)--($(D)!.5 ! (C)$); 
		\draw ($(A)!.5 ! (B)$) node[left]{$16$  m};
		 \draw (-2,0) node[above]{$2$ m};
		\draw  (2,0) node[above]{$2$ m};
		\draw  (0,8) node[above]{$6$ m};
	\end{tikzpicture}}
	 \shortans{23,1}
	\loigiai{ 
	\immini{Chọn hệ  tọa độ $Oxy$ như hình vẽ, đơn vị trên hệ trục là mét.\\
		Khi đó, hypebol có phương  trình chính tắc là $\dfrac{x^2}{a^2}-\dfrac{y^2}{b^2}=1$, $(a>0$, $b>0)$.\\
		Ta thấy hypebol đi qua hai điểm $P(1; 0)$, $D(3; 8)$ nên $$\heva{&\dfrac{1^2}{a^2}-\dfrac{0^2}{b^2}=1\\&\dfrac{3^2}{a^2}-\dfrac{8^2}{b^2}=1} \Rightarrow a^2=1;  b^2=8.$$
		Vậy phương  trình hypebol là $\dfrac{x^2}{1}-\dfrac{y^2}{8}=1 \Rightarrow y= \pm \sqrt{8 x^2-8}$.\\
		Diện tích mảnh đất hình chữ nhật $ABCD$ là  $S=6\cdot 16=96 \mathrm{~m}^2$.\\
		Diện tích  trồng hoa (hai phần đất được giới hạn bởi đường cong hypebol và chiều dài mảnh đất) là $$S_1=4 \displaystyle\int_1^{3} \sqrt{8 x^2-8}~\mathrm{d}x \approx 38 ~\text{ m}^2.$$
		Diện tích phần lát  đá là $S_2=S-S_1 \approx 58 \mathrm{~m}^2$.\\
		Tổng chi phí mà ông A bỏ  ra để hoàn thành dự định của mình là
		$$0{,}15 \cdot S_1+0,3 \cdot S_2 \approx 23{,}1 \text{ triệu đồng}.$$}{\begin{tikzpicture}[scale=0.7,line join=round,line cap=round,>=stealth]
		 \tikzset{every node/.style={scale=1}}
		\def\xmin{-4.5}\def\xmax{4.5}\def\ymin{-11}\def\ymax{11} 
		\draw[->] (\xmin-0.2,0)--(\xmax+0.2,0) node[below]{$x$}; 
		 \draw[->] (0,\ymin-0.2)--(0,\ymax+0.2) node[right]{$y$};
		 \draw (0,0) node[below left]{$O$};
		\foreach \x in {-4,-2,2,4}\draw  (\x,0.1)--(\x,-0.1) node[below]{$\x$};
		\foreach \x in {-1}\draw  (\x,0.1)--(\x,-0.1);
		\foreach \y in {-10,-6,-4,-2,2,4,6,10}\draw (0.1,\y)--(-0.1,\y) node[left]{$\y$}; 
		\foreach \y in {-8}\draw (0.1,\y)--(-0.1,\y) node[below  left]{$\y$};
		\foreach \y in {8}\draw (0.1,\y)--(-0.1,\y) node[above left]{$\y$}; 
		\clip (\xmin,\ymin)  rectangle (\xmax,\ymax);
		 \draw[smooth,samples=200,domain=1:\xmax] plot (\x,{sqrt (8*(\x)^(2.0)-8)});
		\draw[smooth,samples=200,domain=1:\xmax] plot  (\x,{-sqrt (8*(\x)^(2.0)-8)});
		\draw[smooth,samples=50,domain=\xmin:-1] plot  (\x,{sqrt (8*(\x)^(2.0)-8)});
		\draw[smooth,samples=50,domain=\xmin:-1]  plot (\x,{-sqrt (8*(\x)^(2.0)-8)});
		\coordinate [label=left:$B$](B)  at (-3,8);
		\coordinate [label=left:$A$](A) at  (-3,-8);
		\coordinate [label=right:$C$](C)  at (3,-8);
		\coordinate [label=right:$D$](D) at  (3,8);
		\coordinate [label=above left:$P$](P) at (1,0); 
		\foreach  \diem in {A,B,C,D,P} \fill (\diem) circle (1pt);
		\draw[smooth](A)--(B)--(D)--(C)--(A); 
	\end{tikzpicture}}
	 }
\end{bt}

% nguồn: Phan2_Lop_12_Dai_so_Chuong_4
\begin{bt}%[2D4V3-2]%[ma:20f74572]
	Một viên gạch hoa hình vuông cạnh $4$ dm. Người thiết kế đã sử dụng bốn đường parabol có chung đỉnh tại tâm  viên gạch để tạo ra bốn cánh hoa (được tô màu như hình vẽ dưới). Giả sử để hoàn thiện sơn và phủ bóng viên gạch thì chi phí phần cánh hoa (tô màu) là $400$ nghìn/m$^2$, chi phí
	phần còn lại (màu trắng) là $300$ nghìn/m$^2$. Tính chi  phí (đơn vị nghìn đồng) để sơn và phủ bóng cả viên gạch (kết quả làm tròn đến hàng đơn vị).
	\begin{center}
		\begin{tikzpicture}[line cap=butt,line join=miter,>=stealth]
				\tikzset{declare function={xmin=-3;  xmax=3; ymin=-3; ymax=3;
					f (\x)=(\x)^2/2; 
					 g (\x)=-(\x)^2/2;
					}, 
					 smooth,samples=450
				 }
				\fill[cyan!60,opacity=0.6] 
					 plot[domain=0:2] (\x,{f (\x)}) --
					 plot[domain=2:0,variable=\y] ({\y^2/2},\y) -- cycle;
				\fill[cyan!60,opacity=0.6] 
					plot[domain=-2:0] (\x,{f  (\x)}) --
					plot[domain=0:2,variable=\y] ({-\y^2/2},\y)  -- cycle;;
				\fill[cyan!60,opacity=0.6] 
					plot[domain=0:2] (\x,{g (\x)}) -- 
					plot[domain=-2:0,variable=\y] ({-\y^2/2},\y) -- cycle; 
				\fill[cyan!60,opacity=0.6] 
					 plot[domain=-2:0] (\x,{g (\x)}) --
					plot[domain=-2:0,variable=\y] ({\y^2/2},\y) --  cycle;
				\clip (xmin,ymin)  rectangle (xmax,ymax);
				\draw  plot[domain=-2:2] (\x,{f (\x)});
				\draw[violet] plot[domain=-2:2] (\x,{g  (\x)});
				 \draw[blue,domain=0:2] plot (\x,{sqrt (2*\x)});
				\draw[blue,domain=0:2] plot  (\x,{-sqrt (2*\x)});
				\draw[red,domain=-2:0] plot  (\x,{sqrt (-2*\x)});
				\draw[red,domain=-2:0] plot (\x,{-sqrt (-2*\x)}); 
				\draw (-2,2) rectangle  (2,-2);
			\end{tikzpicture}
	\end{center}
	\shortans{53} 
	 \loigiai{
		\begin{center}
			\begin{tikzpicture}[line cap=butt,line join=miter,>=stealth]
					\tikzset{declare function={xmin=-3; xmax=3;  ymin=-3; ymax=3;
						f  (\x)=(\x)^2/2;
						g  (\x)=-(\x)^2/2;
						}, 
						smooth,samples=450 
					} 
					\draw[->] (xmin,0)--(xmax,0) node[shift={(-100:7pt)},font=\normalsize]{$ x  $};
					 \draw[->] (0,ymin)--(0,ymax) node[shift={(190:7pt)},font=\normalsize]{$ y $};
					\fill  (0,0) node[shift={(225:7pt)},font=\normalsize]{$ O $};
					 \fill[cyan!60,opacity=0.6]
						plot[domain=0:2]  (\x,{f (\x)}) --
						 plot[domain=2:0,variable=\y] ({\y^2/2},\y) -- cycle;
					\fill[cyan!60,opacity=0.6] 
						plot[domain=-2:0] (\x,{f (\x)}) -- 
						plot[domain=0:2,variable=\y]  ({-\y^2/2},\y) -- cycle;;
					 \fill[cyan!60,opacity=0.6]
						plot[domain=0:2] (\x,{g (\x)}) -- 
						plot[domain=-2:0,variable=\y] ({-\y^2/2},\y) -- cycle; 
					\fill[cyan!60,opacity=0.6] 
						 plot[domain=-2:0] (\x,{g (\x)}) --
						plot[domain=-2:0,variable=\y] ({\y^2/2},\y) --  cycle;
					\clip (xmin,ymin) rectangle (xmax,ymax); 
					\draw plot[domain=-2:2]  (\x,{f (\x)});
					 \draw[violet] plot[domain=-2:2] (\x,{g (\x)});
					\draw[blue,domain=0:2] plot  (\x,{sqrt (2*\x)});
					\draw[blue,domain=0:2] plot  (\x,{-sqrt (2*\x)});
					\draw[red,domain=-2:0] plot (\x,{sqrt (-2*\x)}); 
					\draw[red,domain=-2:0] plot (\x,{-sqrt (-2*\x)}); 
					\draw (-2,2) rectangle (2,-2); 
				\end{tikzpicture}
		\end{center}
		Chọn hệ tọa  độ như hình vẽ ($1$ đơn vị trên trục bằng $1$ dm), các cánh hoa tạo bởi các đường parabol có phương trình $y=\dfrac{x^2}{2}$, $y=-\dfrac{x^2}{2}$, $x=-\dfrac{y^2}{2}$, $x=\dfrac{y^2}{2}$.\\
		Diện tích một cánh hoa (nằm trong góc phần tư thứ nhất) bằng diện tích hình phẳng giới hạn  bởi hai đồ thị hàm số $y=\dfrac{x^2}{2}$, $y=\sqrt{2x}$ và hai đường thẳng $x=0$, $x=2$.\\
		Do đó diện tích một cánh hoa bằng $\displaystyle\int_0^2\left(\sqrt{2x}-\dfrac{x^2}{2}\right) \mathrm{d}x=\left(\dfrac{2\sqrt{2}}{3} \sqrt{x^3}-\dfrac{x^3}{6}\right)\Bigg|_0  ^2=\dfrac{4}{3}$ $\left(\mathrm{dm}^2\right)$.\\
		Suy ra diện tích phần cánh hoa được  tô đen là $\dfrac{4}{3} \cdot 4=\dfrac{16}{3}\left(\mathrm{dm}^2\right)=\dfrac{4}{75}\left(\mathrm{m}^2\right)$.\\
		Diện  tích phần màu trắng còn lại là $4\cdot 4-\dfrac{16}{3}=\dfrac{32}{3}\left(\mathrm{dm}^2\right)=\dfrac{8}{75}\left(\mathrm{m}^2\right)$.\\
		Số tiền cần dùng là $\dfrac{4}{75} \cdot  400+\dfrac{8}{75} \cdot 300=\dfrac{160}{3} \approx 53$ nghìn đồng.
	} 
\end{bt}

% nguồn: Phan2_Lop_12_Dai_so_Chuong_4
\begin{bt}%[2D4V3-2]%[ma:36a7fa54]
	Cô Hạnh đổ bê tông một đường đi trong vườn (phần được tô màu) với kích thước được cho trong hình sau. Biết rằng đường cong $AB$ được cho bởi đồ thị của một hàm số liên tục và đường cong $DC$ nhận được từ đường cong $AB$ bằng cách tịnh tiến theo phương thẳng đứng lên phía trên $2$ m. Ngoài ra, cô Hạnh quyết định đổ lớp bê tông dày $15$ cm và giá tiền $1$ m$^3$ bê tông là $1\,080\,000$ đồng. Tính số tiền cô Hạnh cần dùng để đổ bê tông con đường đó  (đơn vị nghìn đồng).
	\begin{center}
		\begin{tikzpicture}[line join = round,line cap = round,>=stealth,font=\footnotesize,scale=0.5]
				\fill[cyan!50,smooth] plot[domain=0:10] (\x,{0.008*(\x)^3-0.12*(\x)^2+4})-- (10,2) -- plot[domain=10:0] (\x,{0.008*(\x)^3-0.12*(\x)^2+6}) -- (0,4); 
				 \draw (0,0) rectangle (10,6)
					plot[domain=0:10] (\x,{0.008*(\x)^3-0.12*(\x)^2+4})-- (10,2) -- plot[domain=10:0] (\x,{0.008*(\x)^3-0.12*(\x)^2+6}) --  (0,4);
				\path 
					(0,4)  coordinate (A)
					 (0,6) coordinate (D)
					 (10,2) coordinate (C)
					(10,0) coordinate (B); 
				\foreach \x/\g in {A/180,B/0,C/0,D/180} \fill(\x) circle (1pt)+(\g:3mm) node{$\x$}; 
				\draw[<->] (0,-0.5) --  (10,-0.5);
				\draw (5,-0.5) node[below]  {$10$ m};
				\draw[<->] (-1,0)  -- (-1,6);
				\draw (-1,3) node[left]  {$6$ m};
				\draw[<->] (11,0) -- (11,2); 
				\draw (11,1)  node[right] {$2$ m};
			\end{tikzpicture}
	\end{center}
	 \shortans{3240}
	\loigiai{ 
	\immini{ 
		Xét hệ  trục toạ độ $Oxy$ như hình vẽ bên (đơn vị trên mỗi trục toạ độ là mét).\\
		Trong hệ trục toạ độ đó, đường cong $AB$ được cho bởi đồ thị của hàm số liên tục $y=f(x)$, $0\leq x\leq 10$. Suy ra, đường cong $DC$ được cho bởi đồ thị của hàm  số liên tục $y=f(x)+2$, $0\leq x\leq 10$.\\
		Áp dụng công  thức tính diện tích ta có
		 }{
		\begin{tikzpicture}[line join = round,line cap = round,>=stealth,font=\footnotesize,scale=0.5]
				 \fill[cyan!50,smooth] plot[domain=0:10] (\x,{0.008*(\x)^3-0.12*(\x)^2+4})-- (10,2) -- plot[domain=10:0] (\x,{0.008*(\x)^3-0.12*(\x)^2+6}) -- (0,4);
				\draw (0,0) rectangle (10,6); 
				 \draw[->] (0,0) -- (12,0) node[below] {$x$};
				\draw[->]  (0,0) -- (0,7) node[left] {$y$};
				\draw  (0,0) node[left] {$O$}
					 plot[domain=0:10] (\x,{0.008*(\x)^3-0.12*(\x)^2+4})-- (10,2) -- plot[domain=10:0] (\x,{0.008*(\x)^3-0.12*(\x)^2+6}) -- (0,4);
				 \path
					(0,4)  coordinate (A)
					(0,6) coordinate  (D)
					(10,2) coordinate  (C)
					(10,0) coordinate (B); 
				\draw  (A)--(B) node[midway,sloped,below]{$y=f (x)$};
				\draw (D)--(C)  node[midway,sloped,above]{$y=f (x)+2$};
				\foreach \x/\g in {A/180,B/45,C/0,D/180} \fill(\x) circle (1pt)+(\g:3mm)  node{$\x$};
				\draw (A) node[below right]  {$4$};
				 \draw (D) node[above right] {$6$};
				\draw (B) node[below]  {$10$};
			\end{tikzpicture}
		} 
		\noindent 
		\begin{align*}
			S_{\text{ hình phẳng được  tô màu }ABCD }=\displaystyle\int_{0}^{10} \left[f(x)+2-f(x)\right] \mathrm{\,d}x=\displaystyle\int_{0}^{10} 2 \mathrm{\,d}x=2x\bigg|_0^{10}=20\text{ m}^2.
		\end{align*}
		Do đó, thể tích khối bê tông dùng để đổ con đường  là $20\cdot 0{,}15=3$ m$^3$.\\
		Vậy số  tiền cô Hạnh cần dùng để đổ bê tông con đường đó là
		\begin{align*}
			3 \cdot 1\,080\,000=3\,240\,000=3\,240 \text{ (nghìn đồng)}. 
		\end{align*}
	 }
\end{bt}

% nguồn: Phan2_Lop_12_Dai_so_Chuong_4
\begin{bt}%[2D4V3-2]%[ma:2fcf89bf]
	\immini{Trong hình vẽ bên, đoạn $AD$ được chia làm $3$ bởi các điểm $B$ và $C$ sao cho $AB=BC=CD=2$. Ba nửa đường tròn có bán kính bằng $1$ là $\wideparen{AEB}$, $\wideparen{BFC}$ và $\wideparen{CGD}$ có đường kính tương ứng là $AB$, $BC$ và $CD$. Các điểm $E$, $F$, $G$ lần lượt là tiếp điểm của tiếp tuyến chung $EG$ với $3$ nửa đường tròn. Một đường tròn tâm $F$, bán kính bằng $2$. Diện tích miền bên trong đường tròn tâm $F$ và bên ngoài $3$ nửa đường tròn (miền gạch sọc) bằng bao nhiêu? (\textit{Kết quả làm tròn đến hàng phần mười}).}{\begin{tikzpicture}[scale=1,line join=round,line cap=round,>=stealth]
		\tikzset{every node/.style={scale=1}} 
		\draw [fill=blue!40!white](0,1) circle (2); \coordinate [label=above:$F$](F)  at (0,1);
		\coordinate [label=below:$B$](B) at (-1,0); 
		\coordinate [label=below:$A$](A) at (-3,0); 
		\coordinate  [label=below:$C$](C) at (1,0);
		\coordinate  [label=below:$D$](D) at (3,0);
		\coordinate (I) at (-2,0); 
		 \coordinate [label=above left:$E$](E) at (-2,1);
		\coordinate [label=above right:$G$](G) at (2,1); 
		 \pgfmathsetmacro\a{-1*sqrt (3)}
		\coordinate (K)  at (\a,0);
		\coordinate  (L) at (-\a,0);
		\fill[fill=white]plot[domain=-3:-1](\x,{sqrt  (1-((\x)+2)^(2))})--plot[domain=-1:-3](\x,0)--cycle;
		 \fill[fill=white]plot[domain=1:-1](\x,{sqrt (1-((\x)-2+2)^(2.0))})--plot[domain=-1:1](\x,0)--cycle;
		\fill[fill=white]plot[domain=1:3](\x,{sqrt (1-((\x)-4+2)^(2.0))})--plot[domain=3:1](\x,0)--cycle; 
		\draw  (A) arc (180:0:1);
		 \draw (B) arc (180:0:1);
		\draw (C) arc (180:0:1); 
		\foreach \diem in {A,B,C,D,E,F,G} \fill  (\diem) circle (1pt);
		 \draw[smooth](A)--(D);
	\end{tikzpicture}
	 }
	\shortans{9,6} 
	 \loigiai{
		\begin{center}
			\begin{tikzpicture}[scale=1.5,line join=round,line cap=round,>=stealth]
					\tikzset{every  node/.style={scale=1}}
					\coordinate [label=above  left:$F$](F) at (0,1);
					 \draw [pattern=north east lines](F) circle (2);
					\coordinate [label=below:$B$](B) at (-1,0); 
					 \coordinate [label=below:$A$](A) at (-3,0);
					\coordinate [label=below:$C$](C) at  (1,0);
					\coordinate [label=below:$D$](D) at (3,0); 
					 \coordinate (I) at (-2,0);
					\coordinate [label=above left:$E$](E) at  (-2,1);
					\coordinate [label=above  right:$G$](G) at (2,1);
					\coordinate  [label=below:$I$](I) at (-2,0);
					\pgfmathsetmacro\a{-1*sqrt  (3)}
					 \fill[fill=white]plot[domain=-3:-1](\x,{sqrt (1-((\x)+2)^(2))})--plot[domain=-1:-3](\x,0)--cycle;
					\fill[fill=white]plot[domain=1:-1](\x,{sqrt (1-((\x)-2+2)^(2.0))})--plot[domain=-1:1](\x,0)--cycle; 
					 \fill[fill=white]plot[domain=1:3](\x,{sqrt (1-((\x)-4+2)^(2.0))})--plot[domain=3:1](\x,0)--cycle;
					\draw (A)  arc (180:0:1);
					 \draw (B) arc (180:0:1);
					\draw (C) arc (180:0:1); 
					\coordinate  [label=above right:$K$](K) at (\a,0);
					 \coordinate (L) at (-\a,0);
					\foreach \diem in  {A,B,C,D,E,F,G,I,K,L} \fill (\diem) circle (1pt);
					 \draw[smooth](A)--(D);
					 \def\xmin{-4}\def\xmax{4}\def\ymin{-1}\def\ymax{3.5}
					 \draw[->] (\xmin-0.2,0)--(\xmax+0.2,0) node[below]{$x$};
					\draw[->] (0,\ymin-0.2)--(0,\ymax+0.2) node[right]{$y$}; 
					\draw (0,0) node[below  left]{$O$};
					 \draw[smooth](E)--(-2,0)
						 (G)--(2,0);
					\draw[dashed](E)--(G); 
					\draw (1,1)  node[above]{$y=1$};
					 \draw[dashed,smooth,samples=200,domain=-2:\a] plot (\x,{0-sqrt (4-(\x)^(2.0))+1});
					 \draw[dashed,smooth,samples=200,
						domain=2:-\a] plot  (\x,{0-sqrt (4-(\x)^(2.0))+1});
				\end{tikzpicture}
		\end{center}
		Diện tích hình tròn tâm $F$, bán  bán kính bằng $2$ là $S_F=\pi \cdot 2^2=4 \pi$.\\
		Diện tích nửa hình tròn bán kính bằng $1$ là $S_N=\dfrac{1}{2} \cdot \pi \cdot  1^2=\dfrac{\pi}{2}$.\\
		 Diện tích góc phần tư $\widehat{IBE}$ là $S_I=\dfrac{1}{4} \pi \cdot 1^2=\dfrac{\pi}{4}$.\\
		Dựng hệ trục toạ độ $Oxy$ như hình vẽ.\\ 
		Toạ độ các điểm $F(0; 1)$,  $I(-2; 0)$.
		Phương trình đường tròn  tâm $F$, bán kính bằng $2$ là $x^2+(y-1)^2=4$.\\Suy ra
		cung nhỏ $\wideparen{EK}$ nằm trên đồ  thị $y=1-\sqrt{4-x^2}$.\\
		Hoành độ giao điểm $K$ là nghiệm của phương trình  $1-\sqrt{4-x^2}=0$, $(x<0)$.\\
		Ta có $1-\sqrt{4-x^2}=0 \Leftrightarrow 4-x^2=1  \Rightarrow x=-\sqrt{3}$ (vì $x<0$).\\
		Diện  tích tam giác cong $I E K$ là $S_C=\displaystyle\int_{-2}^{-\sqrt{3}}\left(1-\sqrt{4-x^2}\right)~\mathrm{d}x$.\\
		Diện tích miền tô  đậm trong hình là
		$$S=S_F-S_N-2\left(S_I-S_C\right)=4  \pi-\dfrac{\pi}{2}-2\left(\dfrac{\pi}{4}-\displaystyle\int_{-2}^{-\sqrt{3}}\left(1-\sqrt{4-x^2}\right)~\mathrm{d}x\right) \approx 9{,}6.$$
	 }
\end{bt}

% nguồn: Phan2_Lop_12_Dai_so_Chuong_4
\begin{bt}%[2D4V3-2]%[ma:404c3169]
	Ông A làm một logo bằng tấm nhựa phẳng, có hình dạng là một hình có trục đối xứng. Biết đường viền $AE$, $CF$ ở hai bên là hai nhánh của một parabol và phần lõm $EIF$ phía dưới đáy cũng là một parabol, hai nhánh phía trên là hai đoạn thẳng như hình bên dưới. Cho biết $AC=6$ dm, $AB=3\sqrt{2}$ dm, $EF=2$ dm và khoảng  cách giữa $AC$ và $EF$ bằng $8$ dm . Hỏi diện tích của logo đó bằng bao nhiêu dm$^2$. Viết kết quả dưới dạng số thập phân. Viết kết quả làm tròn đến hàng phần mười.
	\begin{center}
		\begin{tikzpicture}[scale=0.8,font=\footnotesize,line join=round,line cap=round,>=stealth]
				\draw (0,0)  circle (1pt) node[below left]{2 dm};
				\path  (-3,8) coordinate (A)
					(0,5) coordinate  (B)
					(3,8) coordinate  (C)
					(-1,0) coordinate  (E)
					 (1,0) coordinate (F)
					 (0,1) coordinate (I);
				\clip (-4,0) rectangle (4.5,9); 
				 \draw
					 plot[domain=-3:3,smooth,samples=200] (\x,{(\x)^2-1});
				\draw[fill=cyan,opacity=0.4] (A)--plot[domain=-3:3]  (\x,{(\x)^2-1})--(C)--(B)--cycle;
				\draw plot[domain=-1:1,smooth,samples=200] (\x,{1-(\x)^2}); 
				 \draw[fill=white] (E)--plot[domain=-1:1,smooth,samples=200] (\x,{1-(\x)^2})--(F)--cycle;
				 \draw (A)--(B) node[midway,above right]{$3\sqrt{2}$ dm}--(C);
				\foreach \d/\g in {A/120,B/90,C/60,F/30,E/150,I/120} 
				\draw  (\d) circle (1pt) +(\g:3mm) node{$\d$};
				\draw[<->]  (4,0)--(4,8) node[midway,left]{8 dm};
				 \draw[<->] (A)--(C) node[midway,below]{6 dm};
				\draw[<->] (B)--(I) node[midway,right]{4  dm};
			\end{tikzpicture}
	\end{center}
	\shortans{24,3} 
	\loigiai{ 
		Xét hệ tọa độ $Oxy$  như hình vẽ.
		\begin{center}
			\begin{tikzpicture}[scale=0.8,font=\footnotesize,line join=round,line cap=round,>=stealth]
					 \draw[->] (-4,0)--(5,0) node[above]{$x$};
					 \draw[->] (0,-0.5)--(0,9) node[left]{$y$};
					\draw (0,0) circle (1pt) node[below  left]{$O$};
					\path (-3,8)  coordinate (A)
						(0,5) coordinate  (B)
						(3,8) coordinate  (C)
						 (-1,0) coordinate (E)
						 (1,0) coordinate (F);
					\clip (-4,0) rectangle  (4.5,9);
					 \draw
						 plot[domain=-3:3,smooth,samples=200] (\x,{(\x)^2-1});
					\draw[fill=cyan,opacity=0.4] (A)--plot[domain=-3:3]  (\x,{(\x)^2-1})--(C)--(B)--cycle;
					\draw  plot[domain=-1:1,smooth,samples=200] (\x,{1-(\x)^2});
					 \draw[fill=white] (E)--plot[domain=-1:1,smooth,samples=200] (\x,{1-(\x)^2})--(F)--cycle;
					\draw (A)--(B)  node[midway,above,rotate=-45]{$3\sqrt{2}$ dm}--(C);
					 \foreach \d/\g in {A/120,B/90,C/60,F/30,E/150,I/120}
					\draw (\d) circle (1pt) +(\g:3mm) node{$\d$}; 
					\draw[<->] (4,0)--(4,8) node[midway,left]{8  dm};
					 \draw[<->] (A)--(C) node[midway,below]{6 dm};
					\draw[<->] (B)--(I) node[midway,right]{4 dm}; 
				\end{tikzpicture}
		\end{center}
		Khi đó nửa bên phải trục tung là hình phẳng $(H)$ giới  hạn bởi các đường $y=x^2-1$, $y=x+5$, $y=-x^2+1$, $x=0$.\\
		Ta có $S=2S_H=2\left(\displaystyle\int_0^1\left[x+5-\left(1-x^2\right)\right]\mathrm{\,d}x+\displaystyle\int_1^3\left[x+5-\left(x^2-1\right)\right]\mathrm{\,d}x\right) \approx 24{,}3$  dm$^2$.
	 }
\end{bt}

% nguồn: Phan2_Lop_12_Dai_so_Chuong_4
\begin{bt}%[2D4V3-2]%[ma:3bc8cf64]
	\immini{Một nghệ sĩ đang thực hiện vẽ tranh tường $3$D nghệ thuật để trang trí nội thất. Mỗi bức tranh gồm một hình lục giác đều có cạnh bằng $30$ cm , trên mỗi cạnh của hình lục giác đều có một cánh hoa hình parabol, đỉnh của parabol cách cạnh lục giác đều $30$ cm và nằm phía ngoài hình lục giác, đường parabol đó đi qua hai đầu mút của mỗi cạnh (xem hình bên). Biết tổng giá vẽ tranh $3$D (bao gồm công và vật tư) là $500\,000$ đồng $/ \mathrm{m}^2$. Hỏi nghệ sĩ nhận được bao nhiêu tiền sau khi hoàn thành bức tranh? (\textit{Làm tròn đến nghìn} đồng \textit{gần nhất}).}{\begin{tikzpicture}[scale=2]
		\pgfmathsetmacro{\R}{1} 
		\pgfmathsetmacro{\Angle}{30} 
		\pgfmathsetmacro{\PetalLength}{0.8} 
		\pgfmathsetmacro{\PetalWidth}{0.3} 
		\coordinate  (O) at (0,0);
		 \foreach \i in {1,...,6}{
			\pgfmathsetmacro{\Ang}{\Angle  + (\i-1)*60}
			\coordinate  (V\i) at (\Ang:\R);
		} 
		\foreach \i in {1,...,6}{ 
			 \pgfmathsetmacro{\j}{ifthenelse (\i==6,1,\i+1)}
			\fill[orange!85!yellow] (O) -- (V\i) -- (V\j) -- cycle; 
		} 
		\foreach  \i in {1,...,6}{
			\pgfmathsetmacro{\j}{ifthenelse  (\i==6,1,\i+1)}
			\pgfmathsetmacro{\MidAngle}{\Angle  + (\i-1)*60 + 30}
			\coordinate (M) at  (\MidAngle:\R);
			\coordinate  (E) at (\MidAngle:\R+\PetalLength);
			\pgfmathsetmacro{\Color}{ifthenelse (mod  (\i,2)==1,"red!80!black","blue!80!black")}
			 \pgfmathsetmacro{\CtrlAngleA}{\MidAngle-90}
			 \pgfmathsetmacro{\CtrlAngleB}{\MidAngle+90}
			 \coordinate (C1) at (\MidAngle-30:\R+\PetalWidth/2);
			\coordinate (C2)  at (\MidAngle+30:\R+\PetalWidth/2);
			 \fill[\Color] (V\i)
			..  controls (\MidAngle-10:\R+\PetalWidth*1.5) and (\MidAngle-10:\R+\PetalLength).. (E)
			.. controls (\MidAngle+10:\R+\PetalLength) and (\MidAngle+10:\R+\PetalWidth*1.5).. (V\j) 
			 -- cycle;
		 }
		 \foreach \i in {1,...,6}{
			 \pgfmathsetmacro{\j}{ifthenelse (\i==6,1,\i+1)}
			\pgfmathsetmacro{\MidAngle}{\Angle  + (\i-1)*60 + 30}
			\coordinate (E) at  (\MidAngle:\R+\PetalLength);
			\draw[line  width=2pt] (V\i)
			.. controls (\MidAngle-10:\R+\PetalWidth*1.5) and (\MidAngle-10:\R+\PetalLength).. (E) 
			 .. controls (\MidAngle+10:\R+\PetalLength) and (\MidAngle+10:\R+\PetalWidth*1.5).. (V\j);
		} 
		\draw[line width=2pt] (V1) -- (V2) -- (V3) -- (V4) -- (V5)  -- (V6) -- cycle;
		\foreach  \i in {1,...,6}{
			 \draw[line width=0.2pt] (O) -- (V\i);
		 }
	\end{tikzpicture}
	 }
	 \shortans{297}
	\loigiai{Diện tích lục giác là diện tích  của $6$ tam giác đều cạnh $0{,}3 \mathrm{~m}$ là $S_1=6 \cdot \dfrac{(0{,}3)^2 \sqrt{3}}{4}=\dfrac{27 \sqrt{3}}{200}~\left(\mathrm{m}^2\right)$.\\
		Diện  tích $6$ cánh hoa là diện tích $6$ parabol $d=0{,}3$ m; $h=0{,}3$ m là
		$$S_2=6 \cdot \dfrac{2}{3} h  \cdot d=6 \cdot \dfrac{2}{3} \cdot 0{,}3 \cdot 0{,}3=\dfrac{9}{25}\left(\text{m}^2\right).$$
		Tổng diện tích bức  tranh $S=S_1+S_2$.\\
		Tổng giá tiền để vẽ bức tranh là $$\left(\dfrac{9}{25}+\dfrac{27 \sqrt{3}}{200}\right) \cdot 500 \approx 297 ~(\text{nghìn đồng}).$$ 
	} 
\end{bt}

% nguồn: Phan2_Lop_12_Dai_so_Chuong_4
\begin{bt}%[2D4V3-2]%[ma:799b8bc8]
	Một bức tường hình chữ nhật $ABCD$ có  kích thước lần lượt là $6$ m và $4$ m được bạn Minh trang trí bằng cách vẽ hai đồ thị hàm số $f(x)=a^x$ $(0<a\ne1)$ và $g(x)=\log_b x$ $(0<b\ne1)$ đối xứng nhau qua đường thẳng $y=x$. Bức tường được chia làm 3 phần (tham khảo hình vẽ).
	\begin{center}
		\begin{tikzpicture}[>=stealth,line join=round,line cap=round,scale=1.2]
				\fill[blue!60] 
					plot[domain=-3:1](\x,{2^(\x)})--(1,2) 
					 --(-3,2)--(-3,0.125)--cycle;
				 \fill[green!60]
					plot[domain=0.25:3](\x,{ln  (\x)/ln (2)})--(3,1.5849)--(3,-2)--(0.25,-2)--cycle;
				 \fill[yellow!40]
					plot[domain=0.25:3](\x,{ln  (\x)/ln (2)})--(3,1.5849)--(3,2)--(1,2)--plot[domain=1:-3](\x,{2^(\x)})--(-3,0.125)--(-3,-2)--(0.25,-2)--cycle;
				\draw[->](-4,0)--(4,0); 
				\draw[->](0,-3)--(0,3); 
				 \draw (-3,-2) rectangle (3,2);
				\foreach \x in  {-3,-2,-1,1,2,3}\draw (\x,0) circle (1pt) node[below]{$\x$};
				 \foreach \y in {-2,-1,1,2}\draw (0,\y) circle (1pt) node[left]{$\y$};
				 \draw (4,0) node[above]{$x$}
					 (0,3)
					 node[left]{$y$}
					(0,0) node[below 
					left]{$O$}; 
				 \draw (-2,1) node[above] {$H_1$};
				\draw  (-2,-1) node[below] {$H_2$};
				\draw (2,-1)  node {$H_3$};
			\end{tikzpicture}
	\end{center}
	Phần $H_1$ được sơn màu xanh da trời, phần $H_2$  sơn màu vàng và phần $H_3$ được sơn màu xanh lá cây.
	Bạn Minh cần  mua các hộp sơn mà mỗi hộp chỉ sơn được một màu tương ứng với các màu mà bạn Minh định sơn.
	Biết rằng mỗi hộp sơn chỉ sơn được tối đa $3$ m$^2$ tường, giá một hộp sơn màu xanh da trời là $120\,000$ đồng, giá một hộp sơn màu xanh lá cây là  $140\,000$ đồng, giá một hộp sơn màu vàng là $160\,000$ đồng.
	Cửa hàng chỉ bán số các hộp sơn là các  số nguyên dương.
	Bạn Minh cần bao nhiêu triệu đồng để sơn bức tường  đó?
	 \shortans{1,3}
	 \loigiai{
		\begin{center}
			\begin{tikzpicture}[>=stealth,line join=round,line cap=round,scale=1.2]
					 \fill[blue!60]
						plot[domain=-3:1](\x,{2^(\x)})--(1,2) 
						 --(-3,2)--(-3,0.125)--cycle;
					\fill[green!60] 
						 plot[domain=0.25:3](\x,{ln (\x)/ln (2)})--(3,1.5849)--(3,-2)--(0.25,-2)--cycle;
					 \fill[yellow!40]
						plot[domain=0.25:3](\x,{ln  (\x)/ln (2)})--(3,1.5849)--(3,2)--(1,2)--plot[domain=1:-3](\x,{2^(\x)})--(-3,0.125)--(-3,-2)--(0.25,-2)--cycle;
					\draw[smooth]plot[domain=-3:1.2](\x,{2^(\x)}) node[right]{$y=a^x$}; 
					 \draw[smooth]plot[domain=0.2:3.3](\x,{ln (\x)/ln (2)}) node[right]{$y=\log_ax$};
					\draw[->](-4,0)--(4,0); 
					 \draw[->](0,-3)--(0,3);
					\draw (-3,-2)  rectangle (3,2);
					\foreach \x in {-3,-2,-1,1,2,3}\draw (\x,0) circle  (1pt) node[below]{$\x$};
					\foreach \y in {-2,-1,1,2}\draw (0,\y) circle (1pt) node[left]{$\y$}; 
					 \draw (4,0) node[above]{$x$}
						(0,3) 
						node[left]{$y$} 
						(0,0)  node[below
						 left]{$O$};
					 \draw (-2,1) node[above] {$H_1$};
					\draw (-2,-1) node[below]  {$H_2$};
					 \draw (2,-1) node {$H_3$};
				\end{tikzpicture}
		\end{center}
		Đồ thị hàm số $f(x)=a^x$ $(0<a\ne1)$ và đồ  thị hàm số $g(x)=\log_b x$ $(0<b\ne1)$ đối xứng nhau qua đường thẳng $y=x$ nên $a=b$.\\
		Ta thấy đồ thị hàm số $g(x)=\log_b x$  đi qua điểm $B(3;2)$ nên suy ra $b=\sqrt{3}$.\\
		Vậy  bạn Minh đã vẽ hai đồ thị của hàm số $f(x)=\sqrt{3}^x$ và $g(x)=\log_{\sqrt{3}} x$.\\
		Ta  có
		$$ 
		\sqrt{3}^x=2 \Leftrightarrow x=\log_{\sqrt{3}}2 \Rightarrow S_{H_1} = \displaystyle\int_{-\log_{\sqrt{3}}2}^{2} \left(2 -  \sqrt{3}^x\right)\mathrm{d}x \approx 5{,}23 .$$
		Do đó cần mua 2 hộp sơn để sơn phần $H_1$ với số tiền là $2.0{,}12=0{,}24$  (triệu đồng).\\
		$$ 
		\log_{\sqrt{3}}x = -2 \Leftrightarrow x = \dfrac{1}{3} \Rightarrow S_{H_3} = \displaystyle\int_{\frac{1}{3}}^{3} \left(\log_{\sqrt{3}}x +  2\right)\mathrm{d}x \approx 7{,}15.$$
		Do đó cần mua 3 hộp sơn để sơn phần $H_3$ với số tiền là $3\cdot0{,}14=0{,}42$  (triệu đồng).\\
		$$ 
		S_{H_2} = S_{HCN} - S_{H_1} - S_{H_3} \approx 6\cdot 4 - 5{,}23 - 7{,}15 \Rightarrow  S_{H_2} \approx 11{,}62 .$$
		Do đó cần mua 4 hộp sơn để  sơn phần $H_2$ với số tiền là $4\cdot 0{,}16=0{,}64$ (triệu đồng).\\
		Vậy số tiền bạn Minh cần dùng để sơn bức tường đó là $0{,}24 + 0{,}42 + 0{,}64 = 1{,}3$  (triệu đồng).
	 }
\end{bt}

% nguồn: Phan2_Lop_12_Dai_so_Chuong_4
\begin{bt}%[2D4V3-2]%[ma:15bc4b3c]
	 \immini{
	Một cánh cổng của một tòa nhà có dạng parabol gồm hai phần: phần hai cánh cửa hình chữ nhật $CDEF$, còn lại là phần xiên hoa trang trí. Biết rằng $GH=4$ m, $AB=4$ m và $AC=BD=0{,}9$ m.  Diện tích phần cổng làm xiên hoa trang trí (kết quả làm tròn đến hàng phần~trăm) bằng bao nhiêu mét vuông?
	}{\begin{tikzpicture}[smooth,samples=100,scale=0.8,>=stealth,font=\small]
		 \path
			(0,0)  coordinate (H)
			(0,4) coordinate (G) 
			(-2,0) coordinate  (A)
			(2,0) coordinate  (B)
			(-1.1,0) coordinate (C) 
			(1.1,0) coordinate (D) 
			($(-1.1,-1.1^2+4)$) coordinate  (F)
			 ($(1.1,-1.1^2+4)$) coordinate (E);
		\draw[domain=-2:2] plot  (\x,{-(\x)^2+4});
		\draw  (H) circle (1pt) node[shift={(-90:4mm)}]{$H$};
		\fill (A) circle  (1pt) node[shift={(-90:4mm)}]{$A$};
		 \fill (B) circle (1pt) node[shift={(-90:4mm)}]{$B$};
		 \fill (C) circle (1pt) node[shift={(-90:4mm)}]{$C$};
		\fill (D)  circle (1pt) node[shift={(-90:4mm)}]{$D$};
		\fill (E) circle (1pt)  node[shift={(0:4mm)}]{$E$};
		 \fill (F) circle (1pt) node[shift={(180:4mm)}]{$F$};
		\fill (G) circle  (1pt) node[shift={(90:4mm)}]{$G$};
		\draw[pattern=north  west lines] (D)--(E)--(F)--(C)--cycle;
		 \draw[dashed] (H)--(G);
		 \draw (A)--(B);
	\end{tikzpicture}
	 }
	 \shortans{4,53}
	\loigiai{\begin{center}
	\begin{tikzpicture}[smooth,samples=100,>=stealth,font=\small]
			 \draw[->] (-3,0)--(3,0) node[shift={(-90:4mm)}]{$x$};
			 \draw[->] (0,-1)--(0,4.5) node[shift={(0:4mm)}]{$y$};
			\path 
				(0,0) coordinate  (H)
				(0,4)  coordinate (G)
				(-2,0)  coordinate (A)
				(2,0)  coordinate (B)
				(-1.1,0) coordinate (C) 
				(1.1,0) coordinate (D) 
				($(-1.1,-1.1^2+4)$) coordinate (F) 
				 ($(1.1,-1.1^2+4)$) coordinate (E);
			\draw[domain=-2:2] plot (\x,{-(\x)^2+4}); 
			\draw (H) circle  (1pt) node[shift={(-60:4mm)}]{$H$};
			\fill (A) circle (1pt)  node[shift={(-90:4mm)}]{$A$} node[shift={(150:4mm)}]{$-2$};
			\fill (B) circle (1pt) node[shift={(-90:4mm)}]{$B$} node[shift={(30:4mm)}]{$2$}; 
			\fill (C) circle (1pt)  node[shift={(-90:4mm)}]{$C$} node[shift={(30:5mm)}]{$-1{,}1$};
			\fill (D) circle  (1pt) node[shift={(-90:4mm)}]{$D$} node[shift={(135:4mm)}]{$1{,}1$};
			 \fill (E) circle (1pt) node[shift={(0:4mm)}]{$E$} node[shift={(30:14mm)}]{$f (x)=-x^2+4$};
			\fill (F) circle (1pt)  node[shift={(180:4mm)}]{$F$};
			\fill (G) circle  (1pt) node[shift={(30:4mm)}]{$G$} node[shift={(120:4mm)}]{$4$};
			\fill (0,2.79) circle (1pt) node[shift={(-30:5mm)}]{$2{,}79$}; 
			\draw  (D)--(E)--(F)--(C)--cycle;
			\draw (A)--(B); 
		\end{tikzpicture}
\end{center}
Chọn hệ trục tọa độ $Oxy$ sao cho $AB$ trùng $Ox$, $H$ trùng với $O$. Khi đó parabol có đỉnh là $G(0; 4)$ đi  qua điểm $A(-2; 0)$ và điểm $B(2; 0)$.\\
Suy ra hàm số parabol có dạng $y = f(x) =  - x^2 + 4$.\\
Diện tích cánh cổng là $$S = \displaystyle \int_{-2}^{2} \left(-x^2 + 4\right) \mathrm{\,d}x = \left(-\dfrac{x^3}{3} + 4x\right)\bigg|_{-2}^{4}  = \dfrac{32}{3} \left(\text{m}^2\right).$$
Chiều cao của cửa là hình chữ nhật $CDEF$ có $CF = DE = f(1{,}1) = 2{,}79$; chiều rộng của cửa $CD =  4-2\cdot 0{,}9 = 2{,}2$ (m).\\
Diện tích phần cánh cửa là $$S_{CDEF} =  CD\cdot CF = 2{,}79 \cdot 2{,}2 = 6{,}138 \left(\text{m}^2\right).$$
Diện tích phần xiên hoa trang trí là $$S_{xh} = S - S_{CDEF} = \dfrac{32}{3} - 6{,}138 \approx  4{,}53 \left(\text{m}^2\right).$$
} 
\end{bt}

% nguồn: Phan2_Lop_12_Dai_so_Chuong_4
\begin{bt}%[2D4V3-2]%[ma:51b7215a]
	\immini{Bác Bình muốn làm cửa rào sắt hình parabol có chiều cao từ mặt đất đến đỉnh là $3$ m, chiều rộng tiếp giáp với mặt đất là $2$ m (được mô phỏng bởi hình vẽ bên dưới). Giá $1$ m$^2$ vật tư và công làm là $1{,}4$ triệu đồng, các chi phí khác không đáng kể. Vậy để làm cửa rào sắt đó thì số tiền Bác Bình phải trả là bao~nhiêu triệu đồng (làm tròn kết quả đến hàng phần chục)?}{\begin{tikzpicture}[line join=round,line cap=round,>=stealth,font=\scriptsize]
		\tikzset{every node/.style={scale=1}} 
		\path (-1,0) coordinate  (A)
			(1,0)  coordinate (B)
			(0,3)  coordinate (C);
		\begin{scope}
			\clip (-1.5,-1) rectangle  (1.5,3.5);
			\fill[blue,thin,pattern=north east lines](-1,0) --plot[samples=100,domain=-1:1,smooth,variable=\x]  (\x,{-3*(\x)^2+3})--(1,0);
			 \draw plot[samples=100,domain=-1:1,smooth,variable=\x] (\x,{-3*(\x)^2+3})
				 (-1,0)--(1,0);
		\end{scope}
		\path  (A)--(B)--([turn]-90:0.35) coordinate (xt)
			 ($(A)-(B)+(xt)$) coordinate (yt);
		\draw[>=stealth,|<->|] (xt)--(yt) node[fill=white,inner sep=0pt,midway,sloped]{$2\  \mathrm{m}$};
		\path ($(C)+(0:1.2)$) coordinate  (xt)
			 ($(B)+(0:0.2)$) coordinate (yt);
		\draw[dash pattern=on 2pt off  2pt] (C)--(xt)
			 (B)--(yt);
		\draw[>=stealth,|<->|] (xt)--(yt)  node[fill=white,inner sep=0pt,midway,sloped]{$3\ \mathrm{m}$};
	\end{tikzpicture}}
	\shortans{5,6} 
	\loigiai{\immini{Do tính đối xứng của parabol  nên chọn hệ trục tọa độ $Oxy$ sao cho parabol có đỉnh $I\in Oy$.\\
		Gọi phương trình parabol $(P)\colon  y=ax^2+bx+c$ $(a\ne 0)$.\\
		Ta có $\heva{& I(0;3)\in (P) \\ & A(-1;0)\in (P) \\ & B(1;0)\in (P) }\Rightarrow \heva{& 3=c \\ & a-b+c=0 \\  & a+b+c=0} \Leftrightarrow \heva{& a=-3 \\ & b=0 \\ & c=3.}$\\
		Do  đó $(P)\colon y=-3x^2+3$.\\
		Diện tích cửa  rào hình parabol là
		\allowdisplaybreaks 
		\begin{eqnarray*}
			&S=&\displaystyle\int_{-1}^1 (-3x^2+3)\mathrm{\,d}x\\ 
			&=&2\displaystyle\int_0^1 (-3x^2+3)\mathrm{\,d}x\\ 
			 &=&2\cdot \left(-x^3+3x\right)\big|_0^1 \\
			&=&4\,  \left(\text{m}^2\right).
		\end{eqnarray*}
		Vậy số tiền Bác Bình phải trả là $4\cdot 1{,}4=5{,}6$ triệu đồng.}{\begin{tikzpicture}[line join=round,line cap=round,>=stealth,font=\scriptsize]
		\tikzset{every node/.style={scale=0.9}} 
		\draw[->] (-2,0)--(2.5,0) node[below left] {$x$}; 
		\draw[->] (0,-1.1)--(0,4) node[below left] {$y$}; 
		\draw (0,0) node [above  left] {$O$};
		\begin{scope}
			\clip  (-1.5,-1) rectangle (1.5,3.5);
			\fill[blue,thin,pattern=north east lines](-1,0)  --plot[samples=100,domain=-1:1,smooth,variable=\x] (\x,{-3*(\x)^2+3})--(1,0);
			\draw plot[samples=100,domain=-1:1,smooth,variable=\x] (\x,{-3*(\x)^2+3}); 
		\end{scope}
		\draw (1,0) node[below]{$B  (1; 0)$}
			(-1,0) node[below]{$A (-1; 0)$} 
			(0,3) node[above  right]{$I (0; 3)$};
	\end{tikzpicture}}
	} 
\end{bt}

% nguồn: Phan2_Lop_12_Dai_so_Chuong_4
\begin{bt}%[2D4V3-2]%[ma:38d27128]
	Một cái cổng có kích thước như hình bên dưới.  Vòm cổng có dạng một Parabol. Tính diện tích hai cánh cửa cổng \textit{(kết quả làm tròn đến hàng phần chục)}. \\
	 \centerline{
	\begin{tikzpicture}[scale=1,>=stealth,line join=round,line cap=round,font=\footnotesize]
			 \draw [smooth,samples=100,fill=brown!80!black!80] (-2.5,0)--(2.5,0)--++(0,2)--plot [domain=2.5:-2.5] (\x,{-2/25*(\x)^2+2.5})--cycle;
			\draw [white]  (0,0)--(0,2.5);
			 \foreach \x/\y in {2.575/-.05,-3.575/-.05}{
				\fill [red] (\x,\y)--++(1,0)--++(90:2.2)--++(-1,0)--cycle; 
				\fill [pattern=bricks,pattern  color=white] (\x,\y)--++(1,0)--++(90:2.2)--++(-1,0)--cycle;
				\fill [gray!40]  (\x,\y)++(-.05,2.2) rectangle ++(1.1,.2);
				\fill  [black!80] (\x,\y)++(.25,2.4) rectangle ++(.5,.05)
				(\x,\y)++(.45,2.4) rectangle  ++(.1,.1);
				\fill [orange!80] (\x,\y)++(.5,2.75) circle  (.25);
			 }
			\fill  (2.5,.4) rectangle ++(.075,.4)
				(2.5,1.3)  rectangle ++(.075,.4)
				(-2.5,.4) rectangle ++(-.075,.4) 
				(-2.5,1.3) rectangle ++(-.075,.4); 
			\draw (0,2.5)--++(4.5,0) 
				 (0,0)--++(5,0)
				(2.5,2)--++(1.2,0); 
			\draw [<->] (4.5,0)--++(90:2.5) node [midway,right] {$2,5$ m}; 
			\draw [<->] (3.7,0)--++(90:2)  node [midway,right] {$2$ m};
			\draw [<->] (-2.5,-.3)--(2.5,-.3) node [midway,fill=white,inner sep=.5pt] {$5$ m}; 
		\end{tikzpicture}
	 }
	\shortans{11,7} 
	 \loigiai{
		Chọn  hệ tọa độ $Oxy$ như hình vẽ bên dưới (đơn vị trên mỗi trục tọa độ là $1$ m). \\
		 \centerline{
		\begin{tikzpicture}[scale=1,>=stealth,line join=round,line cap=round,font=\footnotesize]
				\draw  [smooth,samples=100,fill=brown!80!black!80] (-2.5,0)--(2.5,0)--++(0,2)--plot [domain=2.5:-2.5] (\x,{-2/25*(\x)^2+2.5})--cycle;
				 \draw [white] (0,0)--(0,2.5);
				\foreach \x/\y in {2.575/-.05,-3.575/-.05}{ 
					\fill [red]  (\x,\y)--++(1,0)--++(90:2.2)--++(-1,0)--cycle;
					\fill [pattern=bricks,pattern color=white]  (\x,\y)--++(1,0)--++(90:2.2)--++(-1,0)--cycle;
					\fill [gray!40] (\x,\y)++(-.05,2.2) rectangle  ++(1.1,.2);
					\fill [black!80] (\x,\y)++(.25,2.4) rectangle  ++(.5,.05)
					(\x,\y)++(.45,2.4)  rectangle ++(.1,.1);
					\fill  [orange!80] (\x,\y)++(.5,2.75) circle (.25);
				} 
				\fill (2.5,.4) rectangle  ++(.075,.4)
					(2.5,1.3) rectangle  ++(.075,.4)
					 (-2.5,.4) rectangle ++(-.075,.4)
					(-2.5,1.3) rectangle ++(-.075,.4); 
				\draw [->]  (-4,2.5)--(4,2.5) node [above] {$x$};
				\draw  [->] (0,-.5)--(0,3) node [left] {$y$};
				 \draw [dashed] (2.5,2.5)--(2.5,2)--(0,2)
					 (-2.5,2.5)--(-2.5,2);
				\draw (0,2.5) circle  (1pt) node [below left] {$O$}
					(2.5,2.5) circle  (1pt) node [above] {$2{,}5$}
					(-2.5,2.5) circle (1pt)  node [above] {$-2{,}5$}
					(0,0) circle (1pt)  node [below left] {$-2{,}5$}
					(0,2) circle (1pt)  node [below right] {$-0{,}5$};
			\end{tikzpicture}
		 } \\
		Diện tích hai cánh cổng bằng diện tích phần hình phẳng giới hạn bởi Parabol, đồ  thị $y=-2{,}5$ và hai đường thẳng $x=-2{,}5$, $x=2{,}5$. \\
		Vì Parabol có đỉnh $O(0;0)$  nên phương trình có dạng $y=ax^2$. \\
		Parabol  đi qua điểm $(2;-0{,}5)$ nên
		 $$-0{,}5=a \cdot (2{,}5)^2 \Leftrightarrow a=-\dfrac{2}{25}.$$
		Như vậy Parabol có  phương trình $y=-\dfrac{2}{25}x^2$. \\
		 Diện tích hai cánh cổng là
		 $$S=\displaystyle\int_{-2{,}5}^{2{,}5}\left|-\dfrac{2}{25}x^2+2{,}5\right|\mathrm{\,d}{x}=\dfrac{35}{3} \approx 11{,}7 \; \left(\text{m}^2\right).$$
	 }
\end{bt}

% nguồn: Phan2_Lop_12_Dai_so_Chuong_4
\begin{bt}%[2D4V3-2]%[ma:bb56c187]
	Một biển quảng cáo có dạng hình vuông $ABCD$ cạnh  bằng $4$ m và $I$ là trung điểm của đoạn thẳng $CD$. Trên tấm biển đó có đường parabol đỉnh $I$ đi qua $A$, $B$ và cắt đường chéo $BD$ tại $M$ \textit{(tham khảo hình vẽ)}
	\begin{center}
		\begin{tikzpicture}[scale=1,font=\footnotesize,line join=round,line cap=round,>=stealth]
				\path 
					 (0,0) coordinate (A)
					(4,0)  coordinate (B)
					 (4,4) coordinate (C)
					(0,4) coordinate (D) 
					(1,3) coordinate (M) 
					($(D)!1/2!(C)$)  coordinate (I);
				\draw(A)--(B)--(C)--(D)--(A); 
				 \draw(B)--(D);
				 \draw[smooth,samples=100,domain=0:4] plot (\x,{-1*(\x)^2+4*(\x)});
				 \fill[pattern=crosshatch,pattern color=gray] plot[domain=1:4](\x,{-1*(\x)^2+4*(\x)})--(B);
				\fill[pattern=dots] (A)--plot[domain=0:1](\x,{-1*(\x)^2+4*(\x)})--(M)--(B); 
				\draw  (1.5,1) node[above]{$S_2$}
					(2.5,2.5) node[above]{$S_1$}; 
				\foreach \x/\g in {A/-120,B/-90,C/0,D/180,I/90,M/180}  \fill (\x) circle (1pt)+(\g:0.3) node{$\x$};
			\end{tikzpicture}
	\end{center}
	Chi phí sơn phần tô hình kẻ chéo(có diện tích $S_1$) là $200\,000$ đồng/$\mathrm{m}^2$, chi phí sơn phần tô chấm( có  diện tích $S_2$) là $150\,000$ đồng /$\mathrm{m}^2$ và phần còn lại là $120\,000$ đồng/$\mathrm{m}^2$. Số tiền cần chi trả để sơn tấm biển quảng cáo là bao nhiêu nghìn đồng?
	 \shortans{2465}
	\loigiai{ 
		Gắn hệ trục tọa độ $Oxy$ như hình  vẽ.
		\begin{center}
			\begin{tikzpicture}[scale=1,font=\footnotesize,line join=round,line cap=round,>=stealth]
					\path 
						 (0,0) coordinate (A)
						(4,0) coordinate  (B)
						(4,4) coordinate  (C)
						(0,4) coordinate (D) 
						(1,3)  coordinate (M)
						($(D)!1/2!(C)$) coordinate (I); 
					\draw[->]  (-0.5,0) -- (5,0) node[below]{$x$};
					\draw[->] (0,-0.5) --  (0,5) node[left]{$y$};
					 \draw(A)--(B)--(C)--(D)--(A);
					 \draw(B)--(D);
					\draw[smooth,samples=100,domain=0:4] plot (\x,{-1*(\x)^2+4*(\x)}); 
					\fill[pattern=crosshatch,pattern color=gray]  plot[domain=1:4](\x,{-1*(\x)^2+4*(\x)})--(B);
					\fill[pattern=dots]  (A)--plot[domain=0:1](\x,{-1*(\x)^2+4*(\x)})--(M)--(B);
					 \draw (1.5,1) node[above]{$S_2$}
						(2.5,2.5) node[above]{$S_1$}; 
					\foreach \x/\g in {A/-120,B/-90,C/0,D/180,I/90,M/180} \fill (\x)  circle (1pt)+(\g:0.3) node{$\x$};
				\end{tikzpicture}
		\end{center}
		Khi đó $BD\colon y=-x+4$ và Parabol $(P)\colon y=ax^2+bx+c$ đi qua ba điểm $A\equiv O(0;0)$, $B(4;0)$, $I(2;4)$ nên Parabol có dạng $y=-x^2+4x$.\\ 
		$M(1;3)$ là  giao điểm của $BD$ và $(P)$.\\
		 Suy ra $S_1=\displaystyle\int_{1}^{4} \left|(-x^2+4x)-(-x+4)\right| \mathrm{\,d}x=\dfrac{9}{2}\left(\mathrm{m}^2\right)$ và $S_2=\displaystyle\int_{0}^{4} |-x^2+4x| \mathrm{\,d}x-\dfrac{9}{2}=\dfrac{37}{6}\left(\mathrm{m}^2\right)$.\\
		Diện tích  phần còn lại là $S_3=4^2-\displaystyle\int_{0}^{4} |-x^2+4x| \mathrm{\,d}x=\dfrac{16}{3}$.\\
		Số tiền  cần chi trả để sơn tấm biển quảng cáo là
		\begin{eqnarray*}
			&&200\,000\cdot S_1+150\,000\cdot  S_2+120\,000\cdot S_3 \\
			&=&200\,000\cdot \dfrac{9}{2}+150\,000\cdot \dfrac{37}{6}+120\,000\cdot \dfrac{16}{3}\\ 
			&=&2\,465\,000 \text{ đồng}. 
		\end{eqnarray*}
	} 
\end{bt}

% nguồn: Phan2_Lop_12_Dai_so_Chuong_4
\begin{bt}%[2D4V3-2]%[ma:9e9e193a]
	\immini[thm]{Hưởng ứng chính sách  hiện đại hóa nông thôn, người dân ở khu phố $A$ đồng lòng cùng nhau góp tiền đổ bê tông một đường đi trong khu phố (phần được tô đậm trong hình vẽ). Biết rằng khi chọn hệ trục tọa độ $Oxy$ với đơn vị độ dài trên mỗi trục tọa độ là $10$ m, các đường cong $AB$, $CD$ là mép đường được cho bởi đồ thị hàm số $f(x)=\dfrac{1}{32}x^2-\dfrac{3}{8}x+\dfrac{5}{2}$ và $g(x)=\dfrac{1}{32}x^2-\dfrac{5}{8}x+1$, đồng thời lớp bê tông được đổ dày 16 cm và giá tiền $1$ m$^3$ bê tông là $1\,080\,000$ đồng. Tính số tiền (đơn vị triệu đồng) cần dùng để đổ bê tông con đường đó (làm tròn kết quả đến hàng đơn vị).}
	{\begin{tikzpicture}[scale=0.6,font=\footnotesize,line join=round,line cap=round,>=stealth]
		\def\xmin{-2} 
		\def\xmax{6} 
		\def\ymin{0} 
		\def\ymax{6} 
		\draw[teal,->] (\xmin-1,0)--(\xmax+1,0) node[below] {$x$}; 
		\draw[teal,->] (0,\ymin-1)--(0,\ymax+1) node[left] {$y$}; 
		\draw[teal] (0,0)  node [below right] {$O$};
		\newcommand{\f}[1]{0.06*#1^3 -  0.1*#1^2 - 0.56*#1+0.62}
		 \draw[teal,blue,domain=-2.2:5.3,smooth] plot (\x,{\f{\x}});
		\newcommand{\g}[1]{0.06*#1^3 -  0.04*#1^2 - 0.51*#1+1.52}
		 \draw[teal,blue,domain=-0.5:3,smooth]
			plot  (\x,{\g{\x}});
		 \draw[teal,smooth,samples=200] plot[domain=-2.3:5.3](\x,{0.06*(\x)^3-0.04*\x^2-0.51*\x+1.52});
		\coordinate  (C) at (-1,{\f{-1}});
		\draw[teal,red,fill=red] (C) circle (2pt); 
		\node[below left] at  (C) {$C$};
		 \coordinate (A) at (-1,{\g{-1}});
		 \draw[teal,red,fill=red] (A) circle (2pt);
		 \node[above left] at (A) {$A$};
		\coordinate (B)  at (5,{\g{5}});
		\draw[teal,red,fill=red] (B)  circle (2pt);
		 \node[above left] at (B) {$B$};
		\coordinate  (D) at (5,{\f{5}});
		\draw[teal,red,fill=red]  (D) circle (2pt);
		\node[below  right] at (D) {$D$};
		 \draw[teal,dashed] (-1,0) node[below]{$-2$}--(A);
		 \draw[teal,dashed] (5,0) node[below]{$7$}--(B);
		 \fill[smooth,pattern=north east lines](C)-- plot[domain=-1:5](\x,{\f{\x}})--(B)--plot[domain=5:-1](\x,{\g{\x}})--(A)--cycle;
	\end{tikzpicture}
	} 
	\shortans{253} 
	 \loigiai{
		Diện tích hình phẳng được  giới hạn bởi hai đường cong $AB$, $CD$ là $\displaystyle\int_{-2}^{7} \left(\dfrac{1}{4}x+1\right)\mathrm{\,d}x =\dfrac{117}{8}$.\\
		Vì đơn vị trên mỗi trục là $10$ m nên diện tích phần đổ bê tông thực tế là $\dfrac{117}{8} \cdot 100  =\dfrac{2\,925}{2}\, \text{ m}^2.$\\
		Thể tích của bê tông để đổ con đường đó  là
		$0{,}16  \cdot \dfrac{2\,925}{2} = 0{,}08 \cdot 2\,925 = 234$ m$^3$.\\
		 Số tiền cần dùng để đổ bê tông là
		$1\,080\,000 \cdot 234  = 252\,720\,000 \approx 253$ triệu đồng.
	} 
\end{bt}

% nguồn: Phan2_Lop_12_Dai_so_Chuong_4
\begin{bt}%[2D4V3-2]%[ma:cd612b0c]
	Mái của một nhà thi đấu rộng $20$\,m và dài $60$\,m được kéo theo dạng vòm parabol. Hỏi thể tích không khí  tăng thêm là bao nhiêu mét khối nếu chọn chiều cao đỉnh mái $10$\,m thay vì phương án ban đầu $8$\,m?
	\begin{center}
		\begin{tikzpicture}[scale=0.5,line join=round,line cap=round,>=stealth]
				 \tikzset{declare function={f (\x)=-3/16*(\x)^2+1.5*(\x); g (\x)=-0.25*(\x)^2+2*(\x);
					xI=4; yI=4; xmin=xI-4.5; xmax=xI+4.5; 
					ymin=yI-0.5; ymax=yI+2;  a=7.8;},
					smooth,samples=100 
				 }
				 \draw[->] (xmin-0.5,0)--(0,0) node[below left]{$O$} --(xmax+1,0) node[below left] {$x$};
				 \draw[->] (0,-1.5)--(0,ymax) node[below right] {$y$};
				\fill (0,0)  circle (1pt);
				\draw (8,0)--($(8,0)+(50:a)$)  coordinate (A) node[pos=0.5,right]{$60$};
				\fill[red!15] (2,3) --++(50:a) plot[domain=7:13] (\x,{g (\x-5)+6})--(8,0)--plot[domain=8:2] (\x,{g (\x)})--cycle; 
				\draw[dashed] (0,3)--(4,3) 
					(0,4)--(4,4); 
				\draw[gray](2,3) --++(50:a); 
				\draw[red,domain=0:8] plot (\x,{f  (\x)});
				 \draw[blue,domain=0:8] plot (\x,{g (\x)}) plot[domain=7:13] (\x,{g (\x-5)+6});
				\fill 
					(4,0) circle (2pt)  node[shift={(-90:7pt)}]{$10$}
					(8,0) circle (2pt)  node[shift={(-90:7pt)}]{$20$}
					(4,3) circle (2pt) 
					(4,4) circle  (2pt)
					(0,4) circle (2pt) node[shift={(180:7pt)}]{$10$} 
					(0,3) circle  (2pt) node[shift={(180:7pt)}]{$8$};
			\end{tikzpicture}
	\end{center}
	\shortans{1600} 
	\loigiai{ 
		Giả sử (trong hệ trục  như hình)
		\begin{itemize}
			\item Với mái  cao $10$\,m, parabol có dạng $f(x)=a x^2+bx+c$. Ta có
			 \[\heva{&f(0)=0\\&f(10)=10\\&f(20)=0}\Rightarrow\heva{&c=0\\&10^2a+10b+c=10\\&20^2a+20b+c=0}\Rightarrow\heva{&a=-\dfrac{1}{10}\\&b=2\\&c=0.}\]
			 Khi đó $f(x)=-\dfrac{1}{10}x^2+2x$.
			\item Với mái cao $8$\,m, parabol có dạng $g(x)=m x^2+nx+p$.  Ta có
			\[\heva{&g(0)=0\\&g(10)=8\\&g(20)=0}\Rightarrow\heva{&p=0\\&10^2m+10n+p=8\\&20^2m+20n+p=0}\Rightarrow\heva{&m=-\dfrac{2}{25}\\&n=\dfrac{8}{5}\\&p=0.}\] 
			 Khi đó $g(x)=-\dfrac{2}{25}x^2+\dfrac{8}{5}x$.
		\end{itemize}
		Diện tích phần chênh lệch trong mặt  cắt ngang là\\ $S=\int_0^{20}\big|f(x)-g(x)\big|\mathrm{d}x=\int_0^{20}\big|-\dfrac{1}{50}x^2+\dfrac{2}{5}x\big|\mathrm{d}x=\dfrac{80}{3}$\,m$^2$.\\
		Vậy thể tích  tăng thêm là $\dfrac{80}{3}\cdot 60=1\,600$\,m$^3$.
	} 
\end{bt}

% nguồn: Phan2_Lop_12_Dai_so_Chuong_4
\begin{bt}%[2D4V3-2]%[ma:2a86e62e]
	\immini{Trên cửa sổ có dạng hình chữ nhật của ngôi nhà một doanh nghiệp kinh doanh hải sản, hoạ sĩ cần thiết kế logo hình con cá. Logo là hình phẳng giới hạn bởi hai parabol với các kích thước được cho trong hình bên (đơn vị trên mỗi trục toạ độ là dm). Hoạ sĩ cần tính diện tích của logo để báo giá cho doanh nghiệp đó trước khi kí hợp đồng. Diện tích của logo bằng $x$ dm$^2$. Tìm $x$  (làm tròn kết quả đến hàng đơn vị).}{
	\begin{tikzpicture}[scale=0.7,font=\footnotesize,line join=round,line cap=round,>=stealth]
			\fill[yellow!50,domain=-4:4] plot (\x,{-(\x)^2/8+2}) -- plot  (\x,{3*(\x)^2/16-3});
			\fill[yellow!50,domain=-5:-4]  (-5,0)-- plot (\x,{3*(\x)^2/16-3})--cycle;
			\fill[yellow!50,domain=-5:-4] (-5,0)--  plot (\x,{-(\x)^2/8+2})--cycle;
			\draw[->] (-6,0)--(5,0) node[below] {$x$}; 
			\draw[->] (0,-4)--(0,3) node[left] {$y$}; 
			 \draw[domain=-5:4,samples=100] plot (\x,{-(\x)^2/8+2});
			 \draw[domain=-5:4,samples=100] plot (\x,{3*(\x)^2/16-3});
			\draw (-5,2) rectangle (4,-3); 
			 \draw (3,0.4) circle (3pt);
			\coordinate (A) at  (-5,{-(-5)^2/8+2});
			\coordinate (B) at (-5,{3*(-5)^2/16-3}); 
			 \coordinate (C) at (-4,0);
			\coordinate (D)  at (0,2);
			\coordinate  (E) at (4,0);
			\coordinate (G) at (0,-3); 
			\coordinate (O)  at (0,0);
			\coordinate (-5) at (-5,0); 
			\coordinate (4) at (4,0); 
			\foreach \x/\y  in {A/180,B/180,C/-90,D/45,E/-45,G/45,O/135,-5/135,4/45}{
				 \draw (\x) circle (1pt)
				($(\x)+(\y:4mm)$)  node {$\x$};
			 }
			 \draw[<->] (4.7,0)--(4.7,2) node[midway,right] {$2$ dm};
			\draw[<->] (4.7,0)--(4.7,-3) node[midway,right]  {$3$ dm};
			\draw[<->] (-5,-3.5)--(-4,-3.5) node[midway,below]  {$1$ dm};
			 \draw[<->] (-4,-3.5)--(0,-3.5) node[midway,below] {$4$ dm};
			\draw[<->] (0,-3.5)--(4,-3.5) node[midway,below]  {$4$ dm};
		\end{tikzpicture}}
	 \shortans{28}
	 \loigiai{
		Parabol $y=f(x)$ có dạng $f(x)=ax^2+bx+c$ $(a\neq  0)$.\\
		Parabol nhận $Oy$  làm trục đối xứng nên $b=0$.\\
		Vì đồ thị đi qua điểm $(4;0)$ và $(0;2)\Rightarrow  a=-\dfrac{1}{8}$ và $c=2$.\\
		 $\Rightarrow f(x)=-\dfrac{1}{8}x^2+2$.\\
		Parabol $y=g(x)$ có dạng $f(x)=a'x^2+b'x+c'$ $(a'\neq 0)$.\\ 
		Parabol nhận $Oy$ làm trục  đối xứng nên $b'=0$.\\
		Vì đồ thị  đi qua điểm $(4;0)$ và $(0;4)\Rightarrow a=\dfrac{3}{16}$ và $c=-3$.\\
		$\Rightarrow g(x)=\dfrac{3}{16}x^2-3$.\\ 
		Diện tích  của logo là $S=\displaystyle \int _{-5}^4\left|f(x)-g(x)\right|\mathrm{\,d}x=\displaystyle \int _{-5}^4\left|5-\dfrac{5}{16}x^2\right|\mathrm{\,d}x\approx 28$ (dm$^2$).
	 }
\end{bt}

% nguồn: Phan2_Lop_12_Dai_so_Chuong_4
\begin{bt}%[2D4V3-2]%[ma:38c563e8]
	Ông Vượng mới khai phá được một mảnh đất hình chữ nhật, nhà nước chưa cấp sổ nên ông chưa biết chính xác diện tích mảnh đất là bao nhiêu. Mảnh đất của ông Vượng nằm ở một vị trí thuận lợi để trồng trọt vì có một dòng suối nhỏ chảy qua với hình dáng một parabol, dòng suối nhỏ này chia mảnh đất ra làm hai phần có diện tích $S_1$, $S_2$  ($S_1 > S_2$). Mảnh đất có diện tích $S_2$ được xem như hình phẳng giới hạn bởi parabol cùng hai tiếp tuyến vuông góc của parabol đó. Vào vụ Hè thu, ông Vượng quyết định trồng lúa trên phần đất có diện tích $S_1$ và trồng ớt trên phần đất có diện tích $S_2$. Dự kiến lợi nhuận mang lại từ việc trồng lúa là $30$ nghìn/m$^2$ và lợi nhuận từ việc trồng ớt là $40$ nghìn/m$^2$ (trong một vụ mùa). Ông quyết định dựng hệ trục $Oxy$ như hình vẽ với gốc $O$ trùng với điểm cực trị của dòng suối dạng parabol, đơn vị trên mỗi trục là $100$ mét. Tổng lợi nhuận theo dự kiến của ông Vượng sau vụ Hè thu này (làm tròn đến hàng đơn vị của triệu đồng) là bao nhiêu? Biết rằng diện tích con suối không đáng kể.
	\begin{center}
		\begin{tikzpicture}[scale=3,font=\footnotesize,line join=round,line cap=round,>=stealth]
				\draw[teal,->] (-.5,0)  -- (1.5,0) node[below] {$x$};
				 \draw[teal,->] (0,-.5) -- (0,1.5) node[left] {$y$};
				\coordinate[label=right:$A$] (A) at (1,1); 
				 \coordinate (O) at (0,0);
				\coordinate[label=above:$B$] (B)  at (-.25,0.06);
				\coordinate[label=above:$C$] (C) at  (0.375,1.31);
				\coordinate[label=below:$E$] (E) at (0.375,-0.25); 
				\draw[teal]  (0,0) node[below left]{$O$}
					(B)--(C)--(A); 
				 \clip (-.5,-.5) rectangle (1.5,1.5);
				 \draw[teal,samples=150,smooth,domain=-5:5] plot (\x,{(\x)^2});
				\draw[teal,samples=150,smooth,domain=-5:5] plot  (\x,{2*(\x)-1});
				 \draw[teal,samples=150,smooth,domain=-5:5] plot (\x,{-0.5*(\x)-1/16});
				\draw[teal,pattern = north  west lines,opacity=.5,pattern color = black] plot[domain=-5:5] (B)--(C)--(A)--plot[domain=1:-.25](\x,{(\x)^2});
				\draw[teal,pattern = north east lines,opacity=.5,pattern color =  black] plot[domain=-5:5] (E)--(B)--plot[domain=-.25:1](\x,{(\x)^2})--(E);
				\foreach \diem in {A,B,C,E,O}\fill (\diem)  circle (.25pt);
				\fill[white] (.4,.6) circle  (3pt) node{$S_1$}
					(.3,-.06)  circle (3pt) node{$S_2$}
					(1.25,1.3) node{$d_1$} 
					(.8,-.35)  node{$d_2$};
			\end{tikzpicture}
	\end{center}
	\shortans{309} 
	 \loigiai{
		\begin{center}
			\begin{tikzpicture}[scale=3,font=\footnotesize,line join=round,line cap=round,>=stealth]
					\draw[teal,->] (-.5,0) -- (1.5,0) node[below] {$x$}; 
					\draw[teal,->] (0,-.5) -- (0,1.5) node[left] {$y$}; 
					 \coordinate[label=right:$A$] (A) at (1,1);
					\coordinate (O) at  (0,0);
					 \coordinate[label=above:$B$] (B) at (-.25,0.06);
					\coordinate[label=above:$C$] (C) at  (0.375,1.31);
					\coordinate[label=below:$E$] (E) at (0.375,-0.25); 
					\draw[teal] (0,0) node[below  left]{$O$}
						(B)--(C)--(A); 
					 \clip (-.5,-.5) rectangle (1.5,1.5);
					 \draw[teal,samples=150,smooth,domain=-5:5] plot (\x,{(\x)^2});
					\draw[teal,samples=150,smooth,domain=-5:1.25] plot  (\x,{2*(\x)-1});
					\draw[teal,samples=150,smooth,domain=-5:5] plot (\x,{-0.5*(\x)-1/16}); 
					\foreach \diem in {A,B,C,E,O}\fill (\diem)  circle (.25pt);
					 \draw[teal] (.4,.6) node{$S_1$}
						(.3,-.06) node{$S_2$} 
						(1.25,1.3) node{$d_1$} 
						(.8,-.35)  node{$d_2$};
					\draw[teal,dashed] (1,0) node[below]{1}|-(0,1)  node[left]{1};
				\end{tikzpicture}
		\end{center}
		Đầu tiên ta đặt $A\left(a; a^2\right)$, $B\left(b; b^2\right)$, $ a > 0$, $b < 0$ là hai tiếp điểm ứng với hai tiếp tuyến  vuông góc của parabol $(P)$.\\
		Gọi $d_1$, $ d_2$  lần lượt là các tiếp tuyến của đồ thị $(C)$ tại $A$, $ B$.\\
		Khi đó $d_1\colon y = 2a x - a^2 $; $d_2\colon y =  2b x - b^2$.\\
		Do $d_1 \perp d_2$ nên  $2a \cdot 2b = -1 \Rightarrow b = -\dfrac{1}{4a} \Rightarrow B\left(-\dfrac{1}{4a};\dfrac{1}{16a^2}\right)$.\\
		Khi đó $d_2 \colon y = -\dfrac{x}{2a} -  \dfrac{1}{16a^2}$.\\
		Gọi $E = d_2 \cap d_1$, suy ra $E\left(\dfrac{4a^2 - 1}{8a};\ -\dfrac{1}{4}\right)$; $EA = \dfrac{\sqrt{(4a + 1)^3}}{8a}$;  $ EB = \dfrac{\sqrt{(4a + 1)^3}}{16a^2}$. \\
		Diện tích mảnh đất $S = EA \cdot EB  = \dfrac{(4a+1)^3}{128a^3}$.\\
		Ta có $EA = 2EB \Rightarrow a = 1 \Rightarrow  S = \dfrac{125}{128}$.\\
		Ta có 
		$d_1\colon y = 2x - 1$; $  d_2\colon y = -\dfrac{x}{2} - \dfrac{1}{16}$
		và  $A(1; 1)$, $ B\left(-\dfrac{1}{4}; \dfrac{1}{16}\right)$, $ E\left(\dfrac{3}{8}; -\dfrac{1}{4}\right)$.\\
		 Ta có
		$S_2 =\displaystyle\int_{-\frac{1}{4}}^{\frac{3}{8}} \left[x^2 - \left(-\dfrac{x}{2} - \dfrac{1}{16} \right)\right] \mathrm{\,d}x  +\displaystyle\int_{\frac{3}{8}}^{1} \left[x^2 - (2x - 1)\right] \mathrm{\,d}x = \dfrac{125}{768}$
		và $S_1 = S - S_2 = \dfrac{625}{768}$.\\ 
		Tổng số tiền  thu được của ông Vượng sau vụ hè thu là
		 \[
		\dfrac{625}{768} \cdot 10^2 \cdot 30 +  \dfrac{125}{768} \cdot 10^2 \cdot 40 \approx 309\,244 \text{ nghìn đồng} \approx {309} \text{ triệu đồng}.
		 \]
	} 
\end{bt}

% nguồn: Phan2_Lop_12_Dai_so_Chuong_4
\begin{bt}%[2D4V3-2]%[ma:b452b9a9]
	Khu vực trung tâm một quảng trường có dạng hình tròn đường kính $AB$ bằng $10$ m. Người ta trang trí khu vực này bằng hai đường Parabol đối xứng nhau qua $AB$, nằm trong hình tròn, đi qua các điểm $A$, $B$ và có đỉnh cách mép hình tròn  $1$ m. Phần giới hạn bởi $2$ parabol được trồng hoa với chi phí $200$ nghìn đồng $1$ mét vuông, phần còn lại được lát gốm sứ với chi phí $800$ nghìn đồng $1$ mét vuông. Tính tổng chi phí (triệu đồng) để hoàn thành khu vực này (kết quả làm tròn đến hàng phần chục)?
	\begin{center}
		\begin{tikzpicture}[scale=1,font=\footnotesize,line join=round,line cap=round,>=stealth,x=0.6cm,y=0.6cm,declare function = {f (\x)=(4/25)*(\x)^2-4; g (\x)=-(4/25)*(\x)^2+4;}]
				 \def \xmin{-5.8};
				 \def \xmax{5.8};
				\def  \ymin{-5.1};
				 \def \ymax{5.1};
				\clip (\xmin+0.1,\ymin+0.1) rectangle  (\xmax-0.1,\ymax-0.1);
				\draw[smooth,samples=100,domain=-5:5] plot (\x,{f (\x)}) plot  (\x,{g (\x)});
				\draw (0,0) circle  (5);
				\draw (-5,0)-- (5,0); 
				\foreach \x/\y/\l/\g in {-5/0/A/135,5/0/B/45}{\fill (\x,\y) circle (1pt)+(\g:0.5)  node{$\l$};}
			\end{tikzpicture}
	\end{center}
	 \shortans{30,8}
	 \loigiai{
		Xét hệ trục  độ như hình vẽ
		\begin{center}
			\begin{tikzpicture}[scale=1,font=\footnotesize,line join=round,line cap=round,>=stealth,x=0.6cm,y=0.6cm,declare function = {f (\x)=(4/25)*(\x)^2-4; g (\x)=-(4/25)*(\x)^2+4;}]
					\def  \xmin{-6.5};
					\def  \xmax{7.6};
					 \def \ymin{-5.8};
					\def  \ymax{6.5};
					\draw[->] (\xmin,0.) -- (\xmax,0.) node[anchor=north] {$x$}; 
					\draw[->]  (0.,\ymin) -- (0.,\ymax) node[anchor=west] {$y$};
					\clip (\xmin+0.1,\ymin+0.1)  rectangle (\xmax-0.1,\ymax-0.1);
					\draw[smooth,samples=100,domain=-5:5] plot (\x,{f (\x)}) plot (\x,{g (\x)}); 
					\draw  (0,0) circle (5);
					\draw[dashed] (,) node{} |- (,) node{}; 
					\foreach \x/\y/\l/\g in {0/0/O/225,0/4/4/135,0/5/5/135,-5/0/A/135,5/0/B/45,0/-4/-4/225,0/-5/-5/225}{\fill (\x,\y) circle  (1pt)+(\g:0.5) node{$\l$};}
				\end{tikzpicture}
		\end{center}
		Phương  trình đường tròn là $x^2+y^2=25$.\\
		 Phương trình của Parabol có bề lõm hướng lên là $y=\dfrac{4}{25}x^2-4$.\\
		 Diện tích phần trồng hoa giới hạn bới $2$ parabol là
		\[S_1=2\int_{-5}^5\left|\left(\dfrac{4}{25}x^2-4\right)\right|\mathrm{\,d}x=4\int_0^5\left|\left(\dfrac{4}{25}x^2-4\right)\right|\mathrm{\,d}x=\dfrac{160}{3}\  (\mathrm{m}^2).\]
		Diện tích toàn bộ phần hình tròn  là $S_2=25\pi$ (m$^2$).\\
		Diện tích phần còn lại để trang trí gốm sứ là  $S=S_2-S_1=25\pi-\dfrac{160}{3}$ (m$^2$).\\
		Vậy tổng chi phí để làm khu vực trung tâm quảng trường là 
		\[\dfrac{160}{3}\cdot200+\left(25\pi-\dfrac{160}{3}\right)\cdot 800\approx 30{,}83\, \text{  (triệu đồng)}.\]
	 }
\end{bt}

% nguồn: Phan2_Lop_12_Dai_so_Chuong_4
\begin{bt}%[2D4V3-2]%[ma:7c45a123]
	Một người nghệ sĩ đã vẽ hình chiếc nơ theo một cách khác lạ so với các nhà thiết kế. Anh ta vẽ hình chữ nhật $ABCD$ tâm $O$ có chiều dài bằng $4$\,dm, chiều rộng bằng $2$\,dm. Chiếc nơ hình $(H)$ nằm bên trong hình chữ nhật sao cho khi kẻ tia $Ot$ bất kì cắt $(H)$ và cạnh hình chữ nhật  lần lượt tại $M$ và $N$ thì $MN=1$\,dm. Tính diện tích chiếc nơ hình $(H)$ đó theo dm$^2$ (làm tròn kết quả đến hàng phần mười).
	\begin{center}
		\begin{tikzpicture}[>=stealth,line join=round,line cap=round,font=\footnotesize,scale=1]
				\def\a{6}\def\b{3} 
				\draw[teal] (0,0)  coordinate (C) rectangle +(\a,\b) coordinate (A);
				\path (0,\b) coordinate (B) 
					 (\a,0) coordinate (D)
					($(B)!0.5!(D)$) coordinate  (O);
				\draw[teal,name  path=cong] (O).. controls +(80:0.5cm) and +(-170:0.5cm)..($(O)!0.5!(A)$)
				to[bend right=0.5cm]($(O)!0.5!(D)$).. controls +(170:0.5cm) and +(-80:0.5cm)..(O);
			\draw[teal] (O)..  controls +(100:0.5cm) and +(-10:0.5cm)..($(O)!0.5!(B)$)
			to[bend left=0.5cm]($(O)!0.5!(C)$).. controls +(20:0.5cm) and +(-100:0.5cm)..(O);
		\path (A)++(-1,0) coordinate  (N);
		\path[name  path=ON] (O)--(N);
		\path[name  path=ab] (A)--(B);
		\path[name path=circleO]  (O) circle (1cm);
		\path[name intersections={of= ON and cong}] coordinate (M) at  (intersection-2);
		\draw[teal] (O)--(N)--($(O)!1.4!(N)$)  node[above left]{$t$};
		\foreach  \diem/\goc in {A/0,B/180,C/180,D/0,O/180,N/135,M/90}
		\fill (\diem)  circle (1pt) ++(\goc:0.3) node{$\diem$};
		\end{tikzpicture}
	\end{center}
	 \shortans{1,52}
	\loigiai{ 
		Xét hình vẽ và các kí hiệu  như sau
		\begin{center}
			\begin{tikzpicture}[>=stealth,line join=round,line cap=round,font=\footnotesize,scale=1]
					 \def\a{6}\def\b{3}
					\draw[teal] (0,0) coordinate (C) rectangle +(\a,\b)  coordinate (A);
					\path  (0,\b) coordinate (B)
						(\a,0) coordinate (D) 
						($(B)!0.5!(D)$) coordinate (O); 
					 \draw[teal,->] (O)--+(0:4) node[shift=(-90:0.2)]{$x$};
					 \draw[teal,->] (O)--+(90:3) node[shift=(180:0.2)]{$y$};
					\draw[teal,name path=cong] (O).. controls +(80:0.5cm) and  +(-170:0.5cm)..($(O)!0.5!(A)$)
					to[bend right=0.5cm]($(O)!0.5!(D)$).. controls +(170:0.5cm) and +(-80:0.5cm)..(O);
				\draw[teal] (O).. controls +(100:0.5cm) and  +(-10:0.5cm)..($(O)!0.5!(B)$)
				to[bend left=0.5cm]($(O)!0.5!(C)$).. controls +(20:0.5cm) and +(-100:0.5cm)..(O);
			\path 
				(A)++(-1,0) coordinate  (N)
				(A)++(0,-0.5) coordinate  (N')
				($(A)!0.5!(D)$)  coordinate (H)
				 ($(A)!0.5!(B)$) coordinate (K);
			\path[name path=ON]  (O)--(N);
			\path[name path=ON']  (O)--(N');
			\path[name  path=ab] (A)--(B);
			\path[name path=circleO]  (O) circle (1cm);
			\path[name intersections={of= ON and cong}] coordinate (M)  at (intersection-2);
			\path[name intersections={of= ON' and  cong}] coordinate (M') at (intersection-2);
			\draw[teal] (O)--(N)--($(O)!1.4!(N)$) node[above left]{$t$} 
				(O)--(N')--($(O)!1.4!(N')$)  node[above left]{$t'$};
			\draw[teal] pic["\scriptsize $\theta$",draw,angle eccentricity=2,angle radius=3mm]{angle =  N--O--K};
			\draw[teal] pic["\scriptsize $\varphi$",draw,angle eccentricity=2,angle radius=3mm]{angle  = H--O--N'};
			 \foreach \diem/\goc in {A/90,B/90,C/180,D/0,O/180,N/135,M/90,H/30,K/30}
			 \fill (\diem) circle (1pt) ++(\goc:0.3) node{$\diem$};
			 \foreach \diem/\goc in {M'/45,N'/45}
			\fill (\diem) circle (1pt)  ++(\goc:0.4) node{$\diem$};
			\end{tikzpicture}
		\end{center}
		Gọi $\varphi=(Ox,Ot')$ thì $\cos\varphi=\dfrac{OH}{ON'}=\dfrac{2}{r_{\varphi}+1}\Rightarrow r_{\varphi}=\dfrac{2}{\cos\varphi}-1$; với $r_{\varphi}=OM'$ quay quanh gốc $O$ khi $0< \varphi  < \widehat{HOA}$.\\
		Gọi $\theta=(Oy,Ot)$ thì $\cos\theta=\dfrac{OK}{ON}=\dfrac{1}{r_{\theta}+1}\Rightarrow r_{\theta}=\dfrac{1}{\cos\theta}-1$; với $r_{\theta}=OM$ quay quanh gốc $O$ khi $0< \theta < \widehat{AOK}$.\\ 
		Ta  có
		$\tan\widehat{AOH}=\dfrac{AH}{OH}=\dfrac{1}{2}\Rightarrow \widehat{AOH}=arc\tan0{,}5$.\\ 
		Và  $\tan\widehat{AOK}=\dfrac{AK}{OK}=2\Rightarrow \widehat{AOK}=arc\tan2$.\\
		Do đó diện tích cần  tính là
		\[S=4\cdot  \dfrac{1}{2}\left[\displaystyle\int^{arc\tan0{,}5}_0{{\left(\dfrac{2}{\cos\varphi}-1\right)}^2}\mathrm{d}\varphi+\displaystyle\int^{arc\tan2}_0{{\left(\dfrac{1}{\cos\theta}-1\right)}^2}\mathrm{d}\theta \right]\approx 1{,}52\,(\text{dm}^2).\]
	 }
\end{bt}

% nguồn: Phan2_Lop_12_Dai_so_Chuong_4
\begin{bt}%[2D4V3-2]%[ma:6cbfd012]
	Trong Công viên Toán học có những mảnh đất mang hình dáng khác nhau. Mỗi mảnh được trồng một loài hoa và nó được tạo thành bởi  một trong những đường cong đẹp trong toán học. Ở đó có một mảnh đất mang tên Bernoulli, nó được tạo thành từ đường Lemmiscate có phương trình trong hệ tọa độ Oxy là $16y^2=x^2 (25-x^2)$ như hình vẽ bên dưới (đơn vị trên trục là mét). Biết chi phí để trồng hoa trên mảnh đất Bernoulli là $150$ nghìn đồng/m². Số tiền cần dùng để trồng hoa trên mảnh đất này bằng bao nhiêu nghìn đồng?
	\begin{center}
		\begin{tikzpicture}[line join=round,line cap=round,>=stealth,scale=0.6]
				\tikzset{every node/.style={scale=0.9}} 
				\draw[->] (-6,0)--(6,0) node[below  left] {$x$};
				\draw[->]  (0,-5.1)--(0,5.1) node[below left] {$y$};
				\draw (0,0) node  [below left] {$O$};
				\draw[blue,samples=200,domain=0:5,smooth,variable=\x] plot  (\x,{((\x)*sqrt (25-(\x)^2))/4});
				\draw[blue,samples=200,domain=-5:0,smooth,variable=\x] plot  (\x,{-((\x)*sqrt (25-(\x)^2))/4});
				\draw[blue,samples=200,domain=-5:0,smooth,variable=\x] plot (\x,{((\x)*sqrt (25-(\x)^2))/4}); 
				\draw[blue,samples=200,domain=0:5,smooth,variable=\x]  plot (\x,{-((\x)*sqrt (25-(\x)^2))/4});
			\end{tikzpicture}
	\end{center}
	\shortans{6250} 
	\loigiai{Ta  có $y=0 \Leftrightarrow x^2 (25-x^2)=0 \Leftrightarrow \hoac{& x=0 \\& x=\pm 5.}$\\
		Suy ra đường Lemmiscate cắt trục hoành tại ba điểm là $A(-5;0)$, $O(0;0)$ và $B(5;0)$.\\ 
		Từ phương trình $16y^2=x^2  (25-x^2) \Rightarrow y^2=\dfrac{1}{16}x^2 (25-x^2)$\\
		Suy ra  $y=\pm \dfrac{1}{4}x\sqrt{25-x^2}$.\\
		Xét góc phần tư thứ nhất, đường Lemmiscate có  phương trình là $y=\dfrac{1}{4}x\sqrt{25-x^2}$ với $x \in [0;5]$.\\
		Khi  đó \[S_{1}=\displaystyle\int_0^5 \dfrac{1}{4}x\sqrt{25-x^2} \mathrm{\,d}x=\dfrac{125}{12}
		\Rightarrow S=4S_{1}=\dfrac{125}{3} (\text{m}^2).\] 
	Số tiền cần dùng  để trồng hoa trên mảnh đất này là: $\dfrac{125}{3} \cdot 150=6\,250$ nghìn đồng.}
\end{bt}

% nguồn: Phan2_Lop_12_Dai_so_Chuong_4
\begin{bt}%[2D4V3-2]%[ma:1aef5603]
	 \immini[thm]{
	Trang trí một sân hình chữ nhật kích thước $28$ m $\times 16$ m, trong đó hai Parabol $\left(P_1\right)$ đối xứng với $\left(P_2\right)$ qua đường thẳng đi qua hai trung điểm của chiều dài sân (hình vẽ), khoảng cách giữa hai đỉnh Parabol bằng  $4$ m. Chi phí trang trí cho phần tô đậm là $180$ ngàn đồng trên một mét vuông, phần trắng là $160$ ngàn đồng trên một mét vuông. Tổng chi phí trang trí cho sân là bao nhiêu triệu đồng (làm tròn kết quả đến hàng phần mười)?
	 }{
	\begin{tikzpicture}[scale=0.7,font=\footnotesize,line join=round,line cap=round,>=stealth,xscale=0.3,yscale=0.3,rotate=90]
			\def\xmin{-10} \def\xmax{10} 
			 \def\ymin{-16} \def\ymax{16}
			\def\f[#1]{1/4*(#1)^2-2} 
			 \def\g[#1]{-1/4*(#1)^2+2}
			 \fill[gray!25] plot[domain=-8:8](\x,{\f[\x]})--plot[domain=8:-8](\x,{\g[\x]});
			\clip (\xmin+0.1,\ymin+0.1) rectangle  (\xmax-0.1,\ymax-0.1);
			\draw[teal] (-8,-14)  rectangle (8,14);
			\draw[teal,smooth,samples=300,domain=-8:8]  plot (\x,{\f[\x]});
			\draw[teal,smooth,samples=300,domain=-8:8] plot (\x,{\g[\x]}); 
			\draw[teal,->] (3,7) node  [below] {$\left(P_1\right)$}--(5,6);
			\draw[teal,->]  (3,-7) node [below] {$\left(P_2\right)$}--(5,-6);
		\end{tikzpicture}
	} 
	 \shortans{74,4}
	\loigiai{ 
	 \immini{
		Đặt hình chữ nhật vào hệ tọa độ $Oxy$ sao  cho tâm hình chữ nhật trùng với gốc tọa độ, chiều dài của hình chữ nhật vuông góc với trục $Ox$.\\
		Ta có Parabol $(P_1)\colon y=ax^2+bx+c$ đi qua $(0;-2)$,  $(-8;14)$, $(8;14)$.\\
		Suy ra $\heva{&  c=-2\\& 64a-8b=16\\& 64a+8b=16} \Leftrightarrow \heva{& c=-2\\& b=0\\& a=\dfrac{1}{4}.}$\\
		Do  đó $(P_1)\colon y=\dfrac{1}{4}x^2-2$; $(P_2)\colon y=-\dfrac{1}{4}x^2+2$.\\
		Gọi $S_1$, $S_2$ lần lượt là diện tích phần tô đậm  và diện tích phần trắng.\\
		 Lại có $S_1=2\displaystyle\int_{-2}^{-2\sqrt{2}} \left|\dfrac{1}{2}x^2-4\right| \mathrm{\,d}x + \displaystyle\int_{-2\sqrt{2}}^{2\sqrt{2}} \left|\dfrac{1}{2}x^2-4\right| \mathrm{\,d}x$.\\
		Suy  ra $S_2=16\cdot 28-S_1$.\\
		Vậy tổng chi phí trang trí là $180\,000\cdot S_1+160\,000\cdot S_2\approx 74{,}4$  (triệu đồng).
		 }{
		\begin{tikzpicture}[scale=0.7,font=\footnotesize,line join=round,line cap=round,>=stealth,xscale=0.35,yscale=0.3]
				\def\xmin{-10} \def\xmax{10} 
				\def\ymin{-16}  \def\ymax{16}
				\def\f[#1]{1/4*(#1)^2-2} 
				 \def\g[#1]{-1/4*(#1)^2+2}
				\fill[gray!25]  plot[domain=-8:8](\x,{\f[\x]})--plot[domain=8:-8](\x,{\g[\x]});
				\draw[teal,->]  (\xmin,0)--(\xmax,0) node [below]{$x$};
				 \draw[teal,->] (0,\ymin)--(0,\ymax) node [left]{$y$};
				\fill  (0,0) circle (2.5pt)
					($(0,0)+(-135:12mm)$)  node{$O$};
				 \clip (\xmin+0.1,\ymin+0.1) rectangle (\xmax-0.1,\ymax-0.1);
				\draw[teal]  (-8,-14) rectangle (8,14);
				\draw[teal,smooth,samples=300,domain=-8:8]  plot (\x,{\f[\x]});
				 \draw[teal,smooth,samples=300,domain=-8:8] plot (\x,{\g[\x]});
				\foreach \x/\y/\g/\d in{-8/0/135/-8,8/0/45/8,0/-14/45/-14,0/-2/-45/-2,0/2/45/2,0/14/-45/14} 
				\fill (\x,\y)  circle (2.5pt)
					($(\x,\y)+(\g:16mm)$)  node{$\d$};
			\end{tikzpicture}
		 }
	 }
\end{bt}

% nguồn: Phan2_Lop_12_Dai_so_Chuong_4
\begin{bt}%[2D4V3-2]%[ma:8efd8ce5]
	\immini{Một mảnh đất hình chữ nhật có chiều dài $60$ m, chiều rộng $20$ m. Người ta muốn trồng hoa ở hai đầu của mảnh đất hai hình bằng nhau giới hạn bởi hai  đường parabol có hai đỉnh cách nhau $40$ m và nhận đường trung bình của hình chữ nhật làm trục đối xứng (như hình vẽ bên dưới). Phần còn lại của mảnh đất người ta lát gạch. Tính diện tích phần lát gạch (m$^2$) (kết quả làm tròn đến hàng đơn vị).
	\shortans{933} 
	 }
	{\begin{tikzpicture}[line cap=butt,line join=miter,>=stealth,xscale=.2,yscale=.08,font=\normalsize]
		 \tikzset{declare function={xmin=-5; xmax=30; ymin=-5; ymax=70;
			 f (\x)=-.1*(\x)^2+2*(\x);
			g  (\x)=.1*(\x)^2-2*(\x)+60;
			}, 
			smooth,samples=450 
		} 
		\draw[->] (xmin,0)--(xmax,0) node[shift={(-100:7pt)}]{$  x $};
		\draw[->] (0,ymin)--(0,ymax) node[shift={(190:7pt)}]{$  y $};
		 \fill (0,0) node[shift={(225:10pt)}]{$ O $};
		 \draw (0,60)--(20,60)--(20,0);
		\draw[<->,dashed](22,0)--(22,60) node[pos=0.5,sloped,below]{$60$m}; 
		 \draw[<->,dashed](10,10)--(10,50) node[pos=0.5,sloped,above]{$40$m};
		\draw[<->,dashed](0,65)--(20,65)  node[pos=0.5,sloped,above]{$20$m};
		\draw  (0,60) node[left]{$A$} circle (1.2pt)
			(20,60)  node[above right]{$B$} circle (1.2pt)
			 (20,0) node[below]{$C$} circle (1.2pt)
			(10,10)  node[above right]{$F$} circle (1.2pt)
			(10,50) node[below right]{$E$} circle (1.2pt); 
		\fill[pattern=north west  lines]plot[domain=0:20] (\x,{f (\x)})--(0,0)--(20,0);
		\fill[pattern=north west lines]plot[domain=0:20] (\x,{g (\x)})--(0,60)--(20,60); 
		 \clip (xmin,ymin) rectangle (xmax,ymax);
		\draw  plot[domain=0:20] (\x,{f (\x)});
		\draw[violet] plot[domain=0:20] (\x,{g (\x)}); 
	\end{tikzpicture}}
	\loigiai{Gán  hệ trục tọa độ như hình vẽ
		\begin{center}
			\begin{tikzpicture}[line cap=butt,line join=miter,>=stealth,xscale=.2,yscale=.08,font=\normalsize]
					\tikzset{declare  function={xmin=-5; xmax=30; ymin=-5; ymax=70;
						f  (\x)=-.1*(\x)^2+2*(\x);
						g  (\x)=.1*(\x)^2-2*(\x)+60;
						 },
						smooth,samples=450 
					} 
					\draw[->]  (xmin,0)--(xmax,0) node[shift={(-100:7pt)}]{$ x $};
					\draw[->] (0,ymin)--(0,ymax) node[shift={(190:7pt)}]{$ y  $};
					\fill (0,0) node[shift={(225:10pt)}]{$  O $};
					\draw  (0,60)--(20,60)--(20,0);
					\draw[<->,dashed](22,0)--(22,60) node[pos=0.5,sloped,below]{$60$m}; 
					\draw[<->,dashed](10,10)--(10,50) node[pos=0.5,sloped,above]{$40$m}; 
					\draw[<->,dashed](0,65)--(20,65)  node[pos=0.5,sloped,above]{$20$m};
					\draw  (0,60) node[left]{$A$} circle (1.2pt)
						(20,60) node[above  right]{$B$} circle (1.2pt)
						(20,0) node[below]{$C$} circle  (1.2pt)
						(10,10) node[above right]{$F$} circle  (1.2pt)
						(10,50) node[below right]{$E$} circle  (1.2pt);
					\fill[pattern=north west lines]plot[domain=0:20] (\x,{f  (\x)})--(0,0)--(20,0);
					\fill[pattern=north west lines]plot[domain=0:20] (\x,{g (\x)})--(0,60)--(20,60); 
					\clip (xmin,ymin) rectangle (xmax,ymax); 
					\draw plot[domain=0:20] (\x,{f  (\x)});
					\draw[violet] plot[domain=0:20]  (\x,{g (\x)});
				\end{tikzpicture}
		\end{center}
		Tọa độ các điểm $O(0;0)$, $A(0;60)$, $B(20;60)$, $C(20;0)$; $E(10;50)$; $F(10;10)$.\\ 
		 Parabol $(P)$ có đỉnh $F(10;10)$, có phương trình dạng $y=a(x-10)^2+10$.\\
		$(P)$ đi qua $C(20;0)$ $\Leftrightarrow a\cdot(20-10)^2+10=0\Leftrightarrow  a=-\dfrac{1}{10}$.\\
		Vậy $(P)\colon  y=-\dfrac{1}{10} (x-10)^2+10$.\\
		Diện tích trồng cỏ là $S=2\displaystyle\displaystyle\int\limits_0^{20}\left[-\dfrac{1}{10} (x-10)^2+10\right]\mathrm{\,d}x=\dfrac{800}{3}$~(m$^2$).\\ 
		Diện tích lát gạch  $S_g=20\cdot 60-S=1200-\dfrac{800}{3}=\dfrac{2800}{3}\approx 933$~(m$^2$).\\
	 }
\end{bt}

% nguồn: Phan2_Lop_12_Dai_so_Chuong_4
\begin{bt}%[2D4V3-2]%[ma:5348d1d2]
	Trong cơ khí chế tạo, một chi tiết máy hình đĩa tròn có dạng như hình vẽ, nhận $AB$ và $CD$ làm các trục đối xứng. Người ta cần phủ sơn cả hai mặt của chi tiết. Biết rằng đường tròn lớn có bán kính $5$ dm, các đường tròn nhỏ đều có bán kính bằng  $2$ dm, $AB=CD=4$ dm và chi phí sơn là $82.000$ đồng/m$^2$. Chi phí để sơn hoàn thiện chi tiết máy bằng bao nhiêu nghìn đồng (kết quả làm tròn đến hàng đơn vị).
	\begin{center}
		\begin{tikzpicture}[scale=0.7,font=\footnotesize,line join=round,line cap=round,>=stealth]
				\filldraw[fill=white,draw] (0,0)  circle (4.985cm);
				\clip (0,0) circle  (5.01cm);
				\filldraw[fill=white,draw] (4,0)  circle (2cm);
				\filldraw[fill=white,draw] (-4,0) circle  (2cm);
				\filldraw[fill=white,draw] (0,4)  circle (2cm);
				 \filldraw[fill=white,draw] (0,-4) circle (2cm);
				\draw[dashed](-2,0) node[above right]{$A$}--(2,0) node[above right]{$B$}; 
				 \draw[dashed](0,-2) node[above right]{$C$}--(0,2) node[above right]{$D$};
				\fill (2,0) circle (2pt); 
				\fill (-2,0) circle  (2pt);
				\fill (0,2)  circle (2pt);
				\fill (0,-2) circle  (2pt);
			\end{tikzpicture}
	\end{center}
	\shortans{65} 
	\loigiai{Ta dựng hệ trục toạ độ  $Oxy$ như sau:
		\begin{center}
			\begin{tikzpicture}[scale=0.7,font=\footnotesize,line join=round,line cap=round,>=stealth]
					\filldraw[fill=white,draw] (0,0)  circle (4.985cm);
					\clip (0,0)  circle (5.01cm);
					\filldraw[fill=white,draw] (4,0) circle (2cm); 
					\filldraw[fill=white,draw] (-4,0) circle  (2cm);
					\filldraw[fill=white,draw] (0,4) circle (2cm); 
					 \filldraw[fill=white,draw] (0,-4) circle (2cm);
					\draw[dashed](-2,0) node[above  right]{$A$}--(2,0) node[above right]{$B$};
					\draw[dashed](0,-2)  node[above right]{$C$}--(0,2) node[above right]{$D$};
					\fill  (2,0) circle (2pt);
					\fill  (-2,0) circle (2pt);
					\fill (0,2) circle  (2pt);
					\fill  (0,-2) circle (2pt);
					 \fill (0,0) circle (2pt);
					\node  at (-0.3,-0.3){$O$};
					\fill (4,0) circle (2pt); 
					\draw[->](-5,0)--(4,0)  node[below]{$I_1$}--(4.8,0) node[below]{$x$};
					\draw[->](0,-5)--(0,4.8)  node[right]{$y$};
					\node  at (1.8,-0.25){$2$};
				\end{tikzpicture}
		\end{center}
		\begin{itemize}
			\item Vì các đường tròn nhỏ đều có bán kính là $2$ dm và nhận  $AB$, $CD$ là trục đối xứng nên ta chỉ cần xét $1$ đường tròn là được.
			\item Ta có: $OA=OB=\dfrac{AB}{2}=2$ dm. 
			\item Xét đường tròn $(S_1)$ có bán kính $R_1=2$ do đó $(S_1)$ có tâm $I_1(4,0)$. Phương trình  đường tròn $(S_1)$:
			$$(S_1):(x-4)^2+y^2=4 \Leftrightarrow y^2=4-(x-4)^2  \Leftrightarrow y=\pm \sqrt{4-(x-4)^2}.$$
			\item Vì đường tròn đối xứng qua trục $Ox$ nên ta chỉ cần xét phần dương. Tức  là: $$y= \sqrt{4-(x-4)^2}.$$
			\item Xét đường tròn  $(S)$ có bán kính $R=5$ và tâm $I(0,0)$. Phương trình đường tròn $(S)$:
			$$(S):x^2+y^2=25 \Leftrightarrow y^2=25-x^2 \Leftrightarrow y=\pm \sqrt{25-x^2},$$ và tương tự ta cũng chỉ cần  xét phần dương: $y=\sqrt{25-x^2}$.
			\item Phương trình hoành độ giao điểm  của hai đường tròn $(S)$ và $(S_1)$:
			$$\sqrt{25-x^2} = \sqrt{4-(x-4)^2} \Leftrightarrow 25-x^2 =4-(x-4)^2 \Leftrightarrow  8x=37 \Leftrightarrow x =4{,}625.$$
			\item  Ta sẽ tính diện tích phần bị cắt đi, ta sẽ chia thành $2$ phần như sau:
		\subitem Diện tích được giới hạn trên đoạn $[2;4{,}625]$, nửa đường  tròn: $y=\sqrt{4-(x-4)^2}$ và đường thẳng $y=0$.
		\subitem Diện tích được giới hạn trên đoạn $[4{,}625;5]$, nửa đường tròn: $y=\sqrt{25-x^2}$ và đường thẳng  $y=0$.
			\item Như đã nói trên ta chỉ xét phần dương của đường tròn nên khi tính toán ta tính thêm phần âm bằng cách nhân kết quả cho $2$ và có $4$ phần bị cắt  như nhau nên ta nhân kết quả cho $4$. Ta tính:
			$$S_{\text{ bị cắt đi}}=4\cdot 2 \cdot \left(\int_{2}^{4{,}625}\sqrt{4-(x-4)^2} \mathrm{\,d}x+\int_{4{,}625}^{5}\sqrt{25-x^2} \mathrm{\,d}x\right) \approx 38{,}8\text{ $($dm}^2).$$ 
			\item Diện tích $1$ mặt của chi tiết máy: $S_{\text{ chi tiết máy}}=S_{\text{ đĩa tròn}}-S_{\text{ bị cắt đi}}=5^2\cdot \pi - 38{,}8 \approx 39{,}74$. 
			\item Vì người ta cần sơn $2$ mặt  nên tổng diện tích là: $2\cdot 39{,}74=79{,}48$ dm$^2$ = $0{,}7948$ m$^2$.
			\item Số tiền để sơn là: $T=0{,}7948 \cdot 82.000\approx  65.174$ (đồng) $\approx 65$ (nghìn đồng).
	\end{itemize}}
\end{bt}

% nguồn: Phan2_Lop_12_Dai_so_Chuong_4
\begin{bt}%[2D4V3-2]%[ma:d3d0bfa6]
	Dọc bờ nước có một đường xe đạp. Trên một đoạn dài $20$\,m, người ta muốn thay bằng tuyến mới đi vòng một cái giếng (đường kính $2$\,m). Các đoạn nối phải trơn (chuyển tiếp êm), đoạn đường vòng qua giếng là phần đồ thị hàm số $f(x)=ax^4+bx^2+c$, nhận trục đối  xứng là đường thẳng đi qua tâm của giếng và vuông góc với đoạn đường.
	\begin{center}
		\begin{tikzpicture}[scale=1,line join=round,line cap=round,>=stealth]
				\tikzset{declare  function={a=-1/16; b=0.5; c=-1; f (\x)=a*(\x)^4+b*(\x)^2+c;
					xI=0; yI=c; xmin=xI-2.5;  xmax=xI+2.5;
					 ymin=yI-0.5; ymax=yI+2;},
					 smooth,samples=100
				 }
				 \tikzset{Icon-cano/.pic={
				\fill[ball color=brown!80!cyan!80!white](-4.65,0.08)..controls+(-23:2.5)  and+(-165:3.5)..(4.6,0)..controls+(15:0.15) and+(-15:0.5)..(3.9,0.2)..controls+(-172:5) and+(15:2)..(-4.65,0.08);
				\fill[ball  color=gray!10!brown] [red](-4.65,0.08)--+(-25:0.4)..controls+(10:2) and++(-172:5)..(3.8,0)--++(110:0.2)..controls+(-172:5) and+(15:2)..(-4.65,0.08);
				\fill[ball  color=gray!90!cyan](-4.65,0.08)..controls+(-23:2.5) and+(-165:3.5)..(4.6,0)..controls+(15:0.05) and++(65:1.2)..(4.2,-1)..controls++(-152:0.75) and+(13:2)..(-1,-2.5)..controls++(-167:1.1) and+(-50:3.5)..(-4.65,0.08);
				\fill[ball color=gray!90!cyan](2,0)  coordinate (e)--++(0:0.4)--++(88:0.65)--++(175:0.3) coordinate (f)--cycle;
					 }}
				 \fill[pattern=crosshatch] (-2,0)--plot[domain=-2:2] (\x,{f (\x)});
				\fill[cyan!10] (-4,0.1) rectangle (3.5,2); 
				 \draw[fill= cyan!80](0,-0.5) circle (7pt);
				\draw[red,line width=2pt,domain=-2:2] plot (\x,{f (\x)}); 
				\draw[gray,line  width=2pt](-4,0)--(-2,0)
					(2,0)--(3.5,0); 
				\draw  (-2,0)--(2,0)
					 (-4,0.1)--(3.5,0.1);
				 \path (-2.5,1) pic[scale=0.15]{Icon-cano}
					(2.6,1.3) pic[scale=0.12]{Icon-cano}; 
				\draw[|<->|](-2,0.2)--(2,0.2) node[pos=0.5,above]{$20$\,m}; 
				\draw[|<->|](2,-1)--(2,0)  node[pos=0.5,right]{$6$\,m};
				\draw[dashed,line width=0.3pt] (0,-1)--(2,-1); 
			\end{tikzpicture}
	\end{center}
	Diện tích phần tô màu cần lát  giữa bờ và tuyến mới là bao nhiêu mét vuông (làm tròn kết quả
	đến hàng đơn  vị)
	\shortans{61} 
	 \loigiai{
		Từ giả thiết ta chọn mô  hình hàm số $f(x)=a x^4+bx^2+c$ có đồ thị trong hệ truc $O xy$ như hình vẽ
		\begin{center}
			\begin{tikzpicture}[scale=1,line join=round,line cap=round,>=stealth]
					\tikzset{declare function={a=-1/16;  b=0.5; c=-1; f (\x)=a*(\x)^4+b*(\x)^2+c;
						xI=0; yI=c; xmin=xI-2.5;  xmax=xI+2.5;
						ymin=yI-0.5;  ymax=yI+2;},
						 smooth,samples=100
					 }
					\draw[->]  (xmin-0.5,0)--(0,0) node[below left]{$O$} --(xmax+1,0) node[below left] {$x$};
					\draw[->] (0,-1.5)--(0,ymax) node[below  right] {$y$};
					\fill (0,0) circle (1pt); 
					\draw[red,domain=-2.45:2.45]  plot (\x,{f (\x)});
					\fill (-2,0)  circle (1pt) node[shift={(90:7pt)}]{$-10$}
						(2,0)  circle (1pt) node[shift={(90:7pt)}]{$10$}
						(0,-1)  circle (1pt) node[shift={(-35:7pt)}]{$-6$};
				\end{tikzpicture}
		\end{center}
		Ta có $f'(x)=4a x^3+2bx$.\\ Từ hình vẽ ta có hệ phương trình  sau:
		\[\heva{&f(0)=-6\\&f(10)=0\\&f'(10)=0}\Rightarrow\heva{&c=-6\\&10^4  a+10^2b+c=0\\&4\cdot 10^3a+2\cdot 10b=0}\Rightarrow\heva{&a=-\dfrac{3}{5\,000}\\&b=\dfrac{3}{25}\\&c=-6.}\]
		 Do đó $f(x)=-\dfrac{3}{5\,000}x^4+\dfrac{3}{25}x^2-6$.\\
		Diện tích phần lát kể cả diện tích miệng  giếng là \[S_1=\int_{-10}^{10}\big|f(x)\big|\mathrm{d}x=2\int_{0}^{10}\big|-\dfrac{3}{5\,000}x^4+\dfrac{3}{25}x^2-6\big|\mathrm{d}x=64\text{ m}^2.\]
		 Diện tích miệng giếng là $S_2=\pi\cdot 1=\pi$\,m$^2$.\\
		Vậy diện tích phần lát là $S_1-S_2=64-\pi\approx  61$\,m$^2$.
	 }
\end{bt}

% nguồn: Phan2_Lop_12_Dai_so_Chuong_4
\begin{bt}%[2D4V3-2]%[ma:ee0750a7]
	Cho ba hàm số $y=f(x),\,y=g(x)$ và $y=h(x)$ có đồ thị như hình vẽ. Diện tích mỗi miền được ghi bằng số trong hình. Tính $\displaystyle\int_1^6{[h(x)-g(x)]\mathrm{\,d}x}+\displaystyle\int_0^6{[f(x)-g(x)]\mathrm{\,d}x}$. 
	\begin{center}
		\begin{tikzpicture}[scale=1,font=\footnotesize,line join=round,line cap=round,>=stealth]
				 \def\a{2} \def\b{-3} \def\c{1}
				\def\xmin{-1}  \def\xmax{7}
				\def\ymin{-1.} \def\ymax{4} 
				\draw[->]  (\xmin,0)--(\xmax,0) node [below]{$x$};
				\draw[->] (0,\ymin)--(0,\ymax) node  [right]{$y$};
				 \fill (0,0) circle (1.0pt) node[below left]{$O$};
				 \clip (\xmin+0.1,\ymin+0.1) rectangle (7.5,\ymax-0.1);
				\draw[smooth,samples=300,domain=0:7] plot  (\x,{0.02*(\x)^3-0.35*(\x)^2+1.84*(\x)});
				\draw[smooth,samples=300,domain=0:6.7] plot  (\x,{-0.05*(\x)^3+0.56*(\x)^2-1.94*(\x)+3});
				\draw[smooth,samples=300,domain=0:7] plot (\x,{0.01*(\x)^3-0.16*(\x)^2+0.36*(\x)+3}); 
				\fill (0,1) circle  (1.0pt) node[left]{$1$}
					(1,0) circle (1.0pt)  node[below]{$1$}
					(3,0)  circle (1.0pt) node[below]{$3$}
					(6,0)  circle (1.0pt) node[below]{$6$};
				 \draw[dashed] (1,0)--(1,1.5)
					(3,0)--(3,2.9) 
					(6,0)--(6,2.9); 
				 \fill[blue,smooth,pattern=north west lines,opacity=0.5]
					plot[domain=0:1](\x,{0.02*(\x)^3-0.35*(\x)^2+1.84*(\x)}) 
					 -- plot[domain=1:0] (\x,{-0.05*(\x)^3+0.56*(\x)^2-1.94*(\x)+3}) -- cycle;
				 \fill[blue,smooth,red!30,opacity=0.5]
					plot[domain=1:3](\x,{0.02*(\x)^3-0.35*(\x)^2+1.84*(\x)}) 
					--  plot[domain=3:1] (\x,{-0.05*(\x)^3+0.56*(\x)^2-1.94*(\x)+3}) -- cycle;
				\fill[blue,smooth,pattern=north east  lines,opacity=0.5]
					 plot[domain=3:6](\x,{0.02*(\x)^3-0.35*(\x)^2+1.84*(\x)})
					-- plot[domain=6:3] (\x,{0.01*(\x)^3-0.16*(\x)^2+0.36*(\x)+3}) --  cycle;
				 \fill[blue,smooth,pattern=fivepointed stars,opacity=0.5]
					plot[domain=3:6](\x,{-0.05*(\x)^3+0.56*(\x)^2-1.94*(\x)+3}) 
					-- plot[domain=6:3]  (\x,{0.01*(\x)^3-0.16*(\x)^2+0.36*(\x)+3}) -- cycle;
				\fill (1,0)  circle (1.0pt) node[below]{$1$}
					 (3,0) circle (1.0pt) node[below]{$3$}
					(6,0) circle (1.0pt)  node[below]{$6$}
					(0.5,1.5) circle (0pt) node{$5$} 
					(2,1.6) circle (0pt) node{$6$} 
					(1.2,2.5) circle  (0pt) node{$8$}
					(4.5,1.6)  circle (0pt) node{$4$}
					(5.4,2.5) circle (0pt) node{$3$} 
					(6.7,3) circle (0pt) node{$y=g (x)$} 
					(6.7,1.7) circle (0pt)  node{$y=f (x)$}
					 (6.7,.8) circle (0pt) node{$y=h (x)$};
			\end{tikzpicture}
	\end{center}
	\shortans{-3} 
	 \loigiai{
		 Ta có
		\begin{align*}
			&\displaystyle\int_1^6{[h(x)-g(x)]\mathrm{\,d}x}+\displaystyle\int_0^6{[f(x)-g(x)]\mathrm{\,d}x}  \\
			 =&\displaystyle\int_1^6{[h(x)-g(x)]\mathrm{\,d}x}+\displaystyle\int_0^3{[f(x)-g(x)]\mathrm{\,d}x}+\displaystyle\int_3^6{[f(x)-g(x)]\mathrm{\,d}x}\\
			 =&-( 6+4+3)+( 5+8)-3=-3.
		\end{align*}
	} 
\end{bt}

% nguồn: Phan2_Lop_12_Dai_so_Chuong_4
\begin{bt}%[2D4V3-2]%[ma:1b81cdad]
	\immini{Một vật trang trí bằng gỗ có dạng hình vuông cạnh $6$ dm. Người ta tạo hoa văn bởi bốn đường parabol có chung đỉnh là tâm của hình vuông và mỗi parabol đi qua hai đỉnh của hình vuông (như hình vẽ). Biết phần tô đậm trong hình được  sơn với giá $3$ nghìn đồng/dm$^2$ và phần còn lại được sơn với giá $2$ nghìn đồng/dm$^2$. Tính số tiền để sơn vật trang trí trên (đơn vị nghìn đồng). }
	{\begin{tikzpicture}[scale=0.5,font=\footnotesize,line join=round,line cap=round,>=stealth]
		\def\a{1/3} 
		\draw[samples=150,smooth,domain=-3:3] plot  (\x,{\a*(\x)^2});
		\draw[samples=150,smooth,domain=-3:3]  plot (\x,{-\a*(\x)^2});
		\draw  (-3,-3)--(3,-3)--(3,3)--(-3,3)--(-3,-3);
		 \draw[pattern =north west lines,opacity=.5,pattern color =black] plot[domain=-3:3](\x,{\a*(\x)^2})--(3,3)--(-3,3)--cycle;
		\draw[pattern =north west lines,opacity=.5,pattern  color =black] plot[domain=-3:3](\x,{-\a*(\x)^2})--(3,-3)--(-3,-3)--cycle;
	\end{tikzpicture}}
	 \shortans{84}
	\loigiai{ 
		Xét hệ trục tọa độ $Oxy$ như hình vẽ. 
		\begin{center}
			\begin{tikzpicture}[scale=0.5,font=\footnotesize,line join=round,line cap=round,>=stealth]
					 \def\a{1/3}
					\draw[samples=150,smooth,domain=-3:3]  plot (\x,{\a*(\x)^2});
					\draw[samples=150,smooth,domain=-3:3] plot (\x,{-\a*(\x)^2}); 
					\draw (-3,-3)--(3,-3)--(3,3)--(-3,3)--(-3,-3); 
					\draw[pattern =north west lines,opacity=.5,pattern color =black]  plot[domain=-3:3](\x,{\a*(\x)^2})--(3,3)--(-3,3)--cycle;
					\draw[pattern  =north west lines,opacity=.5,pattern color =black] plot[domain=-3:3](\x,{-\a*(\x)^2})--(3,-3)--(-3,-3)--cycle;
					\draw[->] (-3.5,0)--(5,0) node[below]{$x$}; 
					 \draw[->] (0,-3.5)--(0,5) node[left]{$y$};
					\draw  (0,0) node[below left]{$O$}
						(3,0) node[below right]{$3$} 
						(0,3) node[above  left]{$3$};
				\end{tikzpicture}
		\end{center}
		Gọi nhánh parabol phía trên  là $y=ax^2$ ($a\ne 0$).\\
		Do parabol đi qua điểm $(3;3)$,  suy ra $a=\dfrac{1}{3}$.\\
		 Vậy parabol cần tìm là $y=\dfrac{1}{3}x^2$.\\
		Diện tích của phần được tô đậm là $S_{\text{ tô đậm}}=2\displaystyle  \int_{-3}^{3}\left( \dfrac{1}{3}x^2\right) \mathrm{\,d}x=12$ (dm$^2$).\\
		Diện tích của vật trang trí là $36$  dm$^2$. \\
		Diện tích của phần không tô đậm  là $S_{\text{ không tô đậm}}=36-12=24$ (dm$^2$).\\
		Số  tiền để sơn vật trang trí là $12\cdot3 +24\cdot2=84$ (nghìn đồng).
	 }
\end{bt}

% nguồn: Phan2_Lop_12_Dai_so_Chuong_4
\begin{bt}%[2D4V3-2]%[ma:e87fe5cc]
	Sau khi đọc xong truyện ngắn rất cảm động \lq\lq Chiếc lá cuối cùng\rq\rq \,của nhà  văn O.Henry, bạn An muốn thiết kế hình ảnh một chiếc lá để gửi đến chị em Johnsy và Sue. Chiếc lá là một hình phẳng được giới hạn bởi hai phần parabol giống nhau. Biết một parabol có đỉnh là điểm $O$ và trục đối xứng của hình phẳng này là một đường thẳng tạo với phương ngang góc $60^{\circ}$. Với các kích thước được cho như hình vẽ, hãy tính diện tích của chiếc lá. (\textit{Làm tròn kết quả đến hàng đơn vị của} cm$^2$)
	\begin{center}
		\begin{tikzpicture}[scale=0.7,line join=round,line cap=round,>=stealth]
				\tikzset{every node/.style={scale=1}} 
				\def\xmin{-4}\def\xmax{5}\def\ymin{-3}\def\ymax{7} 
				\coordinate  (A) at (2,5);
				\coordinate (C)  at (2,0);
				\pgfmathsetmacro\a{sqrt  (3)}
				\coordinate (B) at  (-5*\a/3+2,0);
				\foreach \diem  in {A,C,B} \fill (\diem) circle (1pt);
				\draw[smooth]($(A)! 1.5 ! (B)$)--($(B)!  1.25 ! (A)$);
				\draw[dashed]($(A)! 1.25 ! (C)$)--($(C)!  1.25 ! (A)$)
					 ($(A)+(-1,0)$)--($(A)+(2,0)$)
					 (0,-1.25)--(0,6.25)
					(-3,0)--(4,0); 
				 \draw[<->](0,-1)--(2,-1);
				\draw (1,-1) node  [below]{$20$ cm};
				\draw[<->](3,0)--(3,5); 
				\draw (3,2.5) node [right]{$50$ cm}; 
				\draw (0,0) node[below right]{$O$}; 
				 \clip (\xmin,\ymin) rectangle (\xmax,\ymax);
				\draw[smooth,samples=200,domain=-2+.8*\a:2] plot (\x,{(5/4)*((\x)^2)}); 
				\draw[xshift=-38,yshift=22,rotate=120,smooth,samples=200,domain=-2+.8*\a:2] plot  (\x,{-(5/4)*((\x)^2)});
			\end{tikzpicture}
	\end{center}
	 \shortans{745}
	 \loigiai{
	\immini{Chọn hệ trục $Oxy$  như hình vẽ.\\
		Vì parabol có  đỉnh là điểm $O(0; 0)$ và đi qua điểm $A(20; 50)$ nên phương trình parabol là $y=\dfrac{1}{8} x^2$.\\
		Xét  tam giác $ABC$, ta có $\cot 60^{\circ}=\dfrac{BC}{AC} \Leftrightarrow BC=\dfrac{50 \sqrt{3}}{3}$ cm.\\
		Suy ra tọa độ điểm $B\left(-\dfrac{50 \sqrt{3}}{3}+20;  0\right)$.\\
		Phương  trình đường thẳng $AB\colon y=\sqrt{3} x+50-20 \sqrt{3}$.\\
		Xét phương trình hoành độ giao điểm của parabol $y=\dfrac{1}{8} x^2$ và đường thẳng $AB$, ta có $$\sqrt{3} x+50-20 \sqrt{3}=\dfrac{1}{8} x^2 \Leftrightarrow\heva{&x=20 \\& x=-20+8 \sqrt{3}.}$$}{\begin{tikzpicture}[scale=0.7,line join=round,line cap=round,>=stealth]
		\tikzset{every node/.style={scale=1}} 
		\def\xmin{-4}\def\xmax{5}\def\ymin{-3}\def\ymax{7} 
		\coordinate [label=below  right:$A$](A) at (2,5);
		\coordinate [label=below right:$C$](C) at (2,0); 
		\pgfmathsetmacro\a{sqrt  (3)}
		\coordinate [label=below:$B$](B)  at (-5*\a/3+2,0);
		 \foreach \diem in {A,C,B} \fill (\diem) circle (1pt);
		\draw[smooth]($(A)! 1.5 ! (B)$)--($(B)!  1.25 ! (A)$);
		\draw[dashed]($(A)! 1.25  ! (C)$)--($(C)! 1.25 ! (A)$)
			($(A)+(-1,0)$)--($(A)+(2,0)$); 
		\draw[<->](0,-1)--(2,-1); 
		\draw (1,-1) node [below]{$20$ cm}; 
		\draw[<->](3,0)--(3,5); 
		\draw (3,2.5) node  [right]{$50$ cm};
		 \draw[->] (\xmin-0.2,0)--(\xmax+0.2,0) node[below]{$x$};
		\draw[->]  (0,\ymin-0.2)--(0,\ymax+0.2) node[right]{$y$};
		\draw (0,0) node[below right]{$O$}; 
		\clip  (\xmin,\ymin) rectangle (\xmax,\ymax);
		\draw[smooth,samples=200,domain=-2+.8*\a:2] plot (\x,{(5/4)*((\x)^2)}); 
		 \draw[xshift=-38,yshift=22,rotate=120,smooth,samples=200,domain=-2+.8*\a:2] plot (\x,{-(5/4)*((\x)^2)});
	\end{tikzpicture}}
		Vậy diện tích của chiếc lá  là
		$S=2 \cdot \displaystyle\int_{-20+8 \sqrt{3}}^{20}\left(\sqrt{3}  x+50-20 \sqrt{3}-\dfrac{1}{8} x^2\right) \mathrm{d} x \approx 745~\left(\mathrm{cm}^2\right)$.
	 }
\end{bt}

% nguồn: Phan2_Lop_12_Dai_so_Chuong_4
\begin{bt}%[2D4V3-2]%[ma:9c32c219]
	\immini[thm]{Kiến trúc sư thiết kế một con đường để chia khu đất hình chữ nhật $ABCD$ với $AB=40$\,m, $AD=60$\,m thành hai phần (hai lề đường là các đường cong $HEK$ và $PQIR$). Trong đó phần giới hạn bởi đường cong $HEK$ và các đoạn $AH$, $AB$, $BK$ là sân chơi, phần giới hạn bởi đường cong $PQIR$ và các đoạn $PD$, $DC$, $CR$ để trồng hoa (tham khảo hình vẽ). Nếu gắn một hệ trục tọa độ vuông góc có trục hoành, trục tung lần lượt cùng phương với các đường thẳng $AB$, $AD$ thì đường cong $PQIR$ là một phần của đồ thị hàm số bậc ba $y=f(x)$. Đường cong $HEK$ nhận được bằng cách tịnh tiến đường cong $PQIR$ theo phương thẳng đứng lên phía trên $10$\, m. Biết $AP=30$\,m, $BR=20$\,m, điểm $Q$ cách các cạnh $AB$, $AD$ lần lượt $25$\,m và $20$\,m, điểm $I$ cách các cạnh  $AB$, $AD$ lần lượt $25$\,m và $35$\, m. Gọi $S_1$ là diện tích phần sân chơi và $S_2$ là diện tích phần trồng hoa. Tính $\dfrac{S_1}{S_2}$ (làm tròn kết quả đến hàng phần trăm).}
	{\begin{tikzpicture}[scale=1,font=\footnotesize,line join=round,line cap=round,>=stealth]
		\def\a{1/7} 
		\def\b{-6/7} 
		\def\c{39/28} 
		\def\d{3} 
		\def\e{4} 
		\def\xmin{-.5}  \def\xmax{5}
		\def\ymin{-.5} \def\ymax{7} 
		\draw (0,0) node[below  left]{$D$};
		\draw[samples=100,smooth,domain=0:4]  plot (\x,{\a*(\x)^3+(\b)*(\x)^2+(\c)*(\x)+(\d)});
		\draw[samples=100,smooth,domain=0:4] plot (\x,{\a*(\x)^3+(\b)*(\x)^2+(\c)*(\x)+(\e)}); 
		\coordinate[label=left:$A$] (A) at (0,6); 
		\coordinate[label=above:$B$] (B) at (4,6); 
		\coordinate[label=below:$C$] (C) at  (4,0);
		\coordinate[label=left:$P$]  (P) at (0,3);
		\coordinate[label=left:$H$] (H)  at (0,4);
		\coordinate[label=below:$Q$] (Q)  at (2,3.5);
		 \coordinate[label=above:$E$] (E) at (2,4.5);
		 \coordinate[label=below:$I$] (I) at (3.5,3.5);
		\coordinate[label=right:$R$]  (R) at (4,4);
		\coordinate[label=right:$K$] (K)  at (4,5);
		\draw  (C)--(B)--(A)--(0,0)--(C);
		\foreach  \diem in {A,B,C,H,P,E,Q,I,K,R}\fill (\diem) circle (1pt);
		\fill[pattern  = north west lines] plot[domain=0:4](\x,{\a*(\x)^3+(\b)*(\x)^2+(\c)*(\x)+(\d)}) -- (K)-- plot[domain=4:0](\x,{\a*(\x)^3+(\b)*(\x)^2+(\c)*(\x)+(\e)}) --(P);
	\end{tikzpicture}}
	\shortans{0,43} 
	 \loigiai{
		Chọn hệ trục tọa độ $Oxy$ sao cho trục $Ox$ trùng với cạnh $AB$, trục $Oy$  trùng với cạnh $AD$.
		\begin{center}
			\begin{tikzpicture}[scale=1,font=\footnotesize,line join=round,line cap=round,>=stealth]
					 \def\a{1/7}
					\def\b{-6/7} 
					 \def\c{39/28}
					\def\d{3} 
					 \def\e{4}
					\def\xmin{-.5}  \def\xmax{6}
					\def\ymin{0} \def\ymax{7} 
					\draw[->]  (\xmin,6)--(\xmax,6) node [below]{$x$};
					 \draw[->] (0,\ymin)--(0,\ymax) node [left]{$y$};
					\draw (0,0) node[below left]{$D$}; 
					 \draw[samples=100,smooth,domain=0:4] plot (\x,{\a*(\x)^3+(\b)*(\x)^2+(\c)*(\x)+(\d)});
					 \draw[samples=100,smooth,domain=0:4] plot (\x,{\a*(\x)^3+(\b)*(\x)^2+(\c)*(\x)+(\e)});
					\coordinate[label=below left:$A$] (A) at (0,6); 
					\coordinate[label=above:$B$] (B)  at (4,6);
					 \coordinate[label=below:$C$] (C) at (4,0);
					\coordinate[label=left:$P$] (P)  at (0,3);
					 \coordinate[label=left:$H$] (H) at (0,4);
					\coordinate[label=below:$Q$] (Q) at  (2,3.5);
					\coordinate[label=above:$E$]  (E) at (2,4.5);
					\coordinate[label=below:$I$]  (I) at (3.5,3.5);
					\coordinate[label=right:$R$] (R)  at (4,4);
					\coordinate[label=right:$K$] (K) at  (4,5);
					\draw (0,0)--(C)--(B)--(A) 
						(0,0)  circle (1pt)
						 (2,5.5) node{$S_1$}
						(2,2)  node{$S_2$};
					 \foreach \diem in {A,B,C,H,P,E,Q,I,K,R}\fill (\diem) circle (1pt);
					\fill[pattern = north west lines]  plot[domain=0:4](\x,{\a*(\x)^3+(\b)*(\x)^2+(\c)*(\x)+(\d)}) -- (K)-- plot[domain=4:0](\x,{\a*(\x)^3+(\b)*(\x)^2+(\c)*(\x)+(\e)}) --(P);
				\end{tikzpicture}
		\end{center}
		Khi đó ta có $A(0; 0)$, $P(0;-30)$,  $Q(20;-25)$, $I(35;-25)$, $R(40;-20)$.\\
		Gọi đường cong $PQIR$ có phương  trình là $y=f(x)=a x^3+b x^2+c x+d$.\\ Từ đó ta có hệ phương trình $\heva{&d=-30\\&8\,000a+400b+20c+d=-25\\
		 &42\,875a+1\,225b+35c+d=-25\\&64\,000a+1\,600b+40c+d=-20}\Leftrightarrow\heva{&a=\dfrac{1}{700}\\&b=-\dfrac{3}{35} \\&c=\dfrac{39}{28} \\&d=-30.}$\\
		Ta suy  ra đường cong $PQIR$ là đồ thị hàm số $y=f(x)=\dfrac{1}{700} x^3-\dfrac{3}{35} x^2+\dfrac{39}{28} x-30$.\\
		Và đường cong $HEK$ là  đồ thị hàm số $y=g(x)=\dfrac{1}{700} x^3-\dfrac{3}{35} x^2+\dfrac{39}{28} x-20$.\\
		 Diện tích phần sân chơi $S_1=\displaystyle\int_0^{40}[-g(x)] \mathrm{\,d}x=\displaystyle\int_0^{40}\left[-\left(\dfrac{1}{700} x^3-\dfrac{3}{35} x^2+\dfrac{39}{28} x-20\right)\right] \mathrm{\,d}x=600$\,m$^2$.\\
		Diện tích phần trồng hoa là  $S_2=\displaystyle\int_0^{40}\left[f(x)+60\right] \mathrm{\,d}x=\displaystyle\int_0^{40}\left[\dfrac{1}{700} x^3-\dfrac{3}{35} x^2+\dfrac{39}{28} x-30+60\right] \mathrm{\,d}x=1\,400$\,m$^2$.\\
		Vậy  tỉ số $\dfrac{S_1}{S_2}=\dfrac{600}{1\,400}=\dfrac{3}{7} \approx 0{,}43$.
	 }
\end{bt}

% nguồn: Phan2_Lop_12_Dai_so_Chuong_4
\begin{bt}%[2D4V3-2]%[ma:c0bdd321]
	 \immini{
	Ông An xây dựng một sân bóng đá mini hình chữ nhật có chiều rộng $30$ m và  chiều dài $50$ m. Để giảm bớt kinh phí cho việc trồng cỏ nhân tạo, ông An chia sân bóng ra làm hai phần (tô màu và không tô màu) như hình vẽ.
	\begin{itemize}
		\item Phần tô màu gồm hai miền diện tích  bằng nhau và đường cong $AIB$ là một parabol có đỉnh $I$.
		\item Phần tô màu được trồng cỏ nhân tạo với giá $130$ nghìn đồng/m$^2$ và phần còn lại được trồng cỏ nhân tạo  với giá $90$ nghìn đồng/m$^2$.
	\end{itemize}
	Hỏi số tiền ông An phải trả để trồng cỏ nhân tạo cho sân bóng là bao~nhiêu triệu đồng? 
	 \shortans{151}
	 }{
	\begin{tikzpicture}[line join=round,line cap=round,>=stealth]
			\def\a{1.5} \def\b{5}  \def\c{1.5}
			\path 
				(180:\a) coordinate  (A)
				 (0:\a) coordinate (B)
				($(A)+(90:\b)$) coordinate  (D)
				($(B)+(D)-(A)$) coordinate (C) 
				($(A)!1/5!(D)$) coordinate  (M)
				($(A)!4/5!(D)$) coordinate (N) 
				($(B)!1/5!(C)$)  coordinate (P)
				($(B)!4/5!(C)$) coordinate  (Q)
				($(M)!1/2!(P)$) coordinate  (I)
				($(N)!1/2!(Q)$) coordinate  (J);
			\draw  (A)--(B)--(C)--(D)--cycle;
			 \draw[dashed] (M)--(P)
				(N)--(Q); 
			\begin{scope}
				\clip (-1.5,0) rectangle (1.5,5); 
				 \draw[samples=100,domain=-1.5:1.5,smooth,variable=\x] plot (\x,{(4/9)*(\x)^2+4});
				\fill[pattern=north west lines,opacity=0.8,draw,thin]  (-1.5,5)--plot[domain=-1.5:1.5](\x,{(4/9)*(\x)^2+4})--(1.5,5)--cycle;
				 \draw[samples=100,domain=-1.5:1.5,smooth,variable=\x] plot (\x,{-(4/9)*(\x)^2+1});
				 \fill[pattern=north west lines,opacity=0.8,draw,thin] (-1.5,0)--plot[domain=-1.5:1.5](\x,{-(4/9)*(\x)^2+1})--(1.5,0)--cycle;
			\end{scope}
			\foreach \x/\y in{A/-150,B/-50,I/90}{\draw (\x) circle (1pt)++(\y:0.35) node{$\x$};} 
			\draw[|<->|,transform canvas={yshift=-9pt}]  (A)--(B) node[midway,below]{$30$ m};
			\draw[|<->|,transform canvas={xshift=-9pt}]  (A)--(D) node[midway,above,rotate=90]{$50$ m};
			\draw[|<->|,transform canvas={xshift=9pt}] (C)--(Q) node[midway,below,rotate=90]{$10$  m};
			\draw[|<->|,transform canvas={yshift=-9pt}] (N)--(J) node[midway,below]{$15$ m}; 
		\end{tikzpicture}
	 }
	 \loigiai{
	 \immini{
		Diện tích sân  bóng là $S=50\cdot 30=1500$ m$^2$.\\
		Đặt hệ trục toạ độ như  hình vẽ.\\
		Phương trình đường parabol đi qua  $A$, $I$, $B$ là $y=-\dfrac{2}{45}x^2+10$.\\
		Diện tích hình phẳng  giới hạn bởi đường cong $AIB$ và đường thẳng $AB$ là
		$$S=\displaystyle\displaystyle\displaystyle\int_{-1{,}5}^{1{,}5}\left(-\dfrac{2}{45}x^2+10\right)\,dx=200\,(m^2).$$ 
		Số tiền trải cỏ nhân tạo cho sân bóng  là
		\begin{eqnarray*}
			T&=&2\cdot 200\cdot 130\,000+(1500-2\cdot 200)\cdot  90\,000\\
			 &=&151\,000\,000\,(\text{đồng})=151\,(\text{triệu đồng}).
		\end{eqnarray*}
		 }
		{\begin{tikzpicture}[line join=round,line cap=round,>=stealth]
		\def\a{1.5} \def\b{5}  \def\c{1.5}
		\path 
			(0:0)  coordinate (O)
			(180:\a) coordinate (A) 
			 (0:\a) coordinate (B)
			($(A)+(90:\b)$) coordinate (D) 
			($(B)+(D)-(A)$)  coordinate (C)
			($(A)!1/5!(D)$) coordinate (M) 
			 ($(A)!4/5!(D)$) coordinate (N)
			($(B)!1/5!(C)$)  coordinate (P)
			($(B)!4/5!(C)$)  coordinate (Q)
			($(M)!1/2!(P)$) coordinate  (I)
			($(N)!1/2!(Q)$)  coordinate (J);
		\draw  (A)--(B)--(C)--(D)--cycle;
		 \draw[dashed] (M)--(P)
			(N)--(Q); 
		\begin{scope}
			\clip (-1.5,0) rectangle (1.5,5); 
			 \draw[samples=100,domain=-1.5:1.5,smooth,variable=\x] plot (\x,{(4/9)*(\x)^2+4});
			\fill[pattern=north west  lines,opacity=0.8,draw,thin] (-1.5,5)--plot[domain=-1.5:1.5](\x,{(4/9)*(\x)^2+4})--(1.5,5)--cycle;
			\draw[samples=100,domain=-1.5:1.5,smooth,variable=\x]  plot (\x,{-(4/9)*(\x)^2+1});
			\fill[pattern=north west lines,opacity=0.8,draw,thin]  (-1.5,0)--plot[domain=-1.5:1.5](\x,{-(4/9)*(\x)^2+1})--(1.5,0)--cycle;
		\end{scope}
		\draw[->]($(B)!1.3!(A)$)--($(A)!1.3!(B)$) node[below]{$x$}; 
		\draw[->](-90:0.5)--(90:5.5) node[left]{$y$}; 
		\foreach \x/\y in{A/150,B/30,I/30,O/-45}{\draw  (\x) circle (1pt)++(\y:0.35) node{$\x$};}
		\path 
			(-1.5,0) node[below]{$-15$} 
			 (1.5,0) node[below]{$15$}
			(0,1) node[above left]{$10$}; 
	\end{tikzpicture}}
	} 
\end{bt}

% nguồn: Phan2_Lop_12_Dai_so_Chuong_4
\begin{bt}%[2D4V3-2]%[ma:cd059bcd]
	 Một cửa vòm có dạng hình parabol được lắp các tấm kính hình tròn đường kính $1$ m và các tấm kính hình vuông có cạnh $1$ m như hình vẽ. Phần còn lại cửa được sơn màu trang trí với mức giá $1{,}2$ triệu đồng/m$^2$. Chi phí sơm màu là bao nhiêu triệu đồng (kết quả làm tròn đến hàng phần chục)?
	\begin{center}
		\begin{tikzpicture}[scale=1,font=\footnotesize,line join=round,line cap=round,>=stealth]
				 \def\r{0.5}
				 \def\s{1}
				\begin{scope}
					 \clip (-3.5,0) rectangle (3.5,4);
					\fill[gray!60]  plot[domain=-3.5:3.5,smooth] (\x,{-0.5*(\x)^2+4});
				\end{scope}
				\node[below left] at (-2.5,0) {\textcolor{blue}{A}}; 
				\node[below right] at (2.5,0) {\textcolor{blue}{B}}; 
				\draw[fill=white] (-2,0.5)  circle (\r);
				\draw[fill=white] (2,0.5) circle (\r); 
				 \draw[fill=white] (0,3.5) circle (\r);
				\draw[fill=white]  (-1,2) rectangle ++(1,1);
				\draw[fill=white] (0,2) rectangle ++(1,1); 
				\foreach \x  in {-2,-1,0,1}{
					\draw[fill=white]  (\x,1) rectangle ++(1,1);
				 }
			\end{tikzpicture}
	\end{center}
	\shortans{8,1} 
	\loigiai{ 
		Chọn hệ trục tọa độ  $Oxy$ như hình vẽ.
		\begin{center}
			\begin{tikzpicture}[scale=1,font=\footnotesize,line join=round,line cap=round,>=stealth]
					\def\r{0.5} 
					\def\s{1} 
					\draw[->]  (-3.5,0)--(3.5,0) node[below]{$x$};
					\draw[->] (0,-0.5)--(0,5)  node[right]{$y$};
					\draw (0,0)  node[below left]{$O$};
					\begin{scope}
						 \clip (-3.5,-0.3) rectangle (3.5,4);
						\draw plot[domain=-3:3,smooth] (\x,{-0.5*(\x)^2+4}); 
					\end{scope}
					 \node[below left] at (-2.5,0) {\textcolor{blue}{A}};
					 \node[below right] at (2.5,0) {\textcolor{blue}{B}};
					\node[below left]  at (-2,2.5) {\textcolor{blue}{C}};
					\node[below right] at (2,2.5)  {\textcolor{blue}{D}};
					\draw  (-2,0.5) circle (\r);
					 \draw (2,0.5) circle (\r);
					\draw (0,3.5)  circle (\r);
					\draw (-1,2) rectangle  ++(1,1);
					 \draw (0,2) rectangle ++(1,1);
					\foreach \x in {-2,-1,0,1}{ 
						 \draw (\x,1) rectangle ++(1,1);
					} 
					\foreach  \x in {-2,-1,0,1}{
						\draw (\x,0)  rectangle ++(1,1);
					} 
					\draw  (0,4) node[above left]{$4$};
					\draw (-2,0)  node[below]{$-2$};
					 \draw (2,0) node[below]{$2$};
				\end{tikzpicture}
		\end{center}
		Gọi phương trình parabol có dạng  $y=ax^2+bx+c$.\\
		Do parabol đi qua các điểm $(0;4)$,  $(2;2)$ và $(-2;2)$ nên
		\begin{eqnarray*}
			 &&\heva{&4a+2b+c=2\\&4a-2b+c=2\\&c=4}\\
			&\Leftrightarrow&  \heva{&a=-\dfrac{1}{2}\\&b=0\\&c=4}\\
			 &\Rightarrow&y=-\dfrac{1}{2}x^2+4.
		\end{eqnarray*}
		Parabol cắt trục hoành tại hai  điểm $A$ và $B$ có hoành độ là nghiệm của phương trình
		\[-\dfrac{1}{2}x^2+4=0\Leftrightarrow  \hoac{&x=2\sqrt{2}\\&x=-2\sqrt{2}.}\]
		Diện tích vòm cửa là  $S=\displaystyle\int ^{2\sqrt{2}}_{-2\sqrt{2}}\left(-\dfrac{1}{2}x^2+4\right)\mathrm{\,d}x=\left(-\dfrac{x^3}{6}+4x\right)\left|^{2\sqrt{2}}_{-2\sqrt{2}}\right.=\dfrac{32\sqrt{2}}{3}$.\\
		Diện tích các tấm  kính bằng $6+3\pi\left(\dfrac{1}{2}\right)^2$.\\
		Suy ra diện tích phần  còn lại là $S=\dfrac{32\sqrt{2}}{3}-6-\dfrac{3\pi}{4}$.\\
		Giá tiền sơn là $\left(\dfrac{32\sqrt{2}}{3}-6-\dfrac{3\pi}{4}\right)\cdot 1{,}2\approx 8{,}1$ (triệu đồng). 
	 }
\end{bt}

% nguồn: Phan2_Lop_12_Dai_so_Chuong_4
\begin{bt}%[2D4V3-2]%[ma:1d8c5d7b]
	\immini{ 
	Kiến trúc sư thiết kế một khu sinh hoạt cộng đồng có dạng hình chữ nhật với chiều rộng và chiều dài lần lượt là $60$ m và $80$ m. Trong đó, phần được tô màu đậm là sân chơi, phần còn lại để trồng hoa. Mỗi phần trồng hoa có đường biên cong là một phần của parabol với đỉnh thuộc một trục đối xứng của hình chữ nhật và khoảng cách từ đỉnh đó đến trung điểm cạnh tương ứng của hình chữ nhật bằng $20$ m (xem hình minh họa). Diện  tích của phần sân chơi là bao nhiêu mét vuông?}
	{ 
	\begin{tikzpicture}[line join=round,line cap=round,xscale=1.5,yscale=0.6,>=stealth,scale=0.7]
			\draw[<->]  (0,2.5)--(0,4.75) node [midway,right] {\small 20m};
			\draw[<->] (0,-2.5)--(0,-4.75) node [midway,right] {\small 20m}; 
			\draw[fill=gray!30]  plot[samples=200,domain=-1.5:1.5,smooth,variable=\x] (\x,{(\x)^2+2.5})--plot[samples=200,domain=1.5:-1.5,smooth,variable=\x] (\x,{-(\x)^2-2.5})--cycle;
			\draw (-1.5,4.75)--(1.5,4.75) 
				 (-1.5,-4.75)--(1.5,-4.75);
		\end{tikzpicture}
	} 
	\shortans{3200} 
	 \loigiai{
		Đặt hệ trục tọa  độ $Oxy$ như hình vẽ.
		\begin{center}
			\begin{tikzpicture}[line join=round,line cap=round,xscale=1.5,yscale=0.6,>=stealth,scale=0.7]
					 \path
						(0,0) coordinate (O') 
						(0,-2.5) coordinate  (I)
						(-1.5,-4.75) coordinate (O) 
						(-1.5,4.75) coordinate (C) 
						(1.5,-4.75)  coordinate (A)
						 ($(O)!(I)!(C)$) coordinate (y);
					\path (y)  node[left]{$20$}
						(C)  node[left]{$80$}
						(0,-4.75) node[below]{$30$} 
						 (A) node[below]{$60$};
					\draw[<->]  (0,2.5)--(0,4.75) node [midway,right] {\small 20m};
					 \draw[<->] (0,-2.5)--(0,-4.75) node [midway,right] {\small 20m};
					\draw[->] (C)--($(C)+(90:2)$) node[left]{$y$}; 
					 \draw[->] (A)--($(A)+(0:1)$) node[below]{$x$};
					 \draw[fill=gray!30] plot[samples=200,domain=-1.5:1.5,smooth,variable=\x] (\x,{(\x)^2+2.5})--plot[samples=200,domain=1.5:-1.5,smooth,variable=\x] (\x,{-(\x)^2-2.5})--cycle;
					 \draw[dashed] (I)--(y);
					\draw  (-1.5,4.75)--(1.5,4.75)
						(-1.5,-4.75)--(1.5,-4.75); 
					 \foreach \t/\g in {O/200,I/90,A/40}{
						 \draw[fill=white] (\t) circle (1pt) node[shift={(\g:7pt)},font=\scriptsize]{$ \t $};
					} 
				\end{tikzpicture}
		\end{center}
		Nhận thấy hai phần  trồng hoa dạng đường parabol có diện tích như nhau. Vậy ta xét đường parabol ở bồn hoa phía dưới có dạng $(P)\colon y=a x^2+b x+c$ với $a \neq 0$.
		Ta có đồ thị $(P)$ đi qua ba điểm là $O(0; 0)$, $I(30; 20)$ và $A(60; 0)$. Thay  vào ta thu được hệ phương trình
		 \[
		 \heva{&c=0\\&900a+30b+c=20\\&3\,600a+60b+c=0} \Leftrightarrow \heva{&a=-\dfrac{1}{45} \\&b=\dfrac{4}{3} \\&c=0.}\]
		Như vậy, ta có phương trình parabol đó  là $y=-\dfrac{1}{45} x^2+\dfrac{4}{3} x$ và diện tích của phần trồng hoa phía dưới là
		 $$
		S_{\text{ hoa}}=\displaystyle\int_0^{60} \left\lvert-\dfrac{1}{45} x^2+\dfrac{4}{3}x\right\vert \mathrm{d}  x=\displaystyle\int_0^{60}\left(-\dfrac{1}{45} x^2+\dfrac{4}{3} x\right)\mathrm{d}x=800~(\mathrm{m}^2).
		$$ 
		 Khi đó, diện tích của phần sân chơi là $S=60\cdot 80-2S_{\text{ hoa}}=3\,200$ (m$^2$).
	 }
\end{bt}

% nguồn: Phan2_Lop_12_Dai_so_Chuong_4
\begin{bt}%[2D4V3-2]%[ma:d7edada3]
	\immini{ 
	Cô Hà đổ bê tông  một đường đi trong vườn (phần được tô màu) với kích thước được cho trong hình vẽ bên. Biết rằng đường cong $AB$ được cho bởi đồ thị của một hàm số liên tục $y=f(x)$ và đường cong $DC$ có phương trình $y=f(x)+2$. Cô Hà đổ lớp bê tông dày $15$ cm để làm con đường thì thể tích khối bê tông cần đổ là bao nhiêu mét khối?
	} 
	 {
	\begin{tikzpicture}[line join = round,line cap = round,>=stealth,font=\footnotesize,scale=0.5]
			\fill[cyan!50,smooth] plot[domain=0:10] (\x,{0.008*(\x)^3-0.12*(\x)^2+4})-- (10,2)  -- plot[domain=10:0] (\x,{0.008*(\x)^3-0.12*(\x)^2+6}) -- (0,4);
			\draw (0,0) rectangle (10,6); 
			 \path
				(0,4) coordinate  (A)
				(0,6)  coordinate (D)
				 (10,2) coordinate (C)
				 (10,0) coordinate (B);
			\foreach \x/\g in  {A/180,B/0,C/0,D/180} \fill(\x) circle (1pt)+(\g:3mm) node{$\x$};
			\draw[<->] (0,-0.5) -- (10,-0.5); 
			 \draw (5,-0.5) node[below] {$10$ m};
			 \draw[<->] (-1,0) -- (-1,6);
			\draw  (-1,3) node[left] {$6$ m};
			 \draw[<->] (11,0) -- (11,2);
			\draw (11,1)  node[right] {$2$ m};
		\end{tikzpicture}
	 }
	 \shortans{3}
	\loigiai{ 
	 \immini
	 {
	Xét  hệ trục toạ độ $Oxy$ như hình vẽ bên (đơn vị trên mỗi trục toạ độ là mét).\\
	Trong hệ trục toạ độ đó, đường cong $AB$  được cho bởi đồ thị của hàm số liên tục $y=f(x)$, $0\leq x\leq 10$. Suy ra, đường cong $DC$ được cho bởi đồ thị của hàm số liên tục $y=f(x)+2$, $0\leq x\leq 10$.\\
	Áp dụng công thức tính diện  tích ta có
	} 
	 {
	\begin{tikzpicture}[line join = round,line cap = round,>=stealth,font=\footnotesize,scale=0.5]
			 \fill[cyan!50,smooth] plot[domain=0:10] (\x,{0.008*(\x)^3-0.12*(\x)^2+4})-- (10,2) -- plot[domain=10:0] (\x,{0.008*(\x)^3-0.12*(\x)^2+6}) -- (0,4);
			\draw (0,0) rectangle (10,6); 
			\draw[->] (0,0) -- (12,0) node[below]  {$x$};
			 \draw[->] (0,0) -- (0,7) node[left] {$y$};
			 \draw (0,0) node[left] {$O$};
			 \path
				 (0,4) coordinate (A)
				(0,6) coordinate  (D)
				(10,2) coordinate  (C)
				(10,0)  coordinate (B);
			\draw (A)--(B) node[midway,sloped,below]{$y=f  (x)$};
			\draw (D)--(C)  node[midway,sloped,above]{$y=f (x)+2$};
			\foreach \x/\g in {A/180,B/45,C/0,D/180} \fill(\x)  circle (1pt)+(\g:5mm) node{$\x$};
			\draw (A) node[below right] {$4$}; 
			 \draw (D) node[above right] {$6$};
			\draw  (B) node[below] {$10$};
		\end{tikzpicture}
	 }
	Gọi  $S$ là diện tích hình phần được giới hạn bởi đồ thị hàm số $y=f(x)$, $y=f(x)+2$, đường thẳng $x=0$ và $x=10$.
	\begin{align*}
		S=\displaystyle\int_{0}^{10} \left[f(x)+2-f(x)\right]  \mathrm{\,d}x=\displaystyle\int_{0}^{10} 2 \mathrm{\,d}x=2x\Big|_0^{10}=20\text{ \,m}^2.
	\end{align*}
	Do đó, thể tích  khối bê tông dùng để đổ con đường là $20\cdot 0{,}15=3$\,m$^3$.
	} 
\end{bt}

% nguồn: Phan2_Lop_12_Dai_so_Chuong_4
\begin{bt}%[2D4V3-2]%[ma:2cacabaf]
	Kiến trúc sư thiết kế một khu vui chơi có dạng hình chữ nhật với chiều rộng $30$ m và chiều  dài $50$ m. Trong đó, phần tô màu đậm trồng cỏ nhung Nhật, phần còn lại trải thảm cỏ nhân tạo. Mỗi phần trồng cỏ nhung Nhật có đường biên cong là một phần của parabol với đỉnh $I$ thuộc trục đối xứng của hình chữ nhật và khoảng cách từ đỉnh $I$ đến trung điểm của cạnh $AB$ bằng $10$ m (xem hình minh họa). Biết giá cỏ nhung Nhật là $140\,000$ đồng/m$^2$, giá cỏ nhân tạo là $100\,000$ đồng/m$^2$. Tính tổng số tiền để trồng và trải thảm cỏ hết khu vui chơi theo đơn vị triệu đồng.\\
	 \shortans{166}
	\begin{center}
		\begin{tikzpicture}[>=stealth,line join=round,line cap=round,scale=0.15]
				 \coordinate (A) at (0,30);
				\coordinate (B) at (0,0); 
				\coordinate (C) at  (50,0);
				\coordinate  (D) at (50,30);
				 \draw (A)--(B)--(C)--(D)--cycle;
				 \draw[dashed] (10,0)--(10,30);
				\draw[dashed]  (40,0)--(40,30);
				\fill[gray!60,domain=0:30,variable=\y] 
					plot ({10 - (2/45)*(\y - 15)^2},{\y}) -- (0,30) -- (0,0) --  cycle;
				 \fill[gray!60,domain=0:30,variable=\y]
					plot  ({40 + (2/45)*(\y - 15)^2},{\y}) -- (50,30) -- (50,0) -- cycle;
				\foreach \p/\pos in {A/above left,B/below  left,C/below right,D/above right}{
					\fill (\p)  circle (5pt);
					\node[\pos] at (\p)  {\p};
				 }
				\fill (10,15)  circle (5pt);
				 \node[left] at (10,15) {$I$};
				\fill  (40,15) circle (5pt);
				 \node[right] at (40,15) {$J$};
				\draw[<->] (-2,0) -- (-2,30) node[midway,left] {$30\,\text{  m}$};
				\draw[<->] (0,-4) --  (50,-4) node[midway,below] {$50\,\text{ m}$};
				\draw[<->] (0,33)  -- (10,33) node[midway,above] {$10\,\text{ m}$};
			\end{tikzpicture}
	\end{center}
	\loigiai{ 
		Xét hệ trục tọa độ  $Oxy$ như hình vẽ.
		\begin{center}
			\begin{tikzpicture}[>=stealth,line join=round,line cap=round,scale=1.2]
					 \fill[gray!60,domain=-2:2,variable=\x]
						(-2,1) -- plot ({\x},{0.25*\x*\x}) -- (2,1) --  cycle;
					\draw[domain=-2:2,variable=\x] plot  ({\x},{0.25*\x*\x});
					\draw  (-2,1)--(2,1);
					\draw[dashed] (-2,1)--(-2,0); 
					\draw[dashed] (2,1)--(2,0); 
					\fill (0,0) circle (1.2pt) node[below left] {$O$}; 
					\fill (-2,1) circle (1.2pt)  node[left] {$D$};
					\fill (2,1) circle (1.2pt) node[right] {$C$}; 
					\draw[->] (-2.8,0) -- (2.8,0)  node[below right] {$x$};
					 \draw[->] (0,-0.3) -- (0,2) node[left] {$y$};
				\end{tikzpicture}
		\end{center}
		Gọi phương trình parabol $(P)$ có dạng $y = ax^2 +  bx + c$.\\
		Khi đó $(P)$ đi qua gốc $O$, điểm $C(15;10)$  và $D(-15;10)$.\\
		Ta có hệ phương  trình:
		 $$
		 \heva{
		&c =  0\\
		&225a + 15b  + c = 10\\
		 &225a - 15b + c = 10
		} 
		\Rightarrow 
		\heva{ 
		&a  = \dfrac{2}{45}\\
		 &b = 0\\
		&c  = 0.
		} 
		$$ 
		Vậy $(P)\colon y =  \dfrac{2}{45}x^2$.\\
		Diện tích hình chữ nhật là $S  = 50.30 = 1\,500\, (\text{m}^2)$.\\
		Diện tích trồng cỏ nhung Nhật  là
		$$ 
		S_1 = 2\displaystyle\int_{-15}^{15} \left(10 - \dfrac{2}{45}x^2\right)\mathrm{d}x  = 400 \,(\text{m}^2).$$
		Diện tích trải  thảm nhân tạo là
		$$ 
		S_2 = S - S_1 = 1\,500 - 400  = 1100 \,(\text{m}^2).$$
		Tổng số tiền để trồng và trải thảm cỏ  hết khu vui chơi là
		$$400\cdot 140\,000 + 1100\cdot 100\,000  = 166\,000\,000 \, (\text{đồng}) =166 \,(\text{triệu đồng}).$$
	} 
\end{bt}

% nguồn: Phan2_Lop_12_Dai_so_Chuong_4
\begin{bt}%[2D4V3-2]%[ma:2df4b06b]
	\immini{Để trang trí một ô hình chữ nhật với kích thước $80 \mathrm{~cm} \times 60 \mathrm{~cm}$ trên tường căn phòng của mình, anh Bình đã tô màu $2$ hoạ tiết giới hạn bởi các parabol $(P)$, parabol $(C)$ và các cạnh của ô chữ nhật bằng màu vàng (kí hiệu sọc nghiêng) và màu xanh (vùng tô đậm) (như hình vẽ). Biết rằng các parabol $(P)$ và $(C)$ đối xứng nhau qua đường thẳng $AB$ và diện tích phần màu vàng bằng $\dfrac{1600}{3} \mathrm{~cm}^2$. Tính tổng số tiền (\textit{đơn vị:} triệu đồng) mà anh Bình cần phải trả để mua sơn trang trí. Biết mỗi lọ sơn, sơn được tối đa $30 \mathrm{~cm}^2$. Cho biết giá của một lọ sơn vàng là $50\,000$ đồng, giá của một lọ sơn xanh là $30\,000$ đồng (\textit{biết rằng mỗi lọ sơn không được bán lẻ}).}{\begin{tikzpicture}[scale=0.7,line join=round,line cap=round,>=stealth]
		\tikzset{every  node/.style={scale=1}}
		 \def\xmin{-1}\def\xmax{9}\def\ymin{-1}\def\ymax{8}
		\fill (0,0) circle  (1pt);
		\coordinate [label=below right:$B$](B) at  (1,0);
		\coordinate  [label=above left:$A$](A) at (7,6);
		\coordinate  [label=left:$C$](C) at (0,6);
		\draw[smooth,samples=200,domain=0:4] plot (\x,{.5*((\x)^2)-2*\x+6}); 
		\fill[pattern=north  east lines]plot[domain=0:4](\x,{.5*((\x)^2)-2*\x+6})--plot[domain=4:0](\x,{6})--cycle;
		\foreach \diem in  {A,B,C} \fill (\diem) circle (1pt);
		\draw[smooth](C)--(0,0)--(8,0)--(8,6)--(C) 
			 (A)--(B);
		 \draw[<->](0,6.5)--(4,6.5);
		\draw (2,6.5) node[above]{$40$  cm};
		\draw[<->](7,6.5)--(8,6.5); 
		 \draw (7.5,6.5) node[above]{$10$ cm};
		\draw[<->](0,-.5)--(1,-.5); 
		\draw (.5,-.5) node[below]{$10$  cm};
		 \draw (4,4.5) node[above]{$(P)$};
		\draw (5.5,2.25) node[above]{$(C)$}; 
		\draw [samples=50,rotate  around={-90.:(5.,1.)},xshift=5.cm,yshift=1.cm,
			domain=-2.45:1)] plot (\x,{(\x)^2/2/1.0}); 
		\fill[color=gray!150!black,rotate around={-90.:(5.,1.)},xshift=5.cm,yshift=1.cm, 
			 domain=-2.45:1)]plot[domain=-2.5:1](\x,{(\x)^2/2/1.0})--plot[domain=1:-2.5](\x,{3})--cycle;
	\end{tikzpicture}}
	\shortans{667} 
	 \loigiai{
		\begin{center}
			\begin{tikzpicture}[scale=0.8,line join=round,line cap=round,>=stealth]
					 \tikzset{every node/.style={scale=1}}
					 \def\xmin{-1}\def\xmax{9}\def\ymin{-1}\def\ymax{8}
					\draw[->]  (0,0)--(\xmax+0.2,0) node[below]{$x$};
					 \draw[->] (0,0)--(0,\ymax+0.2) node[right]{$y$};
					\draw (0,0) node[below left]{$O$}; \fill (0,0) circle (1pt); 
					 \foreach \x in {8}\draw (\x,0.1)--(\x,-0.1) node[below]{$80$};
					\foreach \y in {5}\draw (0.1,\y)--(-0.1,\y) node[left]{$10+10\sqrt{16}$}; 
					\coordinate [label=below right:$B$](B)  at (1,0);
					\coordinate  [label=above left:$A$](A) at (7,6);
					 \coordinate [label=left:$C$](C) at (0,6);
					\coordinate [label=below  left:$I$](I) at (2,4);
					\coordinate  [label=left:$I'$](I') at (5,1);
					\coordinate [label=right:$H$](H)  at (3.5,2.5);
					\coordinate  [label=below:$C'$](C') at (7,-1);
					 \coordinate (J) at ($(I)+(0,2)$);
					\draw[smooth,samples=200,domain=0:4] plot (\x,{.5*((\x)^2)-2*\x+6}); 
					\fill[pattern=north east  lines]plot[domain=0:4](\x,{.5*((\x)^2)-2*\x+6})--plot[domain=4:0](\x,{6})--cycle;
					\foreach \diem in {A,B,C,I,I',H,C',J} \fill (\diem)  circle (1pt);
					\draw[dashed](0,5)--(8,5); 
					\draw[smooth](C)--(8,6)--(8,0) 
						 (A)--(B)
						(I)--(I') 
						(I)--(J); 
					 \draw[<->](0,6.5)--(4,6.5);
					\draw ($(I)!.5 !  (J)$) node[left]{$20$};
					\draw (2,6.5)  node[above]{$40$ cm};
					\draw[<->](7,6.5)--(8,6.5); 
					\draw  (7.5,6.5) node[above]{$10$ cm};
					\draw[<->](0,-.5)--(1,-.5); 
					\draw  (4,4.5) node[above]{$(P)$};
					\draw (5.5,2.25)  node[above]{$(C)$};
					\draw  (.5,-.5) node[below]{$10$ cm};
					\draw [samples=50,rotate  around={-90.:(5.,1.)},xshift=5.cm,yshift=1.cm,
						domain=-2.45:1)]  plot (\x,{(\x)^2/2/1.0});
					\fill[color=gray!150!black,rotate  around={-90.:(5.,1.)},xshift=5.cm,yshift=1.cm,
						 domain=-2.45:1)]plot[domain=-2.5:1](\x,{(\x)^2/2/1.0})--plot[domain=1:-2.5](\x,{3})--cycle;
				\end{tikzpicture}
		\end{center}
		Chọn hệ trục $Oxy$ như hình vẽ. 
		$AB\colon  y=a x+b$ đi qua $A(70; 60)$, $B(10; 0)$.\\
		Suy ra $\heva{&70 a+b=60  \\ &10 a+b=0} \Rightarrow\heva{&a=1 \\& b=-10}\Rightarrow AB\colon y=x-10$.\\ Ta có
		$S_{\text {vàng }}=\dfrac{2}{3} \cdot D \cdot h$ trong đó $D$ là khoảng cách  giữa hai điểm của parabol nằm trên $O x$ và $h$ là khoảng cách từ đỉnh xuống trục $O x$.\\
		Suy ra $h=\dfrac{3  S}{2 D}=\dfrac{3 \cdot \dfrac{1600}{3}}{2 \cdot 40}=20$.\\
		Gọi $I$ là đỉnh của parabol $(P)\colon y=a x^2+b x+c$ đi qua $C(0; 60)$, có đỉnh  $I(20; 40)$.\\
		Suy ra $\heva{&c=60 \\& 400 a+20 b+c=40 \\& -\dfrac{b}{2 a}=20} \Rightarrow\heva{&a=\dfrac{1}{20} \\& b=-2 \\& c=60}\Rightarrow(P)\colon y=\dfrac{1}{20} x^2-2 x+60$.\\  Mặt khác
		$I'$ là  đỉnh của $(C)$
		$I I'$ đi  qua $I(20; 40)$ và $\perp AB$.
		\\Suy ra phương trình $I I'\colon  1(x-20)+1(y-40)=0 \Leftrightarrow x+y-60=0$.\\
		 Gọi $H=AB \cap I I'$ thì tọa độ $H$ thỏa mãn hệ phương trình
		$$\heva{&x-y-10=0 \\& x+y-60=0} \Leftrightarrow\heva{&x=35 \\&  y=25} \Rightarrow H(35; 25).$$
		$H$ là trung điểm của $I I'$ suy ra $\heva{&x_{I'}=2 x_H-x_I=50 \\& y_{I'}=2  y_H-y_I=10} \Rightarrow I'(50; 10)$.\\
		$C'$ là điểm đối xứng của  $C$ qua $H$ nên $C'(70;-10)$.\\
		Ta có $(C)\colon x=m y^2+n  y+p$ đi qua $C'(70;-10)$ và có đỉnh $I'(50; 10)$ suy ra
		$$\heva{&100 m-10  n+p=70 \\& 20 m+n=0 \\& 100 m+10 n+p=50} \Rightarrow\heva{&m=\dfrac{1}{20} \\& n=-1 \\& p=55}\Rightarrow (C)\colon x=\dfrac{1}{20} y^2-y+55.$$
		Phương trình tung độ giao điểm của đường thẳng $x=80$ và $(C)$ là 
		 \allowdisplaybreaks
		\begin{eqnarray*}
			80&=&\dfrac{1}{20} y^2-y+55 \\\Rightarrow  y&=&10+10 \sqrt{6}\\
			S_{(C)}&=&\displaystyle\int_{0}^{10+10 \sqrt{6}}\left(80-\left(\dfrac{1}{20} y^2-y+55\right)\right) \mathrm{d} y \approx 773\left(\mathrm{~cm}^2\right). 
		\end{eqnarray*}
		Số hộp sơn vàng cần là $\dfrac{773}{30} \approx 26$ (hộp sơn), số hộp sơn xanh cần là  $\dfrac{\dfrac{1600}{3}}{30} \approx 18$ (hộp sơn).\\
		Số tiền cần là 
		$26 \cdot  30000+18 \cdot 50000=1680000$ (đồng) $=1{,}68$ (triệu đồng).
	 }
\end{bt}

% nguồn: Phan2_Lop_12_Dai_so_Chuong_4
\begin{bt}%[2D4V3-2]%[ma:4c66c9ad]
	Một viên gạch hình vuông có cạnh bằng $40$  cm. Người ta dùng bốn đường parabol có chung đỉnh tại tâm của viên gạch để tạo ra bốn cánh hoa (được tô màu như hình vẽ).
	\begin{center}
		\begin{tikzpicture}[line join=round,line cap=round,>=stealth,font=\footnotesize,scale=1]
				\def\a{3.8} 
				 \path (0:0) coordinate (O)
					 ++(45:\a*0.7) coordinate (A)
					++(-90:\a) coordinate  (B)
					++(180:\a) coordinate  (C)
					 ++(90:\a) coordinate (D);
				 \draw (A)--(B)--(C)--(D)--(A);
				\begin{scope}
					\clip  plot[domain=0:\a/2](\x,{2/(\a)*(\x)^2})--plot[domain=(\a)/2:0]({2/(\a)*(\x)^2},\x)--(0,0);
					\fill[blue!20] plot[domain=(\a)/2:0]({2/(\a)*(\x)^2},\x)--plot[domain=0:\a/2](\x,{2/(\a)*(\x)^2}); 
				\end{scope}
				\begin{scope}
					\clip plot[domain=0:-\a/2](\x,{2/(\a)*(\x)^2})--plot[domain=(\a)/2:0]({-2/(\a)*(\x)^2},\x)--(0,0); 
					\fill[blue!20]  plot[domain=(\a)/2:0]({-2/(\a)*(\x)^2},\x)--plot[domain=0:-\a/2](\x,{2/(\a)*(\x)^2});
				\end{scope}
				\begin{scope}
					\clip  plot[domain=-\a/2:0](\x,{-2/(\a)*(\x)^2})--plot[domain=0:-\a/2]({-2/(\a)*(\x)^2},\x)--(0,0);
					 \fill[blue!20] plot[domain=-\a/2:0](\x,{-2/(\a)*(\x)^2})--plot[domain=0:-\a/2]({-2/(\a)*(\x)^2},\x)--(0,0);
				\end{scope}
				\begin{scope}
					 \clip plot[domain=\a/2:0](\x,{-2/(\a)*(\x)^2})--plot[domain=0:-\a/2]({2/(\a)*(\x)^2},\x)--(0,0);
					\fill[blue!20]  plot[domain=\a/2:0](\x,{-2/(\a)*(\x)^2})--plot[domain=0:-\a/2]({2/(\a)*(\x)^2},\x)--(0,0);
				\end{scope}
				 \draw[smooth,blue] plot[domain=-(\a)/2:(\a)/2](\x,{2/(\a)*(\x)^2});
				 \draw[smooth,blue] plot[domain=-(\a)/2:(\a)/2](\x,{-2/(\a)*(\x)^2});
				\draw[smooth,blue] plot[domain=-(\a)/2:(\a)/2]({2/(\a)*(\x)^2},\x); 
				\draw[smooth,blue]  plot[domain=-(\a)/2:(\a)/2]({-2/(\a)*(\x)^2},\x);
			\end{tikzpicture}
	\end{center}
	Hỏi diện  tích phần cánh hoa được tô màu của viên gạch bằng bao nhiêu centimét vuông? (kết quả làm
	tròn đến hàng đơn  vị)
	 \shortans{533}
	 \loigiai{
		Chọn hệ  trục tọa độ $Oxy$ với $O$ trùng với tâm của viên gạch như hình vẽ.
		\begin{center}
			\begin{tikzpicture}[line join=round,line cap=round,>=stealth,font=\footnotesize,scale=1]
					\def\a{3.8} 
					\path (0:0) coordinate  (O)
						++(45:\a*0.7)  coordinate (A)
						++(-90:\a) coordinate  (B)
						++(180:\a) coordinate  (C)
						++(90:\a) coordinate (D); 
					\draw  (A)--(B)--(C)--(D)--(A);
					\begin{scope}
						\clip plot[domain=0:\a/2](\x,{2/(\a)*(\x)^2})--plot[domain=(\a)/2:0]({2/(\a)*(\x)^2},\x)--(0,0); 
						 \fill[blue!20] plot[domain=(\a)/2:0]({2/(\a)*(\x)^2},\x)--plot[domain=0:\a/2](\x,{2/(\a)*(\x)^2});
					\end{scope}
					\begin{scope}
						 \clip plot[domain=0:-\a/2](\x,{2/(\a)*(\x)^2})--plot[domain=(\a)/2:0]({-2/(\a)*(\x)^2},\x)--(0,0);
						\fill[blue!20]  plot[domain=(\a)/2:0]({-2/(\a)*(\x)^2},\x)--plot[domain=0:-\a/2](\x,{2/(\a)*(\x)^2});
					\end{scope}
					\begin{scope}
						 \clip plot[domain=-\a/2:0](\x,{-2/(\a)*(\x)^2})--plot[domain=0:-\a/2]({-2/(\a)*(\x)^2},\x)--(0,0);
						 \fill[blue!20] plot[domain=-\a/2:0](\x,{-2/(\a)*(\x)^2})--plot[domain=0:-\a/2]({-2/(\a)*(\x)^2},\x)--(0,0);
					\end{scope}
					\begin{scope}
						 \clip plot[domain=\a/2:0](\x,{-2/(\a)*(\x)^2})--plot[domain=0:-\a/2]({2/(\a)*(\x)^2},\x)--(0,0);
						 \fill[blue!20] plot[domain=\a/2:0](\x,{-2/(\a)*(\x)^2})--plot[domain=0:-\a/2]({2/(\a)*(\x)^2},\x)--(0,0);
					\end{scope}
					\draw[smooth,red] plot[domain=-(\a)/2:(\a)/2](\x,{2/(\a)*(\x)^2}) node  [above right]{$(P_1)$};
					\draw[smooth,blue] plot[domain=-(\a)/2:(\a)/2](\x,{-2/(\a)*(\x)^2}); 
					\draw[smooth,violet]  plot[domain=(\a)/2:-(\a)/2]({2/(\a)*(\x)^2},\x) node [below right]{$(P_2)$};;
					\draw[smooth,blue]  plot[domain=-(\a)/2:(\a)/2]({-2/(\a)*(\x)^2},\x);
					\draw (O) node[below left] {$O$}; 
					\draw[->] (-\a/2-0.8,0)--(\a/2+0.8,0) node [below] {$x$}; 
					\draw[->] (0,-\a/2-0.8)--(0,\a/2+0.8) node [right] {$y$}; 
					\draw  (\a/2,0) node [above right]{$20$};
					\draw (-\a/2,0) node  [above left]{$-20$};
					 \draw (0,\a/2) node [above right]{$20$};
					\draw (0,-\a/2) node  [below right]{$-20$};
				\end{tikzpicture}
		\end{center}
		Gọi $(P_1) \colon y=ax^2+bx+c$ là parabol đi qua điểm có tọa  độ lần lượt là $(0;0)$; $(20;20)$; $(-20;20)$.\\
		Do đó ta có hệ phương  trình
		$$\heva{&0^2a+0b+c=0\\&20^2 a+20b+c=20\\&(-20)^2a-20b+c=20} \Leftrightarrow \heva{&c=0\\&400a+20b+c=20\\&400a-20b+c=20}  \Leftrightarrow \heva{&a=\dfrac{1}{20}\\&b=0\\&c=0.}$$
		Vậy  $(P_1) \colon y=\dfrac{1}{20}x^2$.\\
		Gọi $(P_2) \colon x=dy^2+ey+f$ là parabol đi qua các điểm  có tọa độ lần lượt là $(0;0)$; $(20;20)$; $(20;-20)$.\\
		Do đó ta có hệ phương trình 
		$$\heva{&0^2d+0e+f=0\\&20^2  d+20e+f=20\\&(-20)^2d-20e+f=20} \Leftrightarrow \heva{&d=0\\&400d+20e+f=20\\&400d-20e+f=20} \Leftrightarrow \heva{&d=\dfrac{1}{20}\\&e=0\\&f=0.}$$
		Vậy $(P_2) \colon x=\dfrac{1}{20}y^2 \Rightarrow  \hoac{&y=\sqrt{20x}\\&y=-\sqrt{20x}.}$ \\
		Diện tích một  cánh hoa là phần hình phẳng giới hạn bởi đồ thị $(P_1)$, đồ thị của hàm số $y=\sqrt{20x}$ và hai đường thẳng $x=0$, $x=20$ là
		$$S=\displaystyle \int_{0}^{20} \left| \dfrac{1}{20}x^2-\sqrt{20x}\right|=\dfrac{400}{3} \; (\text{cm}^2).$$ 
		Diện tích phần cánh hoa được tô màu của viên gạch là $4S=4 \cdot \dfrac{400}{3}=\dfrac{1600}{2} \approx  533$ (cm$^2$).
	} 
\end{bt}

% nguồn: Phan2_Lop_12_Dai_so_Chuong_4
\begin{bt}%[2D4V3-2]%[ma:a1fff78f]
	Người ta cần làm một cái cửa nhà hình parabol có chiều cao từ mặt đất đến đỉnh là $2{,}25$ mét, chiều rộng tiếp giáp  với mặt đất là $3$ mét. Giá thi công trọn gói mỗi mét vuông là $1\,500\,000$ đồng. Vậy số tiền làm cửa phải trả là bao nhiêu (đơn vị: triệu đồng)?
	\shortans{6,75} 
	 \loigiai{
		Gắn parabol $(P)\colon y=a{x^2}+bx+c$ vào hệ trục tọa độ sao cho $(P)$ đi qua $O\left(0; 0\right)$, $A\left(3; 0\right)$ và $B\left(1{,}5;  2{,}25\right)$.\\
		Ta  được
		\allowdisplaybreaks 
		\begin{eqnarray*}
			\heva{&c=0\\&9a+3b+c=0\\&\dfrac{9}{4}a+\dfrac{3}{2}b+c=\dfrac{9}{4}}\Leftrightarrow\heva{&a=-1\\&b=3\\&c=0.} 
		\end{eqnarray*}
		Suy ra $(P)\colon y=-x^2+3x$ nên diện tích cửa khi đó là:  \[S=\displaystyle\int_0^3\left|-x^2+3x\right|\,\mathrm{d}x=\dfrac{9}{2}.\]
		Vậy số tiền làm  cửa là $\dfrac{9}{2}\cdot1{,}5=6{,}75$ (triệu đồng).
	} 
\end{bt}

% nguồn: Phan2_Lop_12_Dai_so_Chuong_4
\begin{bt}%[2D4V3-2]%[ma:ab12e351]
	 \immini[thm]{Một công ty có ý định thiết kế một logo hình vuông có độ dài nửa đường chéo bằng $4$. Biểu tượng $4$ chiếc lá (được tô màu) được tạo thành bởi các đường cong đối xứng với nhau qua tâm của hình vuông và qua các đường chéo. Một trong số các đường cong ở nửa bên phải của logo là một phần của đồ thị hàm số bậc ba dạng $y = ax^3 + bx^2 - x$ với hệ số $a<0$. Để kỷ niệm ngày thành lập $2/3$, công ty thiết kế để tỉ số diện tích được tô màu so với phần không được tô màu bằng $\dfrac{2}{3}$. Tính $2a + 2b$.
	}{\begin{tikzpicture}[line join=round,line cap=round,>=stealth]
		\foreach \i  in {0,1,2,3}{
			 \fill[cyan,smooth,rotate=90*\i] plot [domain=0:3](\x,{-(\x)*(\x-3)^2/8})--plot [domain=3:0](\x,{(\x)*(\x-3)^2/8})--cycle;
			 \draw (0,0)--(90*\i:3)--(90*\i+90:3);
		 }
	\end{tikzpicture}
	 }
	\shortans{0,8} 
	 \loigiai{
		 Ta có nửa đường chéo hình vuông có độ dài là $4$, cạnh hình vuông sẽ là $4\sqrt{2}$ và diện tích hình vuông là $32$, khi đó ta có được diện tích phần tô màu là $\dfrac{64}{5}$.\\
		Gọi $f(x) = ax^3 +  bx^2 - x$ là hàm số bậc ba biểu diễn đường cong trên logo.\\
		Ta có $x = 4$ là nghiệm của phương trình nên $64a + 16b -  4 = 0 \Leftrightarrow 4a+b=1$.\quad (1)\\
		Do đó phương trình $f(x) = 0$ sẽ có các nghiệm là $0$,  $4$.\\
		Khi đó diện tích hình phẳng $S$  giới hạn bởi các đường $y = f(x)$, trục $Ox$, đường thẳng $x = 0$, $x = 4$ là
		\begin{eqnarray*}
			S&= & \displaystyle \int_{0}^{4}  |ax^3+bx^2-x|\mathrm{\,d}x = -\displaystyle \int_{0}^{4} (ax^3+bx^2-x)\mathrm{\,d}x\\
			&= & -\left(\dfrac{ax^4}{4}+\dfrac{bx^3}{3}-\dfrac{x^2}{2}\right)\Bigg|_0^4 = -64a-\dfrac{64}{3}b+8. 
		\end{eqnarray*}
		Mà $S = \dfrac{1}{8}\cdot \dfrac{64}{5} = \dfrac{8}{5}$ nên $-64a-\dfrac{64}{3}b+8  = \dfrac{8}{5} \Leftrightarrow 64a+\dfrac{64}{3}b=\dfrac{32}{5}$.\quad (2)\\
		Từ $(1)$ và $(2)$,  ta có $\heva{&4a+b=1\\&64a+\dfrac{64}{3}b=\dfrac{32}{5}} \Leftrightarrow \heva{&a = -\dfrac{1}{20}\\&b=\dfrac{9}{20}} \Rightarrow 2a+2b = \dfrac{4}{5}$.
	 }
\end{bt}

% nguồn: Phan2_Lop_12_Dai_so_Chuong_4
\begin{bt}%[2D4V3-2]%[ma:b72c00da]
	 Người ta lát gạch trang trí một mảnh sân hình chữ nhật có kích thước $14\mathrm{~m} \times 12\mathrm{~m}$ như hình vẽ bên dưới, trong đó $(P_1)$, $(P_2)$ là hai parabol đối xứng trục với nhau qua trục đối xứng vuông góc với chiều dài của mảnh sân, $(C)$ là đường tròn có tâm trùng với tâm của mảnh sân và lần lượt có duy nhất một điểm chung với các parabol đó (tham khảo hình vẽ biết phần gạch đậm là phần lát gạch). Chi phí cho phần lát gạch là $240$ nghìn đồng một mét vuông. Trong trường hợp hình tròn $(C)$ có diện tích lớn nhất thì chi phí lát gạch là bao nhiêu triệu đồng? (kết quả làm tròn tới hàng phần chục).
	\begin{center}
		\begin{tikzpicture}[scale=1,font=\footnotesize,line join=round,line cap=round,>=stealth,x=0.5cm,y=0.5cm]
				\tikzset{declare function={f  (\x)=0.25*(\x)^2-2;}}
				\fill[pattern=bricks] (-7,-6)  rectangle (7,6);
				\draw[fill=white] 
					plot[smooth,samples=150,domain=-6:6] ({f  (\x)},\x)
					plot[smooth,samples=150,domain=-6:6]  ({-f (\x)},\x)
					plot[smooth,samples=150,domain=-sqrt (8):sqrt (8)] ({f  (\x)},\x)--
					plot[smooth,samples=150,domain=sqrt  (8):-sqrt (8)] ({-f (\x)},\x)
					(0,0)  circle (2);
				\draw (6,5)  node{$(P_2)$}
					(-6,5) node{$(P_1)$} 
					(2.8,0) node{$(C)$}; 
				 \draw (-7,-6) rectangle (7,6);
				\draw  [<->] (-7,-7)--(7,-7) node[below,midway]{$14$\,m};
				\draw [<->] (8,-6)--(8,6)  node[right,midway]{$12$\,m};
			\end{tikzpicture}
	\end{center}
	\shortans{5361} 
	 \loigiai{
		\begin{center}
			\begin{tikzpicture}[scale=1,font=\footnotesize,line join=round,line cap=round,>=stealth,x=0.5cm,y=0.5cm]
					\tikzset{declare function={f (\x)=0.25*(\x)^2-2;}} 
					 \fill[pattern=dots] (-6,-7) rectangle (6,7);
					 \draw[fill=white]
						plot[smooth,samples=150,domain=-6:6] (\x,{f (\x)}) 
						plot[smooth,samples=150,domain=-6:6] (\x,{-f  (\x)})
						plot[smooth,samples=150,domain=sqrt (8):-sqrt (8)] (\x,{f (\x)})-- 
						plot[smooth,samples=150,domain=-sqrt (8):sqrt (8)]  (\x,{-f (\x)})
						(0,0) circle (2); 
					\draw[->]  (-8,0)--(-6,0) node[above left]{$-6$}--(6,0) node[above right]{$6$}--(8,0) node[above]{$x$};
					\draw[->] (0,-9)--(0,-7) node[below left]{$-7$}--(0,0) node[above left]{$O$}--(0,7) node[below left]{$7$}--(0,9)  node[above]{$y$};
					\draw  (-6,-7) rectangle (6,7);
					 \fill (6,0) circle (1pt)
						 (-6,0) circle (1pt)
						(0,-7) circle  (1pt)
						(0,7)  circle (1pt)
						(0,0) circle  (1pt)
						(-6,7) circle  (1pt)
						(0,-2)  circle (1pt);
					\draw  (-4.5,6) node{$(P_2)$}
						(-4.5,-6)  node{$(P_1)$}
						 (1,2.5) node{$(C)$}
						(-6,7) node[above  left]{$A$}
						(0,-2)  node[below right]{$K$} node[below left]{$k$};
				\end{tikzpicture}
		\end{center}
		Chọn hệ trục tọa độ $Oxy$ như hình vẽ với $M(0; m)~(m  < 0)$ là đỉnh của parabol $\left(P_1\right)$.\\
		Khi đó $\left(P_1\right)\colon y=\dfrac{7-m}{36}  x^2+m$ và $(C)\colon x^2+y^2=m^2$.\\
		 Để $\left(P_1\right)$, $(C)$ có một điểm chung duy nhất thì phương trình sau có nghiệm duy nhất.
		\[x^2+\left(\dfrac{7-m}{36} x^2+m\right)^2=m^2\Leftrightarrow x^2\left[\left(\dfrac{7-m}{36}\right)^2x^2+\dfrac{-m^2+7m+18}{18}\right]=0.\] 
		Yêu cầu bài toán $\Leftrightarrow-m^2+7m+18\geq 0\Leftrightarrow-2\leq m \leq  9$.\\
		Khi đó, bán kính lớn nhất của ($C$) khi $m=-2\Rightarrow  r=2$.\\
		 Hoành độ giao điểm của $\left(P_1\right)\colon y=\dfrac{1}{4} x^2-2$ và trục hoành là $x=\pm 2\sqrt{2}$.\\
		Diện tích phần lát gạch là $S=4\displaystyle\int_{2\sqrt{2}}^6\left(\dfrac{1}{4} x^2-2\right)\mathrm{d}  x+\pi r^2=\dfrac{18+8\sqrt{2}}{3}+4\pi$.\\
		Số tiền  lát gạch là $240S=5\,361$ nghìn đồng.
	} 
\end{bt}

% nguồn: Phan2_Lop_12_Dai_so_Chuong_4
\begin{bt}%[2D4V3-2]%[ma:04dc8d05]
	Để tham gia lễ hội hóa trang, Bạn An dự định làm một chiếc mặt nạ nửa mặt bằng chất liệu giấy cứng. Hình dạng của chiếc mặt nạ được bạn thiết kế trên mặt phẳng tọa độ $Oxy$ là phần hình phẳng giới hạn bởi hai đường parabol $\left(P_1\right)$, $\left(P_2\right)$ lần lượt có đỉnh là gốc tọa độ $O$ và điểm $I(0; 12)$, cùng nhận trục $Oy$ làm trục đối xứng và cùng đi qua điểm $M(15; 18)$. Mỗi đơn vị trên các trục  tọa độ có độ dài $1$ cm. Sau đó, bạn vẽ hai hình thoi bằng nhau có độ dài các đường chéo là $2\sqrt{2}$ cm và $4\sqrt{2}$ cm để khoét làm mắt (minh họa như hình vẽ dưới đây)
	\begin{center}
		\begin{tikzpicture}[line cap=butt,line join=miter,>=stealth]
				 \tikzset{declare function={xmin=-4; xmax=5; ymin=-1; ymax=5;
					f  (\x)=0.25*(\x)^2;
					 g (\x)=0.06*(\x)^2+3;
					}, 
					 smooth,samples=450
				} 
				\path (0,0) coordinate (A) 
					(-2,2.56) coordinate (E) 
					(-1.2,1.34) coordinate (G) 
					 (2,2.56) coordinate (F)
					(1.2,1.34)  coordinate (H)
					($(E)!.5!(G)$) coordinate  (I)
					($(F)!.5!(H)$)  coordinate (J)
					($(I)!.5!90:(G)$) coordinate (P) 
					($(I)!.5!-90:(G)$) coordinate (Q) 
					 ($(J)!.5!90:(H)$) coordinate (M)
					 ($(J)!.5!-90:(H)$) coordinate (N);
				\draw[->] (xmin,0)--(xmax,0)  node[shift={(-100:7pt)},font=\normalsize]{$ x $};
				 \draw[->] (0,ymin)--(0,ymax) node[shift={(190:7pt)},font=\normalsize]{$ y $};
				\fill (0,0) node[shift={(225:7pt)},font=\normalsize]{$  O $};
				 \path (0,3) node[above left]{$12$};
				\draw[dashed]  (4,0) node[below]{$15$}--(4,4)--(0,4) node[left]{$18$};
				\fill (0,3) circle (1pt);  \fill (4,4) circle (1pt); \fill (-4,4) circle (1pt);
				 \fill[gray!60,opacity=0.6]
					plot[domain=-4:4] (\x,{f (\x)})  --
					plot[domain=4:-4] (\x,{g  (\x)});
				\draw (E)--(P)--(G)--(Q)--cycle  (F)--(M)--(H)--(N)--cycle;
				\fill[white] (E)--(P)--(G)--(Q)--cycle (F)--(M)--(H)--(N)--cycle; 
				 \clip (xmin,ymin) rectangle (xmax,ymax);
				 \draw plot[domain=-4:4] (\x,{f (\x)});
				\draw  plot[domain=-4:4] (\x,{g (\x)});
			\end{tikzpicture}
	\end{center}
	Công đoạn cuối cùng là tô màu xám cho một mặt của mặt nạ. Tính diện tích cần tô màu theo đơn vị  cm$^2$.
	\shortans{224} 
	\loigiai{ 
		Gọi parabol có đỉnh $O$ có phương trình là $y=a x^2~(a > 0)$, ta có parabol đi qua điểm $M(15; 18)$ nên ta có $15^2a=18\Leftrightarrow a=\dfrac{18}{225}$, phương trình parabol là $y=\dfrac{18}{225}  x^2$.\\
		Gọi parabol có  đỉnh là $I$ có phương trình dạng $y=b x^2+12$, parabol đi qua điểm $M$ nên ta có phương trình $15^2b+12=18\Leftrightarrow b=\dfrac{6}{225}$, phương trình parabol là $y=\dfrac{6}{225} x^2$.\\
		Diện tích mặt nạ là $S_1=2\displaystyle\int_0^{15}\left[\dfrac{6}{225}  x^2+12-\dfrac{18}{225} x^2\right] \mathrm{d} x=240\left(\mathrm{cm}^2\right)$.\\
		Diện tích hai  hình thoi làm mắt là $S_2=2\cdot \dfrac{1}{2} 2\sqrt{2} \cdot 4\sqrt{2}=16\left(\mathrm{cm}^2\right)$.\\
		Diện tích cần sơn mặt nạ là $S=S_1-S_2=240-16=224\left(\mathrm{cm}^2\right)$. 
	} 
\end{bt}

% nguồn: Phan2_Lop_12_Dai_so_Chuong_4
\begin{bt}%[2D4V3-2]%[ma:7f955689]
	Mac Fish muốn làm logo trong suốt dán lên cửa sổ. Nhà thiết kế đưa bản vẽ như hình. (đơn vị trên mỗi trục toạ độ là decimét). Diện tích logo là bao nhiêu decimét vuông?  (Kết quả làm tròn đến hàng đơn vị).
	\begin{center}
		\begin{tikzpicture}[scale=0.7,font=\footnotesize,line join=round,line cap=round,>=stealth,declare function={f (\x)=-1/8*(\x)^2+2;
					g (\x)=3/16*(\x)^2-3;}] 
				\begin{scope}
					 \clip plot[samples=300,smooth,domain=-5:4] (\x,{f (\x)})--plot[samples=300,smooth,domain=4:-5] (\x,{g (\x)});
					 \fill[cyan] (-5,-3) rectangle (4,2);
					\draw (4,0)  circle (2);
					\draw [semithick] (2.8,0.5) circle (2pt); 
				\end{scope}
				\draw[samples=300,smooth,domain=-5:4] plot (\x,{f (\x)}) plot (\x,{g (\x)}); 
				\draw[->]  (-6,0) -- (5,0) node[right] {$x$};
				\draw[->]  (0,-3.5) -- (0,2.5) node[above] {$y$};
				 \draw (-5,-3) rectangle (4,2);
				\draw[<->]  (-5,-3.5)--(-4,-3.5) node[below,pos=.5]{$1$ dm};
				 \draw[<->] (-4,-3.5)--(0,-3.5) node[below,pos=.5]{$4$ dm};
				\draw[<->]  (0,-3.5)--(4,-3.5) node[below,pos=.5]{$4$ dm};
				\draw[<->] (-5.5,-3)--(-5.5,0) node[left,pos=.5]{$3$  dm};
				 \draw[<->] (-5.5,0)--(-5.5,2) node[left,pos=.5]{$2$ dm};
				\fill  (0,0) circle (1pt) node[below left]{$O$}
					 (3,-2.3) node{$y=g (x)$}
					(3.1,1.6) node{$y=f  (x)$}
					 (0,2) node[above left=-.1]{$A$}
					(-4,0) node[below]{$C$} 
					(4,0)  node[above right=-.1]{$B$}
					(0,-3) node[below  left=-.1]{$M$};
			\end{tikzpicture}
	\end{center}
	\shortans{28} 
	\loigiai{ 
		Dựa vào đồ thị, ta có diện tích logo là diện tích được  giới hạn bởi hai parabol $y=f(x)$, $y=g(x)$ và hai đường thẳng $x=4$, $x=-5$.\\
		Quan sát đồ  thị, ta thấy\\
		Parabol $y=f(x)=ax^2+bx+c$  đi qua ba điểm $A(0;2)$, $B(4;0)$, $C(-4;0)$
		\[\Rightarrow  \heva{
		&  c=2 \\
		 & 16a+4b+c=0 \\
		& 16a-4b+c=0}\Leftrightarrow \heva{ 
		 & c=2 \\
		& b=0  \\
		 & a=-\dfrac{1}{8}}
		\Rightarrow f(x)=-\dfrac{1}{8}x^2+2.\] 
		Parabol $y=g(x)=cx^2+\mathrm{\,d}x+e$ đi qua  ba điểm $M(0;-3)$, $B(4;0)$, $C(-4;0)$
		\[\Rightarrow  \heva{
		& e=-3  \\
		&  16c-4d-3=0 \\
		& 16c-4b-3=0} 
		 \Leftrightarrow \heva{
		&  e=-3 \\
		& b=\dfrac{3}{16}  \\
		&  d=0} \Rightarrow f(x)=\dfrac{3}{16}x^2-3.\]
		Vậy diện tích logo là 
		\[\displaystyle\int_{-5}^4{[f(x)-g(x)]\mathrm{\,d}x}=\displaystyle\int_{-5}^4{\left| \dfrac{5}{16}x^2-5  \right|\mathrm{\,d}x}=\dfrac{1345}{48}\approx 28\,\left(\text{dm}^2\right).
		 \]
	 }
\end{bt}

% nguồn: Phan2_Lop_12_Dai_so_Chuong_4
\begin{bt}%[2D4C3-2]%[ma:658ed871]
	\immini{ Vườn hoa của một trường học có hình dạng được giới hạn bởi một đường elip có bốn đỉnh $A$, $B$, $C$, $D$ và hai đường parabol có các đỉnh lần lượt là $E$, $F$ (phần tô đậm của hình vẽ). Hai đường parabol có cùng trục đối xứng $AB$, đối xứng nhau qua trục $CD$, hai parabol cắt elip tại các điểm $M$, $N$, $P$, $Q$. Biết $AB=8$ m, $CD=6$ m, $MN=PQ=3\sqrt{3}$ m, $EF=2$ m. Chi phí để trồng hoa trên vườn là $300\,000/\mathrm{m}^2$. Hỏi số tiền trồng hoa  cho cả vườn là bao nhiêu triệu đồng. Kết quả làm tròn đến hàng phần trăm.}{
	\begin{tikzpicture}[scale=0.6]
			\pgfmathsetmacro{\ycoord}{3*sqrt  (3)/2}
			\draw (0,0) ellipse (4cm  and 3cm);
			\fill[pink!50] plot[domain=\ycoord:-\ycoord,samples=100] ({-(4/27)*(\x)^2-1},\x)--plot[domain=-2:2,samples=100] (\x,{-sqrt  (9-(9/16)*(\x)^2)})
				 -- plot[domain=-\ycoord:\ycoord,samples=100] ({(4/27)*(\x)^2+1},\x)
				-- plot[domain=2:-2,samples=100] (\x,{sqrt  (9-(9/16)*(\x)^2)});
			\draw  plot[domain=-\ycoord:\ycoord,samples=50] ({-(4/27)*(\x)^2-1},\x);
			 \draw plot[domain=-\ycoord:\ycoord,samples=50] ({(4/27)*(\x)^2+1},\x);
			\fill (-4,0) circle (1pt); \node[left] at (-4,0)  {$A$};
			\fill (4,0) circle (1pt);  \node[right] at (4,0) {$B$};
			\fill (0,3) circle (1pt); \node[above]  at (0,3) {$C$};
			\fill (0,-3) circle (1pt); \node[below] at (0,-3)  {$D$};
			\fill (-1,0) circle (1pt); \node[below  left] at (-1,0) {$E$};
			\fill (1,0) circle (1pt); \node[below right] at (1,0) {$F$}; 
			\fill (-2,\ycoord) circle (1.5pt); \node[above left] at  (-2,\ycoord) {M};
			\fill  (-2,-\ycoord) circle (1.5pt); \node[below left] at (-2,-\ycoord) {N};
			 \fill (2,\ycoord) circle (1.5pt); \node[above right] at (2,\ycoord) {P};
			\fill (2,-\ycoord) circle (1.5pt); \node[below right]  at (2,-\ycoord) {Q};
			\draw[dashed] (-4,0)  -- (4,0);
			\draw[dashed] (0,-3) -- (0,3); 
			\draw[dashed] (2,-\ycoord) -- (2,\ycoord); 
		\end{tikzpicture}
	} 
	\shortans{4,81} 
	 \loigiai{
		Chọn hệ trục tọa độ với gốc tọa độ là trung điểm của $EF$ (xem hình vẽ).\begin{center}
		\begin{tikzpicture}[scale=0.6]
				\pgfmathsetmacro{\ycoord}{3*sqrt  (3)/2}
				\draw[->] (-6,0) --  (6,0) node[right] {$x$};
				\draw[->]  (0,-4.5) -- (0,4.5) node[above] {$y$};
				\draw (0,0) ellipse (4cm  and 3cm);
				\fill[pink!50] plot[domain=\ycoord:-\ycoord,samples=100]  ({-(4/27)*(\x)^2-1},\x)--plot[domain=-2:2,samples=100] (\x,{-sqrt (9-(9/16)*(\x)^2)})
					-- plot[domain=-\ycoord:\ycoord,samples=100] ({(4/27)*(\x)^2+1},\x) 
					-- plot[domain=2:-2,samples=100]  (\x,{sqrt (9-(9/16)*(\x)^2)});
				 \draw plot[domain=-\ycoord:\ycoord,samples=50] ({-(4/27)*(\x)^2-1},\x);
				 \draw plot[domain=-\ycoord:\ycoord,samples=50] ({(4/27)*(\x)^2+1},\x);
				\fill (-4,0) circle (1pt); \node[below left] at (-4,0)  {$A$};
				\fill (4,0) circle (1pt); \node[below right]  at (4,0) {$B$};
				 \fill (0,3) circle (1pt); \node[above left] at (0,3) {$C$};
				\fill (0,-3) circle (1pt); \node[below  left] at (0,-3) {$D$};
				 \fill (-1,0) circle (1pt); \node[below left] at (-1,0) {$E$};
				\fill (1,0) circle (1pt); \node[below right] at (1,0) {$F$}; 
				\fill (-2,\ycoord) circle (1.5pt); \node[above  left] at (-2,\ycoord) {M};
				\fill (-2,-\ycoord) circle (1.5pt); \node[below left]  at (-2,-\ycoord) {N};
				\fill (2,\ycoord) circle (1.5pt);  \node[above right] at (2,\ycoord) {P};
				\fill (2,-\ycoord) circle (1.5pt); \node[below right] at  (2,-\ycoord) {Q};
				\draw[dashed] (-4,0) -- (4,0); 
				\draw[dashed]  (0,-3) -- (0,3);
				\draw[dashed] (2,-\ycoord)  -- (2,\ycoord);
			\end{tikzpicture}
	\end{center}
	Khi đó $(E)$ có độ dài trục lớn $AB=2a=8 \Rightarrow a=4$.\\ 
	Độ dài trục bé $CD=2b=6 \Rightarrow  b=3$.
	 \\
	Phương trình $(E)$ là $\dfrac{x^2}{16}+\dfrac{y^2}{9}=1.$\\ Parabol có  phương trình $x=my^2+1$ (vì điểm $F(1;0)$).\\
	Do $PQ=3\sqrt{3}  \Rightarrow P\left(x_P;\dfrac{3\sqrt{3}}{2} \right)$ $(x_P > 0)$, mà $P \in(E)$ nên ta có $$\dfrac{x_P^2}{16}+\dfrac{\left(\dfrac{3\sqrt{3}}{2} \right)^2}{9}=1 \Rightarrow x_P=2.$$
	Mặt khác Parabol đi qua $P$ nên  $$2=a\cdot \left(\dfrac{3\sqrt{3}}{2} \right)^2+1 \Rightarrow a=\dfrac{4}{27}.$$
	Xét phần hình phẳng thuộc góc phần tư thứ nhất khi đó đường cong  $(E)$ có phương trình $x=4\sqrt{1-\dfrac{y^2}{9}}$.\\
	Parabol có phương trình $x = \dfrac{4}{27}y^2 +  1$.\\
	Diện tích hình phẳng giới hạn bởi các đường $y = 0$, $y = 3\sqrt{3}$, $x = 4\sqrt{1-\dfrac{y^2}{9}}$ và $x = \dfrac{8}{27}y^2$ là $$S_1 = \displaystyle\int_{0}^{3\sqrt{3}}\left( 4\sqrt{1-\dfrac{y^2}{9}} - \dfrac{4y^2}{27} -  1 \right)\mathrm{\,d}y.$$
	Diện tích elip là $S  = \pi ab = 12\pi$.\\
Diện tích phần tô đậm là $$S_T  = S-4\cdot S_1 \approx 16{,}03~\mathrm{m}^2.$$ Vậy $$T = 300\,000\cdot S_T \approx 4{,}81~\text{ triệu đồng}.$$}
\end{bt}

% nguồn: Phan2_Lop_12_Dai_so_Chuong_4
\begin{bt}%[2D4C3-2]%[ma:b1d086e5]
	\immini{ \lq\lq Bài toán tính diện tích hình càng cua\rq\rq. Từ bảng gỗ hình vuông $ABCD$ độ dài $12$ cm, người ta chia hình vuông này thành $9$ hình vuông nhỏ bằng nhau rồi vẽ một hình càng cua được giới $24$ hạn bởi các cung phần tư của các đường tròn tâm $A$, $E$, $F$, $G$ (xem hình vẽ). Tính diện  tích hình càng cua theo đơn vị $\mathrm{cm}^2$ (phần tô đậm trong hình vẽ). Viết kết quả làm tròn đến hàng đơn vị.}{
	\begin{tikzpicture}[scale=0.3]
			\draw (0,0) --  (12,0);
			\draw (0,0) -- (0,12); 
			\draw (0.5,0) -- (0.5,0.5)  -- (0,0.5);
			 \node[below left] at (0,0) {$D$};
			 \node[above left] at (0,12) {$A$};
			\node[above right] at (12,12) {$B$}; 
			\node[below right] at (12,0)  {$C$};
			\node[below] at (8,0) {$E$}; 
			\node[above] at (4,12)  {$G$};
			 \node[below left] at (8,4) {$F$};
			\draw (12,12)  arc (0:-90:12);
			 \draw (4,4) arc (-90:0:8);
			\draw  (0,0) arc (180:90:8);
			\draw  (4,4) arc (180:90:4);
			\fill[red!50] (0,0) arc (180:90:8)--(8,8) arc (90:180:4)--(4,4) arc (-90:0:8)--(12,12) arc (0:-90:12)--cycle; 
			\foreach \x in  {4,8,12}{
				\draw  (\x,0) -- (\x,12);
			 }
			\foreach \y in {4,8,12}{ 
				\draw  (0,\y) -- (12,\y);
			 }
			\foreach  \x in {0,4,8,12}{
				 \foreach \y in {0,4,8,12}{
				 \fill (\x,\y) circle (1pt);
			 }
			} 
			\node[below left] at  (8,4) {$F$};
		\end{tikzpicture}
	 }
	\shortans{37} 
	\loigiai{\begin{center}
		\begin{tikzpicture}[scale=0.3]
				\draw[->] (0,0)  -- (15,0) node[below]{$x$};
				\draw[->] (0,0)  -- (0,15) node[left]{$y$};
				\draw  (0.5,0) -- (0.5,0.5) -- (0,0.5);
				\node[below] at (4,0)  {$4$};
				 \node[below] at (8,0) {$8$};
				\node[below] at (12,0) {$12$}; 
				\node[left]  at (0,4) {$4$};
				\node[left] at (0,8)  {$8$};
				\node[left]  at (0,12) {$12$};
				\node[below left] at (0,0)  {$D$};
				\node[above right] at (0,12)  {$A$};
				 \node[above right] at (12,12) {$B$};
				\node[above  right] at (12,0) {$C$};
				 \node[above right] at (8,0) {$E$};
				 \node[above] at (4,12) {$G$};
				\draw  (12,12) arc (0:-90:12);
				\draw (4,4) arc  (-90:0:8);
				 \draw (0,0) arc (180:90:8);
				\draw (4,4) arc (180:90:4); 
				\fill[red!50] (0,0) arc (180:90:8)--(8,8) arc (90:180:4)--(4,4) arc  (-90:0:8)--(12,12) arc (0:-90:12)--cycle;
				\foreach \x in  {4,8,12}{
					\draw (\x,0) --  (\x,12);
				} 
				\foreach \y in  {4,8,12}{
					\draw  (0,\y) -- (12,\y);
				 }
				\foreach \x in {0,4,8,12}{ 
					\foreach  \y in {0,4,8,12}{
					\fill (\x,\y) circle  (1pt);
				} 
				 }
				\node[below left] at (8,4)  {$F$};
			\end{tikzpicture}
	\end{center} Đặt hệ trục tọa độ như hình vẽ ta có tọa độ các điểm $A(0;12)$; $B(12;12)$; $C(12;0)$; $D(0;0)$; $E(8;0)$; $G(4;12)$; $F(8;4)$.\\
	Phương trình đường tròn tâm $A$ có bán kính $R=12$ là $$(C_A)\colon x^2+(y-12)^2=144 \Rightarrow y=12-\sqrt{144-x^2}.$$ 
	Phương trình đường tròn tâm  $E$ có bán kính $R=8$ là $$(C_E)\colon (x-8)^2+y^2=64 \Rightarrow y=\sqrt{64-(x-8)^2}.$$
	Phương trình đường tròn tâm $G$  có bán kính $R=8$ là $$(C_G)\colon (x-4)^2+(y-12)^2=64 \Rightarrow y=12-\sqrt{64-(x-4)^2}.$$ Phương trình đường tròn tâm $F$ có bán kính $R=4$ là $$(C_F)\colon (x-8)^2 + (y-4)^2 = 16 \Rightarrow y = 4 + \sqrt{16-(x-8)^2}.$$
	Vậy \begin{eqnarray*}
	S  =&& \displaystyle\int_{0}^{4} \left( \sqrt{64-(x-8)^2} - 12 + \sqrt{144-x^2} \right) \mathrm{\,d}x\\
	&+& \displaystyle\int_{4}^{8} \left( \sqrt{64-(x-8)^2} - 4 - \sqrt{16-(x-8)^2} \right) \mathrm{\,d}x\\ 
	 &+& \displaystyle\int_{4}^{12} \left( 12-\sqrt{64-(x-4)^2} - 12 + \sqrt{144-x^2} \right) \mathrm{\,d}x\\
	\approx&& 37. 
\end{eqnarray*} }
\end{bt}

% nguồn: Phan2_Lop_12_Dai_so_Chuong_4
\begin{bt}%[2D4C3-2]%[ma:129f6e84]
	\immini{	Khu vườn nhà ông Ba có dạng hình tròn, bán kính $10$ m. Ông Ba dự định trồng hoa Hồng ở khu vực $S_1$ và hoa Ly ở khu vực hình bán nguyệt $S_2$. Với $S_1$ là phần diện tích giới hạn bởi đường parabol đi qua tâm hình tròn và $S_2$ là phần diện tích giới hạn bởi nửa đường elip không chứa tâm hình tròn (kích thước như hình vẽ). Biết rằng kinh phí trồng hoa Hồng là $100\,000$ đồng/$\mathrm{m}^2$, kinh phí trồng hoa Ly là $150\,000$ đồng/$\mathrm{m}^2$. Hỏi ông Ba phải mất bao nhiêu triệu đồng để trồng hoa lên hai dải đất đó (làm tròn kết quả đến hàng phần mười).}{\begin{tikzpicture}[scale=0.3]
		\coordinate  (O) at (0,0);
		\draw (O)  circle (10);
		\coordinate (S1_left) at  (-8,6);
		\coordinate (S1_right) at (8,6); 
		\fill[red!50] (S1_right) arc (36.87:143.13:10) -- plot[domain=-8:8,samples=50] (\x,{3/32  * \x*\x}) -- cycle;
		\node at (0,7.5)  {$S_1$};
		\draw[<->] (-8,6) -- (8,6) node[midway,below] {$16$ m}; 
		\draw[<->] (-8,0) -- (-8,6) node[midway,right] {$6$ m}; 
		\fill[blue!50] plot[domain=-10:10,samples=50] (\x,{-5*sqrt (1 -  \x*\x/100)}) arc (0:-180:10) -- cycle;
		\node  at (0,-7.5) {$S_2$};
		\draw[<->] (0.5,0) -- (0.5,-5) node[midway,right] {$5$  m};
		 \draw (O) circle (10);
		\draw[domain=-8:8,samples=100] plot (\x,{3/32 * \x*\x}); 
		 \draw[domain=-10:10,samples=100] plot (\x,{-5*sqrt (1 - \x*\x/100)});
		\draw[<->]  (-10,0)--(10,0);
	\end{tikzpicture}}
	\shortans{22,7} 
	 \loigiai{
		Gắn  hệ trục tọa độ như hình vẽ.
		\begin{center}
			\begin{tikzpicture}[scale=0.25]
					\coordinate  (O) at (0,0);
					\draw  (O) circle (10);
					\coordinate (S1_left)  at (-8,6);
					\coordinate (S1_right) at (8,6); 
					\fill[red!50] (S1_right) arc (36.87:143.13:10) -- plot[domain=-8:8,samples=50] (\x,{3/32 *  \x*\x}) -- cycle;
					 \node at (-2,7.5) {$S_1$};
					\fill[blue!50]  plot[domain=-10:10,samples=50] (\x,{-5*sqrt (1 - \x*\x/100)}) arc (0:-180:10) -- cycle;
					\node at (-2,-7.5)  {$S_2$};
					\draw  (O) circle (10);
					\draw[domain=-8:8,samples=100] plot (\x,{3/32 * \x*\x}); 
					\draw[domain=-10:10,samples=100] plot (\x,{-5*sqrt (1 - \x*\x/100)}); 
					\draw[->]  (-15,0)--(15,0) node[below]{$x$};
					\draw[->]  (0,-15)--(0,15) node[left]{$y$};
					\path (0,0) node[below left]{$O$} 
						(-10,0)  node[below left]{$-10$}
						(10,0) node[below  right]{$10$}
						(0,10) node[above  left]{$10$}
						(0,-10)  node[below left]{$-10$}
						(0,-5)  node[above left]{$-5$}
						(0,6) node[above right]{$6$} 
						(8,0)  node[below]{$8$};
					 \draw[dashed] (0,6)--(8,6)--(8,0);
				\end{tikzpicture}
		\end{center}Phương trình parabol đi qua tâm hình tròn là $y = \dfrac{3}{32} x^2$ và phương trình đường tròn là\break $x^2 + y^2 = 100$.\\
		Diện tích dải đất trồng hoa Hồng là diện tích hình phẳng được giới hạn bởi  parabol $(P)$, hình tròn $(C)$ và hai đường thẳng $x = -8$, $x = 8$. Khi đó
		$$S_1 = \displaystyle\int_{-8}^{8} \left(\sqrt{100 - x^2} -  \dfrac{3x^2}{32} \right) \mathrm{\,d}x \Rightarrow S_1 = 108{,}73~\mathrm{m}^2.$$
		Xét nửa hình tròn chứa hình bán nguyệt, ta thấy nửa hình tròn tạo bởi hình bán nguyệt và nửa hình elip với hình elip có độ dài hai bán trục $a = 5$, $b  = 10$ nên $$S_{(E) } = \pi ab = 50 \pi~ \mathrm{m}^2.$$
		Vậy diện tích hình bán  nguyệt là $$S_2 = \dfrac{S_{(C) } - S_{(E)}}{2} = \dfrac{10^2 \pi - 50 \pi}{2} = 25 \pi~ \mathrm{m}^2.$$
	Tổng số  tiền ông Ba cần phải chi để trồng hoa là $$T=100\,000\cdot S_1+150\,000\cdot S_2\approx22{,}7~\text{ triệu đồng}.$$}
\end{bt}

% nguồn: Phan2_Lop_12_Dai_so_Chuong_4
\begin{bt}%[2D4C3-2]%[ma:de52806b]
	\immini{ Vào ngày $8/3$, từ một tấm bìa hình vuông cạnh $8$ (dm), An dùng các đồ thị của hàm số $y=f(x)$ và $x=f(y)$ để vẽ hoa tặng mẹ, nhụy hoa được vẽ bằng đường tròn có tâm là tâm hình vuông và bán kính bằng $2$ (dm). Khi đặt trong hệ tọa độ $Oxy$ thì điểm $A(2;0)$  và đường cong $AB$ là một phần của đồ thị hàm đa thức bậc ba $y=\dfrac{1}{2}x^3+ax^2+bx+c$. Biết rằng tỷ lệ phần tô đậm chiếm $\dfrac{3}{4}$ diện tích hình vuông. Tính $abc$.}{
	\begin{tikzpicture}[scale=0.8]
			\draw  (-4,-4) rectangle (4,4);
			\node[below  left] at (4,0) {$4$};
			\node[below left] at (0,4)  {$4$};
			\fill[teal!30]  (0,0) circle (2);
			 \draw (0,0) circle (2);
			\path[fill=blue,draw] (2,0) arc (0:90:2) 
				.. controls (0.25,3.5) and (3.5,4).. (4,4) 
				.. controls (4,3.5) and (3.5,0.25)..  (2,0);
			 \path[fill=blue,draw] (-2,0) arc (180:90:2)
				.. controls  (-0.25,3.5) and (-3.5,4).. (-4,4)
				.. controls (-4,3.5)  and (-3.5,-0.25).. (-2,0);
			\path[fill=blue,draw] (-2,0) arc (-180:-90:2) 
				.. controls (-0.25,-3.5)  and (-3.5,-4).. (-4,-4)
				.. controls  (-4,-3.5) and (-3.5,-0.25).. (-2,0);
			\path[fill=blue,draw] (2,0) arc  (0:-90:2)
				.. controls (0.25,-3.5)  and (3.5,-4).. (4,-4)
				.. controls (4,-3.5) and (3.5,-0.25)..  (2,0);
			\node[above left,xshift=-2pt] at  (2,0) {$A$};
			\node[above  right] at (4,4) {$B$};
			\node[below left] at (0,0) {$O$}; 
			 \draw[->] (-4.5,0) -- (4.5,0) node[right] {$x$};
			 \draw[->] (0,-4.5) -- (0,4.5) node[above] {$y$};
			\node[below left] at  (2,0) {$2$};
		\end{tikzpicture}}
	 \shortans{72}
	\loigiai{Diện tích hình vuông  của hình vẽ là $8\cdot 8=64$ $(\mathrm{dm}^2)$.\\
		Dựa vào  đồ thị hình vẽ ta có $A(2;0)$; $B(4;4)$.\\
		Lại có $y=\dfrac{1}{2}x^3+ax^2+bx+c$ đi qua $A(2;0)$; $B(4;4)$ nên  ta có $$\heva{& \dfrac{1}{2}\cdot8+a\cdot4+b\cdot 2+c=0 \\ &\dfrac{1}{2}\cdot 64+a\cdot 16+b\cdot 4+c=4} \Leftrightarrow \heva{&4a+2b+c=-4 \\& 16a+4b+c=-28.}\quad (1)$$
		Ta có diện tích phần trắng là $$8\cdot \displaystyle\int_{2}^{4} \left(\dfrac{1}{2}x^3+ax^2+bx+c  \right)\mathrm{\,d}x.$$
		Khi đó diện tích phần in đậm là $$S_{\text{ in đậm}}  = S_{\text{ hv}} - S_{\text{ trắng}} = 64 - 8\cdot \displaystyle\int_{2}^{4} \left( \dfrac{1}{2}x^3 + ax^2 + bx + c \right)\mathrm{\,d}x.$$
		Theo giả thiết ta  có
		\begin{eqnarray*}
			&&\dfrac{S_{\text{ in đậm}}}{S_{\text{ hv}}} = \dfrac{64 - 8\cdot \displaystyle\int_{2}^{4} \left(  \dfrac{1}{2}x^3 + ax^2 + bx + c \right)\mathrm{\,d}x}{64} = \dfrac{3}{4}\\
			&\Rightarrow& \displaystyle\int_{1}^{2} \left( \dfrac{1}{2}x^3 + ax^2 + bx + c \right) \mathrm{\,d}x = 2\\ 
			&\Leftrightarrow& \left[\dfrac{1}{8}x^4 + \dfrac{1}{3}ax^3 + \dfrac{1}{2}bx^2  + cx \right]\Big|_{2}^{4} = 2\\
			&\Leftrightarrow&  30 + \dfrac{56}{3}a + 6b + 2c = 2\quad (2)
		\end{eqnarray*}
		Từ  (1) và (2) suy ra $\heva{& a = -3 \\& b = 6 \\& c = -4.}$\\
	 Vậy $abc = 72$.}
\end{bt}

% nguồn: Phan2_Lop_12_Dai_so_Chuong_4
\begin{bt}%[2D4C3-2]%[ma:c8d8358c]
	Một cửa vòm có hình dạng parabol được  lắp các tấm kính hình tròn đường kính $1$ m và các tấm kính hình vuông có cạnh $1$ m như hình vẽ. Phần còn lại của cửa được sơn màu trang trí với mức giá $1{,}2$ triệu đồng/$\mathrm{m^2}$. Chi phí sơn màu là bao nhiêu triệu đồng? (kết quả làm tròn đến hàng phần chục)
	\begin{center}
		\begin{tikzpicture}[scale=1]
				\def\r{0.5} 
				 \def\s{1}
				\begin{scope}
					\clip (-3.5,0) rectangle (3.5,4); 
					 \fill[gray!60] plot[domain=-3.5:3.5,smooth] (\x,{-0.5*(\x)^2+4});
				\end{scope}
				\node[below left]  at (-2.5,0) {\textcolor{blue}{A}};
				 \node[below right] at (2.5,0) {\textcolor{blue}{B}};
				 \draw[fill=white] (-2,0.5) circle (\r);
				\draw[fill=white] (2,0.5) circle  (\r);
				\draw[fill=white] (0,3.5) circle (\r); 
				\draw[fill=white] (-1,2) rectangle  ++(1,1);
				\draw[fill=white] (0,2) rectangle ++(1,1); 
				 \foreach \x in {-2,-1,0,1}{
					\draw[fill=white] (\x,1)  rectangle ++(1,1);
				} 
			\end{tikzpicture}
	\end{center}
	 \shortans{8,1}
	 \loigiai{
	\begin{center}
		\begin{tikzpicture}[scale=1]
				\def\r{0.5} 
				 \def\s{1}
				\draw[->] (-3.5,0)--(3.5,0) node[below]{$x$}; 
				\draw[->] (0,-0.5)--(0,4.5)  node[right]{$y$};
				\draw (0,0) node[below  left]{$O$};
				\begin{scope}
					\clip (-3.5,-0.3) rectangle  (3.5,4);
					\draw plot[domain=-3.5:3.5,smooth]  (\x,{-0.5*(\x)^2+4});
				\end{scope}
				\node[below left] at (-2.5,0) {\textcolor{blue}{A}}; 
				\node[below  right] at (2.5,0) {\textcolor{blue}{B}};
				\node[below left] at (-2,2.5) {\textcolor{blue}{C}}; 
				 \node[below right] at (2,2.5) {\textcolor{blue}{D}};
				\draw (-2,0.5)  circle (\r);
				\draw (2,0.5) circle  (\r);
				\draw (0,3.5) circle  (\r);
				 \draw (-1,2) rectangle ++(1,1);
				\draw (0,2) rectangle ++(1,1); 
				\foreach \x in  {-2,-1,0,1}{
					 \draw (\x,1) rectangle ++(1,1);
				 }
				\foreach \x  in {-2,-1,0,1}{
					\draw (\x,0)  rectangle ++(1,1);
				} 
			\end{tikzpicture}
	\end{center}
	Gắn hệ trục tọa độ như  hình vẽ suy ra $C(2;2)$, $D(-2;2)$, $E(0;4)$.\\
	Cửa vòm hình parabol có dạng  $y=ax^2+bx+c$ đi qua $E$, $C$, $D$ ta có hệ phương trình
	$$\heva{& 4a+2b+c=2 \\ & 4a-2b+c=2\\&c=4}\Leftrightarrow  \heva{& a=-\dfrac{1}{2}\\ & b=0\\
	& c=4} \Rightarrow  y=-\dfrac{1}{2}x^2+4.$$
	Phương trình hoành độ giao điểm của (P) và trục hoành là $-\dfrac{1}{2}x^2+4=0\Leftrightarrow  \hoac{& x=-2\sqrt{2} \\ & x=2\sqrt{2}.}$\\
	Diện tích cổng vòm parabol  là $\displaystyle\int_{-2\sqrt{2}}^{2\sqrt{2}}{\left(-\dfrac{1}{2}{x^2}+4\right)}\mathrm{d\,}x$.\\
	Diện tích được lắp đặt các tấm kính hình tròn đường  kính 1m và hình vuông cạnh $1$ m là
	 $3\cdot\pi{\left(\dfrac{1}{2}\right)^2}+6\cdot1=\dfrac{3}{4}\pi+6$.\\
	Diện tích của sơn màu  là $\displaystyle\int_{-2\sqrt{2}}^{2\sqrt{2}}{\left(\dfrac{-1}{2}x^2+4\right)}\mathrm{d\,}x-\left(\dfrac{3}{4}\pi+6\right)$.\\
	Vậy chi phí sơn màu là  $1{,}2\cdot\left(\displaystyle\int_{-2\sqrt{2}}^{2\sqrt{2}}{\left(\dfrac{-1}{2}{x^2}+4\right)}dx-\left(\dfrac{3}{4}\pi+6\right)\right)\approx 8{,}1$ triệu đồng.
	 }
\end{bt}

% nguồn: Phan2_Lop_12_Dai_so_Chuong_4
\begin{bt}%[2D4C3-2]%[ma:e6181f1f]
	 \immini
	{ 
	Để chuẩn bị quảng bá sản phẩm, người ta trang trí tấm pano dạng parabol như hình vẽ, biết $OS=8$ m, $AB=6$ m với $O$ là trung điểm của $AB$. Tấm pano được chia  thành ba phần để trang trí với mức chi phí khá nhau: phần trên là phần kẻ sọc giá $100\,000$ đồng/m$^2$, phần giữa hình quạt tâm $O$ bán kính $2$ m được tô đậm giá $200\,000$ đồng/m$^2$, phần còn lại giá $150\,000$ đồng/m$^2$. Tính tổng chi phí để trang trí tấm pano (đơn vị triệu đồng, kết quả làm tròn đến hàng phần trăm).
	} 
	 {
	\begin{tikzpicture}[scale=1,font=\footnotesize,line join=round,line cap=round,>=stealth]
			\def\xmin{-3} \def\xmax{3} 
			 \def\ymin{-1} \def\ymax{5}
			\draw[smooth,samples=300,domain=-2:2]  plot (\x,{-(\x)^2+4});
			\draw[smooth,samples=300,domain=-1.732:1.732] plot (\x,{(4-(\x)^2)^0.5}); 
			 \draw[pattern=north east lines,opacity=0.5] (-1.732,1) plot[domain=-1.732:1.732] (\x,{-(\x)^2+4})--(0,0)--cycle;
			\fill[samples=100,domain=-1.732:1.732,gray!60] plot  (\x,{(4-(\x)^2)^0.5})--(0,0)--cycle;
			\node at (0,0)  [below]{$O$};
			\node at (0,4)  [above]{$S$};
			\draw (-2,0) circle (.5pt) node[left]{\footnotesize  $A$};
			\draw (2,0) circle  (.5pt) node[right]{\footnotesize $B$};
			 \draw (-1.732,1)--(0,0)--(1.732,1)
				 (-2,0)--(2,0);
			 \clip (\xmin+0.1,\ymin+0.1) rectangle (\xmax-0.1,\ymax-0.1);
		\end{tikzpicture}
	 }
	 \shortans{4,44}
	\loigiai{ 
		\begin{center}
			\begin{tikzpicture}[scale=1,font=\footnotesize,line join=round,line cap=round,>=stealth]
					 \def\xmin{-3} \def\xmax{3}
					 \def\ymin{-1} \def\ymax{5}
					\draw[smooth,samples=300,domain=-2:2] plot  (\x,{-(\x)^2+4});
					\draw[smooth,samples=300,domain=-1.732:1.732] plot (\x,{(4-(\x)^2)^0.5}); 
					\draw[pattern=north east lines,opacity=0.5]  (-1.732,1) plot[domain=-1.732:1.732] (\x,{-(\x)^2+4})--(0,0)--cycle;
					\fill[samples=100,domain=-1.732:1.732,gray!60] plot (\x,{(4-(\x)^2)^0.5})--(0,0)--cycle; 
					\draw[->] (\xmin,0)--(\xmax,0) node [below]{$x$}; 
					 \draw[->] (0,\ymin)--(0,\ymax) node [left]{$y$};
					\node  at (0,0) [below right]{$O$};
					\node  at (1.2,3) [right]{$(P)$};
					\node at  (.5,1.5) [right]{$(C)$};
					\node at (0,4)  [above right]{$S$};
					\node at (-1.732,1)  [left]{$M$};
					 \node at (1.732,1) [right]{$N$};
					\draw (-2,0) circle  (.5pt) node[below left]{\footnotesize $A$};
					\draw (2,0)  circle (.5pt) node[below right]{\footnotesize $B$};
					\draw  (-2,0) circle (.5pt) node[above left]{\footnotesize $-3$};
					\draw  (2,0) circle (.5pt) node[above right]{\footnotesize $3$};
					\draw (-1.732,0) circle (.5pt) node[below]{\footnotesize  $-\frac{3\sqrt{55}}{8}$};
					 \draw (1.732,0) circle (.5pt) node[below]{\footnotesize $\frac{3\sqrt{55}}{8}$};
					\draw (0,1) circle (.5pt) node[above  right]{\footnotesize $\dfrac{9}{8}$};
					\draw[dashed]  (-1.732,0)--(-1.732,1)--(1.732,1)--(1.732,0);
					 \draw (-1.732,1)--(0,0)--(1.732,1);
					\clip (\xmin+0.1,\ymin+0.1)  rectangle (\xmax-0.1,\ymax-0.1);
				\end{tikzpicture}
		\end{center}
		Chọn hệ trục tọa độ $ Oxy $ như  hình vẽ.\\
		$(P)$ đi  qua $A(-3;0)$, $B(3;0)$ và $S(0;8)$.\\
		 Suy ra phương trình parabol là $(P)\colon y=-\dfrac{8}{9}x^2+8$.\\
		Rìa của hình quạt là cung tròn của đường  tròn $(C)$ có phương trình
		\[x^2+y^2=9\Rightarrow  y=\pm \sqrt{9-x^2}.\]
		Hoành độ điểm $N$ là nghiệm phương trình 
		\begin{eqnarray*}
			-\dfrac{8}{9}x^2+8&=&\sqrt{9-x^2},\ 0<x<3\\ 
			x&=&\pm \dfrac{3\sqrt{55}}{8}. 
		\end{eqnarray*}
		Ta có  $y_N=-\dfrac{8}{9}x_N^2+8=\dfrac{9}{8}$.\\
		 Suy ra $N\left(\dfrac{3\sqrt{55}}{8};\dfrac{9}{8}\right)$.\\
		Phương  trình $ON\colon y=\dfrac{3}{\sqrt{55}}x$.\\
		Vì tấm pano đối xứng qua trục $Oy$ nên  ta có
		\begin{itemize}
			\item  Diện tích phần gạch sọc là $ S_1=2\displaystyle\int_0^{\tfrac{3\sqrt{55}}{8}}\left(-\dfrac{8}{9}x^2+8-\sqrt{9-x^2}\right)\mathrm{\,d}x\approx 17{,}94$ m$^2$.
			\item Diện tích hình quạt là $ S_2=2\displaystyle\int_0^{\tfrac{3\sqrt{55}}{8}}\left(\sqrt{9-x^2}-\dfrac{3}{\sqrt{55}}x\right)\mathrm{\,d}x\approx  10{,}68$ m$^2$.
			 \item Diện tích phần còn lại là $ S_3=\displaystyle\int_{-3}^{3}\left(-\dfrac{8}{9}x^2+8\right)\mathrm{\,d}x-S_1-S_2\approx 3{,}38$ m$^2$.
		\end{itemize}
		Tổng chi phí  là
		\[T=S_1\cdot 100\,000+ S_2\cdot  200\,000+ S_3\cdot 150\,000\approx4\,440\,000\ \text{ đồng}\approx 4{,}44\ \text{ (triệu đồng)}.\]
	 }
\end{bt}

% nguồn: Phan2_Lop_12_Dai_so_Chuong_4
\begin{bt}%[2D4C3-2]%[ma:485be932]
	\immini{	Từ một miếng tôn hình vuông, người thợ làm chậu cảnh đã tạo mẫu và cắt đi phần hình phẳng không tô đậm trên hình, phần tô đậm được giữ lại làm khuôn quay thành các đôn để đặt các chậu hoa. Tính diện tích hình phẳng đã bị cắt bỏ (đơn vị cm$^2$) biết đường cong trong hình là một parabol có đỉnh nằm trên đường chéo của hình vuông (làm tròn kết quả đến hàng phần mười).}{\begin{tikzpicture}[scale=0.5]
		\fill[teal!60!green] (0,0) rectangle  (12,12);
		\fill[white] (0,0) -- (0,10) -- (2,10)  -- (2,8) -- (4,8) arc (180:270:4) -- (8,4) -- (8,2) -- (10,2) -- (10,0) -- cycle;
		\draw (0,0)  rectangle (12,12);
		\draw  (0,0) -- (0,10) -- (2,10) -- (2,8) -- (4,8) arc (180:270:4) -- (8,4) -- (8,2) -- (10,2) -- (10,0) -- cycle;
		\draw[dashed] (0,0)  -- (6,6);
		\draw[dashed] (0,12) --  (12,0);
		 \node[left] at (0,11) {2 cm};
		\node[below]  at (11,0) {2 cm};
		\node[below] at  (9,2) {2 cm};
		\node[left] at (0,5) {10 cm}; 
		\node[below] at (5,0) {10 cm}; 
		\node[below] at (3,8) {2  cm};
		\node[above  left] at (3.5,3.5) {$5\sqrt{2}$ cm};
	\end{tikzpicture}
	 }
	 \shortans{58,7}
	\loigiai{\begin{center}
		\begin{tikzpicture}[scale=0.5]
				\fill[teal!60!green]  (0,0) rectangle (12,12);
				\fill[white] (0,0) -- (0,10) -- (2,10) -- (2,8) -- (4,8) arc (180:270:4) -- (8,4)  -- (8,2) -- (10,2) -- (10,0) -- cycle;
				\draw (0,0) rectangle  (12,12);
				\draw (0,0) -- (0,10) -- (2,10)  -- (2,8) -- (4,8) arc (180:270:4) -- (8,4) -- (8,2) -- (10,2) -- (10,0) -- cycle;
				\draw[dashed] (0,0) -- (6,6); 
				\draw[dashed] (0,12) --  (12,0);
				\node[left] at (0,11)  {2 cm};
				\node[below] at (11,0) {2  cm};
				\node[below] at  (9,2) {2 cm};
				\node[left] at (0,5) {10  cm};
				\node[below] at (5,0) {10  cm};
				\node[below] at  (3,8) {2 cm};
				\node[above right]  at (6,6) {$O$};
				 \node[above left] at (3.5,3.5) {$5\sqrt{2}$ cm};
				\node[above left] at (-1,-1) {$y$}; 
				 \node[below right] at (13,-1) {$x$};
				\draw[->]  (6,6)--(-1,-1);
				\draw[->] (6,6)--(13,-1); 
			\end{tikzpicture}
	\end{center}Chọn hệ trục như hình vẽ suy ra phương trình parabol có dạng $y=a\left(x-2\sqrt{2}\right) \left(x+2\sqrt{2}\right)$.\\
	Đỉnh parabol có tọa độ  là $\left(0;\sqrt{2}\right)$ thuộc parabol suy ra $a=-\dfrac{1}{4\sqrt{2}}$.\\
	Diện tích  phần parabol là $\dfrac{1}{4\sqrt{2}}\displaystyle\int_{-2\sqrt{2}}^{2\sqrt{2}}(8-x^2) \mathrm{\,d}x=\dfrac{16}{3}$.\\ Diện tích cần tìm là $$S=\dfrac{12^2}{2}-4\cdot2-\dfrac{16}{3}=\dfrac{176}{3}\approx 58{,}7~\mathrm{cm}^2.$$
} 
\end{bt}

% nguồn: Phan2_Lop_12_Dai_so_Chuong_4
\begin{bt}%[2D4C3-2]%[ma:c5afbe5a]
	Một nhóm các kĩ sư muốn xây dựng một cây cầu vòm dàn thép với giá đỡ dưới bằng thép cao cấp có hình dáng là một parabol nối từ 2 cột trụ $A$ và $B$ nằm bên dưới cây cầu. Biết hai cột trụ cách nhau $400$ m, khoảng cách từ trụ $A$ đến cây cầu là $50$  m và $AB$ song song với mặt đường.
	Gắn hệ trục $Oxy$ vào cây cầu với đơn vị trục tọa  độ là $10$ m. Giá đỡ dưới bằng thép là đường cong parabol tạo với hai trục tọa độ các hình phẳng có diện tích $S_1$, $S_2$ như hình vẽ bên. Biết rằng $S_2-2S_1=\dfrac{2200}{21}$. Điểm cao nhất của giá đỡ dưới bằng thép cao cấp cách mặt đường cây cầu bao~nhiêu mét? (làm tròn đến hàng phần mười)
	 \shortans{64,3}
	\loigiai{ 
		Parabol có  dạng $(P)\colon y=ax^2+bx+c$.\\
		Vì $(P)$ đi  qua $(0;-5)$ nên $f(0)=-5 \Rightarrow c=-5$.\\
		 Và $(P)$ đi qua $(40;-5)$ nên $f(40)=-5 \Rightarrow 1600a+40b+c=-5 \Rightarrow 1600a+40b=0 $. \hfill (1)\\
		 Ta có
		 \allowdisplaybreaks
		\begin{eqnarray*}
			S_2-2S_1 &=&  S_2-(S_{OABC}-S_3)=(S_2+S_3)-S_{OABC}\\
			 &=&\displaystyle\int_0^{40} \left( ax^2+bx+c+5\right) \mathrm{\,d}x-5\cdot40\\
			&=&  \dfrac{64000}{3}a+800b+40c.
		\end{eqnarray*}
		Suy ra $\dfrac{64000}{3}a+800b +40c= \dfrac{6400}{21}$. \hfill  $(2)$\\
		Từ $(1)$ và $(2)$ ta giải  được $a=-\dfrac{1}{35}$, $b=\dfrac{8}{7}$. \\
		Suy  ra $(P)\colon y=-\dfrac{1}{35}x^2+\dfrac{8}{7}x-5$.\\
		Do đó $(P)$  có đỉnh $I\left( 20;\dfrac{45}{7}\right) $.\\
		Vậy điểm cao nhất của giá đỡ dưới bằng thép cao cấp cách mặt đường cây cầu $10\cdot  \dfrac{45}{7} \approx 64{,}3$ m.
	 }
\end{bt}

% nguồn: Phan2_Lop_12_Dai_so_Chuong_4
\begin{bt}%[2D4C3-2]%[ma:a64d6b44]
	Bạn An thực hiện thiết kế một logo hình con mắt cho một phòng khám nhãn khoa. Logo là một hình phẳng giới hạn bởi hai parabol $y=f(x)$ và $y=g(x)$ có các kích thước như hình vẽ dưới đây (phần được gạch chéo) và một hình tròn có bán kính bằng $0{,}5$ dm ở giữa là phần con ngươi (phần  được tô màu xanh), đơn~vị trên mỗi trục tọa độ là decimet.
	\begin{center}
		\begin{tikzpicture}[scale=1.0,font=\footnotesize,line join=round,line cap=round,>=stealth];
				\fill[pattern=north east  lines,smooth]plot[domain=-2.449489743:2.449489743] (\x,{-1/4*((\x)^2)+2})-- plot[domain=-2.449489743:2.449489743] (\x,{1/4*((\x)^2)-1})-- cycle;
				 \draw[fill=cyan] (0,0.5) circle (0.5);
				\draw[->] (-3,0) -- (3,0) node[below] {\scriptsize $x$}; 
				\draw[->] (0,-1.5)  -- (0,2.5) node[left] {\scriptsize $y$};
				\draw (0.1,0.1) node[below left]{\scriptsize  $O$};
				 \draw[thin,samples=150,smooth,domain=-2.449489743:2.449489743] plot (\x,{-1/4*((\x)^2)+2});
				\draw[thin,samples=150,smooth,domain=-2.449489743:2.449489743]  plot (\x,{1/4*((\x)^2)-1});
				\draw[dashed] (-2,0) node[below]{$-2$} -- (-2,1)--(2,1)--(2,0) node[below]{$2$}; 
				 \draw (-2,1) circle (1.2pt)
					(0,1) circle (1.2pt) 
					(2,1) circle (1.2pt) 
					(0,-1)  circle (1.2pt)
					 (0,0.5) circle (1.2pt)
					(0,2) circle  (1.2pt);
				\draw (0,2.2)  node[right]{$2$}
					(0,-1.2) node[right]{$-1$} 
					 (0,1.2) node[right]{$1$};
			\end{tikzpicture}
	\end{center}
	Biết rằng chi phí để sơn phần con ngươi  hình tròn màu xanh là $25\,000$ đồng/dm$^2$ và chi phí để sơn phần còn lại màu đen là $20\,000$ đồng/dm$^2$. Chi phí để sơn logo trên là bao nhiêu nghìn đồng? (Kết quả làm tròn đến hàng đơn~vị)
	\shortans{200} 
	\loigiai{ 
		Gọi phương trình parabol $y=f(x)$ là $f(x)=ax^2+bx+c$ với $a\neq 0$, đi  qua các điểm có tọa độ lần lượt là $(2;1)$, $(0;2)$, $(-2;1)$.\\
		Lập hệ phương trình ba ẩn, từ đó  suy ra được $\heva{&a=-\dfrac{1}{4}\\&b=0\\&c=2}$ nên $f(x)=-\dfrac{1}{4}x^2+2$.\\
		Gọi phương trình parabol $y=g(x)$ là $g(x)=mx^2+nx+p$ với $m\neq 0$, đi  qua các điểm có tọa độ lần lượt là $(-2;0)$, $(0;-1)$, $(2;0)$.\\
		Lập hệ phương  trình ba ẩn, từ đó suy ra được $\heva{&m=\dfrac{1}{4}\\&n=0\\&p=-1}$ nên $g(x)=\dfrac{1}{4}x^2-1$.\\
		Phương trình  hoành độ giao điểm $$f(x)=g(x)\Leftrightarrow-\dfrac{1}{4}x^2+2=\dfrac{1}{4}x^2-1\Leftrightarrow \dfrac{1}{2}x^2=3\Leftrightarrow x=\pm \sqrt{6}.$$
		Diện tích hình tròn được  sơn màu xanh là $S_1=\pi r^2=\pi(0{,}5)^2=0{,}25\pi\,\left(\text{dm}^2\right)$.\\
		 Diện tích hình phẳng giới hạn bởi hai đường thẳng $y=f(x)$ và $y=g(x)$ là
		 \[S_2=\displaystyle\int^{\sqrt{6}}_{-\sqrt{6}}|f(x)-g(x)|\,\mathrm{\,d}x=\displaystyle\int^{\sqrt{6}}_{-\sqrt{6}}\left|\left(-\dfrac{1}{4}x^2+2\right)-\left(\dfrac{1}{4}x^2-1\right)\right|\,\mathrm{\,d}x.\]
		 Diện tích phần sơn được tô màu đen là $$S_3=S_2-S_1=\displaystyle\int^{\sqrt{6}}_{-\sqrt{6}}\left|\left(-\dfrac{1}{4}x^2+2\right)-\left(\dfrac{1}{4}x^2-1\right)\right|\mathrm{\,d}x-0{,}25\pi.$$
		Chi phí để sơn logo là 
		 \allowdisplaybreaks
		\begin{eqnarray*}
			 T&=&0{,}25\pi\cdot 25\,000+\left[\displaystyle\int^{\sqrt{6}}_{-\sqrt{6}}\left| \left(-\dfrac{1}{4}x^2+2\right)-\left(\dfrac{1}{4}x^2-1\right)\right|\mathrm{\,d}x-0{,}25\pi\right]\cdot 20\,000\\
			&=&200 \,\text{ nghìn  đồng}.
		\end{eqnarray*}
	} 
\end{bt}

% nguồn: Phan2_Lop_12_Dai_so_Chuong_4
\begin{bt}%[2D4C3-2]%[ma:3a78edc2]
	Một biển quảng cáo có dạng hình vuông $ABCD$ cạnh $AB=4$ m. Trên tấm biển đó có các đường tròn tâm $A$ và đường tròn tâm $B$ cùng bán kính $R=4$ m, hai đường tròn cắt nhau như hình vẽ. Chi  phí để sơn phần gạch chéo là $150\,000$ đồng/m$^2$, chi phí sơn phần màu đen là $100\,000$ đồng/m$^2$ và chi phí để sơn phần còn lại là $250\,000$ đồng/m$^2$.
	\begin{center}
		\begin{tikzpicture}[scale=1.0,font=\footnotesize,line join=round,line cap=round,>=stealth];
				\path 
					(0,0) coordinate  (A)
					(4,0) coordinate  (B)
					(4,4)  coordinate (C)
					(0,4)  coordinate (D);
				\draw  (A)--(B)--(C)--(D)--(A);
				\draw[thin,samples=150,smooth,domain=0:4]  plot (\x,{sqrt (-1*((\x)^2)+16)});
				 \draw[thin,samples=150,smooth,domain=0:4] plot (\x,{sqrt (-1*((\x-4)^2)+16)});
				\draw[dashed] (2,0)--(2,3.464101615); 
				\fill[gray,smooth]plot[domain=0:4] (\x,{sqrt (-1*((\x)^2)+16)})-- plot[domain=4:0] (\x,{sqrt  (-1*((\x-4)^2)+16)})-- cycle;
				\fill[pattern=north east lines,smooth]plot[domain=4:2] (\x,{sqrt (-1*((\x)^2)+16)})-- plot[domain=2:4]  (\x,{0});
				\fill[pattern=north  east lines,smooth]plot[domain=0:2] (\x,{sqrt (-1*((\x-4)^2)+16)})-- plot[domain=2:0] (\x,{0});
				 \foreach \x/\g in{A/-90,B/-90,C/90,D/90}
				\fill(\x)  circle (1pt)
					 ($(\x)+(\g:3mm)$) node{$\x$};
				\draw (2,0) node[below]{$4$m} 
					 (4,2) node[right]{$4$m};
			\end{tikzpicture}
	\end{center}
	Hỏi số tiền để sơn biển quảng cáo theo cách trên (Đơn vị: triệu đồng và kết quả làm tròn đến chữ số thập phân thứ hai sau  dấu phẩy)?
	\shortans{2,20} 
	\loigiai{ 
		Gọi $I$ là giao điểm của 2 cung tròn $AC$  và $BD$.
		\begin{center}
			\begin{tikzpicture}[scale=1.0,font=\footnotesize,line join=round,line cap=round,>=stealth];
					\path 
						(0,0) coordinate  (A)
						(4,0) coordinate  (B)
						(4,4) coordinate (C) 
						(0,4)  coordinate (D);
					\draw  (A)--(B)--(C)--(D)--(A);
					\draw[thin,samples=150,smooth,domain=0:4]  plot (\x,{sqrt (-1*((\x)^2)+16)});
					\draw[thin,samples=150,smooth,domain=0:4] plot (\x,{sqrt  (-1*((\x-4)^2)+16)});
					 \draw[dashed] (2,0)--(2,3.464101615) node[above]{$I$};
					\fill[gray,smooth]plot[domain=0:4] (\x,{sqrt (-1*((\x)^2)+16)})--  plot[domain=4:0] (\x,{sqrt (-1*((\x-4)^2)+16)})-- cycle;
					\fill[pattern=north east lines,smooth]plot[domain=4:2] (\x,{sqrt (-1*((\x)^2)+16)})-- plot[domain=2:4] (\x,{0}); 
					 \fill[pattern=north east lines,smooth]plot[domain=0:2] (\x,{sqrt (-1*((\x-4)^2)+16)})-- plot[domain=2:0] (\x,{0});
					\foreach  \x/\g in{A/-60,B/-90,C/90,D/180}
					 \fill(\x) circle (1pt)
						 ($(\x)+(\g:3mm)$) node{$\x$};
					\draw (2,0) node[below]{$4$m} 
						(4,2) node[right]{$4$m}; 
					\draw[->] (-1,0) -- (5,0) node[below] {\scriptsize $x$}; 
					\draw[->] (0,-1) --  (0,5) node[left] {\scriptsize $y$};
					\draw (0.1,0.1) node[below  left]{\scriptsize $O$};
				\end{tikzpicture}
		\end{center}
		 Chọn hệ trục tọa độ như hình vẽ, khi đó:\\
		Đường tròn tâm $A(0;0)$ bán kính $R=4$ là  $x^2+y^2=16$.\\
		Đường tròn tâm $B(4;0)$ bán  kính $R=4$ là $(x-4)^2+y^2=16$.\\
		Ta có $A(0;0)$, $B(4,0)$, $I(2;2\sqrt{3})$.\\ 
		Xét cung tròn có phương trình $y=\sqrt{16-x^2}$.\\ 
		Phần diện tích gạch chéo: $$S_1=2\cdot  \displaystyle\int^4_2\sqrt{16-x^2}\,\mathrm{\,d}x=\dfrac{16\pi}{3}-4\sqrt{3}.$$
		 Phần diện tích màu đen:
		$$S_2=2\cdot \left(\dfrac{1}{4}\pi\cdot 4^2-\dfrac{16\pi}{3}+4\sqrt{3}\right)=\dfrac{-8\pi}{3}+8\sqrt{3}.$$ 
		Phần  diện tích còn lại: $$S_3=16-\left(\dfrac{16\pi}{3}-4\sqrt{3}+\dfrac{-8\pi}{3}+8\sqrt{3}\right)=16-\dfrac{8\pi}{3}-4\sqrt{3}.$$
		Số tiền để sơn biển quảng cáo:\\ 
		 \allowdisplaybreaks
		$\left(\dfrac{16\pi}{3}-4\sqrt{3}\right)\cdot 150\,000+\left(\dfrac{-8\pi}{3}+8\sqrt{3}\right)\cdot 100\,000+\left(16-\dfrac{8\pi}{3}-4\sqrt{3}\right)\cdot 250\,000\\ 
		 \approx 2{,}2\, \text{ triệu\ đồng}.$
	 }
\end{bt}

% nguồn: Phan2_Lop_12_Dai_so_Chuong_4
\begin{bt}%[2D4C3-2]%[ma:23e35a39]
	Một biển quảng cáo có dạng hình elip với bốn đỉnh $A_1$, $A_2$, $B_1$, $B_2$ như hình vẽ dưới. Biết chi phí để sơn phần tô đậm là $200\,000$ (đồng) và phần còn lại $100\,000$ (đồng). Biết $A_1A_2 = 8$ m, $B_1B_2 = 6$ m và tứ giác $MNPQ$ là hình chữ nhật có $MQ = 3$  m. Hỏi số tiền để sơn theo cách trên (làm tròn đến hàng phần trục, đơn vị triệu đồng) bằng
	\begin{center}
		\begin{tikzpicture}[line join = round,line cap=round,>=stealth]
				 \path
					(-4,0)  coordinate (A1)
					(4,0) coordinate  (A2)
					 (0,-3) coordinate (B1)
					 (0,3) coordinate (B2)
					 ({-2*sqrt (3)},1.5) coordinate (M)
					({2*sqrt (3)},1.5)  coordinate (N)
					 ({-2*sqrt (3)},-1.5) coordinate (Q)
					({2*sqrt (3)},-1.5)  coordinate (P);
				\fill[inner color=cyan!85!blue,outer color=blue,fill opacity=0.45]  (0,0) ellipse ({4} and {3});
				\path[left  color=blue,right color=teal, middle color=cyan]
					(M) -- (Q) -- ($(A1)-(0,1.5)$) --  ($(A1)+(0,1.5)$) -- (M);
				\path[left  color=blue,right color=teal, middle color=cyan]
					(N) -- (P) -- ($(A2)-(0,1.5)$) -- ($(A2)+(0,1.5)$) -- (N); 
				\draw (0,0) ellipse ({4}  and {3})
					(M) -- (N) -- (P) -- (Q) --  (M);
				 \fill
					(A1)  circle (1pt) node[left]{$A_1$}
					(A2)  circle (1pt) node[right]{$A_2$}
					(B1) circle (1pt)  node[below]{$B_1$}
					(B2) circle (1pt) node[above]{$B_2$} 
					(M)  circle (1pt) node[above left]{$M$}
					(N)  circle (1pt) node[above right]{$N$}
					(P)  circle (1pt) node[below right]{$P$}
					 (Q) circle (1pt) node[below left]{$Q$};
			\end{tikzpicture}
	\end{center}
	\shortans{7,3} 
	\loigiai{ 
		\begin{center}
			\begin{tikzpicture}[line join = round,line cap=round,>=stealth,scale=0.7]
					 \path
						(-4,0)  coordinate (A1)
						(4,0)  coordinate (A2)
						 (0,-3) coordinate (B1)
						(0,3) coordinate  (B2)
						({-2*sqrt (3)},1.5)  coordinate (M)
						({2*sqrt (3)},1.5) coordinate  (N)
						 ({-2*sqrt (3)},-1.5) coordinate (Q)
						({2*sqrt (3)},-1.5)  coordinate (P);
					\draw(0,0)  ellipse ({4} and {3});
					\draw 
						(M) -- (Q) -- ($(A1)-(0,1.5)$) -- ($(A1)+(0,1.5)$) -- (M); 
					\draw 
						 (N) -- (P) -- ($(A2)-(0,1.5)$) -- ($(A2)+(0,1.5)$) -- (N);
					\draw[->] 
						(-5,0) -- (5,0)  node[above]{$x$};
					 \draw[->]
						 (0,-4) -- (0,4) node[right]{$y$};
					\draw (0,0) ellipse ({4} and {3}) 
						(M) -- (N) -- (P)  -- (Q) -- (M);
					 \fill
						(A1) circle  (1pt) node[below left]{$A_1$}
						(A2)  circle (1pt) node[below right]{$A_2$}
						(B1) circle (1pt) node[below left]{$B_1$} 
						(B2) circle  (1pt) node[above left]{$B_2$}
						(M)  circle (1pt) node[above left]{$M$}
						(N)  circle (1pt) node[above right]{$N$}
						(P)  circle (1pt) node[below right]{$P$}
						 (Q) circle (1pt) node[below left]{$Q$}
						(0,0) circle  (1pt) node[below left]{$O$};
				\end{tikzpicture}
		\end{center}
		Gọi phương trình chính tắc của elip $(E)$ có dạng  $\dfrac{x^2}{a^2}+\dfrac{y^2}{b^2} = 1$.\\
		Với $\heva{&A_1A_2 = 8 = 2a\\&B_1B_2  = 6 = 2b} \Leftrightarrow \heva{&a=4\\&b=3}$ suy ra $(E) \colon \dfrac{x^2}{16} + \dfrac{y^2}{9} = 1 \Leftrightarrow y = \pm \dfrac{3}{4}\sqrt{16 - x^2}$.\\
		Diện tích của hình elip là $S_{(E)} = \pi a \cdot b = 12\pi$ m$^2$.\\ 
		Vì $MNPQ$ là hình chữ nhật và $MQ =  3 \Rightarrow M \left(x_M; \dfrac{3}{2} \right) \in (E)$ với $x_M < 0$.\\
		Suy ra $\dfrac{x_M^2}{16} +  \dfrac{1}{4} = 1 \Rightarrow x_M^2 = 12 \Rightarrow x_M = -2\sqrt{3}$.\\
		Từ đó  ta được $M\left(-2\sqrt{3}; \dfrac{3}{2} \right)$, tương tự ta được $N\left(3\sqrt{2}; \dfrac{3}{2} \right)$.\\
		Gọi $S_1$; $S_2$ lần lượt là  diện tích phần được tô màu và phần không được tô màu.\\
		Ta có $S_2 = 4 \cdot \dfrac{3}{4}\displaystyle\int^{4}_{2\sqrt{3}} \sqrt{16 - x^2}\mathrm{\,d}x  = 3\displaystyle\int^{4}_{2\sqrt{3}} \sqrt{16 - x^2}\mathrm{\,d}x$.\\
		 Đặt $x = 4\sin{t}$ ta tìm được $S_2 = 4\pi - 6\sqrt{3}$ m$^2$.\\
		Suy ra $S_1 =  S_{(E)} - S_2 = 8\pi + 6\sqrt{3}$.\\
		Gọi $T$ là tổng chi phí.  Khi đó ta có
		$T = \left(4\pi - 6\sqrt{3} \right) \cdot 100 + \left(8\pi + 6\sqrt{3} \right)  \cdot 200 \approx 7\,332\,000$ đồng.\\
		 Vậy chi phí cần để sơn theo cách trên là $7{,}3$ (triệu đồng).
	 }
\end{bt}


%%%%%%%%%%%%%%%% TỰ LUẬN %%%%%%%%%%%%%%%%

% nguồn: Phan2_Lop_12_Dai_so_Chuong_4
\begin{bt}%[2D4H3-2]%[ma:a1e8937e]
	Một vật chuyển động với tốc độ $v(t)=3t+4 \mathrm{(m/s)}$, với  thời gian $t$ tính theo giây, $t\in[0;5]$. Tính quãng đường vật đi được trong khoảng thời gian từ $t=0$ đến $t=5$.
	 \loigiai{
		Quãng đường vật đi được trong khoảng thời gian trên là\\ $s=\displaystyle \int_{0}^{5}v(t)\mathrm{\,d}x=\displaystyle  \int_{0}^{5}(3t+4)\mathrm{\,d}x=\left(\dfrac{3t^2}{2}+4t\right)\bigg|_{0}^{5}=\dfrac{75}{2}+20=\dfrac{115}{2}$.
	 }
\end{bt}

% nguồn: Phan2_Lop_12_Dai_so_Chuong_4
\begin{bt}%[2D4H3-2]%[ma:8fae938d]
	Trong xác suất thống kê thì một biến ngẫu nhiên $X$ được gọi là có hàm mật độ $f(x)$ nếu xác suất để $X$ nhận giá  trị trong $[a;b]$ là diện tích của hình phẳng giới hạn bởi đồ thị hàm số $f(x)$, đường thẳng $y=a$ và $y=b$. Đặc biệt, khi
	 $f(x)=\heva{
	& \lambda \mathrm{e}^{-\lambda x}  &,\, x\ge 0 \\
	 & 0 &,\,x<0}$
	trong đó $\lambda$ là số thực cho trước thì $X$ được gọi là biến ngẫu nhiên có phân phối mũ với tham số $\lambda$. Chẳng hạn, tuổi thọ của một bóng đèn  (đơn vị năm) là một biến ngẫu nhiên $X$ tuân theo phân phối mũ với hệ số $\lambda =\dfrac{1}{5}$. Hãy tính xác suất để một bóng đèn có tuổi thọ lớn hơn $5$ năm và nhỏ hơn $10$ năm.
	\begin{center}
		\begin{tikzpicture}[line join=round,line cap=round,>=stealth,scale=1,font=\footnotesize,y=5cm]
				\tikzset{every node/.style={outer  sep=-0.5mm}}
				\def \xmin{2} \def  \xmax{13}
				\def  \ymin{-.1} \def \ymax{.7}
				\def \fx{e^(-0.2*\x)} 
				\draw[->]  (\xmin,0)--(\xmax,0) node[below left] {\normalsize$x$};
				\foreach \x in  {5,10} \fill (\x,0) circle (1pt) node[below] {$\x$};
				\begin{scope}
					\clip  (\xmin+0.01,\ymin+0.01) rectangle (\xmax-0.01,\ymax-0.01);
					\draw[blue!80!black,samples=150,domain=2:13,smooth,variable=\x] plot  (\x,{\fx});
					\draw[pattern=north  east lines,opacity=0.5] (5,0)--plot[domain=5:10] (\x,{\fx})--(10,0)--cycle;
				\end{scope}
			\end{tikzpicture}
	\end{center}
	\loigiai{ 
		Hàm  mật độ xác suất của $X$ là
		$f(x)=\heva{ 
		& \dfrac{1}{5}\mathrm{e}^{-\tfrac{1}{5}x} &,\, x\ge  0 \\
		 & 0 &,\,x<0.}$ \\
		 Xác suất để một bóng đèn có tuổi thọ lớn hơn $5$ năm và nhỏ hơn $10$ năm là
		\begin{center}
			$P[5\le X\le 10]=\displaystyle\int_{5}^{10}\dfrac{1}{5}\mathrm{e}^{-\tfrac{1}{5}x}\mathrm{\,d}x=\left. -\mathrm{e}^{-\tfrac{1}{5}x} \right|_{5}^{10}=\dfrac{\mathrm{e}-1}{\mathrm{e}^{2}}$. 
		\end{center}
	 }
\end{bt}

% nguồn: Phan2_Lop_12_Dai_so_Chuong_4
\begin{bt}%[2D4V3-2]%[ma:85e6d389]
	Ta biết rằng hàm cầu liên quan đến giá $p$ của một sản phẩm với nhu cầu của người tiêu dùng, hàm cung liên quan đến giá $p$ của sản phẩm với mức độ sẵn sàng cung cấp  sản phẩm của nhà sản xuất. Điểm cắt nhau $(x_0;p_0)$ của đồ thị hàm cầu $p=D(x)$ và đồ thị hàm cung $p=S(x)$ được gọi là điểm cân bằng.\\
	Các nhà kinh tế gọi diện tích của hình giới hạn bởi đồ thị hàm cầu,  đường ngang $p=p_0$ và đường thẳng đứng $x=0$ là thặng dư tiêu dùng. Tương tự, diện tích của hình giới hạn bởi đồ thị của hàm cung, đường nằm ngang $p=p_0$ và đường thẳng đứng $x=0$ được gọi là thặng dư sản xuất, như trong hình bên dưới.
	\begin{flushright}
		(\textit{Theo R. Larson, Brief Calculus: An Applied Approach, 8th edition, Cengage Learning,  2009})
	\end{flushright}
	\begin{center}
		\begin{tikzpicture}[line join=round,line cap=round,>=stealth,scale=1.1,font=\footnotesize]
				\tikzset{every node/.style={outer sep=-0.5mm}} 
				 \def \xmin{-1} \def \xmax{5}
				 \def \ymin{-1} \def \ymax{6}
				\def \fx{1+0.1*(\x)^2} 
				 \def \gx{1+0.05*(\x-9)^2}
				\draw[->] (\xmin,0)--(\xmax,0) node[below  left] {\normalsize$x$};
				\draw[->] (0,\ymin)--(0,\ymax) node[below left]  {\normalsize$y$};
				\draw  (0,0) node[below left] {$O$};
				\fill (0,2.39) circle (1pt)  node[left] {$p_0$};
				\fill (3.73,0)  circle (1pt) node[below] {$x_0$};
				\fill (3.73,2.39) circle (1.25pt) node[right,xshift=2mm]{$(x_0; p_0)$}; 
				\draw (2,1.4)--(2.5,1) node[below]{Hàm cung}; 
				 \draw (1,4.2)--(2,5) node[above]{Hàm cầu};
				\begin{scope}
					\clip  (\xmin+0.01,\ymin+0.01) rectangle (\xmax-0.01,\ymax-0.01);
					\draw[blue!80!black,samples=50,domain=0:\xmax-0.01,smooth,variable=\x] plot  (\x,{\fx});
					\draw[red!80!black,samples=50,domain=0:\xmax-0.01,smooth,variable=\x] plot (\x,{\gx}); 
					\draw[fill=blue!20,opacity=0.5]  (0,1) plot[domain=0:3.73] (\x,{\fx})--(0,2.39)--cycle;
					 \draw[fill=red!20,opacity=0.5] plot[domain=0:3.73] (\x,{\gx})--(0,2.39)--cycle;
				\end{scope}
				\draw (3.73,2.39)--(4,3.5) node[above] {Điểm cân bằng}; 
				\draw[dashed,thin](0,2.39)-|(3.73,0); 
				 \draw (1.6,2.39) node[above]{Thặng dư tiêu dùng} node[below]{Thặng dư sản xuất};
			\end{tikzpicture}
	\end{center}
	Giả  sử hàm cung và hàm cầu của một loại sản phẩm được mô hình hoá bởi:\\
	Hàm cầu: $p=5-0{,}2x$ và hàm cung: $p=1+0{,}02x^2$, trong  đó $x$ là số đơn vị sản phẩm. Tìm thặng dư tiêu dùng và thặng dư sản xuất cho sản phẩm này.
	\loigiai{ 
		Phương trình hoành độ giao điểm của hai hàm là 
		\begin{eqnarray*}
			&&5-0{,}2x=1+0{,}02x^2\\ 
			&\Leftrightarrow&  0{,}02x^2+0{,}2x-4=0\\
			 &\Leftrightarrow&
			\hoac{& x=10  \\ & x=-20.}
		\end{eqnarray*}
		Với $x=10$  thỏa mãn khi đó $x_0=10$; $p_0=5-0{,}2\cdot10=3$.\\
		Thặng dư tiêu dùng là 
		\begin{center}
			$\displaystyle\int_{0}^{10}\left(  5-0{,}2x-3 \right)\mathrm{\,d}x=\int_{0}^{10}\left( 2-0{,}2x \right)\mathrm{\,d}x=10$.
		\end{center}
		Thặng dư sản suất  là
		\begin{center}
			$\displaystyle\int_{0}^{10}\left( 3-1-0{,}02x^2  \right)\mathrm{\,d}x=\int_{0}^{10}\left( 2-0{,}02x^2 \right)\mathrm{\,d}x=\dfrac{40}{3}$.
		\end{center}
	 }
\end{bt}

% nguồn: Phan2_Lop_12_Dai_so_Chuong_4
\begin{bt}%[2D4V3-2]%[ma:e66167fc]
	Cổng của một trại du lịch sinh thái có dạng parabol, lối ra vào cổng là hình chữ nhật như hình vẽ. Biết rằng, lối đi hình chữ nhật $A B C D$ có kích thước $C D=2 \mathrm{~m}$, phía ngoài lối đi hình chữ nhật được trang trí hoa văn cho phù hợp. Chi phí trang trí hoa văn là 240.000 đồng cho một mét vuông. Hỏi số tiền để trang trí hoa văn  của cổng trại du lịch sinh thái là bao nhiêu triệu đồng?
	\loigiai{ 
		Xét hệ trục  toạ độ như hình vẽ
		\begin{center}
			\begin{tikzpicture}[scale=1]
					 \path
						(-1,3) coordinate  (A)
						(1,3)  coordinate (B)
						(1,0) coordinate (C) 
						(-1,0) coordinate  (D);
					\draw[dashed] (-2,0)--(2,0); 
					 \fill[pattern=north east lines] (A)--(B)--(C)--(D)--cycle;
					\foreach \x/\y  in {A/B,B/C,C/D,D/A}
					 \draw (\x)--(\y);
					 \draw[<->] (-2,-.5)--(2,-.5) node[pos=0.5,below]{$4$m};
					 \draw[<->] (-2.5,0)--(-2.5,4) node[pos=0.5,left]{$4$m};
					 \foreach \x/\y in {A/150,B/60,C/-90,D/-90}
					 \draw[fill] (\x) circle (1pt) +(\y:.3) node {$\x$};
					\def  \f (#1){-(#1)^2+4}
					\begin{scope}
						 \draw [samples=200,smooth,domain=-2:2] plot (\x,{\f (\x)}) node [below right=0.5mm]{$ $};
					\end{scope}
				\end{tikzpicture}
		\end{center}
		Parabol có dạng  $y=a x^2+c, a<0$.\\
		 Parabol đi qua các điểm $(0; 4),(2; 0)$ nên có dạng: $y=-x^2+4$.\\
		 Ta có $x_C=1$ nên $y_B=-1^2+4=3$.\\
		Suy ra hình chữ nhật $A B C D$ có kích thước $C D=2 \mathrm{~m}; B C=3 \mathrm{~m}$.\\ 
		Gọi $S_C$ là diện tích  cổng Parabol.\\
		Khi đó, $S_C$ là diện tích hình phẳng giới hạn bởi đồ thị hàm số $y=-x^2+4$, trục hoành và các đường thẳng  $x=-2; x=2$.\\
		Diện tích phần trang trí hoa văn là:\ $S=S_C-S_{A B C D}=\displaystyle\int_{-2}^2\left(-x^2+4\right) \mathrm{d} x-2.3=\dfrac{14}{3}\left(m^2\right)$.\\ Vậy số tiền cho việc trang trí hoa  văn của cổng là: $240000 \cdot \dfrac{14}{3}=1,12$ triệu đồng.
	 }
\end{bt}

% nguồn: Phan2_Lop_12_Dai_so_Chuong_4
\begin{bt}%[2D4V3-2]%[ma:d22b1bee]
	Trên cửa sổ có dạng hình chữ nhật, họa sĩ thiết kế logo hình con cá cho một doanh nghiệp kinh doanh hải sản. Logo là hình phẳng giới hạn bởi hai parabol với các kích thước được cho trong hình sau (đơn  vị trên mỗi trục toạ độ là decimét).
	\begin{center}
		\begin{tikzpicture}[line join=round,line cap=round,>=stealth,scale=.7,font=\footnotesize]
				 \tikzset{every node/.style={outer sep=-0.5mm}}
				 \def \xmin{-6} \def \xmax{6}
				 \def \ymin{-3.5} \def \ymax{3}
				 \def \fx{-0.125*(\x)^2+2}
				\def  \gx{0.1875*(\x)^2-3}
				\draw[->] (\xmin,0)--(\xmax,0) node[below  left] {\normalsize$x$};
				 \draw[->] (0,\ymin)--(0,\ymax) node[below left] {\normalsize$y$};
				\draw  (0,0) node[below left] {$O$};
				\draw  (-5,-3)|-(4,2)|-cycle;
				\foreach \x/\y/\n in {-5/-1.125/A,-5/1.6875/B,-4/0/C,0/2/D,4/0/E,0/-3/G,2/1.4/M,2/-2.25/N,2.25/0.6/F} \coordinate (\n)  at (\x,\y);
				 \foreach \t/\g in {A/180,B/180,C/-90,D/45,E/45,G/45}{
					\fill  (\t) circle (1pt) node[shift={(\g:8pt)}]{$ \t $};
				 }
				\draw[<->] (-5,-4)--(-4,-4)  node[midway,above]{$1$\,dm};
				\draw[<->] (-4,-4)--(0,-4)  node[midway,above]{$4$\,dm};
				\draw[<->] (0,-4)--(4,-4) node[midway,above]{$4$\,dm}; 
				\draw[<->]  (5,0)--(5,2) node[midway,right]{$2$\,dm};
				 \draw[<->] (5,-3)--(5,0) node[midway,right]{$3$\,dm};
				\draw pic[thin,draw,angle radius=2cm] {angle= M--E--N}; 
				\draw  (F) circle (3mm);
				\draw  (F) circle (1.5mm);
				\draw (-2,1.55) node[below right]{$y=f (x)$}; 
				\draw (-3,-1.3) node[above  right]{$y=g (x)$};
				\begin{scope}
					\clip  (\xmin+0.01,\ymin+0.01) rectangle (\xmax-0.01,\ymax-0.01);
					\draw[samples=150,domain=-5:4,smooth,variable=\x] plot  (\x,{\fx});
					\draw[samples=150,domain=-5:4,smooth,variable=\x]  plot (\x,{\gx});
				\end{scope}
			\end{tikzpicture}
	\end{center}
	\begin{listEX}
		\item Lập phương trình các parabol  $y=f(x)$ và $y=g(x)$.
		\item Tính diện tích của  logo.
		\item Logo chỉ cho  phép $50\%$ lượng ánh sáng đi qua. Lượng ánh sáng đi qua toàn bộ cửa sổ sau khi làm logo sẽ giảm bao nhiêu phần trăm (làm tròn kết quả đến hàng phần mười)?
	\end{listEX}
	 \loigiai{
		\begin{listEX}
			\item Từ đồ thị ta thấy parabol $y=f(x)=ax^2+bx+c$ có  đỉnh tại điểm $D(0;2)$ nên $c=2$ và đi qua điểm $E(4;0)$ nên $0=16a+4b+2$.\\
			 Ta có $b=0$ (do parabol đối xứng quá trục tung) nên $16a=-2\Rightarrow a=-\dfrac{1}{8}$ nên
			\begin{center}
				$y=f(x)=-\dfrac{1}{8}x^2+2$. 
			\end{center}
			Parabol  $y=g(x)=mx^2+nx+p$ có đỉnh tại điểm $G(0;-3)$ nên $p=-3$ và đi qua điểm $E(4;0)$ nên $0=16m+4n-3$.\\
			Ta có $n=0$ (do parabol đối xứng  quá trục tung) nên $16a=3\Rightarrow m=\dfrac{3}{16}$ nên
			\begin{center}
				$y=g(x)=\dfrac{3}{16}x^2-3$. 
			\end{center}
			\item Diện tích của  logo là
			\begin{center}
				$S=\displaystyle\int_{-5}^{4}\left| f(x)-g(x) \right|\mathrm{\,d}x=\int_{-5}^{4}\left| -\dfrac{5}{16}x^2+5  \right|\mathrm{\,d}x=\dfrac{1345}{48}$ (dm$^2$).
			\end{center}
			\item Diện tích cửa  sổ là $5\cdot9=45$ (dm$^2$)\\
			Lượng ánh sáng đi qua toàn bộ cửa sổ  sau khi làm logo sẽ giảm
			\begin{center}
				$\dfrac{45-50\%\cdot\dfrac{1345}{48}}{45}\cdot100\approx68{,}86$ ($\%$). 
			\end{center}
		\end{listEX}
	 }
\end{bt}

% nguồn: Phan2_Lop_12_Dai_so_Chuong_4
\begin{bt}%[2D4V3-2]%[ma:f75a10b1]
	Năng lượng gió trên đất liền là một trong những công nghệ năng lượng tái tạo đang được phát triển ở quy mô toàn cầu. Năng lượng gió không trực tiếp phát thải khí nhà kính, không thải ra môi trường các chất ô nhiễm khác, cũng như không tiêu thụ nước để làm mát cho các  nhà máy. Các turbine gió thường có ba cánh quay trên một trục ngang, lấy động năng từ quá trình di chuyển dòng không khí (gió) để chuyển đổi thành điện năng thông qua một máy phát điện được kết nối với lưới điện. Hình thang cong (tô màu vàng) trong hình vẽ mô tả một phần mặt cắt đứng của cánh turbine, được giới hạn bởi các đường thẳng $x=2, x=25$, trục $Ox$ và đồ thị hàm số $$y = f(x) = -\dfrac{1}{800}(x^3 - 33x^2+120x-400).$$
	 Hãy tính diện tích hình thang cong đó.
	\begin{center}
		\begin{tikzpicture}[scale=0.8,font=\footnotesize,line join=round,line cap=round,>=stealth,xscale=0.5,yscale=0.5]
				\def\xt{-2} 
				\def\xp{27} 
				\def\yd{-1} 
				\def\yt{4.5} 
				\begin{scope}
					\clip (\xt,\yd) rectangle (\xp,\yt); 
					\draw[blue,fill=orange](2,0)--(2,0.355).. controls (8,0.355) and (14,4)..  (20,4) node[midway,sloped,above]{\fontsize{12}{10}\selectfont $y=f (x)$}
						.. controls (21,4)  and (22.5,4).. (25,3)--(25,0) node[midway,right]{$3$ m}--cycle;
					\foreach \x in {2,5,10,15,20,25} \draw (\x,-0.1)  node[below]{$\x$} --++(90:0.2);
				\end{scope}
				\draw[->](\xt,0) -- (0,0)  node[below left]{$O$} -- (\xp,0) node[below]{$x$(m)};
				 \draw[->](0,\yd)--(0,\yt) node[right]{$y$(m)};
				\draw[dashed]  (2,0.355)--(0,0.355) node[left]{$0{,}355$}
					(20,0)--(20,4) node[midway,right]{$4$ m}; 
			\end{tikzpicture}
	\end{center}
	\loigiai{Diện tích của  hình thang cong được tô màu vàng là:
		$$\begin{aligned}
		I & =\int_2^{25}-\frac{1}{800}\left(x^3-33 x^2+120 x-400\right) \mathrm{d}  x \\
		& =-\frac{1}{800}\left(\int_2^{25} x^3 \mathrm{~d} x-33 \int_2^{25} x^2 \mathrm{~d} x+120 \int_2^{25}  x \mathrm{~d} x-400 \int_2^{25} \mathrm{~d} x\right) \\
		&  =-\frac{1}{800}\left(\frac{x^4}{4}\bigg|_2 ^{25}-11 x^3\bigg|_2 ^{25}+60 x^2\bigg|_2 ^{25}- 400 x\bigg|_2 ^{25}\right)=\frac{184\,299}{3200}\left(\mathrm{~m}^2\right) .
\end{aligned}$$}
\end{bt}

% nguồn: Phan2_Lop_12_Dai_so_Chuong_4
\begin{bt}%[2D4V3-2]%[ma:7a322d95]
	Đơn đặt hàng của nhà máy cho một loại máy điều hoà không khí là khoảng $6\,000$ chiếc mỗi tuần khi giá là $331$ USD/chiếc và khoảng $8\,000$ chiếc  mỗi tuần khi giá là $303$ USD/chiếc. Hàm cung được cho bởi $p=0{,}0275 x$, trong đó $x$ là số lượng máy điều hoà được bán với giá $p$ USD một chiếc. Tìm thặng dư tiêu dùng và thặng dư sản xuất (giả sử hàm cầu là hàm bậc nhất).
	\loigiai{ 
		Hàm cầu có dạng $p=ax+b$. Do đơn đặt hàng của nhà máy cho một loại máy điều hoà không khí là $6\,000$ chiếc mỗi tuần khi giá là $331$ USD/chiếc và $8\,000$ chiếc  mỗi tuần khi giá là $303$ USD/chiếc nên ta có
		\begin{center}
			 $\heva{& 331 = 6\,000a + b \\
			& 303 =  8\,000a + b}
			\Leftrightarrow  \heva{&a=-0{,}014 \\&b=415.}$
		\end{center}
		Vậy hàm cầu là $p=-0{,}014 x+415$.\\ 
		Xét phương trình $-0,014 x+415=0,0275 x \Leftrightarrow  x=10\,000$ khi đó $p=275$.\\
		Thặng dư tiêu dùng  là
		\begin{center}
			$\displaystyle\int_{0}^{10\,000}\left( -0{,}014x+415-275 \right)\mathrm{\,d}x=700\,000$  (USD).
		\end{center}
		 Thặng dư sản xuất là
		\begin{center}
			$\displaystyle\int_{0}^{10\,000}\left( 275-0{,}0275x \right)\mathrm{\,d}x=1\,375\,000$ (USD). 
		\end{center}
	 }
\end{bt}

% nguồn: Phan2_Lop_12_Dai_so_Chuong_4
\begin{bt}%[2D4V3-2]%[ma:a9ff40a1]
	Mặt cắt của một  cửa hầm có dạng là hình phẳng giới hạn bởi một parabol và đường thẳng nằm ngang như hình sau. Tính diện tích của cửa hầm.
	\begin{center}
		\begin{tikzpicture}[line join=round,line cap=round,>=stealth,scale=.7,font=\footnotesize]
				\tikzset{every  node/.style={outer sep=-0.5mm}}
				 \def \xmin{-4} \def \xmax{4}
				\def \ymin{-.5}  \def \ymax{7}
				 \def \fx{-2/3*(\x)^2+6}
				 \draw[fill=blue!20] (-4,0)|-(4,7)|-cycle;
				 \draw[<->] (-3,-1)--(3,-1) node[midway,above]{$6$\,m};
				\draw[<->] (5,0)--(5,6)  node[midway,right]{$6$\,m};
				\draw[dashed]  (-3,0)--(-3,-1)
					 (3,0)--(3,-1)
					(4,0)--(5,0) 
					 (0,6)--(5,6);
				\begin{scope}
					\clip (\xmin+0.01,\ymin+0.01) rectangle  (\xmax-0.01,\ymax-0.01);
					\draw[samples=150,domain=-3:3,smooth,variable=\x] plot  (\x,{\fx});
					 \draw[fill=cyan,opacity=0.5] plot[domain=-3:3] (\x,{\fx})--cycle;
				\end{scope}
			\end{tikzpicture}
	\end{center}
	 \loigiai{
		Chọn hệ trục tọa độ $Oxy$ như hình vẽ, đơn vị trên mỗi  trục là mét.
		\begin{center}
			\begin{tikzpicture}[line join=round,line cap=round,>=stealth,scale=.7,font=\footnotesize]
					\tikzset{every  node/.style={outer sep=-0.5mm}}
					\def \xmin{-4} \def  \xmax{4}
					\def \ymin{-.75} \def  \ymax{7}
					\def \fx{-2/3*(\x)^2+6} 
					 \draw[->] (\xmin,0)--(\xmax,0) node[below left] {\normalsize$x$};
					\draw[->] (0,\ymin)--(0,\ymax) node[below left] {\normalsize$y$}; 
					 \draw (0,0) node[below left] {$O$};
					\draw (-4,0)|-(4,7)|-cycle; 
					\fill (3,0)  circle (1pt) node[above right]{$3$};
					\fill  (-3,0) circle (1pt) node[above left]{$-3$};
					\fill (0,6) circle (1pt) node[above left]{$6$}; 
					\draw[<->]  (-3,-1)--(3,-1) node[midway,below]{$6$\,m};
					\draw[<->] (5,0)--(5,6)  node[midway,right]{$6$\,m};
					\draw[dashed] (-3,0)--(-3,-1) 
						(3,0)--(3,-1) 
						 (4,0)--(5,0)
						 (0,6)--(5,6);
					\begin{scope}
						\clip (\xmin+0.01,\ymin+0.01) rectangle  (\xmax-0.01,\ymax-0.01);
						\draw[samples=150,domain=-3:3,smooth,variable=\x]  plot (\x,{\fx});
						\draw[fill=cyan,opacity=0.5] plot[domain=-3:3]  (\x,{\fx})--cycle;
					\end{scope}
				\end{tikzpicture}
		\end{center}
		Gọi parabol $(P)\colon y=ax^2+bx+c$.\\ 
		Ta có các điểm $(3;0)$; $(-3;0)$; $(0;6)$  thuộc $(P)$, suy ra
		\begin{align*}
			\heva{&9a+3b+c=0\\&9a-3b+c=0\\&c=6}\Leftrightarrow\heva{&a=-\dfrac{2}{3}\\&b=0\\&c=6.} 
		\end{align*}
		Suy ra parabol là  $y=-\dfrac{2}{3}x^2+6$.\\
		Diện tích  của cửa hầm là
		\begin{center}
			$S=\displaystyle\int_{-3}^{3}\left(  -\dfrac{2}{3}x^2+6 \right)\mathrm{\,d}x=\left.\left( -\dfrac{2}{9}x^3+6x \right)\right|^3_{-3}=24$ (m$^2$).
		\end{center}
	 }
\end{bt}

% nguồn: Phan2_Lop_12_Dai_so_Chuong_4
\begin{bt}%[2D4V3-2]%[ma:9485776f]
	Một công ty quảng cáo muốn làm một bức tranh trang trí như phần $MNEIF$ được tô đậm trong hình vẽ bên dưới ở chính giữa của một bức tường hình chữ nhật $ABCD$ có $BC=6$ m, $CD=12$ m. Biết $MN=4$ m; cung $EIF$ có hình parabol với đỉnh $I$ là trung điểm của cạnh $AB$ và đi qua hai điểm $C$, $D$. Kinh phí làm bức tranh là $1\,200\,000$ đồng/m$^2$. Hỏi công ty đó cần bao nhiêu tiền để làm bức tranh? (\textit{Kết quả làm tròn đến  hàng nghìn}).
	 \loigiai{
		\begin{center}
			\begin{tikzpicture}[scale=0.7,line join=round,line cap=round,>=stealth]

					\fill[pattern=horizontal  lines gray]plot[domain=-2:2](\x,{-1/6*((\x)^2)+0*\x+6})--plot[domain=2:-2](\x,{0})--cycle;
					 \draw[smooth,samples=200
						,domain=-2:2]  plot (\x,{-1/6*((\x)^2)+0*\x+6});
					\coordinate [label=below:$M$](M)  at (-2,0);
					 \coordinate [label=above:$I$](I) at (0,6);
					\coordinate [label=below:$N$](N)  at (2,0);
					\coordinate [label=right:$E$](E)  at (2,16/3);
					\coordinate  [label=left:$F$](F) at (-2,16/3);
					\coordinate [label=below:$C$](C) at (6,0); 
					\coordinate  [label=above:$B$](B) at (6,6);
					\coordinate [label=above:$A$](A) at  (-6,6);
					 \coordinate [label=below:$D$](D) at (-6,0);
					\foreach \diem in {M,N,E,F,A,B,C,D,I} \fill (\diem) circle (1.5pt); 
					 \draw(A)--(B)--(C)--(D)--(A)
						 (F)--(M)
						(E)--(N); 
					\draw[<->,dashed]($(M)- (0,1)$)--($(N)-  (0,1)$);
					\draw  (0,-1) node [below]{$4$ m};
					\draw[<->,dashed]($(A)+ (0,1)$)--($(B)+  (0,1)$);
					 \draw (7,3) node [right]{$6$ m};
					 \draw[<->,dashed]($(B)+ (1,0)$)--($(C)+ (1,0)$);
					 \draw (0,7) node [above]{$12$ m};
				\end{tikzpicture}
		\end{center}Gọi $O$ là trung điểm cạnh $MN$ và trùng với gốc toạ độ, suy ra $M(-2;0)$; $N(2;0)$.\\
		Phương trình parapol đỉnh $I(0;6)$ và đi qua hai điểm $D(-6;0)$; $C(6;0)$  là $$(P)\colon y=-\dfrac{1}{6} x^2+6.$$
		Diện tích hình phẳng giới hạn bởi $(P)\colon y=-\dfrac{1}{6} x^2+6  $; $ y=0$; $x=-2$; $ x=2$ là $S$.\\
		Khi đó $$S=\displaystyle\int_{-2}^2\left|-\dfrac{1}{6} x^2+6\right| \mathrm{\,d} x=\dfrac{208}{9}~\left(\mathrm{m}^2\right).$$ 
		Vậy công ty đó cần số tiền để làm bức tranh là $\dfrac{208}{9}\cdot  1\,200\,000 \approx 27\,733\,000$ (đồng).
	 }
\end{bt}
